\documentclass[11pt]{article}
\usepackage[T1]{fontenc}
\usepackage[utf8]{inputenc}
\usepackage{lmodern}
\usepackage[a4paper,margin=30mm]{geometry}
\usepackage{amsmath,amssymb,amsthm,mathtools}
\usepackage{microtype}
\usepackage{enumitem,booktabs,tabularx,array,longtable,mathrsfs,graphicx,flafter}
\usepackage{pict2e}
\usepackage[numbers,sort&compress]{natbib}
\usepackage{xcolor}
\usepackage{xurl}
\usepackage[hypertexnames=false,colorlinks=true,linkcolor=blue!45!black,citecolor=blue!45!black,urlcolor=blue!45!black]{hyperref}
\usepackage[capitalise,noabbrev,nameinlink]{cleveref}
\ifdefined\pdfinfoomitdate\pdfinfoomitdate=1\fi
\ifdefined\pdftrailerid\pdftrailerid{}\fi
\ifdefined\pdfsuppressptexinfo\pdfsuppressptexinfo=15\fi
\newtheorem{theorem}{Theorem}[section]
\newtheorem{lemma}[theorem]{Lemma}
\newtheorem{proposition}[theorem]{Proposition}
\newtheorem{corollary}[theorem]{Corollary}
\newtheorem{claim}[theorem]{Claim}
\theoremstyle{definition}
\newtheorem{definition}[theorem]{Definition}
\newtheorem{convention}[theorem]{Convention}
\theoremstyle{remark}
\newtheorem{remark}[theorem]{Remark}
\newcommand{\Q}{\mathbb Q}
\newcommand{\Z}{\mathbb Z}
\newcommand{\R}{\mathbb R}
\newcommand{\C}{\mathbb C}
\newcommand{\N}{\mathbb N}
\newcommand{\F}{\mathbb F}
\newcommand{\Qtwo}{\mathbb Q_2}
\newcommand{\Ztwo}{\mathbb Z_2}
\DeclareMathOperator{\Gal}{Gal}
\DeclareMathOperator{\Sel}{Sel}
\DeclareMathOperator{\Sha}{Sha}
\DeclareMathOperator{\corank}{corank}
\DeclareMathOperator{\rank}{rank}
\DeclareMathOperator{\ord}{ord}
\DeclareMathOperator{\length}{length}
\DeclareMathOperator{\Fr}{Fr}
\DeclareMathOperator{\Frob}{Frob}
\DeclareMathOperator{\Cl}{Cl}
\DeclareMathOperator{\Pic}{Pic}
\DeclareMathOperator{\Hom}{Hom}
\DeclareMathOperator{\Ext}{Ext}
\DeclareMathOperator{\Aut}{Aut}
\DeclareMathOperator{\End}{End}
\DeclareMathOperator{\Spec}{Spec}
\DeclareMathOperator{\coker}{coker}
\DeclareMathOperator{\im}{im}
\DeclareMathOperator{\res}{res}
\DeclareMathOperator{\cor}{cor}
\DeclareMathOperator{\GL}{GL}
\DeclareMathOperator{\SL}{SL}
\DeclareMathOperator{\tr}{tr}
\DeclareMathOperator{\diag}{diag}
\DeclareMathOperator{\supp}{supp}
\DeclareMathOperator{\Nrd}{Nrd}
\setlist{topsep=4pt,itemsep=2pt,parsep=0pt}
\setlist[itemize]{label=$\bullet$}
\setcounter{tocdepth}{1}
\emergencystretch=2em
\DeclareMathOperator{\Tr}{Tr}
\DeclareMathOperator{\Nm}{Nm}
\DeclareMathOperator{\adj}{adj}
\DeclareMathOperator{\id}{id}
\theoremstyle{definition}
\newtheorem{hypothesis}[theorem]{Hypothesis}
\newtheorem{assumption}[theorem]{Assumption}
\newtheorem{conjecture}[theorem]{Conjecture}
\newtheorem{example}[theorem]{Example}
\numberwithin{equation}{section}
\title{Goldfeld's analytic density conjecture\\and the $2$-converse for elliptic curves}
\author{OpenAI}
\date{September 23, 2026}
\hypersetup{pdftitle={Goldfeld's analytic density conjecture and the 2-converse for elliptic curves},pdfauthor={OpenAI},pdfsubject={Quadratic twists and the low-corank 2-converse}}
\begin{document}
\maketitle
\begin{abstract}
We prove Goldfeld's analytic density conjecture: for every elliptic curve
$E$ over $\Q$, the quadratic twists of $E$ with analytic rank zero and one
each have density $1/2$ among signed squarefree twist parameters ordered by
absolute value. We also prove the low-corank $2$-converse: if the
$2^\infty$-Selmer corank of $E$ is zero or one, then it equals the analytic
and Mordell--Weil ranks, and the Tate--Shafarevich group is finite.
\end{abstract}
\tableofcontents
\clearpage
\section{Introduction}\label{sec:introduction}

Let $E$ be an elliptic curve over $\Q$. Write
\[
 a(E)=\ord_{s=1}L(E,s),\qquad r(E)=\rank_{\Z}E(\Q),\qquad
 c_2(E)=\corank_{\Ztwo}\Sel_{2^\infty}(E/\Q).
\]
The Selmer group uses the local Kummer conditions at every place.
Its corank includes both the Mordell--Weil rank and any divisible
$2$-primary contribution from the Tate--Shafarevich group.
Modularity supplies the analytic continuation and functional equation
of $L(E,s)$ \citep{BreuilConradDiamondTaylor2001}.

\begin{theorem}[The $2$-converse]\label{thm:main-converse}
For every elliptic curve $E/\Q$ with $c_2(E)\in\{0,1\}$,
\[
 a(E)=r(E)=c_2(E),\qquad \Sha(E/\Q)\text{ is finite}.
\]
\end{theorem}

There is no restriction on reduction type, complex multiplication,
rational torsion, or rational isogenies. The hypothesis concerns the full
$2$-power Selmer corank; a Mordell--Weil rank bound alone does not give it.

For a nonzero squarefree integer $d$, let $E^{(d)}$ denote the quadratic
twist of $E$, and set
\[
 \mathcal D(X)=\{d\in\Z:0<|d|\le X,\ d\text{ squarefree}\}.
\]
Thus positive and negative parameters are counted together and ordered
by absolute value. A residue class can fix the sign of the functional
equation; the half-density statement below concerns the full set
$\mathcal D(X)$.

\begin{theorem}[Goldfeld's analytic density conjecture]\label{thm:main-density}
For every elliptic curve $E/\Q$ and each $j\in\{0,1\}$,
\[
 \lim_{X\to\infty}
 \frac{\#\{d\in\mathcal D(X):a(E^{(d)})=j\}}{\#\mathcal D(X)}
 =\frac12.
\]
Consequently the twists of analytic rank at least two have density zero.
For a density-one set of signed squarefree parameters $d$, the analytic
and Mordell--Weil ranks of $E^{(d)}$ agree and its whole
Tate--Shafarevich group is finite.
\end{theorem}

The theorem follows by combining the pointwise $2$-converse with
Smith's distribution of full $2$-power Selmer coranks
\citep[Theorem~1.1]{Smith2025}. Section~\ref{sec:conclusions} gives
the deduction, including the passage from integer to squarefree
parameters. The main work of this paper is the pointwise converse.

\subsection{Consequences and companion results}

The converse has a finite-descent consequence. Put
\[
 d_2(E)=\dim_{\F_2}\Sel_2(E/\Q)-\dim_{\F_2}E(\Q)[2].
\]

\begin{corollary}[A finite Selmer criterion]\label{cor:finite-descent}
If $d_2(E)\in\{0,1\}$, then
\[
 a(E)=r(E)=c_2(E)=d_2(E),\qquad
 \Sha(E/\Q)\text{ is finite},\qquad \Sha(E/\Q)[2^\infty]=0.
\]
\end{corollary}

Thus a finite $2$-descent can certify the analytic rank and the vanishing
of the entire $2$-primary Tate--Shafarevich group in these two cases.
The proof uses the alternating Cassels--Tate pairing to pass from
$d_2(E)$ to the full Selmer corank, and is given in
Section~\ref{sec:conclusions}.

The density theorem resolves the frequency assertion for analytic ranks
zero and one. The mean analytic-rank assertion requires additional
control: a set of density zero can still contribute to the rank-weighted
average. The separate companion
\citet[Theorems~1.1 and~1.2]{OpenAIMeanAnalyticRank2026} proves the
required analytic tail estimate and, together with
Theorem~\ref{thm:main-density}, obtains mean analytic rank $1/2$
in the same signed squarefree ordering.

A subsequent companion \citep[Theorem~1.1]{OpenAIAllPrimeSelmer2026}
extends the pointwise converse to every prime $p$: full $p$-power Selmer
corank $r\in\{0,1\}$ forces analytic and Mordell--Weil ranks $r$
and finiteness of the whole Tate--Shafarevich group. Its proof uses the
$2$-converse and analytic twist densities established here.

The rank and finiteness assertions are distinct from the
Birch--Swinnerton-Dyer leading-coefficient formula.
The pointwise two-primary formula of
\citet[Theorem~1.1]{OpenAITwoPrimaryBSD2026}, combined with the
Selmer-corank densities in the proof of Theorem~\ref{thm:main-density},
gives the exact two-primary formula on a density-one set of signed
squarefree twists. Their common analytic and Mordell--Weil ranks are
zero and one with density $1/2$ each. This consequence is
\citet[Corollary~16.1]{OpenAITwoPrimaryBSD2026}; its leading-coefficient
assertion concerns the valuation at $2$.

\subsection{Goldfeld's conjecture and analytic nonvanishing}

Birch and Swinnerton-Dyer's comparisons of rational points with central
$L$-values led to the conjectural equality of algebraic and analytic rank
\citep{BirchSwinnertonDyer1965}. The conjecture also predicts finiteness
of $\Sha(E/\Q)$ and a formula for the leading coefficient
\citep[\S1, pp.~415--421]{Tate1966}.
Goldfeld's original 1979 formulation asks for the mean analytic order
in a family parametrized by quadratic-field discriminants
\citep[p.~113, Conjecture~(B)]{Goldfeld1979}.
The density formulation predicts that analytic ranks zero and one each
occur half the time. Our counting convention throughout is the signed
squarefree set $\mathcal D(X)$.

The analytic study of quadratic twists produced both nonzero central
values and simple central zeros. Bump--Friedberg--Hoffstein and
Murty--Murty supplied auxiliary twists with the nonvanishing needed in
Heegner-point arguments; Iwaniec and Perelli--Pomyka\l a developed
quantitative estimates for central derivatives
\citep{BFH1990,MurtyMurty1991,Iwaniec1990,PerelliPomykala1997}.
Ono--Skinner obtained a lower bound of order $X/\log X$ for the
number of nonvanishing central values among twists by fundamental
discriminants $D$ with $0<|D|\le X$
\citep[Corollary~3]{OnoSkinner1998}.
More recently, Kumar--Mallesham--Sharma--Singh sharpened derivative
nonvanishing estimates through second moments, and
Kriz--Nordentoft improved central-value lower bounds for curves with
no rational two-torsion
\citep[Theorem~1]{KumarEtAl2024}
\citep[Theorem~1.1 and Remark~1.2]{KrizNordentoft2025}.
These quantitative results yield many twists without by themselves
giving a positive proportion for every fixed curve.

Positive proportions were obtained in special families. James proved
analytic-rank-zero proportions for particular curves, and Vatsal proved
positive proportions of both analytic ranks zero and one for $X_0(19)$
\citep[Theorem~1]{James1998} \citep[Theorem~II]{Vatsal1998}.
Kriz--Li proved both positive proportions for every rational elliptic
curve admitting a rational $3$-isogeny
\citep[Theorem~1.5]{KrizLi2019}.
In the congruent-number family, Tian's Heegner-point work gives rank
one and odd-order Tate--Shafarevich group in specified subfamilies
\citep{Tian2014}. Genus-period and genus-point congruences, followed by
Selmer statistics, give further analytic density results
\citep{TianYuanZhang2017,SmithCongruent2016}.
These arguments connect coefficient congruences and arithmetic
statistics with analytic nonvanishing, a connection also central to
the present proof.

The statistical input used here has its own progression.
Smith first established higher $2$-power Selmer distributions for curves
with full rational two-torsion and no rational cyclic subgroup of order
four \citep{Smith2017}. His subsequent work treated the irreducible
residual case and further rational two-isogeny cases
\citep{SmithDistributionI,SmithDistributionII}.
The unrestricted theorem gives full Selmer coranks zero and one with
density $1/2$ each for every rational elliptic curve
\citep[Theorem~1.1]{Smith2025}. Its analytic corollary assumes
Birch--Swinnerton-Dyer \citep[Corollary~1.2]{Smith2025}.
We use the Selmer theorem itself and supply the pointwise analytic
implication through Theorem~\ref{thm:main-converse}.

\subsection{The Selmer converse problem}

Gross--Zagier and Kolyvagin established the forward implication:
analytic rank at most one gives equality with the Mordell--Weil rank
and finiteness of the whole Tate--Shafarevich group
\citep{GrossZagier1986,Kolyvagin1990}; see also the combined formulation
in \citet[\S1]{Skinner2020}.
The Gross--Zagier formula relates a central derivative to the height
of a Heegner point, and Kolyvagin's derived classes control the Selmer
group. A converse must produce analytic nonvanishing from Selmer
information.

Cyclotomic Iwasawa theory gave rank-zero converses at good ordinary
odd primes in the work of Skinner--Urban, under residual
irreducibility and ramification hypotheses
\citep{SkinnerUrban2014}.
Skinner and Wei Zhang proved rank-one converse criteria through
Heegner points and derived classes, under their respective reduction,
residual-image and local hypotheses \citep{Skinner2020,Zhang2014}.
Burungale--Castella--Skinner subsequently removed the auxiliary
ramification conditions in the good ordinary, residually irreducible
case for $p>3$ \citep[Theorem~1.2.2 and the discussion following
Remark~1.2.3]{BurungaleCastellaSkinner2025}.
At good supersingular primes $p>3$, Castella--Wan obtained the rank-one
Selmer-corank implication for semistable curves
\citep[Theorem~A]{CastellaWan2024}.

Reducible residual representations call for a different Iwasawa-theoretic
comparison. Castella--Grossi--Lee--Skinner proved the corank-zero and
corank-one converses at good ordinary Eisenstein primes $p>2$, with
local isogeny character neither trivial nor cyclotomic
\citep[Theorem~E]{CastellaGrossiLeeSkinner2022}.
Keller--Yin treated potentially good ordinary reduction at odd
Eisenstein primes without that character exclusion
\citep[Theorem~B]{KellerYin2024}.
These results enlarge the allowed residual representations while
retaining conditions on the prime and its reduction type.

The CM theory includes substantial results at the prime two.
Burungale--Tian proved the rank-one converse at good ordinary primes
$p>3$ \citep{BurungaleTian2020};
Burungale--Castella--Skinner--Tian included the small ordinary primes
under an additional conductor hypothesis
\citep[Theorem~A]{BurungaleCastellaSkinnerTian2022}.
Ressler and Yu obtained the small-prime good-ordinary result without
that conductor condition; see Yu's Theorem~1 and the attribution in
Burungale--Castella--Skinner--Tian, Remark~C
\citep{Yu2021,BurungaleCastellaSkinnerTian2022}.
Burungale--Skinner's combined ordinary formulation includes $p=2$,
and their good supersingular CM theorem also permits $p=2$ under an
additional surjectivity condition on Selmer localization at $p$
\citep[Theorems~A.3 and~A.1]{BurungaleSkinner2024}.

In corank zero, Burungale--Tian's theorem covers every rational CM
elliptic curve at every prime, without a reduction restriction
\citep[Theorem~1.1]{BurungaleTian2026}.
Combined with Smith's unrestricted distribution, it therefore already
gives the analytic-rank-zero half for every CM twist family.
Their Theorem~1.2 gives rank zero for almost all positive squarefree
$n\equiv1,2,3\pmod8$ in the congruent-number family
$ny^2=x^3-x$.
Kriz's preprint states a converse for curves with CM by the full ring
of integers at the unique ramified prime of the CM field, and derives
analytic densities for restricted CM families, including congruent-number
curves \citep[Theorems~1.1 and~10.17]{Kriz2022}.
Its dyadic pointwise statement concerns the CM fields $\Q(i)$ and
$\Q(\sqrt{-2})$.
Together with Smith's distribution theorem, the stated dyadic converse
would yield both half-density conclusions for quadratic twists of every
curve with CM by either of these maximal orders.
Theorem~\ref{thm:main-converse} treats the full $2$-power Selmer corank
for every rational elliptic curve, in both ranks zero and one.

\subsection{Proof strategy}

Fix the curve under consideration. We construct finite binary families
of twists
\[
 h_x=\prod_{q\in\mathcal Q}(q^*)^{\lambda_q(x)},\qquad
 q^*=(-1/q)q,\qquad x\in\F_2^b,
\]
where $\mathcal Q$ is a finite set of distinct good odd primes and the
$\lambda_q$ are linear forms. The parameter at $x=0$ is $1$.
The families have the same local twisting classes at every fixed bad
place. Their nonzero vertices have analytic order zero or one, according
to the corank of the fixed curve. A scalar determinant coordinate
attached to the relevant class and Selmer complex has a nonzero central
specialization there, with uniformly bounded $2$-adic valuation.
Bounded integral interpolation then recovers nonvanishing at the missing
vertex $0$.

At corank zero, the interpolated class is a Beilinson--Kato class and
its central specialization detects $L(E,1)$ \citep{Kato2004}.
At corank one, it is a Heegner class over an imaginary quadratic field
$K=\Q(\sqrt{k})$,
whose height detects
\[
 \left[L(E,s)L(E^{(k)},s)\right]'_{s=1}.
\]
An analytically nonvanishing companion $E^{(k)}$ makes the Selmer
corank over $K$ equal to one. The argument retains integral Kummer
conditions at $2$ and every bad prime. The rational Iwasawa input is
used in its stated scope; the uniform integral estimates at $2$ are
proved in Sections~\ref{sec:heegner-interpolation} and
\ref{sec:cy:interpolation}. The Heegner construction is related to
Howard's Kolyvagin-system framework \citep{Howard2012}; the integral
bounds required here are established with the local conditions and
uniformity stated in this paper.

The finite families are built from modular coefficients. Waldspurger's
formula supplies the even-sign coefficients \citep{Waldspurger1981}.
In odd sign, reductions
of genus Heegner sums give weighted ternary theta coefficients at
increasing finite precisions. A uniform lower bound allows a least
normalized valuation to be selected. In the even case this is a minimum
among the coefficients of one form. In the odd case it is a minimum
over bounded-weight sequences as the precision increases; each finite
collection of tests is then realized at a sufficiently high actual
precision. A test attaining the minimum detects the required analytic
order and gives a uniform upper bound. Weight-two Hecke traces turn these coefficient
tests into contraction and deletion rules for finite configurations
of primes.

There are two exhaustive residual cases. If $E[2]$ is reducible, its
invariant line is a rational point of order two. Graphs encode prescribed
quadratic residue symbols between auxiliary primes. The resulting finite
construction makes two separate coefficient tests nonzero together at
every nonzero binary vertex. The odd argument fixes one companion field
before the dimension grows. In this branch, the even construction first
supplies the lower bound needed to normalize the odd coefficient test.
If $E[2]$ is irreducible, local Selmer equations give matrices over
$\F_2$. Subtracting a correction determined by finitely many local types
makes these matrices symmetric; their nullities still give a lower
bound on the Selmer dimension up to a fixed cost. Removing one selected prime or pair of primes, or retaining
it with either of two allowed local choices, gives three configurations.
Their coefficient tests satisfy the same sum-zero relation as the
corresponding matrix determinants. This relation does not identify the
coefficient test with the determinant.
We group the configurations into blocks and arrange that the matrix
entries between distinct blocks vanish.
A nonzero coefficient test bounds the matrix kernel, and hence the number
of singular blocks in the finite configuration. Fixing a configuration
with a maximal number of such blocks lets us arrange nonzero tests at
every nonzero binary vertex. An additional alternating-matrix construction
supplies a companion with uniformly bounded prime count and Selmer length.

Two parts of the argument have uses beyond this assembly. The
interpolation algebra clears denominators through a bounded complex,
so the required congruence precision is independent of the number of
new primes. The graph construction treats the coefficient tests as polynomials in
prescribed residue symbols and makes separately nonzero tests nonzero
together. Its proof retains squared variables until after multiplying
their highest-degree parts.
Neither construction presupposes the desired nonvanishing at the base.

\paragraph{Organization.}
Section~\ref{sec:preliminaries} records the rank implications and local
conventions. Section~\ref{sec:transfer-criteria} states the two arithmetic
transfer criteria, the lower bounds needed to normalize coefficients,
and auxiliary nonvanishing with prescribed local conditions.
Sections~\ref{cs:coefficient-systems} and~\ref{sec:trace-calculus}
turn analytic nonvanishing into coefficient tests and prime-substitution
identities. Section~\ref{gr:section} develops the graph construction
used in the rational-torsion converse of Section~\ref{pw:section}.
Section~\ref{sec:lm:local-matrices} supplies the Selmer matrices for
the irreducible converse in Section~\ref{sec:sa:saturation}.
The remaining main-text sections prove the arithmetic inputs:
Section~\ref{sec:auxiliary-nonvanishing} establishes auxiliary
nonvanishing, Section~\ref{sec:binary-algebra} develops the common
finite-complex and binary transfer arguments, and
Sections~\ref{sec:cy:interpolation} and~\ref{sec:heegner-interpolation}
construct the cyclotomic and Heegner determinant coordinates and prove
their integral bounds. Finally, Section~\ref{sec:conclusions} assembles the unrestricted
result and its density and finite-descent consequences.

\section{Selmer ranks and local conventions}\label{sec:preliminaries}

We first record the arithmetic implications used to pass between curves
and to interpret the final analytic nonvanishing. They apply to every
elliptic curve over $\Q$.

Throughout, we use Mordell--Weil finite generation and the positive
definiteness of the canonical height modulo torsion over number fields
\citep[Chapter~IV, finite-basis theorem, Theorem~4.7,
Remark~4.12, and Theorem~6.1]{MilneEC2021}.

\begin{lemma}[Ranks and the forward implication]\label{pr:ranks}
For an elliptic curve $A/\Q$ there is an exact sequence
\[
 0\longrightarrow A(\Q)\otimes\Qtwo/\Ztwo
 \longrightarrow\Sel_{2^\infty}(A/\Q)
 \longrightarrow\Sha(A/\Q)[2^\infty]\longrightarrow0.
\]
In particular
\[
 c_2(A)=r(A)+\corank_{\Ztwo}\Sha(A/\Q)[2^\infty].
\]
If $a(A)\le1$, then $r(A)=a(A)$, the whole group $\Sha(A/\Q)$ is
finite, and $c_2(A)=a(A)$.
\end{lemma}

\begin{proof}
The finite Kummer exact sequences are
\[
 0\longrightarrow A(\Q)/2^nA(\Q)
 \longrightarrow\Sel_{2^n}(A/\Q)
 \longrightarrow\Sha(A/\Q)[2^n]\longrightarrow0.
\]
Their direct limit gives the displayed sequence; taking coranks gives
the identity. The last assertion is the Gross--Zagier--Kolyvagin
forward theorem \citep{GrossZagier1986,Kolyvagin1990}; see also its
unrestricted formulation in \citet[\S1]{Skinner2020}.
\end{proof}

\begin{lemma}[Parity and change of curve]\label{pr:invariance}
\label{prop:selmer-parity}
For every elliptic curve $A/\Q$ with root number $\varepsilon(A)$,
\[
 (-1)^{c_2(A)}=\varepsilon(A).
\]
A rational isogeny preserves $a(A)$, $r(A)$ and $c_2(A)$.
For a quadratic field $K=\Q(\sqrt{k})$,
\[
 \corank_{\Ztwo}\Sel_{2^\infty}(A/K)
       =c_2(A)+c_2(A^{(k)}).
\]
\end{lemma}

\begin{proof}
The parity statement is the $p=2$ case of
\citet[Theorem~1.4, equivalently Theorem~4.19]{DokchitserDokchitser2010}.
At $p=2$ this is Monsky's theorem \citep{Monsky1996}.
Isogenies and their duals induce maps respecting every local Kummer
condition; their composites are multiplication by their fixed degree.
After rationalization these maps preserve ranks and coranks. Their
$L$-functions agree by isogeny invariance.

Restriction to $K$ and the two conjugation projections $1+c$ and $1-c$
give maps between $\Sel_{2^\infty}(A/K)$ and the sum of the Selmer
groups of $A$ and $A^{(k)}$ over $\Q$, with composites multiplication
by $2$. They respect local Kummer images, so rationalization gives the
corank identity. An integral direct sum is not needed.
\end{proof}

\begin{lemma}[The residual alternatives]\label{pr:residual-cases}
The representation $E[2]$ is reducible over $\F_2$ if and only if
$E(\Q)[2]\ne0$. In that case it is an extension of two trivial
$\F_2$-representations. Otherwise its image is $C_3$ or $S_3$ in
$\GL_2(\F_2)$. Quadratic twisting leaves this representation unchanged.
\end{lemma}

\begin{proof}
An invariant line over $\F_2$ has a unique nonzero vector, which must
be fixed. Conversely a rational point of order two spans an invariant
line. Both that line and the one-dimensional quotient are trivial,
since $\F_2^\times=\{1\}$. The irreducible subgroups of
$\GL_2(\F_2)\simeq S_3$ are $C_3$ and $S_3$. Finally a quadratic
character reduces to $1$ modulo $2$.
\end{proof}

The distinction here is irreducibility over $\F_2$, rather than
absolute irreducibility: the cyclic cubic case is included.
In the rational-torsion case we may first replace the curve by a
rationally isogenous one with full rational two-torsion, if such a
curve exists. Otherwise, for any chosen rational two-isogeny
$E\to E'$, both $\Q(E[2])$ and $\Q(E'[2])$ are quadratic.
Lemma~\ref{pr:invariance} justifies this choice before any constructions.

\subsection{Finite lengths and twisting data}

All valuations above $2$ are normalized by $v(2)=1$, with $v(0)=+\infty$.
For a finite $\Ztwo$-module $M$, its length is $\log_2\#M$.
The corank of a cofinitely generated discrete $\Ztwo$-module is the
rank of its Pontryagin dual. Its quotient by the maximal divisible
subgroup is finite. This finite quotient must be distinguished from
finiteness of the original group.

A fixed support $S$ is a finite set of primes containing $2$ and
all primes of bad reduction; it is enlarged explicitly when necessary.
Write $N$ for the conductor of $E$.
A local filter specifies a real sign and unit squareclasses at $S$
for a variable odd fundamental discriminant prime to $S$.
The empty discriminant $h=1$ denotes the untwisted curve.
We also write $E^{(h)}$ for twisting by a rational squareclass.
For an odd prime $p$, the signed prime discriminant is
$p^*=(-1/p)p$. Products of distinct $p^*$ are odd fundamental
discriminants, with the empty product interpreted as $1$.

For an odd fundamental discriminant $h$ prime to $2N$, the
coprime-twist root-number formula is
\[
 \varepsilon(E^{(h)})=\varepsilon(E)\chi_h(-N)
\]
\citep[Section~1]{RadziwillSoundararajan2015}.
Thus matching unit squareclasses at the primes dividing $2N$ and
opposite real signs give opposite functional signs: the values of
$\chi_h(N)$ agree, whereas $\chi_h(-1)=\operatorname{sgn}(h)$.
This verifies the sign condition when choosing the auxiliary partners
below.

Put $T=T_2(E)$ and $V=E[2]$. Frobenius is arithmetic, and
$a_p(E)=p+1-\#E(\F_p)$ at a good prime. For $h$ prime to $S$, set
\begin{equation}\label{eq:h:weights}
 t_p=\dim_{\F_2}V^{\Fr_p},\qquad
 w(h)=\frac12\sum_{p\mid |h|}t_p.
\end{equation}
The residual identity, a transposition and a rotation of order three
have $t_p=2,1,0$, respectively. We call these split, simple and root
primes; their weights are $1,1/2,0$. Only types present in the
residual image are used. Quadratic residue symbols written additively
take the values $0,1\in\F_2$, corresponding to $+1,-1$.

Whenever its group is finite, write
\[
 s(h)=\length_{\Ztwo}\Sha(E^{(h)}/\Q)[2^\infty].
\]
This includes every twist already proved to have analytic rank zero
or one. It also includes a base of Selmer corank zero, since its full
$2$-primary Selmer group is then finite. We make no whole-Sha
finiteness assertion at such a base until analytic nonvanishing has
been established.

Let $\Omega_E^+$ be a positive real period and let $\Omega_E^-$ be
the positive magnitude of a nonzero anti-invariant period. Put
\begin{equation}\label{eq:h:algebraic-value}
 \mathcal L(h)=
 \frac{\sqrt{|h|}\,L(E^{(h)},1)}{\Omega_E^{\operatorname{sgn}(h)}}.
\end{equation}
Valuations refer to this algebraic period-normalized value at a chosen
place above $2$. Changes of the fixed period, modular parametrization
or isogeny affect valuations by bounded constants. All parametrizations
send the cusp $\infty$ to the origin.
Replacing Equation~\eqref{eq:h:algebraic-value} by the central value
divided by a N\'eron real period of $E^{(h)}$ changes a finite valuation
by $O_E(1)$: the differential changes by the usual twisting factor,
and there are only finitely many twisting classes at $2$.

Every uniform bound below specifies its fixed data. Constants may
depend on the curve, a chosen fixed twist and finite local preparation;
they may not depend on the subsequent cube dimension or the number of
new primes. At finite precision, the sizes of primes realizing a
specified configuration may depend on that entire configuration.
A limiting construction will only be used through its finite
realizations.

\begin{lemma}[Finite descent and corank]\label{pr:finite-descent}
Let $F$ be the quotient of $\Sha(E/\Q)[2^\infty]$ by its maximal
divisible subgroup. Then
\[
 d_2(E)=c_2(E)+\dim_{\F_2}F[2],
 \qquad \dim_{\F_2}F[2]\text{ is even}.
\]
Thus $d_2(E)\in\{0,1\}$ implies $c_2(E)=d_2(E)$ and $F=0$.
\end{lemma}

\begin{proof}
The finite Kummer sequence at $n=1$ gives
$d_2(E)=r(E)+\dim_{\F_2}\Sha(E/\Q)[2]$.
The divisible subgroup is a direct sum of copies of
$\Qtwo/\Ztwo$ and splits as an abstract abelian subgroup.
Its $2$-torsion dimension is its corank. Lemma~\ref{pr:ranks}
therefore gives the displayed identity.
The Cassels--Tate pairing on $F$ is perfect. For an elliptic curve,
the rational divisor $[O]$ defining its principal polarization makes
the pairing alternating
\citep[Chapter~I, Remark~6.10(b) and Theorem~6.13(a)]{MilneADT2006}.
A finite abelian $2$-group with such a pairing is an orthogonal sum
of pairs of cyclic groups of equal order: choose an element of
maximal order, pair it with an element detecting that order, split
off their nondegenerate plane, and continue. Hence its $2$-rank
is even. A finite $2$-group with zero $2$-rank is zero, proving
the last assertion.
\end{proof}

\section{Arithmetic estimates for finite twist families}
\label{sec:transfer-criteria}

The pointwise converse will follow by placing the given curve at the
zero vertex of a large finite family of twists. We arrange analytic
nonvanishing at every other vertex, together with an upper bound on a
normalized valuation. The estimates in this section explain which
normalization is needed and what must be uniform as the family grows.
Their proofs occupy Sections~\ref{sec:binary-algebra},
\ref{sec:cy:interpolation} and~\ref{sec:heegner-interpolation}.

Fix $E/\Q$ of conductor $N$ and retain the local weights $w(h)$,
finite Sha lengths $s(h)$, and period-normalized values $\mathcal L(h)$
of Section~\ref{sec:preliminaries}. In particular $s(h)$ is defined at
a corank-zero base and at every twist already known to have analytic
rank zero or one. All valuations satisfy $v(2)=1$.

A binary family has the form
\[
 h_x=h_0\prod_{q\in\mathcal Q}(q^*)^{\lambda_q(x)},
 \qquad x\in\F_2^b,\qquad q^*=(-1/q)q,
\]
where the new primes are distinct and prime to $2Nh_0$, and each
$\lambda_q$ is linear. The old factors of $h_0$ remain at every vertex;
no new factor is present at zero. We will hold the required local
squareclasses fixed while $b$ grows. Both the set of primes and their
sizes may depend on $b$.

Here is the common transfer mechanism. For families satisfying the
hypotheses of the transfer theorems below, the arithmetic proofs attach
power series
\[
                 u_x(t)\in\Ztwo[[t]][1/2]
\]
to the vertices, using a common family of cochain models and compatible
trivializations of the determinant lines. The differences of these
coordinates therefore have a fixed meaning. Their constant terms detect
the central value in rank zero or a Heegner point in rank one. If the analytic nonvanishing sought
at the base fails, then $u_0(0)=0$. The nonzero vertices instead satisfy
$v(u_x(0))\le B_1$ with $B_1$ independent of $b$.
The denominator analysis proves that, for every fixed integer $m$ and
sufficiently large $b$, some $x\ne0$ satisfies
\[
                   v\bigl(u_x(0)-u_0(0)\bigr)\ge m.
\]
Taking $m>B_1$ contradicts these two assertions. The work is to obtain
this congruence with a dimension threshold independent of the number
of new primes. Section~\ref{sec:binary-algebra} does so after the new
local blocks have been eliminated from the Selmer complex: a clearing
factor for the remaining bounded complex makes finitely many coefficient
congruences sufficient. Only that clearing factor needs a uniform pole
bound; the numerator may involve the whole prime set.

\subsection{Transfer of a central value}

For a twist with nonzero central value, put
\begin{equation}\label{eq:transfer-normalized-value}
                  D(h)=v(\mathcal L(h))-2w(h)-s(h).
\end{equation}
The subtraction records the local contributions of the ramified primes
and the finite Selmer contribution. In the cyclotomic proof, Kato's class
and the invariant line at the real place form the determinant coordinate
$u_x$. Its central valuation differs from $D(h_x)$ by a fixed bounded
amount. Thus an upper bound on $D$ is exactly the arithmetic input to
the preceding contradiction.

\begin{theorem}[Cyclotomic missing vertex]
\label{thm:cyclotomic-interpolation}
There is a constant $C_E$, depending only on $E$, with the following
property. Fix an odd fundamental discriminant $h_0$, prime to $2N$, with
$c_2(E^{(h_0)})=0$. Suppose that for arbitrarily
large $b$ there is a family
\begin{equation}\label{eq:cy:cube}
 h_x=h_0\prod_{q\in\mathcal Q}(q^*)^{\lambda_q(x)},
 \qquad x\in\F_2^b,\qquad \lambda_q\in(\F_2^b)^*,
\end{equation}
where $\mathcal Q$ consists of distinct primes outside $2Nh_0$,
$q^*=(-1/q)q$, and every $h_x$ has the same sign and the same local
squareclasses at $2N$. Assume that a constant $B$, independent of $b$,
satisfies
\begin{equation}\label{eq:cy:punctured-bound}
 L(E^{(h_x)},1)\ne0,\qquad D(h_x)\le B\qquad(x\ne0).
\end{equation}
If $E[2]$ is irreducible, assume also that some $L$, independent of $b$,
satisfies
\begin{equation}\label{eq:cy:prime-bound}
 v_2(q^2-1)\le L\qquad(q\in\mathcal Q).
\end{equation}
Then $L(E^{(h_0)},1)\ne0$ and
\begin{equation}\label{eq:cy:transferred-bound}
 D(h_0)\le B+C_E.
\end{equation}
The necessary dimension may depend on $h_0$, its finite Selmer group,
and the bound in Equation~\eqref{eq:cy:prime-bound}; $C_E$ does not.
\end{theorem}

The final constant $C_E$ bounds the specialization error. The finite
Selmer group of the base and, in the irreducible case, the bound on
$v_2(q^2-1)$ affect how large a cube is needed, but not this transferred
error. We shall construct the nonzero vertices so that
$v(\mathcal L(h_x))\le2w(h_x)+C$; since $s(h_x)\ge0$, this supplies
the hypothesis on $D(h_x)$.

\subsection{Transfer of a Heegner point}
\label{sec:h:statements}

In corank one the coordinate comes from a Heegner point. The family
must also carry a companion twist whose central value is nonzero.
We first specify the two Heegner constructions whose indices enter
the estimates.

Let $K=\Q(\sqrt{k})$, where $k<0$ is an odd fundamental discriminant and
every prime dividing $2N$ splits in $K$. We omit the two fields with extra
units and, if $E$ has complex multiplication, its CM field. Fix a nonzero modular parametrization $X_0(N)\to E$, translated so
that the cusp $\infty$ maps to the origin, and compatible orientations
of the Heegner points $y_c\in E(H_c)$, where $H_c$ is the ring class
field of conductor $c$. The Manin--Drinfeld theorem makes the images of
cusp differences torsion \citep[Section~1]{Drinfeld1973}. Since the
modular curve and parametrization are fixed, one integer kills all
these images and removes cuspidal ambiguities.

If $\chi_e|_{G_K}$ is a primitive ring class character of conductor $c$,
write $P_c(e)$ for the character \emph{sum} of $y_c$. Equivalently, twist
over $H_c$ and trace to $E^{(e)}(K)$. If this latter group has rank one,
$j(P_c(e))$ is the valuation of its index on the free rank-one lattice.
Passing between this lattice and the rank-one lattice of the contributing
curve among $E^{(e)}(\Q)$ and $E^{(ek)}(\Q)$ changes the index by a bounded
amount: restriction and addition or subtraction of conjugates have
composites multiplication by $2$.

We use two choices of data:
\begin{align*}
 \text{genus data:}\quad& k=h h_*,\quad (h,h_*)=1,
                         \quad e=h_*,\quad c=1,
                         \quad h_*\text{ fixed};\\
 \text{ring data:}\quad& (k,h)=1,\quad e=h,\quad c=|h|.
\end{align*}
The character in genus data is unramified. In ring data, restriction at
each $q\mid h$ has conductor exponent one: this follows by pulling the
quadratic character back under the local norm. Thus these are the asserted
primitive ring class conductors.

\begin{proposition}[Comparison of the two Heegner constructions]
\label{prop:heegner-comparison}
For either set of data,
\begin{equation}\label{eq:h:height}
 \widehat h(P_c(e))
 =C_E\,c\sqrt{|k|}\,
       \left[L(E^{(e)},s)L(E^{(ek)},s)\right]'_{s=1},
\end{equation}
with a fixed nonzero constant determined by the parametrization and height
conventions. In particular, a simple zero of the product gives total
Mordell--Weil rank one and finite Shafarevich--Tate groups.

Suppose $h$ has analytic rank one, and use genus data with a fixed partner
$h_*$ of analytic rank zero and ring data with a possibly varying $k$,
where $L(E^{(hk)},1)\ne0$. Then
\begin{equation}\label{eq:h:index-comparison}
 2j(P_{|h|}(h))-2j(P_1(h_*))
       =v(\mathcal L(hk))+O_{E,h_*}(1).
\end{equation}
Here the two points are formed in their respective quadratic fields.
If instead $h_*$ has analytic rank one and $L(E^{(h)},1)\ne0$, genus
data give
\begin{equation}\label{eq:h:reverse-comparison}
 2j(P_1(h_*))=v(\mathcal L(h))+O_{E,h_*}(1).
\end{equation}
\end{proposition}


The genus construction varies the quadratic field $\Q(\sqrt{hh_*})$
while holding $h_*$ fixed. It is the construction detected by the odd
coefficients below. Interpolation instead uses the ring construction:
within one binary family the field $K$ is common to every vertex.
Equation~\eqref{eq:h:index-comparison} converts a genus-index estimate
and an even central-value estimate into the ring-index estimate needed
for transfer. The field $K$ may change when a new family is chosen.

\begin{theorem}[Heegner interpolation]\label{thm:heegner-interpolation}
Fix $E$ and an initial discriminant $h_0$. Consider families
$h_x$, $x\in\F_2^b$, obtained from $h_0$ by adjoining new prime factors
with linear activation functions on $\F_2^b$. No old factor is removed;
the signed factor belonging to an odd prime $q$ is
$q^*=(-1/q)q$. Suppose the local twisting classes at $2N$ are constant.
For each family let $k$ define an imaginary field satisfying the standing
Heegner conditions, with $(k,h_x)=1$ for every $x$. Each new prime must
either split in $K$ or satisfy $t_q=0$.

Assume that the following bounds are fixed independently of $b$:
\begin{enumerate}[label=\textup{(\roman*)}]
 \item a bound for $\omega(k)$;
 \item a bound for the length of the finite quotient of
       $\Sel_{2^\infty}(E^{(h_0)}/K)$ by its divisible subgroup;
 \item a constant $B$.
\end{enumerate}
Assume this base Selmer group has corank one. Suppose, for every $x\ne0$,
the product $L(E^{(h_x)},s)L(E^{(h_xk)},s)$ has a simple zero and
\begin{equation}\label{eq:h:interpolation-hypothesis}
 2j(P_{|h_x|}(h_x))-4w(h_x)-s(h_x)-s(h_xk)\le B.
\end{equation}
Then, for sufficiently large $b$, the base product has a simple zero.
The dimension threshold depends only on the fixed base and the displayed
bounds, not on the number or sizes of the new primes. The field $K$ may
vary with the family.
\end{theorem}

In these ring data, the relevant normalized quantity is
\[
       2j(P_{|h|}(h))-4w(h)-s(h)-s(hk).
\]
We use the theorem to prove nonvanishing at the base. In our applications,
the genus-index and even central-value estimates will give
$2j(P_{|h_x|}(h_x))\le4w(h_x)+C$ at every nonzero vertex. Subtracting
the nonnegative Sha lengths supplies the required upper bound.

The finite-quotient condition allows two uses of the theorem. The
rational-two-torsion argument fixes $K$ before increasing $b$. The
irreducible argument lets $K$ vary and verifies the following criterion.

\begin{lemma}\label{lem:h:base-selmer-bound}
The finite-quotient bound in Theorem~\ref{thm:heegner-interpolation} holds
if $K$ belongs to a fixed finite collection. It also holds if
$E^{(h_0)}$ is fixed with Selmer corank one, the curves
$E^{(h_0k)}$ have analytic rank zero with bounded $s(h_0k)$, and
$\omega(k)$ is bounded.
\end{lemma}

\subsection{Lower bounds and the coefficient constructions}

The transfer criteria ask for upper bounds at a well-chosen set of
twists. To find those twists, we first need lower bounds on all the
normalized coefficients under consideration. In even sign a least
normalized valuation can then be selected, and a coefficient attaining
it gives both analytic nonvanishing and an upper bound with a constant
independent of the prime support. In odd sign the coefficients come at
increasing finite precision: the minimum is defined through bounded-weight
sequences, and each use is realized on a finite collection of indices.

For the even coefficient construction,
Section~\ref{cs:coefficient-systems} relates twice a coefficient valuation
to $v(\mathcal L(h))$ up to a fixed constant. For the odd construction it
relates the coefficient valuation to the genus index. The following
three lower estimates supply the normalizations, with different roles
in the two residual cases.

When $E[2]$ is reducible, the cyclotomic argument supplies the even
lower estimate directly. By Lemma~\ref{pr:residual-cases}, the
representation is then an extension of two trivial modules.

\begin{theorem}[Cyclotomic lower bound]\label{thm:cyclotomic-bound}
Suppose that $E[2]$ is an extension of two trivial $\F_2$-modules.
There is a constant $C_E$ such that every odd fundamental discriminant
$h$ prime to $2N$ with
$L(E^{(h)},1)\ne0$ satisfies
\begin{equation}\label{eq:cy:lower-bound}
 v(\mathcal L(h))\ge 2w(h)+s(h)-C_E.
\end{equation}
The constant is independent of the number and sizes of the primes
that divide $h$.
\end{theorem}

When $E[2]$ is irreducible, the genus construction supplies the lower
estimate for either sign. With an analytic-rank-zero fixed partner,
it bounds the odd genus index. With an analytic-rank-one fixed
partner, Equation~\eqref{eq:h:reverse-comparison} turns it into the
even central-value bound.

\begin{theorem}[Genus bound]\label{thm:heegner-genus-bound}
In genus data, suppose the product in
Equation~\eqref{eq:h:height} has a simple zero and $E[2]$ is irreducible.
Then, uniformly in $h$,
\begin{equation}\label{eq:h:genus-bound}
 2j(P_1(h_*))\ge 2w(h)+s(h)-O_{E,h_*}(1).
\end{equation}
In proving the uniform estimate, one may omit the finitely many $h$ for
which $K$ meets $\Q(E[2])$ nontrivially.
\end{theorem}

The remaining case is the odd coefficient construction with rational
two-torsion. Here the genus estimate just stated is unavailable. We use
instead the following ring estimate, whose residual representation is
unrestricted.

\begin{theorem}[Ring bound]\label{thm:heegner-ring-bound}
In ring data, suppose the product in
Equation~\eqref{eq:h:height} has a simple zero, every $q\mid h$ splits in
$K$, and $\omega(k)\le a$, where $\omega$ counts distinct prime factors.
Then
\begin{equation}\label{eq:h:ring-bound}
 2j(P_{|h|}(h))\ge
       4w(h)+s(h)+s(hk)-O_{E,a}(1).
\end{equation}
There is no irreducibility assumption on $E[2]$.
\end{theorem}

The passage from this ring estimate to an odd genus estimate is part
of Proposition~\ref{prop:rational-odd-bound}. The even coefficient construction
first provides, for each already rank-one twist $h$ in the positive
filter, a companion parameter $k'$ satisfying the ring-bound hypotheses:
it is negative, odd fundamental, coprime to $2Nh$, a local square at
every prime of $2Nh$, and defines a nonexceptional imaginary field.
Its prime count is bounded, and the companion obeys
$L(E^{(hk')},1)\ne0$ and
\[
       v(\mathcal L(hk'))\le2w(hk')+s(hk')+C.
\]
Substitute this upper bound into
Equation~\eqref{eq:h:index-comparison} and the ring lower bound.
Since $w(hk')=w(h)+w(k')$ and $w(k')$ is bounded, the result is
\[
                   2j(P_1(h_*))\ge2w(h)+s(h)-C'.
\]
Only then is the odd coefficient system normalized in the reducible
branch. Thus the lower-bound dependencies are
\[
 \begin{array}{c@{\quad\longrightarrow\quad}l}
 \text{irreducible genus bound}&\text{even and odd coefficient bounds},\\[2pt]
 \text{reducible cyclotomic bound}&\text{even coefficients, then bounded companions},\\[2pt]
 \text{ring bound and those companions}&\text{reducible odd coefficient bound}.
 \end{array}
\]
The reducible odd normalization is therefore available after the even
companion construction. The all-split hypothesis belongs to the ring
lower bound used here. The Heegner transfer criterion retains its stated
alternative for new primes with $t_q=0$.

\subsection{Auxiliary twists in prescribed local classes}

The coefficient constructions also require an auxiliary twist with
specified local classes. For even sign the number of its prime factors
must remain bounded as the local conditions vary. The following
nonvanishing result supplies both kinds of partner; its proof is given
in Section~\ref{sec:auxiliary-nonvanishing}.

\begin{proposition}[Auxiliary nonvanishing]
\label{prop:rough-nonvanishing}\label{nv:rough}\label{nv:local-witness}
Let $E/\Q$ have conductor $N$, and let $S$ be a finite set of primes
containing the prime divisors of $2N$. Prescribe a real sign
$\sigma\in\{1,-1\}$ and a unit squareclass
$c_p\in\Z_p^\times/\Z_p^{\times2}$ for each $p\in S$.
Suppose that these data are realized by an odd fundamental discriminant,
and that twists in this filter have root number $\varepsilon$.
For every finite set $B$ of primes there are odd fundamental
discriminants $D$ of arbitrarily large absolute value such that
\[
 \operatorname{sign}(D)=\sigma,\qquad [D]_p=c_p\quad(p\in S),\qquad
 p\nmid D\quad(p\in B),
\]
Moreover:
\begin{enumerate}[label=\textup{(\roman*)}]
\item If $\varepsilon=1$, they can be chosen with
$L(E^{(D)},1)\ne0$ and $\omega(|D|)\le5$.
\item If $\varepsilon=-1$, they can be chosen with
$L'(E^{(D)},1)\ne0$, so the central zero is simple.
\end{enumerate}
The bound five in \textup{(i)} is independent of $E,S,B$ and the
prescribed classes. In either case, every prime factor of $D$ can be
required to exceed any previously fixed bound.
\end{proposition}

\section{Coefficients that detect analytic rank}
\label{cs:coefficient-systems}

The interpolation theorems require nonvanishing with a uniform valuation
bound at the nonzero vertices of a cube. We encode both requirements in
one binary coefficient: it is one when a modular coefficient has the
least normalized valuation. Waldspurger's formula supplies the even-sign
coefficients; reduction of genus Heegner sums supplies the odd-sign
coefficients at finite precision. We prove their modular and arithmetic
properties separately before choosing this normalization.

\subsection{Filters and prime weights}

Fix $E/\Q$ and a finite set $S$ containing the primes of $2N$ and the
ramification set of $F=\Q(E[2])$.  Additional fixed primes may be included.
A \emph{filter} consists of a sign $\sigma\in\{1,-1\}$ and unit
squareclasses at the primes in $S$, imposed on positive squarefree integers
$D$ prime to $S$, so that
\[
                     h=\sigma D
\]
is an odd fundamental discriminant and the functional sign of $E^{(h)}$
is constant.  Multiplication by squares of units preserves the local
conditions.  We impose these conditions on all Fourier indices, not only
on squarefree ones.  Projections onto the required unit squareclasses are
finite combinations of quadratic coefficient twists and coprimality
projections; they enlarge the level only at $S$.

For $p\notin S$, the action on $V=E[2]$ has the following three types:
\[
\begin{array}{c|c|c}
\text{type}&\Fr_p|_V&\text{weight}\\ \hline
\text{split}&1&1\\
\text{simple}&\text{an element of order }2&1/2\\
\text{root}&\text{an element of order }3&0.
\end{array}
\]
Thus $w(D)=\frac12\sum_{p\mid D}\dim_{\mathbf F_2}V^{\Fr_p}=w(h)$.
The fractional part of $w(D)$ is constant on a filter: it is determined
by the quadratic discriminant character of the residual representation.
For an arbitrary squarefree set of primes we use the same additive
weight.  Signs are encoded by $p^*=(-1/p)p$, so that $h=\prod_{p\mid D}p^*$.

We use the ordinary theta multiplier cubed in weight $3/2$. A finite
idele unit character and its Dirichlet character are related by the
reciprocal convention: if a finite Schwartz function obeys
$f(ux)=\alpha(u)f(x)$, its scalar classical character on the lower-right
entry is associated to $\alpha^{-1}$. All valuations satisfy $v(2)=1$.

\subsection{The even coefficient system}

\begin{proposition}[Even coefficients]
\label{prop:even-coefficients}\label{co:even}
Fix a nonempty filter with even functional sign.  There are a holomorphic
weight-$3/2$ cusp form
\[
                         g=\sum_{n\geq1}a(n)q^n
\]
of level supported on $S$, after enlarging exponents at those primes,
and finite-order characters $\psi,\eta$ of $2$-power order, with
$\psi/\eta$ quadratic, having the following properties.
\begin{enumerate}[label=\textup{(\roman*)}]
\item The coefficients vanish off the filter and, for $p\notin S$,
\begin{equation}
 a(p^2n)=\bigl(\eta(p)a_p(E)-\psi(p)(-n/p)\bigr)a(n)
                   -p\eta(p)^2a(n/p^2).
\label{cs:hecke-recurrence}
\end{equation}
As usual a coefficient with a nonintegral index is zero.
\item After normalizing at one nonzero squarefree index, every squarefree
coefficient in the filter is a rational number times a root of unity of
$2$-power order.  For a fixed nonzero squarefree index $D_0$,
\begin{equation}
 2v\left(\frac{a(D)}{a(D_0)}\right)
 =v\left(\frac{\mathcal L(\sigma D)}{\mathcal L(\sigma D_0)}\right)
 \quad\text{whenever }a(D)\ne0.
\label{cs:even-ratio}
\end{equation}
Equivalently, there is a constant $c_g$, independent of $D$, such that
\begin{equation}
       2v(a(D))=v\bigl({\cal L}(h)\bigr)+c_g.
\label{cs:even-value}
\end{equation}
In particular $a(D)\ne0$ if and only if $L(E^{(h)},1)\ne0$.
\item After the fixed auxiliary exclusions described below, there is a
constant $C_g$ such that, at every nonzero squarefree coefficient,
\begin{equation}
           v(a(D))\geq w(D)+\tfrac12s(h)-C_g.
\label{cs:even-lower-bound}
\end{equation}
The constants depend on the fixed curve, filter and auxiliary choices,
and not on $D$ or its number of prime factors.
\end{enumerate}
\end{proposition}

\begin{proof}
We first choose the packet member and its local characters. We then
establish the rationality up to phase needed for valuation comparisons,
apply the coefficient-square formula, and deduce the uniform lower bound
from the interpolation estimates.

Use the Shimura packet attached to the integral-weight newform of
$E^{(-\sigma)}$.  In the classical convention its positive Fourier index
$D$ tests the twist by $\chi_{-D}$, hence the curve $E^{(\sigma D)}$.
The packet and Fourier coefficient results of
\citet{Waldspurger1980,Waldspurger1981} apply to the cuspidal packet
orthogonal to unary theta series.  Their local uniqueness and global
multiplicity assertions are summarized in
\citep[Sections~3.1 and~5.1]{Waldspurger1983}.
A nonzero central value in the filter exists by
Proposition~\ref{prop:rough-nonvanishing}.  It selects a globally occurring
packet member with nonzero Fourier functional at that index.  At good
odd primes choose spherical vectors; at infinity choose the holomorphic
vector of weight $3/2$.

We explain the choices at a prime of $S$.  Let $\lambda_D$ be a nonzero
local Fourier functional and start with a vector on which it is nonzero.
Average upper-unipotent translates against the inverse additive character
of frequency $D$.  If the averaging compact subgroup is sufficiently
large, its translate by a unit scale $u$ can contribute only when $u^2$
is sufficiently close to $1$.  The surviving scales are consequently
small neighborhoods of $1$ and $-1$.  Local constancy then permits
projection onto a unit character which matches the forced action of
$-1$, without killing the Fourier value: the contributions of the two
neighborhoods add.  Such a character may be chosen to have $2$-power
order.  Indeed only its parity is prescribed, and the $2$-primary
quotient of the unit group realizes either possible parity.  Upper
averaging at integral frequency gives period one, and lower-unipotent
invariance holds at sufficiently large depth.  These vectors therefore
give a classical form at a level supported on $S$.  The parity conditions
combine correctly by automorphy.  The subsequent squareclass projections
preserve its character.  Write that character as $\psi$.

In the theta-multiplier convention the good eigenvalues of the packet
are $\psi(p)$ times those of $E^{(-\sigma)}$.  Set
$\eta=\psi\chi_{-\sigma}$.  Then $\psi/\eta$ is quadratic and the
weight-$3/2$ Hecke formula gives Equation~\eqref{cs:hecke-recurrence}.

We need more than algebraicity of the coefficients. After normalization
at one squarefree index, their nonzero squarefree valuations must be
integers, so that a strict valuation inequality gains a full factor of
$2$, even when the character field is ramified. We prove this by showing
that every such coefficient is rational up to a phase of $2$-power order.
At the fixed level,
the weight-$3/2$ cusp space with character $\psi$ has a basis over the
field of character values, and coefficient conjugation changes $\psi$
to its conjugate
\citep[Section~2, Lemma~4(3)--(4), pp.~658--659]{Shimura1981CriticalValues}.
The good Hecke equations and the vanishing of coefficients off the
filter are linear equations over $\Q(\psi,\eta)$.  In a finite-dimensional
space, finitely many of these equations cut out their full common
kernel.  The constructed complex vector has a nonzero coefficient at
a squarefree index $D_0$ of the filter, so this coefficient functional
is nonzero on some algebraic vector in that kernel.  Choose such a
vector.  It remains orthogonal to unary theta series: at good primes
their eigenvalues have absolute value $p+1$, whereas the prescribed
eigenvalues have absolute value at most $2\sqrt p$ by Hasse's bound
\citep[Section~2, Lemma~2.1, and Corollary~5.2]{Purkait2012}.

We use \citet[Corollary~2, pp.~483--484]{Waldspurger1981} to compare
Fourier coefficients.  Its hypothesis (H1) holds for every nonzero
classical eigenspace used here, by Proposition~2 of that paper.
The final paragraph of the corollary's proof removes (H2) when the
indices lie in the same local squareclass at every level prime.
Our filters retain this condition at~$2$ as well as at the odd primes.

Let $u$ be an automorphism of $\overline\Q$ acting by an odd power $e$
on the $2$-power characters, and put $\nu=\eta^{(e-1)/2}$.  The
coefficient conjugate $g^u$ and coefficient twist $g\otimes\nu$ have
exactly the same character and good Hecke system, since
\[
 \psi\nu^2=\psi\eta^{e-1}=\psi^e,
                 \qquad \eta\nu^2=\eta^e.
\]
They may be compared at a common level supported on $S$, in the same
good Hecke eigenspace $W$, and both have a nonzero $D_0$-coefficient.
Let $\ell_n:W\to\C$ denote the $n$th Fourier-coefficient functional.
Waldspurger's corollary applies to every $f\in W$ and gives
\[
                    \ell_D(f)^2=C_D\ell_{D_0}(f)^2,
\]
where $C_D$ is independent of $f$.  If $C_D=0$, then $\ell_D=0$.
Otherwise choose $c_D\in\C$ with $c_D^2=C_D$ and factor the identity
as
\[
       (\ell_D-c_D\ell_{D_0})(\ell_D+c_D\ell_{D_0})=0.
\]
The polynomial ring on $W$ is an integral domain, so one factor
vanishes identically.  Thus the two coefficient functionals are
proportional on the whole space, including oldvectors and linear
combinations.  In particular the two forms have the same ratio of
their $D$- and $D_0$-coefficients.

It follows, for $\xi=a(D)/a(D_0)$, that
\[
                         u(\xi)=\nu(D/D_0)\xi.
\]
The local squareclass conditions imply
$\eta(D/D_0)=z^2$ for a value $z$ in the group of values of $\eta$.
Consequently $u(\xi/z)=\xi/z$ for every $u$, so $\xi/z\in\Q$.
This proves the phase assertion.

With our characters and period normalization, its identity of
\emph{coefficient squares} is
\begin{equation}
 \left(\frac{a(D)}{a(D_0)}\right)^2
 =\psi(D/D_0)\,
   \frac{\mathcal L(\sigma D)}{\mathcal L(\sigma D_0)}.
 \label{cs:even-square-ratio}
\end{equation}
Here $\psi(D/D_0)=\psi(D)\psi(D_0)^{-1}$ has valuation zero;
the factor $\sqrt{D/D_0}$ in the classical central-value comparison
is incorporated in $\mathcal L$. The local unit
squareclasses at $S$ agree; at a good prime the spherical squarefree
factor is one at both valuations zero and one
\citep[Section~I.3, pp.~377--378]{Waldspurger1981}.
There is consequently no extra product over the prime factors of $D$.
Taking valuations gives Equation~\eqref{cs:even-ratio}, hence
Equation~\eqref{cs:even-value}. The same identity, compared with the
nonzero coefficient at $D_0$, gives the asserted equivalence of
vanishing. The corresponding rational-square statement for central values
is Vign\'eras's rationality theorem, recalled in
\citet[Section~6.3, Theorem~6]{Waldspurger1983}.

If $E[2]$ is an extension of trivial modules, combine
Equation~\eqref{cs:even-value} with
Theorem~\ref{thm:cyclotomic-bound}.
If $E[2]$ is irreducible, choose a fixed odd fundamental discriminant
$h_*$ of opposite infinite sign, of analytic rank one, whose local
characters agree with those of $h$ at $2N$.  Proposition~\ref{prop:rough-nonvanishing}
supplies such a choice.  Exclude its prime factors from the filter and
choose its support to avoid the finitely many exceptional genus fields.
The reverse genus comparison in Proposition~\ref{prop:heegner-comparison}
and Theorem~\ref{thm:heegner-genus-bound} give the same lower bound.
These are fixed exclusions and constants, proving
Equation~\eqref{cs:even-lower-bound} in both cases.
\end{proof}

\subsection{Odd coefficients from reduction of Heegner traces}

The odd construction combines reduction of Heegner points with the ternary
quadratic forms attached to definite quaternion orders. We use Deuring's
reduction correspondence in the oriented form of Cornut--Jetchev,
together with the Eichler-order and ternary-form description of
Jetchev--Kane \citep{Deuring1941,CornutJetchev2013,JetchevKane2011}.
The weights below retain the genus character and the local orientations,
so that the resulting Fourier coefficients recover the reduced genus sum
at finite precision.

Fix now a nonempty filter of odd functional sign.  Choose an odd fundamental
$h_*$ of opposite infinite sign, with $L(E^{(h_*)},1)\ne0$, whose local
characters agree with those of $h$ at $2N$.  Exclude its primes.  We may
choose $|h_*|>3$ and include in its support a prime outside the
exceptional imaginary quadratic discriminants.  Thus all genus fields
below have only the units $\{1,-1\}$ and avoid the CM field of $E$.
Write
\[
          k=hh_*<0,\qquad K_D=\Q(\sqrt{k}),\qquad
          P_D=P_1(h_*).
\]
The point $P_D$ is the character \emph{sum} of the conductor-one Heegner
point for $K_D$. We fix the modular parametrization $\varphi$ with
$\varphi(\infty)=0$, so the constants in its Hecke and Atkin--Lehner
relations are torsion.  At a simple zero of $E^{(h)}$, its free line is that of
$E^{(h)}$, up to the fixed-degree identifications in
Proposition~\ref{prop:heegner-comparison}.
Below, the product derivative means the derivative at $s=1$ of
$L(E^{(h)},s)L(E^{(h_*)},s)$.

\begin{proposition}[Odd coefficients under a genus-index bound]
\label{prop:odd-coefficients}\label{co:odd}
Assume that a constant $C_*$ satisfies
\begin{equation}
              j(P_D)\geq w(D)+\tfrac12s(h)-C_*
\label{cs:odd-bound-hypothesis}
\end{equation}
for every $h$ in the filter with a simple zero.  There are arbitrarily
large precisions $M$, good odd primes $\ell$ outside the fixed support,
and holomorphic weight-$3/2$ forms $g_M=\sum a_M(n)q^n$, of level
supported on $S\cup\{\ell\}$, with the following properties.
\begin{enumerate}[label=\textup{(\roman*)}]
\item The additional filter at $\ell$ is
\begin{equation}
  (h_*/\ell)=-\sigma,\qquad (h/\ell)=\sigma,
               \qquad \rho_T(\Fr_\ell)\equiv C\pmod{2^M},
\label{cs:detecting-prime}
\end{equation}
where $C$ is the action of a complex conjugation.  In particular
$\ell\equiv-1\pmod{2^M}$.
\item After a fixed harmless decrease in the working precision, the
forms satisfy Equation~\eqref{cs:hecke-recurrence} modulo that precision,
with characters $\eta_M,\psi_M$ of $2$-power order and quadratic ratio.
Their squarefree coefficients have the form of a $2$-power phase times
an integer modulo $2^M$.
\item If a coefficient is visible, meaning $v(a_M(D))<M$, then
$E^{(h)}$ has a simple zero and
\begin{equation}
                    j(P_D)\leq v(a_M(D))+C_1.
\label{cs:odd-detection}
\end{equation}
For every admissible $D$ the coefficient lower bound is
\begin{equation}
 v(a_M(D))\geq
       \min\{M,\,w(D)+\tfrac12s(h)-C_2\}
\label{cs:odd-lower-bound}
\end{equation}
when the product derivative is nonzero; if it is zero, the coefficient
is zero modulo $2^M$.  Here $C_1,C_2$ are independent of $\ell,M,D$.
\item A fixed squarefree witness $D_w$ in the original filter is
admissible at every chosen $\ell$ and has $v(a_M(D_w))$ bounded as
$M\to\infty$.
\end{enumerate}
For irreducible $E[2]$, the hypothesis is
Theorem~\ref{thm:heegner-genus-bound}.  For rational two-torsion, the positive filters used in the pointwise
argument have this bound by Proposition~\ref{prop:rational-odd-bound},
using only the even source. An arbitrary twist is first absorbed into
the fixed curve, so only these positive filters are needed there.
\end{proposition}

\begin{proof}
The proof has three stages. A reduction map detects one fixed nonzero
genus point at arbitrarily high precisions. Weighted theta series then
supply holomorphic forms with the required Hecke recurrence. Finally,
an exact comparison with genus sums transfers point divisibility and
the assumed lower bound to their squarefree coefficients.

\paragraph{The reduction map and a detecting sequence.}
Choose a nonzero integral row $b$ with $bC=\sigma b$.  Under
Equation~\eqref{cs:detecting-prime}, finite-field Kummer theory gives
\[
 \kappa_M:E(\mathbf F_{\ell^2})\longrightarrow
 H^1(\mathbf F_{\ell^2},E[2^M])
       =E[2^M]/(\Fr_\ell^2-1)E[2^M]=E[2^M].
\]
The last identification is evaluation at arithmetic Frobenius squared.
Put $\lambda_M=b\kappa_M$.  This map has Frobenius sign $\sigma$.

Choose $h_w=\sigma D_w$ in the original filter with a simple zero,
using Proposition~\ref{prop:rough-nonvanishing}.  A fixed multiple of
its genus point is a nonzero point on the rational curve $E^{(h_w)}$.
Its division cocycle has open translation image on the kernel of the
Tate action, even after fixed finite auxiliary fields have been
adjoined.  Here is the relevant argument.  The translation image is a
Galois-stable $\Ztwo$-submodule of the Tate module.  Its rational span is
nonzero: otherwise the point's Kummer class would factor through the
Tate image; that image has a central nonidentity scalar, which kills
its rational first cohomology, contradicting Kummer injectivity.
The rational Tate representation is irreducible by
\citet[Chapter~IV, Section~2.1, pp.~IV-9--IV-10]{Serre1968}, since
$\operatorname{End}_{\Q}(E)=\Z$.  This includes curves with geometric
complex multiplication, whose nonreal endomorphisms are not defined
over $\Q$.  The scalar Lie-algebra assertion of
\citet[Section~2, Theorem~3, p.~62]{Bogomolov1981}, followed by local
logarithm and exponential, supplies the required nonidentity scalar.
The nonzero stable span is therefore the whole Tate module.  The
translation image consequently contains a full lattice.  Passing to
a finite-index subgroup preserves this fact.

Take an element acting as $C$ on $T$ and as complex conjugation on the
fixed quadratic character fields, and modify it by such a translation.
On the twisted module for $h_w$ its action is $\sigma C$; at its square
the cocycle includes $1+\sigma C$.  Projection by $b$ is twice the
corresponding nonzero row.  We can therefore arrange nonzero detection
with a fixed finite valuation.  Chebotarev, applied at successive finite
precisions to this fixed compatible Galois element and its division
data, supplies arbitrarily large primes satisfying
Equation~\eqref{cs:detecting-prime}, also with $h=h_w$, and the asserted
bounded detection.  A congruence on the target curve's Frobenius is all
that is required.  The supersingular objects used next are points of
the modular curve, not an additional reduction hypothesis on $E$.

\paragraph{Supersingular values and local weights.}
Let $B$ be the quaternion algebra ramified at $\ell$ and infinity, and
let $R$ be an Eichler order of level $N$.  Supersingular pairs on $X_0(N)$
are parametrized by
\[
                  B^\times\backslash\widehat B^\times/
                  \widehat R^\times.
\]
They are defined over $\mathbf F_{\ell^2}$.  The reduction and Eichler
order description apply to arbitrary $N$ prime to $\ell$; see
\citep[Sections~1.1--1.2 and~2.5]{JetchevKane2011}.
For a representative $z$ write $R_z=B\cap z\widehat Rz^{-1}$ and let
$P_M(z)$ be the reduced modular-parametrization value followed by
$\lambda_M$.  The cusp differences have torsion image by
\citet[Section~1]{Drinfeld1973}.  Multiply the parametrization once by a fixed integer
$L$ killing the cuspidal constants in its Hecke and Atkin--Lehner
relations and the bounded-degree torsion occurring in genus sums.
This choice is independent of $D,\ell,M$.

The resulting $P_M$ has good Brandt eigenvalues $a_p(E)$, level
Atkin--Lehner eigenvalues $w_{p^e}$ for $p^e\Vert N$, and eigenvalue
$\sigma$ for the right uniformizer at $\ell$.  These identities hold
modulo $2^M$.  Choose integer representatives of its values when making
an ordinary characteristic-zero theta series.

On $B^0$, the trace-zero subspace, use the positive definite quadratic
form $Q(x)=-x^2=\Nrd(x)$. Our aim is to make the squarefree coefficient
$a_M(D)$ equal, modulo $2^M$ and up to a phase of $2$-power order, to a
fixed integer multiple of
\[
                 \lambda_M\bigl(\operatorname{red}_{\ell}(L P_D)\bigr).
\]
We will identify its terms with optimal embeddings of
$\mathcal O_{K_D}$. Their local orientations give several copies of the genus sum.
The weights must make those copies contribute equally, while their
unit-scaling law must give a scalar modular character. The weighted
trace calculation below will determine the exact multiplier.

Define a finite compactly supported, locally constant weight
$f=\prod f_p$ as follows; a value not satisfying a listed condition is zero.
\begin{enumerate}[label=\textup{(\alph*)}]
\item Away from $2Nh_*\ell$, require order integrality and put $f_p=1$.
\item At $2$, require $(1+x)/2$ to be integral and $x^2$ to be a split
fundamental unit squareclass.  At $p^e\Vert N$, require the split
fundamental unit type and its order embedding.  The cyclic Eichler
level line selects one of the two eigenroots $r$ of $x$; weight it by a
$2$-power unit character $\eta_p(r)$ with
\[
                   \eta_p(-1)=w_{p^e}\chi_{h_*}(p)^e.
\]
If $p=2$, use the selected summand of $(1+x)/2$, whose eigenvalues
$(1+r)/2$ and $(1-r)/2$ differ by the odd unit $r$. Scaling by an odd
$u$ sends this operator to
$u(1+x)/2+(1-u)/2$ and preserves that summand; its selected root for
$ux$ is $ur$. Thus the weight changes by $\eta_2(u)$, even when the
two eigenvalue parities interchange. If $2\nmid N$, impose only the
stated integrality restriction.
\item At $p\mid h_*$ require integral $x$ with determinant of valuation
one.  Its reduction is nonzero nilpotent.  On that orbit use the sign
which transforms by $(\det u/p)$ under conjugation by
$u\in\GL_2(\mathbf Z_p)$.  This sign exists because the stabilizer of
a nonzero nilpotent has the form $aI+bX$ and determinant $a^2$.
Scaling $x$ by a unit $u$ changes this sign by $(u/p)$.
\item At $\ell$ require integral unit norm and the unramified quadratic
type.  Choose a basis of the trace-zero line in $\mathbf F_{\ell^2}$
and weight the nonzero residue on that line by its Legendre symbol
relative to the basis.
\end{enumerate}
All these conditions may be imposed at finite depth.  They are stable
under unit scaling.  Right order-unit conjugation changes $f$ by the
inverse of $\chi_{h_*}\circ\Nrd$, exactly canceling the norm character
below.  At $\ell$, order units commute on the residue field, while a
uniformizer negates its trace-zero line.  Since $\ell\equiv-1\pmod4$,
the latter weight changes by $-1$.

Write $\alpha_B=\prod_p\alpha_p$ for the finite unit-scaling
character, so that $f(ux)=\alpha_B(u)f(x)$.  Define the classical
Dirichlet character $\eta_B$ by
\[
 \eta_B(d)=\prod_p\alpha_p(d)
 \qquad\text{for integers }d\text{ prime to the support}.
\]
The associated finite idele character restricts to $\alpha_B^{-1}$
on the ramified unit groups, in accordance with our reciprocal
convention.  Both characters are even.  Indeed their value on $-1$ is
\[
 \left(\prod_{p^e\Vert N}w_{p^e}\chi_{h_*}(p)^e\right)
                      \chi_{h_*}(-1)(-1)=1,
\]
by the even functional sign of $E^{(h_*)}$.  The characters have
$2$-power order.

\paragraph{Transformation and good Hecke relation.}
Choose a fixed positive integer $c_0$ divisible by the orders of all
unit groups which occur.  Such a choice is uniform: a unit group of a
rational definite quaternion order has bounded order.  Form
\begin{equation}
 H_M(n)=c_0\sum_z\frac{P_M(z)\chi_{h_*}(\Nrd z)}{|R_z^\times|}
             \sum_{\substack{x\in B^0\\-x^2=n}}f(z^{-1}xz),
 \qquad H_M(\tau)=\sum_{n\geq1}H_M(n)q^n.
\label{cs:weighted-theta}
\end{equation}
The norm character is trivial on rational $B^\times$, since its rational
reduced norms are positive.  The right transformation just checked makes
Equation~\eqref{cs:weighted-theta} independent of the representatives.

For scalar modularity, apply the Weil representation of $(B^0,Q)$
to $f$ and the Gaussian at infinity \citep{Weil1964}.  At a finite prime,
with the additive character of conductor $\Z_p$, the local formulas are
\citep[Section~4.3, pp.~1965--1966]{PalDeVeraPiquero2019}
\[
 \begin{aligned}
 \omega(n(b))f_p(x)&=\mathbf e_p(bQ(x))f_p(x),\\
 \omega(m(u))f_p(x)&=\gamma_p(u)|u|_p^{3/2}f_p(ux),
 &\omega(w)f_p&=\gamma_p\widehat f_p.
 \end{aligned}
\]
Here $n(b)=\left(\begin{smallmatrix}1&b\\0&1\end{smallmatrix}\right)$,
$m(u)=\operatorname{diag}(u,u^{-1})$, and $w$ is the Weyl element.
Integral $b$ fixes $f_p$, since $Q$ is integral on its support.
Choose a lattice of translation invariance for $f_p$; its dual contains
the support of $\widehat f_p$. For sufficiently large $a_p$, the form
$p^{a_p}Q$ is integral on that dual lattice. Thus
$n(p^{a_p}b)$ fixes $\widehat f_p$, and conjugation by $w$ proves
invariance under sufficiently deep lower unipotents. This argument
also works at $2$ and at arbitrary finite character depth.

Compare with the threefold ordinary-theta reference line, using the
same additive character and lifts.  In a presentation $B=(-a,-b)$,
the quadratic form is $aX^2+bY^2+abZ^2$.  The ratio of its diagonal
Weil factor to the reference factor is
\[
                   (a,u)_p(b,u)_p(ab,u)_p=((ab)^2,u)_p=1.
\]
This calculation holds also over $\Q_2$.  The remaining local
diagonal scalar is therefore $\alpha_p(u)$.  Upper unipotents,
diagonal units, and sufficiently deep lower unipotents generate the
required local $\Gamma_0$ subgroup: explicitly, for $a$ a unit,
\[
 \begin{pmatrix}a&b\\c&d\end{pmatrix}
 =n^-(c/a)m(a)n(b/a).
\]
Choose the depths also to serve the ordinary-theta reference line
and the character conductors.  The central ambiguity of the
metaplectic lifts cancels in the comparison, since both representations
have the same odd central action.  At the chosen depth, $a$ is congruent
to $d^{-1}$ modulo the character conductor, so the local right
eigenvalue is $\alpha_p(d)^{-1}$.  Classical
restriction uses the inverse finite component and consequently gives
the character $\eta_B(d)$.  Conjugation by $z$ preserves $Q$ and
commutes with the Weil action, so all summands have this same level
and scalar character.
At good odd primes the lattice is spherical. The Gaussian gives
weight $3/2$, and positive definiteness gives holomorphy at every cusp.
Thus $H_M$ is a classical form of level supported on $2Nh_*\ell$.

At a prime $p\nmid 2Nh_*\ell$, the neighbor relation has Hecke eigenvalue
\[
                         \eta_B(p)\chi_{h_*}(p)a_p(E).
\]
The factor
$\chi_{h_*}(p)$ comes from the reduced norm of a Brandt neighbor and
$\eta_B(p)$ from scaling at the other primes.

The multiplicities can be read directly in a matrix order.  Put $y=px$.
If $y$ is primitive
and $Q(y)$ is divisible by $p^2$, its nonzero nilpotent reduction has a
unique invariant line; in a basis adapted to that line the other
off-diagonal entry is divisible by $p^2$, so exactly one neighbor makes
$y/p$ integral.  If $y/p$ is integral, its reduction has
$1+(-Q(y/p)/p)$ invariant lines, except that a zero reduction has a
further $p$ lines.  These are precisely the middle and last terms of
the weight-$3/2$ Hecke recurrence.  The factors $|R_z^\times|^{-1}$
are the orbit-stabilizer factors making the Brandt transfer valid.
Since the eigenrelation for $P_M$ holds modulo $2^M$, so does this
relation for $H_M$.

Put $d=|h_*|$. The weights force $Q(x)$ to be divisible by $d$.
Before the final filter projections, the series with coefficients
$a_M(n)=H_M(dn)$ is the same weighted theta sum for $Q/d$.
Its ternary determinant changes by $d^{-3}$, whose squareclass is $d$;
hence the scalar character acquires $(d/\cdot)$.
Substitution of $dn$ in the good-prime recurrence leaves
$\eta_B\chi_{h_*}$ unchanged. Now apply the filter projections.
Extraction of these multiples changes the character to
$\psi_M=\eta_B( |h_*|/\cdot)$, while
$\eta_M=\eta_B\chi_{h_*}$ is the good eigenvalue character.
Their quotient is quadratic.  The neighbor relation therefore becomes
Equation~\eqref{cs:hecke-recurrence}.  Every operation has level supported
on $S\cup\{\ell\}$.

\paragraph{The weighted trace identity.}
The form and its recurrence are now constructed. We next identify a
squarefree coefficient with the reduction of the original genus sum,
keeping the orientation multiplicities and unit stabilizers explicit.

For a squarefree $D$ in the filter, put $n=|h_*|D=-k$.
An $x$ counted in Equation~\eqref{cs:weighted-theta} is exactly an
optimal embedding of the full order $\mathcal O_{K_D}$ into $R_z$,
by sending $\sqrt{k}$ to $x$.  At $2$ the generator is
$(1+\sqrt{k})/2$, which is why the integrality condition above is
necessary.  At a split level prime there are two choices of cyclic
level orientation, and at $\ell$ there are two choices of unramified
quadratic orientation.  There is one local orbit at the remaining
primes.  The split-level assertion follows by decomposing a stable
lattice into its two idempotent summands: a cyclic quotient of order
$p^e$ forces all of that quotient into one summand.  The division-prime
assertion follows from the two reduced-norm parity cosets.  At primes outside $2Nh_*\ell$,
stable lattices are the fractional ideals of the quadratic
order, giving the single local orbit.

Fix one orientation tuple and a rational embedding $x_0$.  The
corresponding global optimal-embedding orbits are parametrized by
$\Pic(\mathcal O_{K_D})$.  Explicitly, once one adelic representative
$z_0$ is fixed, the representatives are $z_{\mathfrak a}=a z_0$ with
$a\in\widehat K_D^\times$; changing $a$ by $K_D^\times$ or
$\widehat{\mathcal O}_{K_D}^{\times}$ gives exactly the same class.
Because $a$ commutes with $x_0$, the Schwartz factor is constant along
this class-group orbit.  Its remaining relative weight is
\[
             \chi_{h_*}(\Nrd a)
              =\chi_{h_*}(\operatorname{Norm}_{K_D/\Q}a),
\]
the unramified genus character.  This computation establishes the
weighted identity, rather than merely an unweighted correspondence.

For completeness the stabilizer factor is also exact.  Each orbit of
embeddings in a given $R_z$ has
$|R_z^\times|/|\mathcal O_{K_D}^\times|=|R_z^\times|/2$ elements.
Thus its contribution after division by $|R_z^\times|$ is one half of
the corresponding class-group term. These are oriented embeddings,
not conjugate pairs: $x$ and $-x$ have opposite orientations at $\ell$
and cannot be conjugate by an order unit, whose residue conjugation
on $\mathbf F_{\ell^2}$ is trivial.  CM reduction at the inert prime
identifies these terms with the reductions of the Heegner points and
commutes with ideal action; the cyclic subgroup carries the selected
level orientation.  Fixing the residue-field identification at $\ell$
retains its orientation in the reduction-to-embedding correspondence
of \citet{Deuring1941}; see
\citep[arXiv version~1, Theorem~3.1, Corollary~3.2, and Section~3.6,
Remark~1]{CornutJetchev2013} and
\citep[Sections~2.4--2.6 and~4.1]{JetchevKane2011}.
The former remark explains that forgetting this identification
identifies conjugate embeddings.  Here both orientations remain in
the theta sum.

Choose the initial orientation to agree with the Heegner point defining
$P_D$.  The constant weight on its orbit is a phase
$\zeta_M(D)\in\mu_{2^\infty}$, since all its local weights are units.
Changing orientation at $p^e\Vert N$ multiplies the Schwartz factor by
$w_{p^e}\chi_{h_*}(p)^e$, and multiplies the norm-character-weighted
parametrization value by the same sign.  At $\ell$ both signs are $-1$,
since $\sigma\chi_{h_*}(\ell)=-1$.  Hence every one of the
$2^{\omega(N)+1}$ orientation tuples contributes the same term.
With the fixed multiplier $L$ already incorporated into $P_M$, the
comparison is consequently
\begin{equation}
 a_M(D)=c_0\,2^{\omega(N)}\zeta_M(D)\,
           \lambda_M\bigl(\operatorname{red}_{\ell}(L P_D)\bigr)
                           \pmod{2^M}.
\label{cs:weighted-trace-identity}
\end{equation}
The initial basis at $\ell$ and initial orientation affect only
$\zeta_M(D)$.  Formula~\eqref{cs:weighted-trace-identity} specifies the
normalization using a character sum, with no class-number denominator.
All rational multipliers are fixed independently of $D$ and $\ell$.

If the product derivative is zero, Gross--Zagier makes $P_D$ torsion.
The fixed multiplier kills it: genus sums and their twist comparisons
are defined over biquadratic fields.  The torsion of a fixed elliptic
curve in degree at most four is uniformly bounded, as one sees by
reducing its prime-to-residue-characteristic torsion at two fixed good
primes of distinct residue characteristic; the residue degrees are at
most four.  Thus the coefficient is zero at
the stated precision.  If the derivative is nonzero, write the point
against the free twist line.  Division by $2^{j(P_D)}$ on that line
implies the same divisibility after reduction and $\lambda_M$, up to
the fixed-degree comparison and the fixed multiplier.  Conversely,
a visible image bounds $j(P_D)$ above by its visible valuation plus a
fixed constant.  Equation~\eqref{cs:weighted-trace-identity} proves
Equations~\eqref{cs:odd-detection} and~\eqref{cs:odd-lower-bound}, using
Equation~\eqref{cs:odd-bound-hypothesis}.  Gross--Zagier--Kolyvagin
then gives the simple zero and finite $\Sha$ when the coefficient is
visible.  The detecting construction for $D_w$ supplies the final
assertion.  Fixed denominators or multipliers only reduce $M$ by a
fixed amount; reindexing the precisions gives the formulation stated.
\end{proof}

\subsection{The least normalized depth}

We now minimize the coefficient valuation after subtracting the prime
weight $w(D)$. This normalization is fixed before any prime array is
chosen. Attainment of the minimum will detect analytic rank and bound
the finite Selmer contribution; for the odd source this is interpreted
through the bounded-weight sequences and finite realizations below.

For the even source, first rescale $g$ so that $a(D_0)=1$ at one
nonzero squarefree index. Its nonzero squarefree valuations are then
integers by Proposition~\ref{prop:even-coefficients}. Set
\begin{equation}
              \delta=\min_D\{v(a(D))-w(D)\},\qquad
 A(D)=\begin{cases}
 1,&v(a(D))=\delta+w(D),\\
 0,&\text{otherwise},
 \end{cases}
\label{cs:even-symbol}
\end{equation}
with $A(D)=0$ off the filter.  The minimum exists: the lower bound is
uniform, visible valuations are integers after normalization, the
fractional part of the weight is fixed, and a nonzero witness exists.

For the odd source fix a nonprincipal ultrafilter $\mathcal U$ on its
sequence of precisions. Consider all sequences $D_M$ of admissible indices for
which $w(D_M)$ is bounded and $v(a_M(D_M))$ is visible on an
ultrafilter-large set.  For each such sequence, its normalized depth
has a lower bound in a discrete set.  Let $\delta$ be the least depth
attainable on an ultrafilter-large set.  This least depth exists and is
attained: the uniform lower bound bounds it below, discreteness gives
a least value, and the fixed witness bounds the possible minimum above.
For a bounded-weight admissible sequence $\mathbf D=(D_M)$, define
\[
 A(\mathbf D)=1
 \quad\Longleftrightarrow\quad
 \{M:v(a_M(D_M))=\delta+w(D_M)\}\in\mathcal U,
\]
and set $A(\mathbf D)=0$ otherwise. We suppress the sequence notation
when discussing finite configurations; each application uses their
simultaneous finite realizations, as proved below.

For the odd source, a least weight supporting a unit symbol is minimized
among the same bounded-weight sequences that attain depth $\delta$.
After restricting to an ultrafilter-large set, the weight of such a
sequence may be held constant, since it lies in a bounded subset of
$\tfrac12\Z_{\ge0}$. The set of these constant weights is nonempty and
has a least element $m$. If unit symbols of weight below $m$ occurred
on an ultrafilter-large set of stages, choosing one such index at each
stage would contradict the definition of $m$.

\begin{lemma}[Uniformity at a fixed weight bound]
\label{cs:fixed-weight-uniformity}
For every fixed $W$, on an ultrafilter-large set of sufficiently high
precisions,
\[
                   v(a_M(D))\geq\delta+w(D)
\]
for every admissible index with $w(D)\leq W$, with valuations truncated
at the working precision.  Any finite list of stabilized symbol
assertions of bounded weight is realized simultaneously at such
precisions.  The definition allows arbitrarily many weight-zero root
factors.
\end{lemma}

\begin{proof}
If the first assertion failed on an ultrafilter-large set, choose a
violating $D_M$ at each of those levels.  Their weights are bounded by
$W$ and their normalized depths lie in a fixed bounded discrete range.
After selecting one depth on an ultrafilter-large set, they contradict
the definition of $\delta$.  For a finite collection of stabilized
bits, intersect their ultrafilter-large realization sets.  A bound on
weight, rather than on the number or size of primes, is all that this
argument uses.  At weights beyond the chosen bound, the truncated
lower estimate of Proposition~\ref{prop:odd-coefficients} still
supplies any prescribed bounded-depth divisibility test once the weight
is sufficiently large.  No claim of one finite level valid for all
unbounded weights is made.
\end{proof}

\begin{proposition}[Detection by a unit symbol]
\label{prop:unit-symbol-detection}
For either coefficient source, $A(D)=1$ implies the prescribed analytic
order: zero in the even case and one in the odd case.  At every such
index,
\begin{align}
 v({\cal L}(h))&\leq2w(D)+C &&\text{in the even case},
 \label{cs:even-sharp}\\
 j(P_D)&\leq w(D)+C &&\text{in the odd case},
 \label{cs:odd-sharp}\\
 s(h)&\leq C,\qquad \dim_{\mathbf F_2}\Sel_2(E^{(h)}/\Q)\leq C.
 \label{cs:unit-selmer-bound}
\end{align}
The constants are fixed after choosing the original curve, filter,
coefficient normalization and genus partner.  For the odd source they
are independent of the detecting prime and the working precision, and
of the number and size of active primes.  Statements about its symbols
are used through finite realizations as in
Lemma~\ref{cs:fixed-weight-uniformity}.
\end{proposition}

\begin{proof}
In the even case substitute $v(a(D))=\delta+w(D)$ into
Equations~\eqref{cs:even-value} and~\eqref{cs:even-lower-bound}.
In the odd case use a precision greater than the finitely many tested
depths and apply Equations~\eqref{cs:odd-detection} and
\eqref{cs:odd-lower-bound}.  The constant $\delta$ is fixed and the
comparison constants are uniform, proving the first inequalities and
the bound on $s(h)$.  At these indices the analytic order is established,
so Gross--Zagier--Kolyvagin gives the same algebraic rank and finite
$\Sha$ \citep{GrossZagier1986,Kolyvagin1990}.  The exact sequence for $2$-Selmer then bounds its dimension
by that rank, the fixed bound on rational $2$-torsion, and the
$2$-primary length of $\Sha$.  This proves the final assertion without
making any finiteness assumption at an index where the analytic order
has not yet been detected.
\end{proof}

\section{Trace identities and prime substitutions}
\label{sec:trace-calculus}

We now turn the least-depth coefficients into rules for replacing,
contracting, and deleting primes. The input is a coefficient system
from Section~\ref{cs:coefficient-systems}, with fixed support $S$,
characters $\eta,\psi$, expansion $g=\sum_{n\ge1}a(n)q^n$,
minimum depth $\delta$, and symbol $A(D)\in\F_2$, extended by zero
off its filter. The coefficient lower bounds and the rational leading
units proved there are essential hypotheses. A unit symbol has the
analytic interpretation of Proposition~\ref{prop:unit-symbol-detection}.

Put $F=\Q(E[2])$.  A prime outside $S$ is called split, simple, or a root
according as its action on $E[2]$ is the identity, a transposition, or an
element of order three; its weight is respectively $1$, $1/2$, or $0$.
The same terminology will be used for Galois elements, according to their
image in $\Gal(F/\Q)$.  A nonroot is an element of either of the first two
types.  A finite set of primes and its squarefree product will occasionally
be denoted by the same letter.  Thus $U/p$ means deletion of $p$ from $U$,
not deletion of a prime power from an arbitrary integer.

An oriented Frobenius label in an extension $L$ of $F$ means a
\emph{rational} arithmetic Frobenius lift in $\Gal(L/\Q)$, with
a chosen place and its residual image in $\Gal(F/\Q)$ retained.
At a simple prime this is the lift $g$ of the transposition, not
only the Frobenius $g^2$ of the residue-degree-two place of $F$.

For the odd source, fix the finite operator patterns, their weight bounds,
and all integer divisions required by the intended tests before choosing
an actual stage. The primes realizing those patterns may depend on the
stage. The lemma below distinguishes the sharp estimate on a prescribed
bounded range of squarefree weights from integrality at every Fourier
index. Both are needed to form the integral weight-two trace lattice. At a
chosen stage the detecting prime belongs to $S$, and the level,
coefficient field, and characteristic-zero lift are fixed. Recurrences
hold modulo the working precision before division.

\subsection{Normalized coefficient operators}

Let $U$ be a finite set of nonroot primes outside $S$.  Write
\[
 \mathsf U_p\!\left(\sum_n c(n)q^n\right)=\sum_n c(pn)q^n,
 \qquad
 \mathsf V_p\!\left(\sum_n c(n)q^n\right)=\sum_n c(n)q^{pn},
\]
and define
\begin{equation}
 B_U=2^{-\lceil\delta+w(U)\rceil}
             \prod_{p\in U}(\mathsf U_p-\eta(p)\mathsf V_p)g,
 \qquad
 e_U=\lceil\delta+w(U)\rceil-\delta-w(U).
 \label{tc:normalized-operator}
\end{equation}
The factors commute.  We say that $U$ has matching parity if
$\delta+w(U)$ is integral, and opposite parity otherwise.  Equivalently,
$w(U)$ has, or does not have, the fractional part prescribed by the filter.
Consequently $e_U=0$ or $1/2$, respectively.

These operations increase the level at a newly introduced $p$ only to
exponent one.  In fact, $\mathsf V_p$ changes the character by
$(p/\cdot)$ and raises the level by $p$, while on that raised level
$\mathsf U_p=\mathsf U_{p^2}\mathsf V_p$.  We denote the resulting
character of $B_U$ by $\psi_U$.  At a good prime $q\notin S\cup U$ the
coefficient recurrence remains
\begin{equation}
 B_U[q^2n]
 =\bigl(\eta(q)a_q(E)-\psi_U(q)(-n/q)\bigr)B_U[n]
       -q\eta(q)^2B_U[n/q^2].
 \label{tc:transformed-recurrence}
\end{equation}
For the odd source, Equation~\eqref{tc:transformed-recurrence} is read
modulo $2^{M-\lceil\delta+w(U)\rceil}$ when $M$ is the working
recurrence precision before normalization. As usual, a coefficient with
nonintegral index is zero. The equality of
the good eigenvalue before and after $\mathsf V_p$ can also be seen
directly: its new middle character contains $(p/q)$, which cancels the
factor $(p/q)$ arising from the change of coefficient index.

\begin{lemma}[Coefficient bound and first shell]
\label{tc:shell-bound}
Write an index prime to $S$ uniquely as
\[
             n=Jr m^2,
 \qquad J\subseteq U,
 \qquad r\text{ squarefree},\quad (r,U)=1,
\]
and put $D=(U/J)r$.

For the even source, $B_U$ has integral coefficients and satisfies the
following assertions at every index. For the odd source, fix bounds on
$w(U)$ and on the squarefree weights to be tested, and fix an integer
$K\ge0$. There is a larger cutoff $W$ and an ultrafilter-large set of
arbitrarily high stages at which $B_U$ is integral at every index, the
assertions below hold whenever $w(D)\le W$, and
\[
                  v(B_U[n])\ge K\qquad\text{if }w(D)>W.
\]
These choices are uniform over prime realizations with the prescribed
weight bounds.

The coefficient $B_U[n]$ is zero unless $D$ belongs to the filter.
Within the range just specified, it satisfies
\begin{equation}
 v(B_U[n])\geq
       w(r)+\frac{\#\{p\in J:p\text{ simple}\}}2-e_U.
 \label{tc:coefficient-bound}
\end{equation}
At a split prime $p\in J$, the factor contributed by the difference
operator, after removing its character phase and one factor of $2$, is
\begin{equation}
 c_p(D)=
 \frac{a_p(E)-(\psi/\eta)(p)(-D/p)-1}{2}\pmod 2.
 \label{tc:shell-factor}
\end{equation}
This is independent of the exponent of $p$ in $m$. At a simple prime
$p\in J$, Equation~\eqref{tc:shell-factor} holds for the first shell,
that is, when $p\nmid m$.
\end{lemma}

\begin{proof}
The good-prime recurrence of Section~\ref{cs:coefficient-systems} expresses
every relevant coefficient in terms of the squarefree coefficient $a(D)$.
After removing the factor $\eta(p)^s$, the multipliers for powers $p^{2s}$
are determined by
\[
 b_0=1,\qquad
 b_1=a_p(E)-(\psi/\eta)(p)(-D/p),\qquad
 b_{s+1}=a_p(E)b_s-pb_{s-1}.
\]
If $p\mid D$, the Legendre symbol in $b_1$ is zero.  The recurrences at
different primes commute and their multipliers multiply.

For $p\in U/J$ the difference operator contributes
$b_s-b_{s-1}$, with $b_{-1}=0$.  Here $p\mid D$, so $b_1=a_p(E)$ is
even.  Modulo $2$ the sequence alternates between $1$ and $0$, and every
such difference is odd.  For $p\in J$ we have $p\nmid D$; now $b_0$
and $b_1$ are odd and every $b_s$ is odd.  The difference
$b_{s+1}-b_s$ is therefore even.

For the more precise assertion at a split prime put
$a=a_p(E)$ and $\kappa=(\psi/\eta)(p)(-D/p)\in\{1,-1\}$.
Since $E[2]$ is rational over $\mathbf F_p$,
$a\equiv p+1\pmod4$.  If $\Delta_s=b_{s+1}-b_s$, the recurrence gives
\[
 \Delta_{s+1}=a\Delta_s-p\Delta_{s-1}\quad(s\geq1).
\]
Every $\Delta_s$ is even, whence
$\Delta_{s+1}\equiv\Delta_{s-1}\pmod4$.  Directly,
\[
 \Delta_1-\Delta_0
 =a^2-a\kappa-2a+2\kappa+1-p\equiv0\pmod4.
\]
Thus $\Delta_s/2\equiv(a-\kappa-1)/2\pmod2$ for every $s$.
For a simple prime this last formula is still the defining calculation
at $s=0$, which is all that will be required.

Put $d_U=\lceil\delta+w(U)\rceil$. Whenever the sharp squarefree
bound $v(a(D))\ge\delta+w(D)$ is available, the factors just computed
give at least one further factor of $2$ for each $p\in J$. Hence
\[
 v(B_U[n])\ge \delta+w(D)+\#J-d_U
 =w(r)+\frac{\#\{p\in J:p\text{ simple}\}}2-e_U.
\]
For an allowed squarefree class this expression is integral, since
$\delta+w(D)$ is integral. Its only possible negative value before
this parity condition is imposed is $-1/2$, so it is nonnegative.
For the even source the sharp bound holds at every squarefree index,
which proves all-index integrality and the displayed estimate.

For the odd source, let $u_0$ bound $w(U)$. Every prime in $U$ is a
nonroot, so $\#U\le2u_0$ and the normalizing exponents $d_U$ are
bounded. Proposition~\ref{prop:odd-coefficients} gives a constant
$C$ such that at every admissible squarefree index
\[
                   v(a_M(D))\ge\min\{M,w(D)-C\}.
\]
This includes the case of zero product derivative, when the coefficient
vanishes modulo $2^M$. Choose $W$ above the prescribed tested weights
and so large that $W-C-d_U>K$ for all permitted $U$. By
Lemma~\ref{cs:fixed-weight-uniformity}, on an ultrafilter-large set
of stages the sharp squarefree bound holds simultaneously for every
$D$ with $w(D)\le W$. Take an arbitrarily high stage in this set,
with precision exceeding all the fixed normalizations, the first-shell
tests, and $\delta+W+2u_0$, and such that $M-d_U\ge K$
for every permitted $U$.

The recurrences and operator identities for this source hold modulo
the working precision. Their multipliers are integral, so iteration
through any square exponent leaves errors in the same precision ideal.
After dividing by $2^{d_U}$, the possible error has valuation at least
$M-d_U$. Thus the sharp calculation above is valid throughout the
bounded range. Outside that range the same calculation with the coarse
squarefree bound gives
\[
 v(B_U[n])\ge
 \min\{M-d_U,\,w(D)-C+\#J-d_U\}\ge K.
\]
This proves integrality at every index of the actual form, while
asserting the sharp estimate only on the specified bounded range.
There is no restriction on the number of weight-zero root factors.
Coefficients excluded by the filter are zero. Every removed phase
is a value of $\eta$, hence a unit reducing to $1$ at the chosen
prime above $2$.
\end{proof}

At equality, the normalized leading coefficient is $A(D)$ times the
displayed shell factors. In particular it is zero when $A(D)=0$.
For the odd source this assertion is used for the finite stabilized
symbol evaluations realized at the chosen stage. All omitted-prime
comparisons use the same source $g$ and minimum $\delta$.

In later coefficientwise divisions, choose $K$ to cover the largest
additional division and the desired final divisibility. If a least-weight
argument also requires all lower-weight symbols to vanish, that condition
holds simultaneously on an ultrafilter-large set: a failure would select
a bounded-weight sequence contradicting the stabilized minimum. Intersect
that set with the one in the lemma. The bounded range then supplies the
sharp calculation and any required zero symbols; the coefficients beyond
the cutoff already have all the necessary factors of $2$. Consequently
all coefficientwise divisions and vanishings used below hold on actual
integral forms at one chosen stage for each finite collection of tests.
No stage is required to work for every unbounded configuration.

\subsection{Removing unary coefficients}

The coefficient bound identifies the terms that can survive at the
lowest valuation. For matching $U$ in the rational-torsion case, these
lie on finitely many square classes. Subtracting the corresponding unary
theta series will permit one further integral division by $2$.

Suppose first that $E(\Q)[2]\ne0$.  There are then no root primes.
For matching $U$, define
\begin{equation}
 \begin{split}
 \Theta_J&=
    \sum_{\substack{m>0\\(m,\prod_{s\in S}s)=1}}
       \eta(m)\chi_{-4}(m)^i m\,q^{Jm^2},\\
 H_U&=\frac12\left(B_U-
       \sum_{\substack{J\subseteq U\\J\text{ split only}}}
                    B_U[J]\Theta_J\right),
 \end{split}
 \label{tc:unary-subtraction}
\end{equation}
where $i\in\{0,1\}$ is chosen so that $\eta\chi_{-4}^i$ is odd.
For opposite $U$, set $H_U=B_U$.  Except where a further division is
explicitly made in Lemma~\ref{tc:minimum-transfer}, in the irreducible
case we also set $H_U=B_U$.

\begin{lemma}
\label{tc:unary-integrality}
The form in Equation~\eqref{tc:unary-subtraction} is integral.  Each
$\Theta_J$ is a weight-$3/2$ unary theta form with character
\[
                \eta\chi_{-4}^i(-J/\cdot),
\]
whose new level primes are the primes of $J$, each to exponent one.
Its character differs from $\psi_U$ by a quadratic character.
\end{lemma}

\begin{proof}
In matching parity the right side of
Equation~\eqref{tc:coefficient-bound} can be zero only if $r=1$ and
$J$ consists of split primes.  At these indices
Equation~\eqref{tc:shell-factor}, together with the odd squared-index
multipliers at primes outside $U$, gives
\[
                  B_U[Jm^2]\equiv B_U[J]\eta(m)\pmod2.
\]
The coefficient of $\Theta_J$ at $Jm^2$ has precisely this residue:
the common phase $\eta(m)$ is retained on both sides, while
$\chi_{-4}(m)^i m$ is an odd integer.  Thus this is divisibility by
$2$, not merely by $\mathfrak m$ in a ramified coefficient ring.
Different squarefree $J$ have disjoint such supports.  The numerator in
Equation~\eqref{tc:unary-subtraction} is consequently divisible by $2$
coefficientwise.

Put $\beta=\eta\chi_{-4}^i$, and let $\beta_0$ be its primitive
inducing character, of conductor $r$ supported on $S$.  The unary
theta transformation \citep[Section~2, p.~2]{Purkait2012} gives
\[
 \Theta_J^0=\sum_{m>0}\beta_0(m)m q^{Jm^2}
 \in S_{3/2}\bigl(4r^2J,\beta_0(-4J/\cdot)\bigr).
\]
Removing the terms with $s\mid m$ applies
$1-s\beta_0(s)\mathsf V_{s^2}$, where
$\mathsf V_{s^2}=\mathsf V_s^2$.  Thus
\[
 \Theta_J=\prod_{s\in S}
       \bigl(1-s\beta_0(s)\mathsf V_{s^2}\bigr)\Theta_J^0.
\]
Writing $P_S=\prod_{s\in S}s$, every term in this product has theta
parameter $Jd^2$ for some $d\mid P_S$.  The same transformation formula
places all terms at level $4r^2JP_S^2$, with character
$\beta_0(-4J/\cdot)$ on the units of that level.  This is the displayed
character $\eta\chi_{-4}^i(-J/\cdot)$.  Since $J$ is prime to $S$,
each fresh prime in $J$ still occurs in the level to exponent one.
Finally $\psi_U/\eta$ is quadratic, as is
$\chi_{-4}^i(-J/\cdot)$, proving the last assertion.
\end{proof}

\subsection{Weight-two traces}

We next express coefficient tests as traces of Galois representations.
The trace identities will give relations among the tests; the
level-prime identity will compare a test with the one obtained by
omitting a prime.

Multiply the integral half-integral-weight form under consideration by
\[
                  \theta=1+2\sum_{m\geq1}q^{m^2}.
\]
This produces a weight-two modular form with the same reduced expansion.
Choose a finite extension of $\Qtwo$ containing all coefficient fields
and character values needed at the fixed level, and denote its ring of
integers and maximal ideal by $\mathcal O$ and $\mathfrak m$.

The characteristic-zero decomposition by the Hecke operators away from
the level supplies commuting operators
\[
               t(u),\ d(u),\qquad u\in G_{\Q},
\]
by taking the traces and determinants of the corresponding
two-dimensional Galois representations.  For cuspidal systems these
are the weight-two representations of Deligne; for Eisenstein systems
they are the sums of the two defining characters.  At a good arithmetic
Frobenius, $t(\Fr_q)$ is the ordinary weight-two $T_q$, and the
determinant is the cyclotomic character times the diamond character.
We use the existence and good-prime compatibility of these representations
in their characteristic-zero form
\citep[\S3, (3.9)--(3.12), and Theorem~4.9]{Deligne1969}, and
local-global compatibility
\citep[Theorem~(A), \S0.7]{Carayol1986}.
We take the contragredients of the cohomological representations in these
references, so arithmetic Frobenius has trace $a_q$ and determinant $q$
times the diamond eigenvalue.

Here is a precise integral framework for the reductions that follow.
Take the finite-dimensional sum of these characteristic-zero systems,
retaining their multiplicity spaces, that contains the vector under
consideration and its diamond transforms.  Its integral Fourier lattice
is finite over $\mathcal O$: finitely many coefficients already give an
injective map of this space, and the lattice embeds in the corresponding
finite free $\mathcal O$-module.  Good Hecke operators preserve integral
expansions.  Diamonds do so as well.  For the subtracted vector this last
claim follows from the formula
\[
 \langle u\rangle H_U
 =\psi_U(u)H_U+
 \frac12\sum_J B_U[J]
       \bigl(\psi_U(u)-\eta(u)\chi_{-4}(u)^i(-J/u)\bigr)\Theta_J:
\]
each parenthesis is zero or twice a unit.  The same observation applies
after multiplying by $\theta$.

Traces at arbitrary Galois elements are limits of good Frobenius traces
on this finite lattice.  Indeed, continuity gives a common finite
quotient at any prescribed accuracy, and Chebotarev gives a good prime
in each required conjugacy class.  Its size is unrestricted, and it may
avoid any specified finite list of indices and primes.  This proves
integrality of all the traces and determinants used below, not only of
the good Hecke operators with which the construction began.

Let $\mathcal T$ be the resulting commutative operator algebra on the
reduced orbit of the vector. Equality means equality of every Fourier
coefficient after any further indicated operator products. We retain
this algebra, including its possible nilpotents.

\begin{lemma}[Trace identities and fresh-prime tests]
\label{tc:trace-algebra}\label{co:traces}
In this reduced operator algebra,
\begin{equation}
 d(u)t(v)=t(v),\qquad
 t(u)t(v)=t(uv)+t(uv^{-1}),\qquad t(1)=0.
 \label{tc:trace-identities}
\end{equation}
A coefficient of a product of traces can be tested by replacing its
arguments with Frobenius elements at distinct fresh primes avoiding
the coefficient index.  If the index is $n$ and those primes are
$q_1,\ldots,q_a$, the test is the reduced coefficient at
$nq_1\cdots q_a$.  Determinants may be applied first.
\end{lemma}

\begin{proof}
For a weight-two form $f$, the good-prime formula is
\[
                  (T_qf)[n]=f[qn]+q(\langle q\rangle f)[n/q].
\]
When $q\nmid n$ the second term is absent.  Choosing the successive
test primes distinct and avoiding $n$ proves the coefficient assertion.
Chebotarev and continuity, on the finite lattice just described, allow
these choices for arbitrary trace arguments at any required accuracy.

All the diamond characters are of $2$-power order.  On a form with one
character its determinant therefore reduces to the identity.  For a
subtracted vector the displayed diamond formula leaves, in addition,
only unary contributions $\Theta_J\theta$.  A fresh trace kills each
such contribution modulo $\mathfrak m$: its coefficient test has
index $nq$ with $q\nmid nJ$, and cannot have the squarefree part $J$.
The same argument applies after any further tests.  This proves
$d(u)t(v)=t(v)$.

For two invertible $2$-by-$2$ matrices the characteristic-zero identity
is
\[
 \tr(A)\tr(B)=\tr(AB)+\det(B)\tr(AB^{-1}).
\]
Apply it to every characteristic-zero system and use the determinant
identity just proved.  Finally, $t(1)=2$ before reduction.
\end{proof}

The bounds of Lemma~\ref{tc:shell-bound} are particularly well suited to
these tests: an added prime contributes its nonnegative weight, and
additional root primes contribute no weight.  A zero established by
one of these bounds consequently remains zero after all the further
tests used here.  In the moving-precision case all the preceding
constructions are made at the actual fixed level of the chosen stage.
The required divisibility is uniform at the specified weight cutoff;
continuity is used only in a finite-dimensional space at that stage.

The only level-prime identity we need is the following.  We distinguish
the weight-two operator $\mathsf U_p^{(2)}$ from the earlier
half-integral-weight notation, although both extract $p$-multiples in
their respective Fourier expansions.

\begin{lemma}[Inertia and a level operator]
\label{tc:inertia-identity}\label{co:inertia}
For $p\in U$, choose a place over $p$, an inertia element $\tau_p$
on which the ramified quadratic character is $-1$, and a corresponding
arithmetic Frobenius lift $\phi_p$.  On the characteristic-zero space
under consideration, before reduction, one has
\begin{equation}
 \mathsf U_p^{(2)}
       \bigl(t(g\tau_p\phi_p)-t(g\phi_p)\bigr)
 =d(\phi_p)\bigl(t(g\tau_p)-t(g)\bigr).
 \label{tc:exact-level-identity}
\end{equation}
The identity commutes with further good traces.  After reduction it
therefore evaluates a difference at coefficient $1$ by the corresponding
difference at coefficient $p$.
\end{lemma}

\begin{proof}
At $p$ the summands in question are either old from a level omitting
$p$, or have conductor exponent one and primitive ramified quadratic
diamond character.  For the former the representation is unramified,
so both trace differences are zero, irrespective of the action of
$\mathsf U_p^{(2)}$ on the old multiplicity space.

For the latter, local-global compatibility identifies the local
representation with a principal series with one unramified and one
ramified character; its monodromy is zero.  In a basis of these two
lines write
\[
 \rho(\tau_p)=\begin{pmatrix}1&0\\0&-1\end{pmatrix},
 \qquad
 \rho(\phi_p)=\begin{pmatrix}\lambda&0\\0&\mu\end{pmatrix}.
\]
The level eigenvalue is $\lambda$, the Frobenius eigenvalue on the
unramified line, and $d(\phi_p)=\lambda\mu$.  If the lower-right entry
of $\rho(g)$ is $z$, the two trace differences are $-2z\mu$ and
$-2z$, respectively.  Both sides of
Equation~\eqref{tc:exact-level-identity} are $-2z\lambda\mu$.
The Eisenstein calculation is the identical calculation on its pair
of characters.  Copies old at other primes do not change it.

These are all the possibilities: the new level has exponent one at
$p$, with ramified quadratic determinant, so the Steinberg case with
unramified determinant does not enter.  A unary summand with $p\notin J$
is old from a level omitting $p$; with $p\in J$ it has the same ramified
quadratic local character.  Thus the preceding classification also
covers Equation~\eqref{tc:unary-subtraction}.

Finally apply Equation~\eqref{tc:trace-identities} to remove the
determinant from the right side after reduction.  Taking coefficient
$1$ on the left is taking coefficient $p$ before $\mathsf U_p^{(2)}$.
\end{proof}

\subsection{Rules when there is rational two-torsion}

We now apply the common trace identities to the rational-torsion case.
The coefficient bounds first restrict the Galois data on which a symbol
can depend. The level-prime identity then determines how the symbol
changes when one of those data is varied.

An elementary extension of a number field means a compositum of
quadratic extensions.  Whenever a ramification support is specified,
only extensions unramified outside that support are included.  For
finite support the resulting elementary extension is finite.  Labels
below are Frobenius specifications in these extensions, with a chosen
place when an orientation over $F$ is needed.  They always include the
fixed rational local conditions that specify the filter.

\begin{proposition}[Rational trace and contraction rules]
\label{tc:rational-rules}
Suppose $E(\Q)[2]\ne0$.
\begin{enumerate}[label=\textup{(\roman*)}]
\item For matching $U$, on the subtracted vector $H_U$,
\[
                         t(g)t(u)=0\quad(u\text{ split}).
\]
The trace on split elements factors through the elementary extension
of $F$ supported on $S\cup U$.  For opposite $U$, on $B_U$, every
product of two traces is zero, and traces factor through the rational
elementary extension supported there.
\item If $U\cup\{q\}$ has the parity of the filter, its symbol is
read at coefficient $1$ by $t(\Fr_q)$ on $H_U$ for split $q$, and on
$B_U$ for simple $q$.
\item Fix $p\in U$ and put $g=\Fr_q$.  Replacing $g$ by $g\tau_p$
changes the symbol, additively in $\mathbf F_2$, by
\begin{equation}
                  A\bigl((U/p)\cup\{q'\}\bigr),
       \qquad \Fr_{q'}\sim g\phi_p,
 \label{tc:rational-contraction}
\end{equation}
if $p$ and $q$ are not both simple.  The label of $q'$ is required
only away from the omitted primes and carries their product labels.
The replacement uses the original source with level support
$S\cup(U/p)$: the \emph{whole rational prime} $p$ is omitted,
including both places over it if $p$ splits in $F$. If both endpoints
are simple, the change is zero.
\item A split vertex with trivial label relative to the level with
that vertex omitted has symbol zero.
\end{enumerate}
\end{proposition}

\begin{proof}
For matching $U$, two fresh tests, one of them split, add weight at
least $3/2$ to Equation~\eqref{tc:coefficient-bound}.  The allowed
valuation is integral, so the numerator has valuation at least $2$.
After the division in Equation~\eqref{tc:unary-subtraction} it still
vanishes modulo $\mathfrak m$; fresh tests kill its unary terms.
For opposite $U$, two nonroot tests add weight at least $1$, leaving
at least $1/2$ after subtracting $e_U$, and hence at least $1$ at an
allowed coefficient.  This proves the two product vanishings on the
entire reduced orbit.

The trace identities show why these vanishings give the stated
elementary quotients.  If $u$ is split, apply the identity to
$t(gu)t(u)$ to obtain $t(gu^2)=t(g)$.  The subgroup generated by
squares in $G_F$ has elementary quotient, since a group of exponent
two is abelian.  In the opposite case the same calculation applies
to every $u\in G_\Q$.  Ramification outside the level is absent from
the characteristic-zero representations and remains absent in the
quotient.

For a split $q$, its added weight is $1$: at coefficient $q$ the
division by $2$ in $H_U$ places the coefficient exactly at the
normalized detection depth.  There is no unary term at that index.
For simple $q$, $e_U=1/2$ already gives that depth on $B_U$.
The character phases reduce to $1$, proving (ii).

For (iii), apply Lemma~\ref{tc:inertia-identity} at coefficient $1$,
so that its left side is a coefficient-$p$ test. Replace its trace
arguments by fresh primes with Frobenius $g\tau_p\phi_p$ and
$g\phi_p$. Put $V=U/p$. For $X=U,V$, set
$n_X=\lceil\delta+w(X)\rceil$ and let $m_X$ be one in matching
parity and zero otherwise. After a fresh
trace kills the unary terms, the shell multiplier relative to the
omitted-source test is
\[
 2^{1+n_V-n_U-m_U+m_V}c_p(D'),
 \qquad D'=Vq',\quad \Fr_{q'}\sim g\phi_p.
\]
The exponent is zero for split $p$, and also for simple $p$ with
matching $U$. It is one for simple $p$ with opposite $U$, precisely
the simple--simple case. This checks all integer divisions in the
contraction rule. For simple $p$ only the first shell is used.
The change $\psi_U/\psi_V=(p/\cdot)$ is already included in that
shell calculation; all remaining phases reduce to one.

In the first three cases the shell is therefore $c_p(D')A(D')$.
Both fresh Frobenius elements have the same image in the omitted
source: that source is built from the original modular form, the same
$\eta,\psi,\delta$, and the operators at $V$, so it is unramified at the entire rational prime
$p$. Inertia at either place over $p$ acts trivially on its symbol.
The ramified quadratic coordinate of $\tau_p$ reverses the Legendre
sign in $c_p$, whose additive change is one. The difference of the
two shell values is consequently $A(D')$, proving
Equation~\eqref{tc:rational-contraction}.

If both primes are simple, their residual images are the same
transposition, so $g\phi_p$ is split.  The shell now has positive
valuation after normalization and contributes zero.  Finally (iv)
is (ii) with $t(\Fr_q)=t(1)=0$ in the relevant elementary quotient.
\end{proof}

The contraction rules control changes in existing vertices. The next
calculation gives a way to add two simple vertices while preserving the
symbol.

\begin{proposition}[Insertion of a pair of simple primes]
\label{tc:simple-twins}
Suppose $E(\Q)[2]\ne0$ and $U$ has matching parity.  Let a rational
elementary label, relative to $S$ and all existing split vertices,
be realized by a simple prime $q$ with
\[
                         \#E(\mathbf F_q)\equiv2\pmod4.
\]
Inserting two fresh simple primes with this same label leaves $A(U)$
unchanged.  Their mutual simple-edge label, and their labels against
other simple vertices, do not affect this assertion.
\end{proposition}

\begin{proof}
First test $t(\Fr_q)^2$ on matching $H_U$.  The weight-two Hecke
formula at coefficient $1$ gives, after reduction,
\[
                    H_U[q^2]+q(\langle q\rangle H_U)[1].
\]
Only the $J=1$ unary term can contribute at $1$ or $q^2$.
Write $b=B_U[1]$, $\epsilon=\chi_{-4}(q)$, and
$\kappa=(\psi_U/\eta)(q)$.  Substituting
Equation~\eqref{tc:transformed-recurrence} and the two characters in
Equation~\eqref{tc:unary-subtraction}, the preceding expression is
the reduction of
\begin{equation}
 b\,\frac{\eta(q)}2
 \left(a_q(E)+\kappa(q-\epsilon)
                 -q\epsilon^i(1+\epsilon)\right).
 \label{tc:twin-calculation}
\end{equation}
For $q\equiv3\pmod4$ the parenthesis divided by $2$ is
$a_q(E)/2$ modulo $2$.  For $q\equiv1\pmod4$ it is
$a_q(E)/2+1$.  In either case it equals
$(q+1-a_q(E))/2=\#E(\mathbf F_q)/2$ modulo $2$.
Since $b$ reduces to $A(U)$, this proves that the squared trace
has coefficient $(\#E(\mathbf F_q)/2)A(U)$.

By Lemma~\ref{tc:trace-algebra}, replace the two copies of this
trace by distinct fresh primes realizing the same operator tests.
Their product coefficient reads the symbol with the two primes
inserted.  To pass to any other pair with the labels in the
statement, vary one prime at a time.  After omitting that simple
prime the opposite-parity test of
Proposition~\ref{tc:rational-rules} depends only on rational
elementary labels.  Its changes against other simple primes are
zero by the simple--simple case of the same proposition.  Thus
only the fixed support and existing split labels need be matched.
Fresh intermediate primes can match both old and new finite
specifications, so that no collision of vertices is required.
The assumed odd value of $\#E(\mathbf F_q)/2$ completes the proof.
\end{proof}

\subsection{An elementary character on root elements}

Assume now that $E[2]$ is irreducible.  Its image is either cyclic
of order three or $S_3$.  The reduced trace algebra need not be a
field: the order-two units below can have nontrivial nilpotent
parts.  We therefore retain the algebra, rather than replacing it
by its residue field.

\begin{lemma}[Root character]
\label{tc:root-character}
On matching $B_U$ one has $t(u)=0$ for every nonroot $u$.
On opposite $B_U$ one has $t(u)=0$ for split $u$ and
$t(g)t(u)=0$ for two simple elements.  In either parity,
\[
                         t(u)^2=1\quad(u\text{ a root}).
\]
The traces factor through an extension of $\Gal(F/\Q)$ by an
elementary kernel over $F$.  Moreover, there is a multiplicative
character $\alpha:G_\Q\longrightarrow\mathcal T^\times$ into the
reduced trace algebra, unramified outside $S\cup U$, such that
\[
                  \alpha(u)^2=1\quad\text{for all }u,
 \qquad       t(u)=\alpha(u)\quad(u\text{ a root}).
\]
In particular $\alpha$ factors through a rational elementary
extension with that ramification support.
\end{lemma}

\begin{proof}
The trace vanishings follow from
Equation~\eqref{tc:coefficient-bound}.  In matching parity, one
nonroot test adds at least $1/2$ and hence gives an integral
valuation at least $1$.  In opposite parity, either a split test
or two simple tests add at least $1$, again leaving positive
integral valuation after the loss $e_U=1/2$.

For a root test, use the same fresh prime $q$ twice, avoiding the
coefficient index, rather than two distinct primes.  The weight-two
Hecke formula and Equation~\eqref{tc:transformed-recurrence} give
the identity operator modulo $\mathfrak m$: $a_q(E)$ is odd, and
the middle quadratic-character term cancels the diamond term
modulo $2$.  This proves $t(u)^2=1$ by continuity for every root
$u$.  For completeness, it suffices to check indices prime to the
chosen $q$; each particular index permits such a choice of $q$
with the prescribed finite trace data.

Since split traces vanish in both parities,
$t(gu^2)=t(g)$ for $u\in G_F$.  Thus the traces factor through
the quotient by the squares of $G_F$.  At any fixed accuracy the
trace action is on a finite lattice modulo $\mathfrak m$, so we
may work in a finite quotient
\[
       1\longrightarrow H\longrightarrow\Gamma
          \longrightarrow\Gal(F/\Q)\longrightarrow1
\]
with $H$ elementary abelian.  Choose an order-three lift $R$ of a
rotation.  Such a lift exists because the inverse image of the
order-three subgroup is an extension by a finite $2$-group.
The trace identity with the identity element gives
$t(R)=t(R^{-1})=t(R^2)$.  Applying it to $R,R^2$ gives
$t(R)t(R^2)=t(1)+t(R^{-1})=t(R)$; the left side is $1$.
Hence
\[
                           t(R)=t(R^2)=1.
\]

For $a\in H$ set
\[
                       \alpha(a)=t(aR)=t(aR^2).
\]
The equality follows from $t(a)=0$ and the trace identity with
$R$.  Conjugacy invariance by an element $b\in H$ implies that
$\alpha$ is unchanged on multiplying its argument by
$bR(b)^{-1}$.  It therefore factors through the coinvariants
$H/[H,R]$.  Also it is invariant under conjugation by $R$.
The trace product now gives
\[
 \alpha(a)\alpha(b)
   =t(aR)t(bR)
   =t\bigl(aR(b)R^2\bigr)+t(ab^{-1})
   =\alpha\bigl(aR(b)\bigr)
   =\alpha(ab).
\]
The last equality uses the coinvariant invariance just proved;
the split term is zero.  Since $aR$ is a root,
$\alpha(a)^2=1$.

For cyclic residual image, this already extends to $\Gamma$ by
setting $\alpha(R)=1$.  For $S_3$, choose a transposition lift
$s$ normalizing $\langle R\rangle$.  It can be obtained from any
transposition lift by conjugating its order-three complement by
an element of $H$.  Thus $sRs^{-1}=R^{-1}$ and $s^2\in H$.
Conjugacy of traces shows that $\alpha$ on $H$ is invariant
under $s$.  In both parities we have $t(s)t(sR)=0$, either because
each factor is zero or by the two-simple product rule.  Its trace
identity reads
\[
                       0=\alpha(s^2)+t(R)=\alpha(s^2)+1.
\]
Therefore the extension defined by $\alpha(s)=1$ respects the
relation for $s^2$, as well as conjugation and the relations in
$\langle R\rangle$.  It is the required multiplicative
order-two character, and agrees with the trace on both root
cosets.  At the fixed stage, choose one common finite quotient
through which the traces and residual representation factor, and
pull back the character constructed there.  No canonical choice
across quotients is needed.  Unramifiedness outside the level is preserved.
\end{proof}

\subsection{Irreducible deletion rules and minimal labels}

The root character replaces products of root traces by a single
multiplicative value. We now combine that fact with the level-prime
identity to compute changes at nonroot vertices. The last lemma applies
these rules at the least weight where a symbol is nonzero.

The \emph{total root label} is the product of the rational elementary
labels of all root primes being used.  It includes the fixed-support
conditions and the rational quadratic labels against the active
nonroot vertices.  Its discriminant coordinate is trivial in the
quadratic quotient of $S_3$; for cyclic residual image there is no
such coordinate.

\begin{proposition}[Irreducible substitution rules]
\label{tc:irreducible-rules}
Suppose $E[2]$ is irreducible.
\begin{enumerate}[label=\textup{(\roman*)}]
\item On matching $B_U$, coefficient $1$ of a product of root
traces reads $A$ on the union of $U$ and those roots.  It depends
on the roots only through their total rational elementary label
relative to $S\cup U$.  Any finite collection of roots, including
the empty collection, may be replaced in this test by one virtual
root with that label.
\item Changing the Legendre bit of that root against a split
vertex $p$ has additive derivative equal to deletion of $p$,
with its rational labels away from $p$ transferred to the total
root label.  Changing the root bit against a simple vertex has
derivative zero.
\item To vary a simple vertex $q$, omit it and use opposite
$B_V$, retaining the other root probes.  At another simple
vertex $p$, choose Frobenius lifts $g$ and $\phi_p$ for $q$ and
$p$.  If their transposition images are different, the derivative
under $g\mapsto g\tau_p$ deletes both $p$ and $q$ and transfers
their rational labels to the root total.  If their transposition
images are equal, the derivative is zero.
\item In the test in (iii), variation at a split $p$ is given by
the first-shell factors $c_p$ times symbols of merged
configurations whose weight is one less.  It consequently
vanishes whenever all symbols of that lower weight vanish.
\end{enumerate}
All these assertions preserve every specified label at the other
vertices, and use the chosen place for inertia and Frobenius.
\end{proposition}

\begin{proof}
Root primes have weight zero.  Thus the fresh-prime tests at
coefficient $1$ of matching $B_U$ are precisely the normalized
coefficients defining $A$.  Lemma~\ref{tc:root-character} makes
the product of root traces equal to the value of $\alpha$ on
their total elementary label.  A rotation survives every
compatible elementary specification: an elementary quotient
cannot eliminate the order-three subgroup, and in the $S_3$
case the only compatibility condition is its trivial image in
the quadratic quotient.  Chebotarev therefore supplies the
claimed single virtual root, even for the empty total label.

For (ii), apply Lemma~\ref{tc:inertia-identity} with $g$ a root.
If $p$ is split then $g\phi_p$ is a root.  The coefficient with
$J=\{p\}$ has zero cost in
Equation~\eqref{tc:coefficient-bound}; its first shell is
$c_p$ times the symbol with $p$ deleted.  The latter is
unramified at $p$ and is the same in the two tests, while $c_p$
changes by $1$.  Multiplication of Frobenius labels by $\phi_p$
is exactly the transfer of the deleted vertex's rational labels
to the root.  If $p$ is simple, $g\phi_p$ is simple, and its
weight $1/2$ plus the $1/2$ contributed by $J$ puts the shell
above the testing depth.  The derivative is then zero.

For (iii), $e_V=1/2$.  When the two transpositions differ,
$g\phi_p$ is a rotation, and the $1/2$ contribution from the
simple prime in $J=\{p\}$ exactly cancels $e_V$.  The shell
again equals $c_p$ times the omitted-$p$ symbol, now with both
simple vertices removed and their product represented by a
root.  The Legendre flip proves the asserted derivative.
For equal transpositions their product is split; its additional
weight $1$ forces the shell to vanish.

Finally, if $p$ is split and $g$ simple, then $g\phi_p$ is
simple.  The $p$-shell is at the first-shell depth and is
expressed by Equation~\eqref{tc:shell-factor} with
$D'=(V/p)q'$ and the retained root probes.  Its nonroot weight
is
\[
                 w(V)-1+\frac12
                  =\left(w(V)+\frac12\right)-1.
\]
This is one less than the original weight.  In particular both
shell terms vanish when symbols of that weight vanish.  The
comparison omitting $p$ is unramified there.  Root probes can
be matched on the rational elementary support throughout,
which proves the final label assertion as well.
\end{proof}

At the least nonzero weight, every lower-weight term in the preceding
rules vanishes. This removes the dependence on internal edge labels;
the remaining assertion keeps both the fixed-support labels and the
total root label.

\begin{lemma}[Transfer at the least nonzero weight]
\label{tc:minimum-transfer}
Suppose $E[2]$ is irreducible, and let $m$ be the least weight
of an index on the filter with nonzero symbol.  At weight $m$,
the dependence on each nonroot vertex is only through its
oriented rational Frobenius lift in the Galois group over $\Q$
of the maximal elementary extension of $F$ unramified outside $S$, together with the
requirement on the total root label.  In particular, replacing
the nonroot vertices by fresh vertices with those same
fixed-support labels preserves the symbol, irrespective of
their internal edge labels, provided the root total is matched.
The assertion allows simultaneous changes using disjoint
intermediate configurations.
\end{lemma}

\begin{proof}
Every symbol of weight less than $m$ is zero.  For a simple
vertex, Proposition~\ref{tc:irreducible-rules} therefore makes
its inertia derivatives at every other active nonroot vertex
zero: different transpositions give a lower-weight deletion,
equal transpositions give zero directly, and a split vertex
gives a lower-weight merged configuration.  Its labels against
root vertices are immaterial by the simple case of the root
derivative rule.  These statements hold for conjugates of
inertia too, because the Frobenius and inertia calculation is
made at a chosen place and may be conjugated.  Starting with
the elementary extension over $F$ supported on $S\cup V$,
killing the inertia at the active vertices leaves exactly
fixed-support elementary labels.  This proves the assertion
for a simple vertex.

For a split vertex $q$, let $V$ be the other nonroot vertices,
so $w(V)=m-1$ and $V$ still has matching parity.  We first claim
that $B_V$ is divisible by $2$ integrally.  In
Equation~\eqref{tc:coefficient-bound}, an entry of cost zero
has $r$ consisting only of roots and $J$ consisting only of
split primes.  Its underlying squarefree class has weight
\[
                         w(V)-\#J<m.
\]
Its symbol is zero.  The shell formula at every split prime,
valid at all exponents, therefore adds one more factor of $2$
to the normalized coefficient.  All other allowed entries
already have positive integral cost.  This proves the claim,
including at arbitrarily large coefficient indices.

Use now $H_V=B_V/2$.  Two split traces give zero on $H_V$:
their added weight $2$ in
Equation~\eqref{tc:coefficient-bound} still leaves one factor
of $2$ after this division.  The trace restricted to split
elements consequently factors through the elementary kernel
over $F$, by the same square-invariance argument as before.
Coefficient $1$ with the new split trace and the retained
root probes reads the symbol at weight $m$.

Apply Lemma~\ref{tc:inertia-identity} at $p\in V$ to this
divided vector.  The element $g\phi_p$ has the type of $p$,
because $g$ is split.  Its $p$-shell is $c_p$ times a symbol
of weight $m-1$, whether $p$ is split or simple; for a simple
$p$ only the first shell is used.  For a simple $p$ the normalization can be checked directly:
$J=\{p\}$ and the new simple probe contribute $1/2+1/2=1$ to
Equation~\eqref{tc:coefficient-bound}. Writing
$n_V=\delta+w(V)$, the divided shell is
\[
 \frac{2c_p(D')a(D')}{2^{n_V+1}}
 =c_p(D')\frac{a(D')}{2^{n_V}},\qquad w(D')=w(V)=m-1.
\]
The omitted set $V/p$ has opposite parity, but its ceiling
normalization is still $n_V$. Thus this is exactly its unit-symbol
test, with no fractional loss. Both compared values are zero.
Every such inertia derivative therefore vanishes, leaving only the
fixed-support elementary label.

No character $\alpha$ on this further divided vector is
needed.  Use actual matched root probes in both tests.  At
weight $m$, their changes against a split vertex are already
zero by the split root-deletion rule, because its deletion
has weight $m-1$.  They can therefore be matched while the
nonroot vertex is moved.  Moving vertices successively,
using fresh intermediate copies when necessary, proves the
joint assertion without imposing old internal edge labels.

For a moving coefficient system, minimality is used on a fixed bounded
weight range. A failure of the required lower-weight vanishing on an
ultrafilter-large set of stages would give a bounded-weight sequence
contradicting the stabilized minimum. Choose the actual stage as in
Lemma~\ref{tc:shell-bound}, with the additional division by $2$ included
among its fixed division bounds. Its coarse estimate controls all larger
weights, and its integral square-index recurrences control every shell.
Thus the preceding coefficientwise arguments apply to the actual forms
at that stage. The argument does not require $m$ to be integral.
\end{proof}

\subsection{Summary and finite realization}

We collect the rules in the form used to construct finite prime arrays.
The analytic meaning of a nonzero symbol is unchanged; this section has
established which substitutions preserve it or compute its change.

\begin{theorem}[Symbol rules]
\label{thm:symbol-rules}\label{co:symbol-interface}
For either coefficient source of
Propositions~\ref{prop:even-coefficients} and
\ref{prop:odd-coefficients}, assume its stated coefficient
lower bound and use its normalized symbol and filter.
Then the following rules hold for finite configurations of
primes outside the level support.
\begin{enumerate}[label=\textup{(\roman*)}]
\item With rational two-torsion, the elementary-label
factorizations, contraction except in the simple--simple
case, and isolation by a trivial split label are those of
Proposition~\ref{tc:rational-rules}.  A pair of identical
simple labels realized with $\#E(\mathbf F_q)/2$ odd can be
inserted without changing the symbol, as in
Proposition~\ref{tc:simple-twins}.
\item With irreducible $E[2]$, all root probes combine into
one total rational elementary label.  A root-bit derivative
at a split vertex deletes that vertex, while at a simple
vertex it is zero.  The derivative between two simple
vertices deletes both when their transposition images are
different, and is zero when they are equal.  The remaining
simple--split shell has weight one less.  These are the
rules of Proposition~\ref{tc:irreducible-rules}.
\item At the least weight supporting a nonzero symbol in
the irreducible case, each nonroot may be moved with only
its fixed-support elementary Frobenius label prescribed;
internal edges impose no further condition on the value,
subject to matching the total root label.  This is
Lemma~\ref{tc:minimum-transfer}.
\end{enumerate}
The analytic and finite-Selmer consequences of a nonzero
symbol remain those of
Proposition~\ref{prop:unit-symbol-detection}.  In the odd
source the rules mean finite tests at arbitrarily high stages
chosen in the ultrafilter-large sets of the prescribed stabilization.
\end{theorem}

\begin{proof}
The preceding propositions prove the assertions. For one application,
put the finitely many Galois modules and operator tests in a common
finite quotient at the required precision. Chebotarev then gives
simultaneously fresh Frobenius representatives. The omitted source is
unramified at the whole omitted prime, whereas its ramified rational
quadratic coordinate is exactly the Legendre flip in the shell formula.
These are therefore substitutions with all other labels held fixed.

The arithmetic compatibility of an entire array is addressed in the
following sections. The present assertions concern each finite array;
in the odd case choose a sufficiently high actual stage in the
required ultrafilter-large intersection. Arbitrarily many weight-zero roots cause no additional depth
loss, because the coefficient bound is uniform at a fixed weight and
their traces combine through one total elementary label.

\end{proof}

\section{Prime graphs and simultaneous unit coefficients}\label{gr:section}

The coefficient symbols will now be regarded as functions of finitely
many quadratic residue bits. We first identify the arithmetic data that
those bits record and explain how the trace calculation becomes a
finite difference. This gives polynomials whose derivatives contract
prime labels. We then adjoin auxiliary vertices so that one or two
nonzero coefficient polynomials take the value $1$ at every nonzero
address of a binary cube.

The two-polynomial case is needed for the odd converse: the odd symbol
tests $h$, while an even symbol tests a quadratic companion $kh$.
Separate nonzero polynomials need not have a common unit evaluation;
for example, $z$ and $1+z$ do not. The auxiliary vertices change the
polynomials by controlled contraction operators. The arithmetic input
is the rational contraction and isolation rules of
Theorem~\ref{thm:symbol-rules}, including omission of an entire rational
prime. The same argument applies to any coefficient test with these
rules, such as the restricted even test constructed in
Section~\ref{pw:split-section}.

Encoding quadratic residue symbols by graphs has a precedent in the
congruent-number problem: Feng--Xiong express descent conditions through
prime graphs and the ranks of their binary Laplacians
\citep[\S2, Definitions~2.1 and~2.3 and Lemma~2.2]{FengXiong2004}.
Here the graphs record coefficient-symbol rules, and their oriented
contractions produce simultaneous unit values.

\subsection{Prime labels and their realization}

Let $F$ be $\mathbf Q$ or a quadratic field.  The applications use
$F=\mathbf Q(E[2])$ when this field is quadratic, and also an auxiliary
quadratic field when $E[2]$ is rational.  Fix a finite set $S$ of rational
primes containing $2$ and the ramification support of the fixed data.
Enlarge $S$ so that prime ideals above it generate the class group of $F$,
and let $S_F$ include these places and the infinite places.  Set
\[
 V_S(F)=\{\alpha\in F^\times/F^{\times2}:
       v_{\mathfrak p}(\alpha)\equiv0\pmod2
       \text{ for }\mathfrak p\notin S_F\}.
\]
The valuation sequence expresses this as an extension of the squareclasses
of $S_F$-units by a subgroup of the $2$-torsion of the $S_F$-class group;
in particular, it is finite-dimensional over $\mathbf F_2$
\citep[Theorems~4.4 and~5.11]{MilneANT2020}.
Write $F_S^{(2)}=F(\sqrt{V_S(F)})$.

A \emph{split vertex} is a rational prime $p\notin S$ split in $F$,
with an \emph{orientation} $\mathfrak p\mid p$ when $F$ is quadratic.
Choose a generator $\pi_{\mathfrak p}$ whose divisor outside $S_F$ is
$\mathfrak p$; this is possible by the class-group condition.  Use its
conjugate for the other orientation.  Then
\begin{equation}\label{gr:norm}
 N_{F/\mathbf Q}(\pi_{\mathfrak p})=p u_{\mathfrak p},
 \qquad u_{\mathfrak p}\text{ a signed rational }S\text{-unit}.
\end{equation}
For $F=\mathbf Q$, use the generator $p$.  In the quadratic case an inert
rational prime is a \emph{simple vertex}, with rational generator $p$.
A split vertex has its fixed Frobenius label on $F_S^{(2)}$; a simple
vertex has rational fixed-support labels.  A \emph{type} may also specify
local squareclasses, support valuation parities, real signs, and finite ray
conditions on its generator.  Prescribed types must be compatible and
individually realizable.  A \emph{filter} fixes the total rational labels
and sign of the signed-prime product, where $p^*=(-1/p)p$.

Write quadratic residue symbols additively in $\mathbf F_2$.  Between two
split vertices in the quadratic case there are two independent forward
bits: the residues of $\pi_{\mathfrak p}$ and
$\overline{\pi_{\mathfrak p}}$ at the chosen place of the other vertex.
Hilbert reciprocity and the types determine the reverse bits.
Their sum, with the fixed correction in \eqref{gr:norm}, gives the rational
norm bit.  If an endpoint is simple, or if $F=\mathbf Q$, there is one
rational forward bit.  Vertices, types, and these bits form an
\emph{arithmetic graph}; there is no self-residue datum.

For the coefficient symbols, the vertices are the prime factors of the
Fourier index. Their types fix the support and elementary Frobenius
conditions, while the edge bits record the remaining quadratic residue
tests. Varying an edge leaves those types and the total filter unchanged.
The next calculation explains the resulting contraction rule.

For the standard rational-two-torsion symbol of
Proposition~\ref{tc:rational-rules}, take split primes $p,q$ and let $U$
contain $p$ but not $q$, with matching
parity. Write $V=U/p$ and $g=\Fr_q$. By
Proposition~\ref{tc:rational-rules},
\[
 A(Uq)=\bigl[t(g)H_U\bigr][1].
\]
Changing one residue bit from $p$ to $q$ replaces $g$ by $g\tau_p$,
where $\tau_p$ is the corresponding inertia element. The level identity
\eqref{tc:exact-level-identity} moves the difference to the coefficient
at $p$. Its two first-shell values are
\[
 c_p(D')A(D'),\qquad (c_p(D')+1)A(D'),
 \qquad D'=Vq',\quad \Fr_{q'}\sim g\phi_p.
\]
The omitted source is unramified at the whole prime $p$, so its value
$A(D')$ is the same in the two terms. Their sum is $A(D')$.
Thus, if $z$ is the changed edge bit,
\begin{equation}\label{gr:first-contraction}
 A\big|_{z=0}+A\big|_{z=1}=A(Vq').
\end{equation}
Here $q'$ is a fresh prime carrying the product labels of $p$ and $q$;
it replaces two vertices by one. Its integer value is not the product
$pq$.

In the quadratic field, the two choices of oriented bit correspond to
using $\pi_{\mathfrak p}\pi_{\mathfrak q}$ or
$\pi_{\mathfrak p}\overline{\pi_{\mathfrak q}}$ as a product template,
up to the choice of output orientation. The template adds the support
labels and the external residue bits in the indicated orientations.
There is no internal residue attached to the new vertex. Omitting $p$
removes both places above that rational prime from the source, so a
second difference in the other oriented bit gives zero. Split--simple
contraction has the same interpretation with rational product data;
a simple--simple difference is zero.

To use these formulas on arbitrary arrays of bits, we must know that
the arrays and the product templates have prime representatives, and
that different representatives give the same symbol. The following
lemma supplies both facts.

\begin{lemma}[Finite realization and product labels]\label{gr:realization}
Fix finitely many old vertices and individually realizable new types.
Every choice of independent forward bits is realized by distinct fresh
primes, avoiding any given finite set.  Further fixed local tests may be
included in $S$.

A split cluster may be represented by the product of its oriented
generators, with conjugations prescribed along a contraction tree; its
support types and external bits are the corresponding sums.  A cluster
containing one simple vertex has the corresponding rational product data.
These products may be replaced by fresh prime representatives.  For a
symbol satisfying the omitted-vertex dependence of
Theorem~\ref{thm:symbol-rules}, the graph data determine its value, and
these replacements preserve the contracted symbol.
\end{lemma}

\begin{proof}
The displayed generators account for every relevant radical.  Let $U$
be a finite set of new rational primes and $U_F$ the places above them.
From a squareclass with even valuations outside $S_F\cup U_F$ subtract
the generators at the primes of $U_F$ where its valuation is odd.  The
remainder lies in $V_S(F)$.  At an inert prime its rational generator has
valuation one at the unique place of $F$.  Thus the fixed labels and
residue bits determine every elementary Frobenius test with this
ramification support.

For a new split vertex, start with a template $\pi_0$ of its type, at a
prime ideal $\mathfrak p_0$ away from the old vertices.  Weak approximation
gives $a\in F^\times$ sufficiently close to $1$ at $S_F$, positive at
the specified real places, and with residues at the old prime ideals
chosen so that $a\pi_0$ has the prescribed quadratic residues.  Use a
ray modulus containing deep enough powers of $S_F$, all old prime ideals,
and the real places.  A prime ideal $\mathfrak q$ in the ray class of
$(a)\mathfrak p_0$ satisfies
\[
 \mathfrak q=(ab)\mathfrak p_0,\qquad
 b\equiv1\pmod{\text{ray modulus}}.
\]
Then $ab\pi_0$ has the desired divisor and type.  The prime ideals in
this ray class have positive Dirichlet density
\citep[Chapter~V, Theorem~2.5]{MilneCFT2020}.
For $s>1$, the sum of $(N\mathfrak q)^{-s}$ over prime ideals of
residue degree greater than one is at most $2\sum_p p^{-2s}$, which
stays bounded as $s\downarrow1$.  The corresponding sum over the ray
class diverges.  Hence that class contains infinitely many degree-one
primes.  Outside the fixed ramification support these give split rational
primes, and finite exclusions remove finitely many choices.

At an old inert prime, prescribe instead the norm of the residue in
$\mathbf F_{p^2}^\times$.  Its norm onto $\mathbf F_p^\times$ is
surjective.  For a new simple vertex all conditions are rational
congruences, including the inert class in $F$; the Chinese remainder
theorem and Dirichlet's theorem apply.  Successive selection realizes the
whole graph.  Reciprocity determines reverse bits, since the support
Hilbert symbols are fixed and only the two endpoint residues remain.
Taking norms gives \eqref{gr:norm}, with no extra independent condition.

The same construction applies to a product template.  It preserves its
support data and its residues at every remaining vertex; no internal
residue is imposed on the replacement.  For a simple cluster use rational
congruences for the norm labels.  To check the rational filter in a split
cluster, write $\pi=\prod_i\pi_i$ and $\pi'=ab\pi$ for its replacement.
The multiplier $ab$ is a square unit at $S_F$ and positive at real places.
Consequently
\[
 N\pi=(\prod_i p_i)u,\qquad N\pi'=q u
\]
have the same signed rational $S$-unit $u$: its valuations and sign agree.
Thus $q/\prod_i p_i=N(ab)$ is locally square at the rational support.
The condition at $2$ also identifies $q^*$ with $\prod_i p_i^*$ in the
fixed tests.  This verifies preservation of the norm correction.

Finally, compare two disjoint realizations of the same graph.  Insert
fresh bridge vertices with matching bits to both realizations and replace
vertices one at a time.  The omitted-vertex Frobenius tests agree by the
radical description above, so the symbol values agree.  A preliminary
bridge removes any overlap.  This proves both graph dependence and
invariance under the product replacements.
\end{proof}

Consequently, fixing the vertex types gives a well-defined function
$A_{\mathcal G}$ of the independent forward edge bits: choose any prime
realization of the graph $\mathcal G$ and evaluate its coefficient
symbol. Lemma~\ref{gr:realization} makes the result independent of that
choice. We will write its unique multilinear polynomial as $P$.
Equation~\eqref{gr:first-contraction} says that a derivative of $P$
evaluates the contracted graph.

A split vertex is \emph{primary} if its generator is a square unit at
every finite place of $S_F$ and positive at every real place.  The principal
ray class realizes such types.  Reciprocity makes their fixed radical
labels trivial, and their signed rational contributions are squares in the
fixed tests.  Including stationary vertex primes in those tests makes
primary incidences to them neutral.  Including the support of a quadratic
discriminant $k$ ensures that every primary prime splits in
$\mathbf Q(\sqrt{k})$.

Two symbols may use different fixed-support tests or read only rational
norm labels from this same graph. We call these \emph{compatible
projections} when they forget fixed-label coordinates or take norm data,
preserve primary vertices, and commute with the
omitted-vertex product substitution. In the norm projection either
oriented-bit toggle changes the rational bit once. Thus one physical
edge difference induces the specified difference for each symbol; it
need not give the same polynomial for both symbols.

\begin{definition}[Witnesses and permitted partitions]\label{gr:partitionable}
A \emph{witness} is a finite graph with symbol $1$.  A graph is
\emph{partitionable} if it admits a partition into nonempty clusters, each
containing at most one simple vertex, and oriented contraction trees whose
contracted graph is a witness.  The contraction trees are part of the data;
this definition makes no assertion about a Selmer group.
\end{definition}

For a two-isogenous curve with quadratic field $F'=\mathbf Q(E'[2])$,
a simple vertex is \emph{good} if it is inert in both $F$ and $F'$.
A \emph{donor} is split in $F$ and inert in $F'$.  Attaching a donor
turns a simple vertex split in $F'$ into a good one.  These are conditions
on finite labels; their supply will be established in the arithmetic
application.

\subsection{Forests as finite-difference operators}

For an edge bit $z$, put $\Delta_z f=f|_{z=0}+f|_{z=1}$.
Every function on a finite binary cube has a unique multilinear polynomial
over $\mathbf F_2$; $\Delta_z$ is its formal derivative.  Differences
therefore commute and satisfy $\Delta_z^2=0$.

\begin{lemma}[Oriented forests]\label{gr:forest}
Assume the contraction, isolation, and omitted-prime assertions of
Theorem~\ref{thm:symbol-rules}.  Differentiation along a tree contracts
it, adding product labels and external incidences; the result is zero if
the tree contains two simple vertices.  A collection of edge variables
containing a cycle has zero derivative, including when its orientations
are incompatible.  The two oriented variables on one pair form a cycle
of length two.
\end{lemma}

\begin{proof}
The one-edge rule is the symbol contraction identity, with product
replacement supplied by Lemma~\ref{gr:realization}.  Induction contracts
a tree: split--split mergers remain split, split--simple mergers remain
simple, and merging two simple components gives zero.  A root orientation
determines all relative orientations along the tree.  Its product template
uses support types and external bits, not the differentiated internal bits.

For the two variables on a split pair $p,q$, differentiate first in one
inertia bit over $p$.  The resulting symbol omits the \emph{entire rational
prime} $p$, so its elementary extension is unramified at both places over
$p$.  The derivative in the other oriented bit is zero.  Repeated use of
the same variable gives zero by $\Delta_z^2=0$.

For a longer cycle, contract a spanning path.  The closing edge becomes
internal to one cluster and changes neither its support label nor an
incidence to a remaining vertex.  The omitted-endpoint argument just given
therefore kills its derivative.  Conjugating the product only changes its
orientation; it does not restore the omitted endpoint or create a
self-residue.

For example, after contracting aligned edge $12$, the two bits from its
product to vertex $3$ are
$t_0=z_{13}^0+z_{23}^0$ and $t_1=z_{13}^1+z_{23}^1$, up to fixed type
corrections.  An aligned edge $23$ differentiates $t_0$.  The closing
edge $31$ differentiates either $t_0$ again, giving zero, or $t_1$, giving
the omitted-endpoint zero.  Every longer cycle reduces to these cases.
Further differences commute with this zero derivative.
\end{proof}

Call the original vertices \emph{terminals}; any stationary external
vertices retain their types and incidences. Fix their vertex types and a
realizable baseline graph. Write $z=(z_e)$ for the independent toggles of
the terminal--terminal bits, and let $P(z)$ be the terminal symbol.
A \emph{network} adjoins primary vertices and specified incidences
involving them. It adds no terminal--terminal edge: the original bits
$z$ remain parameters. All unspecified primary incidences are neutral.
Within each network fix a consistent orientation at each terminal.
Different networks may use different alignments.

Write $y=(y_e)$ for the bits on the added incidences, and let $f(z,y)$
be the symbol of the enlarged graph. We shall express $f(z,\mathbf1)$
as a finite-difference operator applied to $P(z)$. The differences in
$y$ first contract the added graph; the forest rule then expresses each
surviving term as differences in the original variables $z$.

For example, join one primary vertex $r$ to two permitted terminals
$a,b$, in a fixed alignment. Expanding in its two added incidences gives
four terms. The term with no difference is zero because $r$ is isolated.
Each one-edge difference contracts $r$ onto one terminal and leaves the
original symbol $P$, so these two terms cancel in $\F_2$. The two-edge
tree contracts the terminal pair and gives $\Delta_{z_{ab}}P$.
Thus this path acts exactly as the terminal difference
$\Delta_{z_{ab}}$. The path has two added edges, but its operator has
contraction degree one. In general, a component meeting $t$ terminals
has terminal contraction degree $t-1$.

The lowest two operator degrees can be read from a matrix. Index the
internal primary vertices by $I$, let $e_{ij}$ be their internal edge
bits, and let $c_a$ be the incidence column from them to terminal $a$.
The grounded Laplacian is the symmetric matrix over $\F_2$ with
\[
 Q_{ij}=e_{ij}\quad(i\ne j),\qquad
 Q_{ii}=\sum_{j\ne i}e_{ij}+\sum_a(c_a)_i.
\]
Equivalently,
\begin{equation}\label{gr:grounding-dictionary}
                         Q\mathbf1=\sum_a c_a.
\end{equation}
Thus prescribed off-diagonal entries and terminal columns realize a
given symmetric $Q$ precisely when this row-sum identity holds. For the
one-primary-vertex path above, $Q=(0)$ and $c_a=c_b=(1)$.

The coefficients below are instances over $\F_2$ of the all-minors
matrix-tree theorem \citep[\S1, Equation~(2)]{Chaiken1982}; we include
the incidence-matrix proof needed here.

\begin{lemma}[Network expansion]\label{gr:network}
A network satisfies $f(z,\mathbf1)=\mathcal N P(z)$, where
$\mathcal N$ is a sum of terminal finite-difference operators indexed by
forests. In each surviving forest term, every component containing an
internal primary vertex meets a terminal. A component meeting $t$
terminals has contraction degree $t-1$;
a component containing two simple terminals acts as zero.  Networks with
disjoint internal vertices and no edges between them multiply as operators
on the same $P$.  In particular, degree greater than $\deg P$ annihilates it.

For the grounded Laplacian $Q$ and terminal columns $c_a$ defined above,
the constant
and the pair-contraction coefficients are
\begin{equation}\label{gr:forest-coefficients}
 \det Q,\qquad c_a^{\mathsf t}\operatorname{adj}(Q)c_b\quad(a\ne b).
\end{equation}
The latter multiplies the terminal difference in the selected alignment;
a prohibited pair multiplies the zero operator.
\end{lemma}

\begin{proof}
For the set $E_0$ of added incidences, keep the terminal variables $z$
fixed and expand in the variables $y$:
\[
 f(z,\mathbf1)=\sum_{T\subseteq E_0}(\Delta_T^y f)(z,\mathbf0).
\]
Lemma~\ref{gr:forest} removes cycles.  A component meeting no terminal
contracts to an isolated primary vertex, so isolation removes it, including
an untouched primary vertex.  In a surviving component the primary labels
and external primary bits are trivial.  Its contracted label and external
incidences are thus exactly those obtained by contracting its terminals
in the original graph.  A terminal tree with the same relative orientations
gives that operator, of degree $t-1$.

Expanding two disjoint networks gives composition: product labels agree
in either order, and any resulting cycle vanishes.  Polynomial
differentiation then proves the degree bound.  This reasoning also verifies
compatibility with projections at every degree.  Norms commute with
oriented products, primary norm corrections are neutral, and either
oriented toggle induces its rational toggle on the original terminal
polynomial.  Repeated projected differences or prohibited simple mergers
give zero.  Forgetting fixed-label coordinates commutes with the same
products.

For the coefficients, write the Laplacian of the added network as
$L=DD^{\mathsf t}$,
where each incidence column has ones at its edge endpoints.  Let $I$ and
$A$ be the internal and terminal sets.  Cauchy--Binet for $Q=L_{I,I}$
counts edge sets whose incidence matrix becomes invertible after deleting
the terminal rows.  They are exactly forests with one terminal in each
component: cycles give dependent columns, a component with no deleted row
gives dependent rows, and a tree with one row deleted is invertible by
leaf removal.  These are the degree-zero terms.

Apply Cauchy--Binet instead to row set $I\cup\{a\}$ and column set
$I\cup\{b\}$.  Each component must contain one root from each of
$A\setminus\{a\}$ and $A\setminus\{b\}$.  Thus $a,b$ lie together
and every other terminal lies in its own component.  The corresponding
minor is
\[
 \det\begin{pmatrix}Q&c_b\\c_a^{\mathsf t}&0\end{pmatrix}
   =c_a^{\mathsf t}\operatorname{adj}(Q)c_b.
\]
The corner is zero because there is no added terminal--terminal edge.  Signs
disappear in characteristic two, and the adjugate identity remains valid
when $Q$ is singular.
\end{proof}

\subsection{Selecting one address}

An \emph{address} is $x\in\mathbf F_2^b$, where $b\ge1$.
A vertex with activation $\lambda\in(\mathbf F_2^b)^*$ is present
when $\lambda(x)=1$.  Fix a nonzero target $x_*$ and an edge bit $z$
joining terminals $a_0,a_1$ active there.  Either all terminals are always
present, or all have nonzero linear activations with at least one terminal
for each nonzero linear form.  For one symbol we need a derivative as the leading term at $x_*$ and
a constant term $1$ elsewhere. For two symbols we must also account for
a derivative allocated to just one factor; the second network includes
a binary parameter for this purpose. We construct these two networks by
prescribing their constant and degree-one terms; the final degree count
will remove every higher term.

Choose a dual basis $\lambda_1,\ldots,\lambda_b$ with
$\lambda_i(x_*)=1$ for every $i$.  Indeed, start with
$\lambda_1(x_*)=1$ and add it to a basis of the annihilator of $x_*$.
Make $b$ groups of three primary vertices, with all three vertices in
group $i$ activated by $\lambda_i$. The physical coordinate vectors index
these actual vertices. In each three-dimensional coordinate space choose
a computational basis $(e_i,u_i,v_i)$ with $e_i=(1,1,1)$. Prescribe a
symmetric form $Q$ by making the $e_i$ orthonormal and orthogonal to all
$u_j,v_j$, and using on the remaining coordinates
\begin{equation}\label{gr:H}
 H(T,B)=\begin{pmatrix}0&T\\T^{\mathsf t}&B\end{pmatrix}.
\end{equation}
Convert this form to its matrix in the physical vertex coordinates,
which we also denote by $Q$. Write $Q_J$ for the restriction to the
physical coordinates in groups $J$.

The column $h_i=Qe_i$ represents the covector dual to $e_i$. Since the
coordinate changes do not mix groups, $h_i$ is supported on group $i$;
it is independent of $T,B$. Attach this incidence column to a terminal
of activation $\lambda_i$, or to any fixed terminal in the stationary
case. When the active groups are $J$, the sum of these anchor columns is
\[
             \sum_{i\in J}h_i|_J=Q_J\mathbf1_J,
\]
since $\mathbf1_J=\sum_{i\in J}e_i$. Use the off-diagonal physical
entries of $Q$ as the internal edge bits. Equation~\eqref{gr:grounding-dictionary}
then gives the prescribed diagonal as well: the anchors make the actual
grounded Laplacian equal to $Q_J$ at every address, even after groups
have disappeared.

Also attach the same column $c$ to $a_0,a_1$.  If their moving activations
$\alpha,\beta$ differ, attach it to a terminal of activation
$\alpha+\beta$ as well.  The additional column sum is always zero:
\[
 \bigl(\alpha(x)+\beta(x)+(\alpha+\beta)(x)\bigr)c=0.
\]
The third terminal is absent at $x_*$.  Equal activations, or two stationary
terminals, require only the first two columns.  Coincident terminal roles
mean addition of columns.  Hence these attachments preserve the grounding
at every restriction.

Write $O_{\ge d}$ for a sum of contraction operators of degrees at least
$d$.  This denotes a filtration, not a numerical estimate.
Figure~\ref{gr:selector-figure} depicts the group-level cycle and path
used in the two programs when $b=4$.

\begin{figure}[tbp]
\centering
\setlength{\unitlength}{1mm}
\begin{picture}(115,36)
\put(15,28){\circle{6}}\put(15,28){\makebox(0,0){$1$}}
\put(40,28){\circle{6}}\put(40,28){\makebox(0,0){$2$}}
\put(40,8){\circle{6}}\put(40,8){\makebox(0,0){$3$}}
\put(15,8){\circle{6}}\put(15,8){\makebox(0,0){$4$}}
\put(18,28){\vector(1,0){19}}\put(40,25){\vector(0,-1){14}}
\put(37,8){\vector(-1,0){19}}\put(15,11){\vector(0,1){14}}
\put(27.5,0){\makebox(0,0){Cycle}}
\put(75,28){\circle{6}}\put(75,28){\makebox(0,0){$1$}}
\put(100,28){\circle{6}}\put(100,28){\makebox(0,0){$2$}}
\put(100,8){\circle{6}}\put(100,8){\makebox(0,0){$3$}}
\put(75,8){\circle{6}}\put(75,8){\makebox(0,0){$4$}}
\put(78,28){\vector(1,0){19}}\put(100,25){\vector(0,-1){14}}
\put(97,8){\vector(-1,0){19}}
\put(87.5,0){\makebox(0,0){Path}}
\end{picture}
\caption{The off-diagonal patterns of $T-I$ for the two programs when
$b=4$.  The vertices represent groups, not the individual primary
vertices of the network.  A proper group restriction breaks the cycle;
the path from $1$ to $4$ survives only when every group is present.}
\label{gr:selector-figure}
\end{figure}

\begin{lemma}[Two address programs]\label{gr:programs}
With these activations and grounding, there are the following networks:
\begin{enumerate}
\item A singular program $\mathcal S$ such that
\[
 \mathcal S_{x_*}=\Delta_z+O_{\ge2},\qquad
 \mathcal S_x=1+O_{\ge1}\quad(x\ne x_*).
\]
\item Two nonsingular programs on the same sites, indexed by
$\varepsilon\in\mathbf F_2$, such that
\[
 \mathcal N_x^{(\varepsilon)}=1+O_{\ge1},\qquad
 \mathcal N_x^{(1)}+\mathcal N_x^{(0)}
 =\begin{cases}\Delta_z+O_{\ge2},&x=x_*,\\
 O_{\ge2},&x\ne x_*.
 \end{cases}
\]
\end{enumerate}
All new sites have nonzero activation and are absent at $0$.
Unspecified incidences, including those to stationary external vertices,
may remain neutral.
\end{lemma}

\begin{proof}
For the first program, take $T=I+P$, with $P$ the cyclic permutation
matrix, and let $B$ be diagonal with exactly one nonzero entry.  Every
proper principal restriction of $P$ is a union of directed paths, hence
nilpotent.  Thus $T_J$ and $H(T_J,B_J)$ are invertible for proper $J$,
including the empty restriction, whose determinant is $1$.

On all groups, $\ker T$ and $\ker T^{\mathsf t}$ are spanned by
$\mathbf1$.  A radical vector of $H$ has $v=t\mathbf1$; summing
$T^{\mathsf t}u+Bv=0$ gives $t=0$, since
$\mathbf1^{\mathsf t}B\mathbf1=1$.  Hence $Q$ has radical generated
by $r=\sum_i u_i$.  Its adjugate is $rr^{\mathsf t}$: it is nonzero
of rank one, symmetry identifies its factors with the kernel, and the
only nonzero scalar in $\mathbf F_2$ is $1$.

Take $c$ dual to $u_1$, so $c^{\mathsf t}r=1$.  All anchors annihilate
$r$.  At $x_*$ the only terminal columns pairing nontrivially with $r$
are therefore those of $a_0,a_1$.  Formula~\eqref{gr:forest-coefficients}
gives constant zero and precisely $\Delta_z$ in degree one.  Every other
address has a proper group subset and constant $1$, since the activation
forms form a basis.  For $b=1$, $T=0$ and $B=(1)$ give the same radical.

For the second program take $T=I+L$, where $L$ is the consecutive lower
shift, and $B_\varepsilon=\varepsilon E_{bb}$.  Every group restriction
of $T$ is invertible, and
\begin{equation}\label{gr:inverse}
 H(T,B)^{-1}=
 \begin{pmatrix}T^{-\mathsf t}BT^{-1}&T^{-\mathsf t}\\
 T^{-1}&0\end{pmatrix}.
\end{equation}
Again take $c$ dual to $u_1$.  On groups $J$, the change in
$c^{\mathsf t}Q_J^{-1}c$ is $((T_J^{-1})_{b1})^2$, interpreted as
zero if an endpoint is absent.  The expansion
$(I+L_J)^{-1}=I+L_J+L_J^2+\cdots$ shows that this is $1$ exactly
when the entire path $1,\ldots,b$ survives, namely at $x_*$.
Pairings with anchors do not change, since $Q_J^{-1}h_i=e_i$ and
$e_i,h_i,c$ are independent of $\varepsilon$.  At the target the
balancing terminal is absent, leaving only the changed pair $a_0,a_1$.
All determinants are $1$.  Lemma~\ref{gr:network} proves the assertion,
also for the single-vertex path $b=1$.
\end{proof}

\subsection{Simultaneous units on a punctured cube}

The preceding programs control only two degrees of their operators.
To make this sufficient, at each address we spend exactly the degree
of a chosen top monomial.  Any additional contraction then vanishes.
For two symbols, the same accounting uses the ordinary product of their
top homogeneous parts. The multiplication is performed before imposing
the Boolean relation $z^2=z$, so one occurrence of $z$ in each factor
remains two contraction requirements.

\begin{theorem}[Address lemma]\label{gr:address}
Assume the contraction, isolation, and whole-prime omission rules of
Proposition~\ref{tc:rational-rules}, or the same rules for another
symbol, together with the realization freedoms of
Lemma~\ref{gr:realization}.  Let $b\ge1$ and
consider either of the following terminal arrangements:
\begin{enumerate}
\item One symbol on a nonempty set of always active terminals, whose
polynomial in the independent pair bits is nonzero.
\item One or two symbols on terminals with nonzero linear activations,
including a terminal for every nonzero linear form.  At each $x\ne0$,
each symbol separately is a nonzero polynomial in the active bits.
The symbols may have stationary external data and may read compatible
projections of a common graph.
\end{enumerate}
At every $x\ne0$, assume that the active terminal data and the stationary
external data obey each required filter.  Then finitely many primary
auxiliaries with nonzero linear activations, and one realizable assignment
of graph bits, make all the symbols $1$ at every $x\ne0$.  No auxiliary
is present at $0$.  All may satisfy fixed splitting and neutral-incidence
conditions included in their primary type.

A nonzero evaluation or a nonzero iterated edge difference suffices for
the polynomial hypothesis.  In particular, a partition as in
Definition~\ref{gr:partitionable} suffices in the stationary case.
\end{theorem}

\begin{proof}
Fix a common realizable baseline for the terminal bits and regard their
toggles as variables; omit variables invisible to a symbol.  A nonzero
iterated difference implies a nonzero polynomial, while every nonzero
polynomial has a top monomial with coefficient $1$.  By
Lemma~\ref{gr:forest}, this also proves the assertion about partitions.

For one symbol, let $P_x$ have degree $d_x$ at $x\ne0$ and choose a
monomial of degree $d_x$ with coefficient $1$.  For each of its variables
adjoin a separate singular program targeted at $x$, all with disjoint
internal vertices.  At address $x$, its own $d_x$ programs have zero
constant, so they consume at least $d_x$ contraction degrees.  Any extra
degree from any program kills $P_x$.  Only the selected degree-one terms
of these programs and the constants of the others survive, giving the
chosen coefficient $1$.  A nonzero constant polynomial needs no program.
Stationary anchors give the same argument for always active terminals.

For two symbols, write their polynomials at $x$ as $P_x,R_x$, of degrees
$d_x,e_x$.  Multiply their top homogeneous parts in the ordinary ring
$\mathbf F_2[z_1,\ldots,z_m]$, before imposing Boolean relations.
This ring is a domain, so choose a coefficient $1$ at a monomial
\begin{equation}\label{gr:chosen-monomial}
 M_x=\prod_z z^{\nu_{xz}},\qquad
 \nu_{xz}\in\{0,1,2\},\qquad
 \sum_z\nu_{xz}=d_x+e_x.
\end{equation}
For each exponent two add one singular program targeted at $x$; for
each exponent one add one nonsingular program, with its own parameter
$\varepsilon_{xz}\in\mathbf F_2$.  Keep all internal sets disjoint.
Let
\[
 \Phi(\varepsilon)=\prod_{x\ne0}
    (\text{final value of }P_x)(\text{final value of }R_x).
\]
We show that $\sum_\varepsilon\Phi(\varepsilon)=1$, which implies
that one choice makes every factor $1$.

Expand the operators by contraction degree.  At its target a singular
program has zero constant in both factors, so it consumes at least one
degree in each.  A nonsingular program has constant $1$ independent of
its parameter.  Summing that parameter cancels a term using only its
constants everywhere; a surviving term must consume at least one degree
somewhere for that program.  The total of these minimum costs is
\begin{equation}\label{gr:budget}
 D=\sum_{x\ne0}(d_x+e_x).
\end{equation}
A nonzero term cannot consume more, since some factor would then be
differentiated beyond its degree.  Every surviving term therefore uses
each minimum cost exactly.

Thus a singular program contributes precisely $\Delta_z$ in each factor
at its target and constants elsewhere.  A nonsingular program contributes
degree one in exactly one factor.  Summing its parameter kills that
contribution away from its target and gives $\Delta_z$ at the target;
every other occurrence contributes the common constant $1$.  In
particular, no higher forest coefficient remains.

At one address put $Z_i=\{z:\nu_{xz}=i\}$ for $i=1,2$, and write
$z_J=\prod_{z\in J}z$.  The surviving allocation sum is
\[
 \sum_{J\subseteq Z_1}
 [z_{Z_2}z_J]P_x^{\rm top}\,
 [z_{Z_2}z_{Z_1\setminus J}]R_x^{\rm top}
   =[M_x](P_x^{\rm top}R_x^{\rm top})=1.
\]
Wrong-degree coefficients are zero.  Since the total degree is $D$
and no factor exceeds its degree, each factor receives its full degree;
its derivative is exactly the displayed coefficient, independent of the
baseline.  Allocations at distinct addresses are independent, so their
product is $1$.  This covers constant factors too.  For instance, if
both top forms are $z$, their ordinary product is $z^2$: a singular
program differentiates each once.  Boolean reduction would lose that
allocation.

Choose parameters with $\Phi=1$ and realize the resulting finite graph
by Lemma~\ref{gr:realization}.  Primary types preserve the filters and
the stipulated neutral incidences and splitting conditions.  Nonzero
linear activations ensure that no new vertex appears at $0$.
\end{proof}

For a stabilized symbol this is a finite construction at each dimension.
First choose $b$, the terminal data, and all networks; their weights and
coefficient tests then have a finite bound.  Select an actual stage where the numerical bounds and the finitely many stabilized symbol
evaluations hold for that range.  The required precision may
grow with $b$.  The unit-valuation constant, however, is the constant of
the coefficient construction and must be independent of $b$; the forest
argument supplies no additional uniformity assertion.

\section{The converse with rational two-torsion}\label{pw:section}

We prove the pointwise converse for curves with $E(\Q)[2]\ne0$.
The finite graph construction supplies binary families with a unit coefficient
at every nonzero address.  Cyclotomic interpolation settles the even case.
For the odd case we first establish the missing lower bound for the moving
coefficient system, then keep one nonvanishing quadratic companion fixed
throughout the interpolation.

Fix such a curve $E/\Q$.  We use the prime weights $w(h)$, normalized values
$\mathcal L(h)$, and finite $2$-primary Shafarevich--Tate lengths $s(h)$
from the preceding sections.  The corank $c_2(E^{(h)})$ is a different
quantity: a corank-zero base has finite $s(h)$ before analytic
nonvanishing is known.  All new primes are odd primes of good reduction.
A filter fixes the real sign and the local unit squareclasses at a finite
support containing $2N$, and may exclude any further fixed finite set.

For a signed parameter $h$, write $A(h)$ for the symbol at the positive
index $|h|$ on its chosen filter.  We will use two bounds attached to a
unit coefficient:
\begin{align}
 A(h)=1\text{ in the even system}\quad&\Longrightarrow\quad
 L(E^{(h)},1)\ne0,\qquad v(\mathcal L(h))\le 2w(h)+C,
 \label{pw:even-unit-interface}\\
 A(h)=1\text{ in the odd system}\quad&\Longrightarrow\quad
 a(E^{(h)})=1,\qquad j(P_h)\le w(h)+C.
 \label{pw:odd-unit-interface}
\end{align}
Here $P_h$ is the genus Heegner sum with a fixed analytic-rank-zero
partner $h_*$, and $j(P_h)$ is its free-line index valuation.
The first bound is available from the even coefficient construction and
Theorem~\ref{thm:cyclotomic-bound}.  We use the second only after
Proposition~\ref{prop:rational-odd-bound} has supplied the arithmetic lower
bound required to normalize the odd system.
The contraction and isolation identities are those of
Theorem~\ref{thm:symbol-rules}; their finite graph consequences were proved
in Section~\ref{gr:section}.

Write $j_h(k)=j(P_{|h|}(h))$ for the ring-class index over $K=\Q(\sqrt{k})$.
Theorem~\ref{thm:heegner-ring-bound} states that, if
$k<0$ is an odd fundamental discriminant with $(k,2Nh)=1$, has at most
$A$ prime factors, is square at $2Nh$, and defines a nonexceptional
imaginary field, and the twists $h,hk$ have analytic
orders one and zero, then
\begin{equation}\label{pw:ring-lower-interface}
 2j_h(k)\ge4w(h)+s(h)+s(hk)-C_{E,A}.
\end{equation}
All coefficient and interpolation constants depend on the fixed curve,
its parametrization, and fixed local data.  A bound on $\omega(k)$ may
enter $C_{E,A}$; the sizes of its primes do not.  The finite networks,
required coefficient precision, and prime sizes may depend on the binary
dimension.  The unit bounds in Equations~\eqref{pw:even-unit-interface}
and~\eqref{pw:odd-unit-interface} may not.

By Lemma~\ref{pr:invariance}, an isogeny preserves analytic rank and
Selmer corank.  If the isogeny class contains a curve with full rational
two-torsion, we use that curve.  Otherwise choose a rational two-isogeny $E\to E'$ and write
$F=\mathbb Q(E[2])$, $F'=\mathbb Q(E'[2])$; both fields are quadratic.
A prime split in $F$ is called split, and one inert in $F$ is called simple.
Their weights are respectively $1$ and $1/2$.  A simple prime is good if it
is also inert in $F'$.  At such a prime
\begin{equation}\label{pw:good-simple}
 \#E(\mathbb F_p)\equiv2\pmod4.
\end{equation}
Indeed there is exactly one nonzero rational two-torsion point in the
reduction.  If the group order were divisible by four, its two-primary
subgroup would contain a rational point of order four.  Its image under the
two-isogeny would supply on $E'$ a rational two-torsion point different from
the dual kernel, contradicting inertness in $F'$.  Thus the good-twin rule
applies.  When $F=F'$, every simple prime is good.

\subsection{Seeds and finite supplies}\label{pw:seeds-section}

The word ``label'' below means an individually realizable fixed-support
elementary Frobenius label, with orientations when required.  It does not
include a self-residue of a prime generator.  All contractions and
comparisons are those of Theorem~\ref{thm:symbol-rules} and
Lemma~\ref{gr:forest}.

\begin{lemma}[Minimum seeds]\label{pw:minimum-seeds}
\begin{enumerate}
\item In an all-split alphabet, suppose a unit symbol has a nonempty
witness.  At the least positive vertex count $r$ of such a witness, all
mutual edge bits are immaterial.  The symbol is determined by the individual
fixed-support labels.  This remains true if the empty configuration is a
unit, provided no contraction to the empty configuration is used.
\item In the quadratic case, suppose the symbol has a nonzero witness,
the alphabet is closed under the allowed contractions, and it contains
the good twin operations.  There are
$r\ge0$ and $s_1\ge1$ such that for every $s\ge s_1$ of the permitted simple
parity, the minimum split count of a unit with $s$ simple vertices is $r$.
At these counts all edge bits are immaterial.  A chosen minimum seed may be
padded by arbitrarily many good twin pairs.
\end{enumerate}
The assertions also hold for a stabilized moving odd symbol, with every
particular finite list of uses made at a sufficiently advanced stage.
\end{lemma}

\begin{proof}
For the first assertion, a derivative in an edge of an $r$-vertex graph is a
symbol on $r-1$ vertices.  If $r\ge2$, this graph remains nonempty and has
value zero by minimality.  If $r=1$, there is no edge.  The omitted-vertex
description then leaves exactly the fixed-support labels.

For the second assertion let $r(s)$ be the minimum split count at simple
count $s$, with value $+\infty$ when no witness exists.  Once a finite value
occurs, good-twin insertion gives $r(s+2)\le r(s)$.  A nonincreasing sequence
of nonnegative integers stabilizes.  Choose $s_1\ge1$ in that range.  A
derivative of an edge with a split endpoint reduces the split count and
preserves the simple count; it is therefore zero at the minimum.  A
simple--simple derivative is zero by the symbol rule.  Twin insertion gives
the last assertion.

For the moving system, take the minimum and the nonincreasing sequence in
the stabilized profile.  Each comparison just used involves bounded vertex
counts.  A failure on an ultrafilter-large set of stages at any one such bound
would produce a bounded-weight sequence contradicting that minimum.  Thus
any fixed finite collection holds on an ultrafilter-large intersection of
stages, from which an arbitrarily high stage may be chosen.  No assertion
about every sufficiently high stage, or about all unbounded counts at one
stage, is needed.
\end{proof}

Recall that a list is partitionable when a permitted contraction forest
has a unit symbol as its output.  Its terminal polynomial is then nonzero,
since the indicated iterated finite difference is nonzero.

\begin{lemma}[Finite label reserves]\label{pw:reserves}
Fix a minimum seed and a finite label alphabet.
\begin{enumerate}
\item In the all-split rational alphabet, a list with the same total label
as the seed and sufficiently many copies of its labels is partitionable.
\item In the quadratic alphabet, retain $r$ split anchors with the seed
labels.  Suppose there are sufficiently many simple vertices of the required
parity and sufficiently many split reserves of each realizable rational
shift.  If $F\ne F'$, also suppose
\begin{equation}\label{pw:donor-balance}
 \begin{aligned}
 \Delta={}&\#\{p:\ p\text{ split in }F\text{ and inert in }F'\}\\
       &-\#\{p:\ p\text{ inert in }F\text{ and split in }F'\}>C_0.
 \end{aligned}
\end{equation}
Here $C_0$ depends only on the anchors and the reserves.  Then all other
split vertices can be attached to simple vertices so that the result is
the minimum seed and good twin pairs.  The twins may match a prescribed
finite enrichment of the rational labels.  The base simple labels may be
specified on any projection on which the required total agrees.
\end{enumerate}
All reserve requirements depend only on the finite alphabets, seed, and
projection.  They do not depend on the original list length.
\end{lemma}

\begin{proof}
In the first case retain the first $r-1$ seed labels and merge all other
vertices to the last cluster.  Sufficient copies ensure that this cluster
is nonempty.  Its label is forced by the total, hence is the last seed
label.  Lemma~\ref{pw:minimum-seeds} makes its unknown edges immaterial.

For the second case, the rational labels of split attachments form the
subgroup with trivial $F$-coordinate.  They act transitively on the simple
coset, so every prescribed change of a simple label has a realizable split
shift.  Reserve sufficiently many representatives of every such shift.
Remove these reserves and the anchors from the donor pool.  If $F\ne F'$,
\eqref{pw:donor-balance}, with $C_0$ exceeding the number removed, leaves a
distinct donor for each bad simple vertex.  Attach these donors; every
simple vertex is now good.  If $F=F'$, this operation was unnecessary.

Pair equal enriched good labels.  By withholding a bounded number of pairs,
leave between $s_1$ and $s_1+C_1$ unpaired vertices, with the correct parity;
$C_1$ depends only on the number of enriched labels.  Use reserved shifts
to turn the excess unpaired vertices into good pairs and to set the first
$s_1-1$ base positions to the desired seed labels.  The seed labels
themselves need not all be good.  Attach every remaining nonanchor split
vertex to the last base position.  Its projected label is forced by the
total, the anchors, the other base positions, and the twin pairs.  It is
the required last seed label.  Each cluster contains one simple vertex,
so every contraction is permitted.  Only a bounded number of reserve
shifts was used.

When incidences to a retained split anchor are needed only to remove twins,
they can be prescribed after the contracted labels are fixed, by the
independent terminal-edge realization.  Fixed stationary vertices instead
enter the finite label enrichment from the outset.  These two uses do not
increase the number of label types with the original list length.
\end{proof}

\begin{lemma}[Neutral packets]\label{pw:packets}
For any of the finite reserve problems above with prescribed rational total
zero, one can choose a finite packet of terminal labels with rational total
zero such that every nonempty union of copies of the packet has the required
supplies. In the quadratic case with $F\ne F'$, its donor surplus may
be made arbitrarily large. Compatible fixed splitting conditions may be
included whenever
the needed individual labels and shifts remain realizable.
\end{lemma}

\begin{proof}
Choose ample copies of every required label, including anchors and shifts.
When the quadratic fields $F,F'$ are distinct, add enough donors to
exceed all the losses from the other labels and the prescribed reserve
margin. When $F=F'$, every simple vertex is already good. Duplicate every
full label in the resulting list.  Duplication gives total zero in every
rational elementary coordinate and preserves any strict donor surplus.
Choose the margins so that after reserving the finitely many anchors in
any one packet, the union of any positive number of packets still meets
the reserve requirements.  All labels are realized by distinct fresh
primes; duplication never means repeating an actual prime factor.
\end{proof}

\subsection{The even converse and its uniform estimate}
\label{pw:lower-section}

\begin{proposition}[Even original base]\label{pw:even-base}
If $c_2(E)=0$, then $a(E)=0$.
\end{proposition}

\begin{proof}
Parity gives even functional sign.  Fix an even symbol on the filter of $1$,
including all auxiliary class-group tests.  If its value at $1$ is a
unit, analytic nonvanishing is already detected.  Otherwise use a
nonempty minimum seed in the full-torsion case, or a stabilized quadratic
seed in the other case.

For each $b$, assign a neutral packet from Lemma~\ref{pw:packets} to
every nonzero linear form on $\mathbb F_2^b$, and activate its primes
by that form.  At every nonzero address at least one packet is active;
their union satisfies the reserve problem and has a nonzero terminal
polynomial.  At zero no terminal is active.  The single-symbol case of
Theorem~\ref{gr:address} adds primary auxiliaries and makes the symbol a
unit at every nonzero address.  All local classes remain those of $1$.
Equation~\eqref{pw:even-unit-interface} gives the uniform bound at all
nonzero vertices.  Theorem~\ref{thm:cyclotomic-interpolation}, applied
at the original corank-zero base, now gives $a(E)=0$.
\end{proof}

\begin{proposition}[Partitionable even estimate]\label{pw:partitionable-even}
Fix an even unit symbol on a prescribed filter.  For every
partitionable parameter $h$ on its filter with $c_2(E^{(h)})=0$,
\begin{equation}\label{pw:partitionable-bound}
 \mathcal L(h)\ne0,\qquad v(\mathcal L(h))-2w(h)-s(h)\le C.
\end{equation}
The constant depends only on the symbol and fixed local data.
\end{proposition}

\begin{proof}
If the parameter has no prime factors, the only contraction is the identity,
so partitionability means that the empty symbol itself is a unit.
Equation~\eqref{pw:even-unit-interface}
then proves the assertion directly.  Otherwise the stationary terminal
set is nonempty.
Keep the prime factors of $h$ stationary.  The iterated edge difference
along the defining contraction forest is $1$, so their terminal polynomial
is nonzero.  The stationary one-symbol case
of Theorem~\ref{gr:address} gives arbitrarily large binary cubes based at
$h$, with only primary new primes, whose other addresses have unit symbol.
Equation~\eqref{pw:even-unit-interface} and $s(h_x)\ge0$ give the uniform
upper bound required by the cyclotomic transfer.  The base corank is zero,
so Theorem~\ref{thm:cyclotomic-interpolation} gives \eqref{pw:partitionable-bound}.  A base-dependent
dimension threshold does not affect the asserted constant.
\end{proof}

\subsection{The lower bound needed for odd coefficients}

The even estimate also supplies the arithmetic hypothesis of
Proposition~\ref{prop:odd-coefficients}. For an already analytic-rank-one
parameter $h$, we choose a companion $k'=k'(h)$ with boundedly many prime
factors and estimate the central value at $hk'$. Comparing the ring and
genus Heegner sums then bounds $j(P_h)$ from below. The companion in this
argument may vary with $h$; the odd coefficient system has not yet been
normalized.

To compare the two fixed isogenous curves $E,E'$, add subscripts to
$s,w,\mathcal L$. At a parameter $h$ already known to have analytic rank
zero, both twists have finite Tate--Shafarevich groups by the forward
theorem. Cassels's isogeny formula for the arithmetic BSD quotient
\citep{Cassels1965}, as explained in
\citet[Theorem~2.1 and the following discussion]{Tate1966}, gives
\begin{equation}\label{pw:isogeny-weight-comparison}
 s_E(h)+2w_E(h)=s_{E'}(h)+2w_{E'}(h)+O_{E,E'}(1).
\end{equation}
Indeed, at an active odd good prime the component-group valuation is
$t_p$ for the corresponding curve, so the active Tamagawa valuations
sum to twice its weight. At the fixed support only finitely many local
twisting classes occur. Differential, period, and torsion ratios are
bounded by the fixed isogeny and its reverse, and the rank-zero regulator
is $1$. These account for the bounded error in
\eqref{pw:isogeny-weight-comparison}. The isogeny identity applies to the
known finite groups and does not assume the BSD value formula.

The two twists have the same $L$-function. Their fixed normalizing periods
therefore give
$v(\mathcal L_E(h))-v(\mathcal L_{E'}(h))=O_{E,E'}(1)$, uniformly in $h$.
Combining this with \eqref{pw:isogeny-weight-comparison} shows that
$v(\mathcal L(h))-2w(h)-s(h)$ changes by a uniformly bounded amount.
Consequently \eqref{pw:partitionable-bound} transfers in either direction
with one constant, even when the choice between the two fixed isogeny
orientations depends on $h$.

\begin{lemma}[A bounded companion for a varying odd parameter]
\label{pw:rough-companion}
Fix a positive odd-sign filter and a negative genus partner $h_*$ with
$L(E^{(h_*)},1)\ne0$, chosen using Proposition~\ref{prop:rough-nonvanishing} and then
excluded from the variable parameters.  There are constants $A,C$ such that every $h$ on this filter of analytic
order one admits a negative odd fundamental discriminant $k'$ for which
\begin{enumerate}
\item $(k',2Nh)=1$, $k'$ is a square at every place of $2Nh$, and
$\mathbb Q(\sqrt{k'})$ is nonexceptional;
\item $\omega(k')\le A$;
\item $L(E^{(hk')},1)\ne0$ and
\begin{equation}\label{pw:rough-bound}
 v(\mathcal L(hk'))\le 2w(hk')+s(hk')+C.
\end{equation}
\end{enumerate}
\end{lemma}

\begin{proof}
Choose the class-group support and all other fixed label tests before
varying $h$.  Prepare an even negative-index symbol for each possible
total type at this support, and for both isogeny orientations if needed.
There are finitely many such preparations.  The positive index $|hk'|$
has nonzero total rational type: its negative-discriminant parity is
already nontrivial among the tests at $2$.  Thus in the full-torsion case
a unit seed has positive count.  In the quadratic case use the stabilized
seed of Lemma~\ref{pw:minimum-seeds}.

When $F\ne F'$, interchanging the isogeny orientation interchanges the two
counts in \eqref{pw:donor-balance}.  Select the orientation with initial
balance at least zero.  The existing donors can then pay for arbitrarily
many existing bad simples.  They do not have to be supplied by new primes.

Let $R$ be the absolute prime-count bound in Proposition~\ref{prop:rough-nonvanishing}.  Insert
in $k'$ a finite dummy supply that meets the anchors, minimum simple count,
and reserve requirements for every one of the finitely many prepared
symbols. When the quadratic fields $F,F'$ are distinct, add donor slack
paying for every reserved loss, the new dummy bad simples, and a further
possible loss of $R$ from the later rough factor. Add a fresh ramified
prime if needed to exclude the exceptional
imaginary fields.  The number of dummy primes is bounded independently
of $h$.  Their labels and initial incidences may be chosen by
Lemma~\ref{gr:realization}.

Absorb $h$ and the dummy product into the curve.  Apply
Proposition~\ref{prop:rough-nonvanishing} to complete $k'$ by a squarefree factor coprime to all
primes already used.  Prescribe its sign and its local unit classes so
that the final $k'$ is an odd negative fundamental discriminant, is a
square at $2Nh$, and has the required original local classes at $2N$.
These conditions are reduced congruences at the level of the absorbed
curve.  Additional fixed primes may be excluded.  The lemma makes the
central value at $hk'$ nonzero with at most $R$ new factors.  Its size
threshold may depend on $h$ and all these congruences; $R$ does not.

No particular resulting unit class at the auxiliary class-group support
was assumed in advance: dummy supplies were prepared for every possible
total there.  For the actual total, Lemma~\ref{pw:reserves} makes $hk'$
partitionable.  Analytic nonvanishing gives Selmer corank zero by the
forward theorem.  Proposition~\ref{pw:partitionable-even} therefore gives
\eqref{pw:rough-bound}, with the isogeny comparison above if the other
orientation was used.  The finite preparation and the bound on dummy and
rough factors give $A,C$ independent of $h$.
\end{proof}

Proposition~\ref{prop:heegner-comparison} compares the genus sum $P_h$
with the ring-class sum over $K=\Q(\sqrt{k})$:
\begin{equation}\label{pw:height-ratio}
 2j_h(k)-2j(P_h)=v(\mathcal L(hk))+O(1)
\end{equation}
whenever the two twists have analytic orders one and zero.
The error is uniform in $h$ for fixed $k$, and uniform in $h,k$ when
$\omega(k)$ is bounded and the prescribed local classes are retained.
Thus the same comparison applies to the bounded companions just constructed.

\begin{proposition}[Uniform lower bound for the odd detector]
\label{pw:odd-lower}\label{prop:rational-odd-bound}
Every analytically rank-one $h$ in the fixed positive filter satisfies
\begin{equation}\label{pw:odd-lower-bound}
 2j(P_h)\ge 2w(h)+s(h)-C.
\end{equation}
In particular $j(P_h)\ge w(h)-C$ uniformly in $h$ and its prime count.
\end{proposition}

\begin{proof}
Choose $k'$ by Lemma~\ref{pw:rough-companion}.  Combine
\eqref{pw:ring-lower-interface}, \eqref{pw:height-ratio}, and
\eqref{pw:rough-bound}.  Since the factors are coprime,
$w(hk')=w(h)+w(k')$, and $w(k')\le\omega(k')\le A$.  Hence
\begin{align*}
 2j(P_h)
 &\ge 4w(h)+s(h)+s(hk')-v(\mathcal L(hk'))-C\\
 &\ge 2w(h)+s(h)-2w(k')-C\\
 &\ge 2w(h)+s(h)-C'.
\end{align*}
Only even symbols and independently nonvanishing companions were used.
Thus this proof precedes, and does not assume, the odd minimum-depth
normalization.
\end{proof}

\subsection{A fixed companion for the odd converse}
\label{pw:fixed-companion-section}

Suppose the functional sign is odd. By
Proposition~\ref{pw:odd-lower}, the odd coefficient system can now be
normalized on the positive filter of $1$.  Fix its old support and genus
partner $h_*$.  Its moving detecting prime $\ell$ has not become an
extra fixed support requirement.

Choose a separate even symbol for negative discriminants square at $2N$,
excluding the fixed odd support, the relevant torsion ramification, and
the exceptional imaginary fields.  Choose a unit $k$ of that symbol.
Then $L(E^{(k)},1)\ne0$.  Fix this $k$ and
$K=\mathbb Q(\sqrt{k})$ for the rest of the construction, before choosing any
cube dimension.  Its fresh ramification makes it independent of the old
fixed label extensions.  At sufficiently advanced stages $\ell$ avoids
$k$.  The intrinsic odd label tests continue to omit the primes of $k$.
When the odd symbol already detects the base, the converse follows directly.
We henceforth consider the remaining case, in which the odd symbol has the
nonempty seed described above.

For the moving odd source, fix its stabilized minimum depth and seed
profiles before constructing packets. A bounded-weight seed has boundedly
many vertices here, since every prime has weight at least $1/2$. For each
fixed field used in the labels, permitting ramification above the one
moving prime $\ell$ adds only boundedly many dimensions to its elementary
squareclass space. Identify these spaces compatibly with their fixed
projections and stabilize the finitely many seed and label profiles.
The reserve bounds of Lemma~\ref{pw:reserves} are uniform over these
profiles, so the packet constructions below have a fixed size. They are
finite tables of labels; their prime representatives will be chosen at
an actual coefficient stage after the networks and their weight bounds
have been specified.

Our remaining task is to construct, for every binary dimension, a family
in which the odd test at $h$ and an even test at $hk$ are simultaneously
units away from the zero address. Every new prime will split in this
single field $K$. For quadratic $F$, the good-twin rule makes the unit
at $k$ a nonzero even test on additional prime configurations. With full
rational two-torsion, the restriction to $K$-split primes can instead
leave only the empty configuration at the original minimum depth. The
next construction supplies a nonempty witness in exactly that case;
its extra division changes the valuation estimate by a fixed constant.

\subsection{A nonempty witness with full rational two-torsion}
\label{pw:split-section}

\begin{lemma}[Recurrence congruences]\label{pw:recurrence}
Let $p$ be odd, $a\equiv p+1\pmod4$, and let
\[
 b_0=1,\quad b_1=a-\kappa,\quad
 b_{s+1}=ab_s-pb_{s-1}.
\]
If $\kappa=\pm1$, then $b_{s+1}-b_s$ is even and
\[
 \frac{b_{s+1}-b_s}{2}\equiv\frac{a-\kappa-1}{2}\pmod2
 \quad(s\ge0).
\]
If $\kappa=0$, put $b_{-1}=0$; then $b_s-b_{s-1}$ is odd for every
$s\ge0$.  Finally, for $\kappa=1$, if $a\equiv p+1\pmod{2^m}$ then
$b_s\equiv p^s\pmod{2^m}$ for every $s\ge0$.
\end{lemma}

\begin{proof}
Modulo $2$ the recurrence is $b_{s+1}=b_{s-1}$.  Its initial values prove
the parity assertions.  For $d_s=b_{s+1}-b_s$, the recurrence for $d_s$
is the same.  Directly,
$d_1-d_0\equiv(p-1)(p-\kappa)\equiv0\pmod4$ when $\kappa=\pm1$.
Thus $d_s/2$ is constant modulo $2$.  The last assertion follows by
substituting the sequence $p^s$ in the recurrence and checking its first
two terms.
\end{proof}

\begin{lemma}[Nonempty $K$-split witness]\label{pw:split-witness}
Suppose $E[2]$ is rational.  Choose a negative unit $k$ of a fixed even
symbol for $E$, prime to all previous fixed support and to the ramification
support of $E[8]$.  Put $K=\mathbb Q(\sqrt{k})$, chosen nonexceptional.
Let $S_0$ contain $2N$, the primes of $k$, and primes generating the class
group of $K$.  On the positive filter of $1$ square at $S_0$, there is an
even unit test on nonempty products of primes split in $K$ with these
properties:
\begin{enumerate}
\item it has a nonempty witness;
\item it satisfies the oriented split contraction and primary-isolation
rules over $K$, for contractions with nonempty output;
\item a unit at $D$ implies
\begin{equation}\label{pw:split-unit-bound}
 L(E^{(kD)},1)\ne0,\qquad v(\mathcal L(kD))\le 2\omega(D)+C_k.
\end{equation}
\end{enumerate}
The proof uses Lemmas~\ref{tc:trace-algebra} and~\ref{tc:inertia-identity}
to construct a further divided coefficient test.
\end{lemma}

\begin{proof}
Let $a_0(n)$ be the coefficients of the original negative-index form for
$E$.  Since $k$ attains its least normalized depth,
$v(a_0(|k|D)/a_0(|k|))\ge\omega(D)$ whenever $kD$ lies on its filter.
Use the ordinary even coefficient form for $E^{(k)}$, normalized by
$a(1)=1$, on the positive filter of $1$ square at $S_0$.
The two exact comparisons in Equation~\eqref{cs:even-ratio} give
\[
 2v(a(D))
 =v\!\left(\frac{\sqrt D\,L(E^{(kD)},1)}{L(E^{(k)},1)}\right)
 =2v\!\left(\frac{a_0(|k|D)}{a_0(|k|)}\right).
\]
Indeed, the period and the factor $\sqrt{|k|}$ cancel in the old ratio,
and the new period cancels in the new ratio.  Enlarging $S_0$ only
restricts the allowed $D$, so the original minimum still gives
\begin{equation}\label{pw:split-basic-depth}
 v(a(D))\ge\omega(D)
\end{equation}
at every allowed squarefree $D$.

If equality holds at a nonempty product of $K$-split primes, use the usual
even bit, projected from oriented $K$-labels to rational labels.  It has
all the required properties by Theorem~\ref{thm:symbol-rules}.  We
therefore assume for the rest of the proof that equality never occurs
there.  By the rational-phase assertion in Proposition~\ref{prop:even-coefficients},
normalization by $a(1)=1$ makes each nonzero squarefree coefficient a
rational number times a root of unity.  Its valuation is an integer,
not just an element of the possibly ramified coefficient field's
discrete value group.  Thus strict inequality implies
\begin{equation}\label{pw:split-extra-depth}
 v(a(D))\ge\omega(D)+1
 \quad\text{for every nonempty allowed all-$K$-split }D.
\end{equation}

Write $g=\sum a(n)q^n$ for the form.  Its characters satisfy
$\psi/\eta=\chi_{-4}$; since the scalar character $\psi$ is even,
$\eta$ is odd.  For a set $U$ of $K$-split
primes put
\[
 B_U=2^{-|U|}\prod_{p\in U}(U_p-\eta(p)V_p)g,
 \qquad
 \Theta_J=\sum_{m>0,\ (m,S_0)=1}\eta(m)m q^{Jm^2}.
\]
All expansions retain the prescribed support sieves.  The ordinary divided
form is
\[
 H_U=\frac12\left(B_U-\sum_{J\subseteq U}B_U[J]\Theta_J\right).
\]
Here $B_U[J]$ denotes its coefficient at the squarefree index $J$.

\smallskip\noindent\emph{The mod-$8$ consequence.}
We first prove
\begin{equation}\label{pw:split-mod8}
 a_p(E^{(k)})\equiv p+1\pmod8
 \quad\text{for every allowed good }p\text{ split in }K.
\end{equation}
Put $f=H_{\{p\}}\theta$ and apply its trace at a Frobenius lift $\phi_p$,
testing coefficient $p^2$.  The terms $B_{\{p\}}\theta$ and
$\Theta_{\{p\}}\theta$ have the same character, whose $p$-part is
primitive quadratic.  Thus they have no $p$-old constituent.  On their
conductor-one ramified principal-series constituents the identity is
\[
 U_p^2t(\phi_p)=U_p^3+d(\phi_p)U_p.
\]
The remaining term $\Theta_\varnothing\theta$ is unraised, and obeys the
ordinary good-prime coefficient identity.  Consequently
\[
 (t(\phi_p)f)[p^2]=f[p^3]+(d(\phi_p)f)[p].
\]
Replace the trace by a fresh split prime $p'$ with the same finite data,
including every support-unit test.  At each rational quadratic character
unramified outside $S_0$, a Frobenius lift at $p$ has the character value
of $p$.  Hence the matched $p'$ has the same value, and $pp'$ has value
$1$.  Equivalently, $p'$ has the same unit squareclasses as $p$ at all
places of $S_0$, so their product is square at all of them, including $2$.
Since $\phi_p\in G_K$, the new prime is also split in $K$.
Its coefficient on the left tests $pp'$; by
\eqref{pw:split-extra-depth} it is zero modulo the maximal ideal after
the two divisions.  To check the determinant term, let $d_1,d_0$ be the
determinant scalars on the raised and unraised character spaces.  Then
\[
 d(\phi_p)f
 =d_1f+\frac{B_{\{p\}}[1]}2(d_1-d_0)
                  \Theta_\varnothing\theta.
\]
The character difference $(d_1-d_0)/2$ is integral.  Also
$f[p]\equiv H_{\{p\}}[p]=0$ modulo the maximal ideal, since
$\Theta_{\{p\}}[p]=1$, whereas
$(\Theta_\varnothing\theta)[p]\equiv\Theta_\varnothing[p]=0$
because $\theta\equiv1$ and $p$ is nonsquare.  Thus
$(d(\phi_p)f)[p]=0$ after reduction.  This cancellation uses the two
character spaces together.

On the right, the indices $p$ and $p^3$ in $B_{\{p\}}$ have underlying
squarefree index $1$; they involve only $a(1),a(p^2),a(p^4)$.  The
$J=\varnothing$ unary term has no coefficient at either nonsquare
index.  Thus no other squarefree coefficient contributes there.
For $a=a_p(E^{(k)})$, the phase-free square coefficients are
$1,a-1,a(a-1)-p$.  Hence the remaining coefficient is, up to a unit,
\[
 H_{\{p\}}[p^3]
 =\frac{a^2-2a-p+1-p(a-2)}4
 =\frac{(a-1)(a-p-1)}4.
\]
This quotient is rational. Its vanishing in the residue field therefore
gives a positive integral $2$-adic valuation, so it is divisible by $2$.
Full rational two-torsion makes $a$ even, so
$a-1$ is odd; this proves \eqref{pw:split-mod8}.  No individual total
filter condition on $p$ was needed: its matched prime $p'$ makes $pp'$
lie on the filter, and the square-index terms on the right come from $1$.

\smallskip\noindent\emph{One further division by $2$.}
Define
\begin{equation}\label{pw:four-divided-form}
 \widetilde H_U=\frac14
 \left(B_U-\sum_{J\subseteq U}B_U[J]\Theta_J\right).
\end{equation}
We verify integrality at every coefficient.  Write a test index uniquely
as $Jrm^2$, where $J\subseteq U$ and $r$ is squarefree and disjoint from
$U$.  The underlying squarefree coefficient is at
$D=(U/J)r$.  Removing the unit phases, a prime of $U/J$ contributes an odd
factor $b_s-b_{s-1}$ with $\kappa=0$, and a prime of $J$ contributes
$b_{s+1}-b_s$ with $\kappa=(D/p)$.  Thus
\begin{equation}\label{pw:shell-valuation}
 v(B_U[Jrm^2])\ge\omega(r)
\end{equation}
by \eqref{pw:split-basic-depth}, with one additional unit of valuation
when $D$ is nonempty and all its primes split in $K$.

If $r\ne1$ and all primes of $D$ split in $K$, the last inequality is
at least $\omega(r)+1\ge2$.  If $D$ contains an inert prime, then $r$
contains at least two: the total rational $K$-character is trivial on
the filter, and all primes of $U$ split.  Then
\eqref{pw:shell-valuation} is again at least $2$.  No unary term has
this squarefree part.  This proves divisibility by $4$ for $r\ne1$.

Suppose $r=1$ and $J\ne U$.  Then $D=U/J$ is nonempty and all split,
so $B_U[J]$ is even.  This cancellation can be checked without dividing
by a leading multiplier, which may be zero.  Factor from both shell
coefficients the common phase and the scalar
$2^{-|U|+|J|}a(U/J)$, whose valuation is at least $1$ by
\eqref{pw:split-extra-depth}.  The remaining factors are the odd
differences at $U/J$, the integral divided differences at $J$, and
the ordinary square-index factors outside $U$.  By
Lemma~\ref{pw:recurrence}, their product is congruent modulo $2$ to
the product of their leading values, even when one of those values
is zero modulo $2$.  Since $m$ is odd, multiplication by $m$ preserves
this congruence.  The common scalar supplies the second factor $2$.
Retaining the phase $\eta(m)$ therefore gives
\[
 B_U[Jm^2]\equiv B_U[J]\eta(m)m\pmod4.
\]
This is exactly the cancellation with $B_U[J]\Theta_J$.

Finally let $J=U$, so $D=1$.  At a prime outside $U$ the recurrence gives
$b_s\equiv p^s\pmod4$.  At a prime of $U$,
\eqref{pw:split-mod8} gives $b_s\equiv p^s\pmod8$, and hence
\[
 \frac{b_{s+1}-b_s}{2}
 \equiv \frac{a_p-2}{2}\,p^s\pmod4.
\]
Multiplying these congruences, again retaining the phases, gives
$B_U[Um^2]\equiv B_U[U]\eta(m)m\pmod4$.  This proves integrality in
the last case and establishes \eqref{pw:four-divided-form}.
Diamonds preserve this lattice: when $J\ne U$, $B_U[J]$ is even and
the quadratic character difference supplies another factor $2$; when
$J=U$, the characters agree exactly.

\smallskip\noindent\emph{The new symbol and its witness.}
Multiply \eqref{pw:four-divided-form} by ordinary theta and use the
integral trace operators.  Two fresh traces in $G_K$ give zero after
reduction at every coefficient.  To see this, start with an arbitrary
test index $Jrm^2$ and choose the two fresh primes outside that index
and the level.  The filter imposes $\chi_K(D)=1$; since $U$ is all split,
the number of inert factors in $r$ is either zero or at least two.
The fresh traces replace $r$ by $rq_1q_2$, with both $q_i$ split.
In the first case the valuation after division by $4$ is at least
$\omega(r)+2+1-2=\omega(r)+1\ge1$.  In the second case it is at least
$\omega(r)+2-2=\omega(r)\ge2$.  The unary terms disappear under these
fresh tests.  Additional fresh $G_K$ traces only improve these estimates,
so vanishing holds on the entire reduced trace orbit.  The trace identity
therefore gives
\[
 t(gu^2)=t(g)\qquad(g,u\in G_K)
\]
on this reduced trace orbit.  Thus the tests factor through the elementary
extension of $K$ unramified off the level.

A single fresh $K$-split trace at coefficient $1$ reads the bit
\begin{equation}\label{pw:deeper-bit}
 \widetilde A(D)=
 \left(2^{-\omega(D)-1}a(D)\text{ with its phase removed}\right)\bmod2
\end{equation}
for a nonempty all-split configuration.  The contraction rule follows
from the same integral local identity as in Lemma~\ref{tc:inertia-identity}.
For completeness its first shell, with $p\in U$ and a fresh trace $q$,
has underlying configuration $D'=(U/p)q$, which is still nonempty, and
its reduction is
\[
 c_p(D')\widetilde A(D'),\qquad
 c_p(D')=\frac{a_p-(D'/p)-1}{2}\bmod2.
\]
Changing the oriented inertia bit changes $c_p$ once.  The omitted-prime
identity gives the product label and removes the entire rational prime
$p$ from the new ramification support.  Consequently its repeated
contraction has precisely the forest interpretation of
Lemma~\ref{gr:forest}.  An omitted-vertex label equal to the identity has
trace zero, which is primary isolation.  No contraction to the empty
configuration has occurred.

It remains to prove that this deeper test has a witness.  Fresh
ramification gives $K\not\subseteq\mathbb Q(E[8])$.  By Chebotarev there
is a good prime $p$ inert in $K$ and trivial on the original $E[8]$.
Then
\[
 p\equiv1\pmod8,\qquad a_p(E^{(k)})\equiv-2\pmod8.
\]
On $\widetilde H_\varnothing$, whose coefficient at $1$ is zero, test
$t(u)^2$ at coefficient $1$, where $u=\operatorname{Fr}_p$.  The square
coefficient and matching unary character give, up to a unit,
\[
 \frac{a_p(E^{(k)})-1-p}{4}\bmod2=1.
\]
The determinant term at coefficient $1$ cancels.  In characteristic two,
$t(u^2)=t(u)^2$, so a single trace in $G_K$ is nonzero.  A fresh prime
realizing this finite trace test is split in $K$.  Moreover $u^2$ is
trivial on every rational quadratic character, so the prime lies on the
required total rational filter.  It is a nonempty witness for
\eqref{pw:deeper-bit}.  Finally the exact coefficient ratio gives
\eqref{pw:split-unit-bound}, with the additional depth at most one.
\end{proof}

\begin{lemma}[Orientation correction]\label{pw:coinvariants}
Let $K$ be imaginary quadratic, let $S_0$ contain $2$ and its ramification
support, and let $V_0$ be the elementary label space over $K$ unramified
off $S_0$.  Complex conjugation acts on $V_0$ by $c$.  If the rational
total of an oriented prime configuration is zero on every rational
quadratic character unramified off $S_0$, then its oriented total belongs
to $(1+c)V_0$.  A reserved supply containing every label allows its
orientations to be flipped to realize any prescribed change in
$(1+c)V_0$.
\end{lemma}

\begin{proof}
Complex conjugation has order two outside $G_K$, so
$G_{\mathbb Q}=G_K\rtimes\langle c\rangle$.  A $c$-invariant quadratic
character of $G_K$ extends by giving $c$ value zero.  At a rational prime
outside $S_0$, inertia lies in $G_K$ and the character is trivial there;
thus its extension remains unramified outside $S_0$.  The assumed total
therefore annihilates $(V_0^*)^c$.

For a finite vector space with involution, the annihilator of
$\ker(1+c^*)=(V_0^*)^c$ is $\operatorname{im}(1+c)$.  This proves the
first assertion.  Flipping a vertex of label $v$ changes its oriented
label from $v$ to $cv$, and hence changes the total by $(1+c)v$.
Reserve one representative for every preimage needed to express the
desired element of the image.  Their independent flips give the second
assertion.  These flips preserve every rational norm label.
\end{proof}

\subsection{Common packets and rank-one interpolation}
\label{pw:assembly-section}

\begin{proposition}[Families with a fixed companion]
\label{pw:fixed-companion-families}
With the fixed odd coefficient system and companion $k$ chosen in
Section~\ref{pw:fixed-companion-section}, suppose the odd symbol at $1$
is zero. There is a constant $C$ with the following property. For every $b\ge1$ there are distinct primes outside
$2Nk$ and the fixed supports of both coefficient tests, with nonzero
linear activations, such that their signed products $h_x$,
$x\in\F_2^b$, satisfy $h_0=1$. Every $h_x$ is positive and has the fixed
local classes of $1$, and every new prime splits in $K=\Q(\sqrt{k})$.
The primes also avoid the detecting prime of the actual odd coefficient
stage. In particular $(h_x,k)=1$ for every $x$. For every $x\ne0$,
\[
 a(E^{(h_x)})=1,\qquad L(E^{(h_xk)},1)\ne0,
\]
and
\[
 j(P_{h_x})\le w(h_x)+C,\qquad
 v(\mathcal L(h_xk))\le 2w(h_x)+C.
\]
The constant $C$ and the discriminant $k$ are independent of $b$.
\end{proposition}

\begin{proof}
We construct common terminal packets, individually split in $K$, such
that the odd terminal polynomial and an even terminal polynomial testing
$hk$ are separately nonzero at each nonzero address.

First suppose $F,F'$ are quadratic.  Keep the original even symbol
evaluated on $hk$, treating the vertices of $k$ as stationary external
data.  Enrich the rational labels of the $h$-vertices by the primes of
$k$.  Their total there and at the even support is required to be zero.
Let $V$ be the resulting rational elementary label space, including the
$F,F'$ and $K$ coordinates and all relations among these characters.
The allowed simple labels form the affine set
$\{v\in V:\chi_F(v)=1,\ \chi_K(v)=0\}$, in additive notation.
The difference of any two of its elements belongs to
$\ker\chi_F\cap\ker\chi_K$, precisely the realizable rational labels
of the permitted split attachments.  Translation by this subgroup is
therefore transitive.  The fields $F,F'$ are unramified outside the old
support, whereas $K$ has a fresh ramified prime; thus $K$ is independent
of their compositum.  When $F\ne F'$, this realizes both the good-simple
pattern $(1,1,0)$ and the donor pattern $(0,1,0)$ in their three
coordinates.  When $F=F'$, all simple labels are good.

For the old odd projection, take a normal closure of its finite label
extension, including the moving prime $\ell$ at any stage realizing the
stabilized profile. It is unramified outside the old support and $\ell$.  Once $\ell$ avoids
$k$, a prime ramified in $K$ remains outside that set.  Hence $K$
has trivial intersection with this normal closure, and every old odd
label can be realized together with splitting in $K$.  This applies
also to oriented split anchors over $F$.  The added stationary $k$
coordinates may have relations, but they are already included in $V$;
they do not impose extra requirements on the omitted odd projection.

Use packets with the required simple supply and reserves in this kernel,
and with sufficiently large donor slack when $F\ne F'$. For the even
test, take no split anchors.
Correct all bad simples to good ones, pair equal enriched labels, and
retain a nonempty good base pair.  Use reserves to match the pair and
the remaining good pairs, and attach every leftover split vertex to the
last base position.  Its label is forced by total zero.  The output is
a collection of good twin pairs, each matching also at every split
vertex of $k$.  Twin insertion into the unit at $k$ proves that this
output has even symbol one.  Thus the even terminal polynomial is
nonzero.

For the odd test, project away the newly added tests at $k$.  The same
packets contain enough split anchors, simple vertices and reserve shifts
to contract to the stabilized odd seed and good pairs, by
Lemma~\ref{pw:reserves}.  Its total on that projection is the prescribed
zero total.  Incidences between the retained anchors and the good pairs
can be set after contraction, by independent edge choices.  The odd
polynomial is therefore nonzero as well.  This odd test does not include
the stationary vertices of $k$.

Now suppose the curve has full rational two-torsion.  Use the even test
of Lemma~\ref{pw:split-witness}, with $K$-oriented labels on its support
$S_0$, and let $v_1,\ldots,v_r$ be a minimum nonempty witness.  Its
oriented total $v_1+\cdots+v_r$ lies in $(1+c)V_0$, by
Lemma~\ref{pw:coinvariants}.  Choose packets containing ample copies of
every needed even label, all odd seed labels on their rational
projection, and orientation-shifting reserves.  In the common compositum,
restriction to the even label extension is surjective, so each even label
has a lift.  Each intrinsic odd label also has a lift fixing $K$: fresh
ramification gives $K\cap L_{\rm odd}=\mathbb Q$, and the moving odd prime
avoids $k$.  Extending these elements to the common compositum gives the
required lifts.  Take sufficiently many lifts for each projection
separately, without prescribing their other projection.  Duplicate each
full label, so the packet has rational total zero.

For the even contraction, retain anchors with labels
$v_1,\ldots,v_{r-1}$.  After reserving these anchors, use orientation
flips to make the total equal to $\sum v_i$; this is possible by
Lemma~\ref{pw:coinvariants}.  Merge all remaining vertices into the
last cluster.  Its label is $v_r$.  The cluster is nonempty, and the
minimum nonempty count makes the output a unit independently of its
mutual edges.  For the odd contraction, use its rational projection.
Every odd seed label lifts to a $K$-split prime because $K$ was freshly
ramified outside its intrinsic support.  Retain all but the last odd
seed label and merge the remainder.  Total zero gives the last label,
so this is also a unit.  Orientation correction is conjugation over
$\mathbb Q$, which preserves every rational quadratic label, including
those at the moving odd prime.  Adding $k$ to the common support has not
added it to the intrinsic odd projection.

In both cases put a sufficient common packet at every nonzero linear
form on $\mathbb F_2^b$.  Each nonzero address has two separately
nonzero terminal polynomials.  Theorem~\ref{gr:address} supplies primary
auxiliaries split in $K$ (and $F$ in the quadratic case), neutral at the
stationary primes of $k$, and makes both symbols one at every nonzero
address.  The auxiliary generators are primary at every support
coordinate used by either projection.  Thus the odd and even graph
identities remain identities on the same common terminal polynomial
data.  All active products $h$ are positive and square at the fixed
tests, and all their primes split in $K$.

For each fixed $b$, the terminal tables and networks now give a finite
list of coefficient tests and a maximum weight. Choose the working
precision after these bounds are fixed. Finer ray neighborhoods affect
prime realization but do not enlarge the elementary label alphabet.
Lemma~\ref{cs:fixed-weight-uniformity} and
Theorem~\ref{thm:symbol-rules} realize the required divisibilities and
symbol identities on an ultrafilter-large intersection of stages.
At an arbitrarily high actual stage in that intersection, realize the
finite graphs by distinct primes. The unit-depth constants and the
companion $k$ were fixed before $b$ was chosen, and therefore remain
uniform in these realizations.

At each nonzero address the odd unit detects analytic rank one and gives
$j(P_h)\le w(h)+C$. The even unit detects $L(E^{(hk)},1)\ne0$ and gives
$v(\mathcal L(hk))\le2w(hk)+C$, or the bound
\eqref{pw:split-unit-bound} in the full-torsion case. The fixed quantity
$2w(k)$ is absorbed into $C$. These are the asserted estimates.
\end{proof}

\begin{proposition}[Odd original base]\label{pw:odd-base}
If $c_2(E)=1$, then $a(E)=1$.
\end{proposition}

\begin{proof}
Parity gives odd functional sign. Proposition~\ref{pw:odd-lower}
therefore permits the odd normalization used above. If its symbol at $1$
is a unit, it already gives $a(E)=1$. Otherwise apply
Proposition~\ref{pw:fixed-companion-families} for arbitrarily large $b$.
At every nonzero address its estimates and
Proposition~\ref{prop:heegner-comparison}, with the fixed $k$, give
\[
 2j_h(k)\le4w(h)+C.
\]
The finite Sha lengths there are nonnegative. Hence
$2j_h(k)-4w(h)-s(h)-s(hk)\le C$, as required for Heegner interpolation.

At address zero, $L(E^{(k)},1)\ne0$ gives $c_2(E^{(k)})=0$ by
Lemma~\ref{pr:ranks}. Lemma~\ref{pr:invariance} gives
\[
 \corank_{\Ztwo}\Sel_{2^\infty}(E/K)
 =c_2(E)+c_2(E^{(k)})=1.
\]
Since $K$ is fixed, the finite quotient of this base Selmer group has a
fixed length. Theorem~\ref{thm:heegner-interpolation} makes
$L(E,s)L(E^{(k)},s)$ have a simple zero. Its second factor is nonzero at
$1$, so $a(E)=1$.
\end{proof}

\begin{theorem}[The rational-two-torsion converse]\label{thm:rational-converse}
Let $E/\Q$ satisfy $E(\Q)[2]\ne0$ and
$c_2(E)=\operatorname{corank}_{\Z_2}\Sel_{2^\infty}(E/\Q)\in\{0,1\}$,
with the local Kummer conditions at every place.  Then
\[
                 \operatorname{ord}_{s=1}L(E,s)=c_2(E).
\]
There is no restriction on the rational isogeny configuration, complex
multiplication, or reduction at $2$ or at the bad primes.
\end{theorem}

\begin{proof}
The initial isogeny choice preserves both quantities and covers both
rational two-torsion configurations.  Proposition~\ref{pw:even-base}
proves the corank-zero case.  In the corank-one case,
Proposition~\ref{prop:rational-odd-bound} supplies the coefficient
normalization, and Proposition~\ref{pw:odd-base} proves analytic rank one
using a companion fixed before the binary dimension grows.
\end{proof}

\section{Selmer matrices and three-state relations}
\label{sec:lm:local-matrices}

Assume that $V=E[2]$ is irreducible, and write
$F=\Q(V)$ and $G=\Gal(F/\Q)\in\{C_3,S_3\}$.
The symbol rules will be useful because a unit symbol bounds the
$2$-Selmer dimension. We now express local Selmer conditions by a
binary matrix and make that bound constrain the number of singular
blocks. The second half of the section develops the elementary matrix
operations needed to apply this constraint.

Write the Weil pairing additively, and choose a basis of $V$ in which
its matrix is
\[
 J=\begin{pmatrix}0&1\\1&0\end{pmatrix}.
\]
Fix a finite set $S$ containing $2$, $3$, the bad primes, all auxiliary
exclusions, and rational primes whose prime ideals generate the class
groups of every subfield of $F$. Local specifications at $S$ will also
include the archimedean places. The detecting prime $\ell$ for an odd
coefficient system is kept separate: all generators below have valuation
zero there. This separation keeps the set of support types finite as
$\ell$ varies. We shall apply the construction to a fixed curve, after
absorbing any initial quadratic twist into $E$.

\subsection{Global classes and local equations}

At a split prime $p$, choose a place of $F$ and put $H_p=1$.
At a simple prime, choose a place with Frobenius a specified
transposition $c$ and put $H_p=\langle c\rangle$.
Set
\[
 F'_p=F^{H_p},\qquad V_p=V^{H_p}.
\]
Thus $V_p$ has dimension two or one, respectively. The chosen place
lies above a degree-one prime $P_p$ of $F'_p$. By the choice of $S$,
there is $\pi_p\in(F'_p)^\times$ whose divisor away from $S$ is $P_p$.
For $x\in V_p$ define
\begin{equation}\label{eq:lm:generator-class}
 C_p(x)=\cor_{F'_p/\Q}(\chi_{\pi_p}x)\in H^1(\Q,V).
\end{equation}
Here $\chi_{\pi_p}$ is the quadratic Kummer character; its product with
$x$ is a cocycle because $H_p$ fixes $x$.

The inertia coordinate of $C_p(x)$ at $p$ is $x$, and all its other
inertia coordinates outside $S$ vanish. Indeed, local
restriction--corestriction has exactly one ramified summand, at the
degree-one place where $\pi_p$ has valuation one. These coordinates
therefore make the classes for distinct $p$ independent. A root prime
has $V^{\Fr_p}=0$ and contributes no coordinates.

Let $h$ be a twist containing the active nonroot primes in its
discriminant. At a good odd prime $p$ ramified in $h$, inertia identifies
the local Kummer condition with $V_p$. To see surjectivity, halve a
Frobenius-fixed two-torsion point over the maximal unramified extension.
The ramified twisting involution negates its half, so the resulting
Kummer inertia is that two-torsion point. The Kummer condition has
dimension $\dim V_p$ by local duality, giving the asserted isomorphism.

Fix arithmetic Frobenius lifts $\phi_p$. A second local coordinate is
unramified evaluation at $\phi_p$, modulo Frobenius coboundaries.
Let $K_p(x)(\phi_p)$ be this coordinate for the Kummer lift of inertia
$x$. It is determined by $E[4]$ and the local twisting squareclass.
Define the block matrix $B$ by
\begin{equation}\label{eq:lm:raw-matrix}
 B_{qp}(y,x)=\langle C_p(x)(\phi_q),y\rangle
 -\mathbf1_{p=q}\langle K_p(x)(\phi_p),y\rangle,
 \qquad x\in V_p,\quad y\in V_q.
\end{equation}
Thus $Bx=0$ is exactly the collection of active-prime Selmer equations
for $\sum_p C_p(x_p)$. Changing the unramified unit part of the twist
adds $J$ to a split diagonal block. On a simple line it changes nothing,
since the Weil pairing restricts to zero on a one-dimensional space.

Record in the \emph{type} $\mathfrak h_p$ the subgroup $H_p$ and the
squareclasses of $\pi_p$ at the places over $S$, including signatures.
Only finitely many types occur. Finer generator neighborhoods will be
used to reproduce self entries, but they are not part of this type
alphabet and do not enter its codimension cost. Global reciprocity gives
\begin{equation}\label{eq:lm:symmetry-defect}
 B_{qp}+B_{pq}^{t}=b(\mathfrak h_q,\mathfrak h_p),
\end{equation}
where $b$ is the sum of the local cup pairings at $S$.
Indeed, at an active prime decompose each class into its Kummer lift
and its unramified difference. Both planes are isotropic, and the
cross pairing is unramified evaluation paired with inertia. The terms
at $p$ and $q$ give the left side; the sum over all places is zero.
Elsewhere outside $S$ the classes are unramified, so the local pairing
vanishes. This includes $\ell$, where the generators have unit valuation.
The same calculation shows that $b(\mathfrak h,\mathfrak h)$ is alternating.

\subsection{Prescribing the matrix entries}

We next prove that the off-diagonal entries and the rational Legendre
bits can be varied independently. The latter bits will control split
diagonals without disturbing the other blocks. The class-field inputs
are ray-class reciprocity, the Hilbert product formula, and Chebotarev
\citep[Chapter~VI, Sections~7--8, and Chapter~VII, Section~13]{Neukirch1999};
see also \citet[Chapter~V, Theorems~3.23, 5.3 and~5.5,
and Example~5.4]{MilneCFT2020}.

\begin{lemma}[Realization of finite matrix prescriptions]
\label{lem:lm:arithmetic-realization}
Fix finitely many individually admissible oriented prime templates,
including generator neighborhoods at $S$, any additional finite places,
and oriented Frobenius on $E[4]$ and fixed quadratic-radical extensions
unramified outside that support. At arbitrarily fresh primes these
templates admit the following simultaneous prescriptions.
\begin{enumerate}[label=\textup{(\roman*)}]
\item Between distinct vertices, one direction of the double-coset
residue bits is free and reciprocity determines the reverse direction.
The available matrix blocks are all binary matrices for two split
vertices in the $S_3$ case, and $J\F_4$ in the $C_3$ case, with the
rotation algebra identified with $\F_4$. For split--simple vertices
the full vector space is available, and for simple--simple vertices
the scalar entry is available.
\item The rational Legendre bit between the absolute primes is
independent of the matrix block, subject to rational reciprocity in
the reverse direction.
\item Repeating a template fixes its simple self entry and its split
self block modulo a scalar multiple of $J$. A sufficiently primary
split template with trivial $E[4]$-Frobenius has self block in $\F_2J$.
\end{enumerate}
Here \emph{primary} means sufficiently close to $1$ at the finite
places of $S$ and totally positive at the real places. New cross-bit
prescriptions are made outside the support. All assertions concern
finite configurations; an additional moving prime $\ell$ may be
included among the prescribed places.
\end{lemma}

\begin{proof}
\emph{Off-diagonal entries.}
Restriction--corestriction in Equation~\eqref{eq:lm:generator-class}
gives, for distinct vertices,
\begin{equation}\label{eq:lm:double-coset}
 B_{qp}(y,x)=
 \sum_{g\in H_q\backslash G/H_p}
 e_{pq}(g)\langle gx,y\rangle.
\end{equation}
The bit $e_{pq}(g)$ is the quadratic residue of $g\pi_p$ at the
oriented $q$-place in $F^{H_q\cap gH_pg^{-1}}$. Indeed, pair local
transfer with the $H_q$-invariant vector $y$; the orbit length
$[H_q:H_q\cap gH_pg^{-1}]$ is precisely the residue degree in the
norm calculation. In this intermediate field, each of $\pi_q$ and
$g\pi_p$ has its off-support divisor at a single prime: its oriented
Frobenius generates the subgroup fixing the generator. Hilbert
reciprocity therefore relates $e_{pq}(g)$ to $e_{qp}(g^{-1})$, with
correction only at $S$.

For a new vertex the independent forward bits are residue conditions
at distinct primes of its own field $F'_p$. Move an initial template
away from the old primes and choose a degree-one prime in the ray class
whose quotient generator has the prescribed neighborhoods at $S$,
real signs, and nonzero residues at those old primes. Multiplying the
template generator by this quotient gives the required new generator.
These conditions also preserve the admissible fixed Frobenius data:
changes by local units at old primes are trivial on the intersection
with extensions unramified outside the fixed support. Chebotarev thus
imposes the ray and fixed-label conditions simultaneously. Prime ideals
of degree greater than one over $\Q$ have density zero among prime
ideals of $F'_p$, so requiring degree one loses no admissible class.
The same argument permits additional finite support, including $\ell$.
Reverse bits are then determined by reciprocity. Rational norm bits at
already chosen root primes may be imposed as additional residue
conditions in the ray modulus: finite-field norm maps are surjective,
and these primes contribute no matrix coordinates.

\emph{Matrix span and the norm bit.}
The norm to $\Q$ writes the rational absolute-prime Legendre bit as
the sum of the double-coset bits, plus a fixed support-unit correction.
For two split vertices the matrices in
Equation~\eqref{eq:lm:double-coset} are $Jg$. They span
$\operatorname{Mat}_2(\F_2)$ for $S_3$ and $J\F_4$ for $C_3$.
A nontrivial rotation satisfies $1+R+R^2=0$: changing these three
bits fixes the matrix and changes their sum. For split--simple
vertices, the three nonzero vectors or dual vectors likewise span
the full space and have the same odd relation. For two simple vertices,
the matching double coset has zero pairing and the other has nonzero
pairing. The first bit changes the norm bit alone. This proves
\textup{(i)} and \textup{(ii)}.

\emph{Self entries.}
The nonidentity double cosets in self-transfer are unramified at the
oriented prime. The identity term is a scalar times the pairing and
is therefore zero on a simple line. At a split vertex the two rotation
terms have fixed sum by reciprocity; changing both alters the block
only by a scalar multiple of $J$. Every other term has an involution
representative $\tau$ exchanging the oriented place with another one.
For a simple vertex this includes a transposition other than $c$
representing the nonmatching self double coset.

Work in $F^\tau$ and put
\[
 \alpha=\pi_p\tau\pi_p,\qquad
 \beta=\pi_p+\tau\pi_p,\qquad
 F=F^\tau(\sqrt{D_\tau}).
\]
The distinct divisors of $\pi_p$ and $\tau\pi_p$ ensure $\beta\ne0$.
If $v_0$ is the place below the oriented prime, the residue symbol of
$\tau\pi_p$ there is $(\alpha,\beta)_{v_0}$: $\alpha$ has valuation
one and $\beta$ is a unit with the required residue. Outside
$S\cup\{v_0\}$, $\alpha$ is a unit. When $\beta$ is a unit, both
$(\alpha,\beta)_v$ and $(-D_\tau,\beta)_v$ vanish. When $\beta$ has
positive valuation, put
$\gamma=(\pi_p-\tau\pi_p)/\sqrt{D_\tau}\in F^\tau$. The identity
\[
 \beta^2-4\alpha=D_\tau\gamma^2
\]
shows that $\alpha$ has residue squareclass $-D_\tau$; the two
symbols are again equal. At $v_0$, $(-D_\tau,\beta)_{v_0}=0$.
Reciprocity applied to both symbols gives, in additive notation,
\begin{equation}\label{eq:lm:self-spin}
 (\alpha,\beta)_{v_0}
 =\sum_{v\in S}\bigl((\alpha,\beta)_v+(-D_\tau,\beta)_v\bigr).
\end{equation}

At finite support places the right side is locally constant in a
sufficiently small neighborhood of the template. At a real place
where $F/F^\tau$ is complex, $\alpha$ and $-D_\tau$ are positive.
At a split real place, opposite signs of $\pi_p$ and $\tau\pi_p$
make the two terms cancel; equal signs determine their sum. Thus
signatures suffice, without archimedean size restrictions.
For a primary template, $\alpha$ is a square at the finite support
places and $\beta$ has the squareclass of $2$. The support sum becomes
$\sum_{v\in S}(-D_\tau,2)_v=0$, since the terms outside $S$ vanish.
The involution terms are therefore zero. Together with the local
Kummer lift determined by $E[4]$, this proves \textup{(iii)}.
\end{proof}

\subsection{Symmetrization and the Selmer bound}

The realization lemma supplies the allowable matrix blocks. We now remove
their fixed-support reciprocity defect and compare the resulting kernels
with actual Selmer classes. The fixed support types, stationary
environment and additional local conditions will account for the loss
in dimension.

Choose a matrix rule $L(\mathfrak h,\mathfrak h')$ in the available
block spans such that
\begin{equation}\label{eq:lm:type-correction}
 L(\mathfrak h,\mathfrak h')+
 L(\mathfrak h',\mathfrak h)^t=b(\mathfrak h,\mathfrak h').
\end{equation}
For distinct types, choose either direction using
Equation~\eqref{eq:lm:symmetry-defect}. For equal types, the right
side is alternating. It is a matrix plus its transpose in the full
span; in the rotation span use $R+R^{-1}=1$ after multiplication by
$J$. On a simple line it is zero. Choose $L=0$ when one type is
primary and trivial at $S$, since then its support pairings vanish.
Define, on every pair including the diagonal,
\[
 M_{qp}=B_{qp}-L(\mathfrak h_q,\mathfrak h_p).
\]
The matrix $M$ is symmetric. Lemma~\ref{lem:lm:arithmetic-realization}
allows zero cross-$M$ entries between prescribed blocks independently
of their internal entries.

A \emph{stationary environment} consists of the fixed coordinates
outside the variable blocks; its dimension is the number of those
coordinates.

\begin{proposition}[Nullity forces Selmer dimension]
\label{prop:selmer-matrix-bound}
Fix the curve, support types, a bound for the number of additional
constrained places, and a bound for the dimension of a stationary
environment. Partition a finite configuration into blocks with zero
cross-$M$ entries, and denote their matrices by $M_i$. There is a
constant depending only on these fixed data such that
\begin{equation}\label{eq:lm:selmer-nullity}
 \dim_{\F_2}\Sel_2(E^{(h)}/\Q)
 \geq\sum_i\dim_{\F_2}\ker M_i-C.
\end{equation}
At a unit coefficient symbol the sum of nullities, and hence the
number of singular blocks, is therefore uniformly bounded. The bound
works for both coefficient sources and is independent of the moving
prime $\ell$.
\end{proposition}

\begin{proof}
On $\bigoplus_i\ker M_i$, impose
\begin{equation}\label{eq:lm:type-sums}
 \sum_{p:\,\mathfrak h_p=\mathfrak h}x_p=0
 \quad\text{for each type }\mathfrak h,
\end{equation}
using the fixed orientation and basis to identify coordinates of the
same type. These equations have codimension at most
$\sum_{\mathfrak h}\dim V_{\mathfrak h}$. They kill the correction
in every row, since that correction is
\[
 \sum_{\mathfrak h}L(\mathfrak h_q,\mathfrak h)
       \sum_{p:\,\mathfrak h_p=\mathfrak h}x_p.
\]
Thus $Bx=0$ at all block positions. Set the environment coordinates
to zero and impose its row conditions, losing at most its dimension.
The Selmer conditions at $S$ and the additional constrained places
have bounded total codimension in local $H^1$. At every remaining
good place the classes satisfy the unramified Kummer condition; at a
root prime local cohomology is zero. Their sums are therefore Selmer
classes, and independent inertia coordinates make the map injective.
This proves Equation~\eqref{eq:lm:selmer-nullity}.

At a unit symbol, Proposition~\ref{prop:unit-symbol-detection} bounds
the Selmer dimension, giving the remaining assertions. The type
constraints use only the fixed $S$; a condition at the single moving
prime contributes a uniformly bounded local dimension.
\end{proof}

\subsection{Three states on two coordinates}

We have obtained two kinds of information from an arithmetic prime
configuration: its coefficient symbol and a matrix whose nullity
bounds its Selmer dimension from below. To compare configurations,
we now vary exactly two coordinates at a time. Both the symbol and
the determinant obey one common three-state relation. They remain
different functions; the relation gives linearity, not their equality.
The bounded number of singular blocks at a unit symbol will be the
additional input used in the next section.

A \emph{site} consists of either one split vertex or an ordered pair
of simple vertices, and always supplies two matrix coordinates.
At a split site use the three states
\[
 \text{absent},\qquad
 \text{present with diagonal }Q,\qquad
 \text{present with diagonal }Q+J.
\]
The rational Legendre bit of the aggregate root selects the two
present diagonals. At a simple-pair site use absence of both vertices
and two present variants differing in their nonmatching cross bit.
Replace the second prime by a prime with the same template, self data,
norm data, and outside cross bits. Its matching bit with the first
prime must also change to keep the rational norm bit fixed; that bit
has zero matrix contribution. Thus the two present matrices again
differ by the symmetric toggle on their coordinate pair. Require all
alternatives to agree against alternatives at other sites. These are
finite simultaneous prescriptions allowed by
Lemma~\ref{lem:lm:arithmetic-realization}.

In comparing states, keep fixed the total parameter label at the fixed
support: the product of the labels of the active sites and the aggregate
roots. The aggregate-root label itself changes when a site is deleted,
absorbing the deleted labels.

\begin{lemma}[The site relations]\label{lem:lm:site-relations}
Fix the total parameter label at the fixed support and the diagonal
choices at the other present sites. In each state choose aggregate
roots restoring this label, compatibly with an order-three residual
Frobenius, and preserving those diagonals. For either coefficient
symbol, its three values at one site sum to zero. The same holds for
$\det M$, with empty determinant one.
Consequently both functions extend linearly to the tensor product of
one copy of $\F_2^2$ per site, whose three nonzero vectors represent
the three states. These factors encode evaluations, not vectors of $V$.
\end{lemma}

\begin{proof}
For a split site this is the root-deletion rule of
Theorem~\ref{thm:symbol-rules}. Deletion transfers the removed
vertex's rational labels to the aggregate root. Signed prime
discriminants preserve the unit twist, and hence the diagonal,
at every remaining prime.

For a simple-pair site, omit the varied second prime from its
opposite-parity trace test. In an extension of $G$ by an elementary
kernel, two lifts of its transposition $c$ are conjugate by the kernel
exactly when their difference lies in the image of $1+c$. Dually,
this means agreement on the $c$-fixed radical squareclasses.
Besides fixed-support radicals, these are generated by products of
the orbit generators $g\pi_p$ over $c$-orbits: their off-support
divisors are independent and conjugation permutes them. Singleton
orbits test matching bits. A two-element orbit has product in $F^c$;
its symbol at the varying degree-one prime is, by the residue norm,
the symbol of one factor in $F$, namely the nonmatching bit.

The variants agree on support radicals and the orbit tests for every
other vertex. At the first prime they change both matching and
nonmatching bits. Modulo the image of $1+c$, their difference is
therefore the product of inertia in the matching transposition
orientation and inertia in the other orientation. Apply these
substitutions successively. The matching-orientation derivative is zero;
the different-transposition derivative deletes the pair by
Theorem~\ref{thm:symbol-rules}. Multiplication
by inertia leaves the residual transposition unchanged. The intermediate
Frobenius test is realized by Chebotarev; it need not be a site variant.
The resulting root has rotation image and transfers the pair's rational
labels to the aggregate. Labels elsewhere are unchanged. At simple
vertices root probes are immaterial; at split vertices the common
rational labels preserve the diagonals. Matched actual probes, or their
aggregate, are available because their discriminant character is trivial.
This proves the symbol relation.

For the determinant, reverse every permutation cycle of length at
least three. The resulting terms cancel in characteristic two, leaving
only involutions. Toggling the symmetric entries $(i,j),(j,i)$ changes
exactly the terms containing $(ij)$, whose sum is the determinant with
coordinates $i,j$ deleted. This includes the split $J$-shift. Finally,
the three-state relation in each factor is precisely linearity on its
three nonzero binary vectors.
\end{proof}

\begin{lemma}[Binary tensor span]\label{lem:binary-tensor-span}
Let $W=\bigotimes_{i=1}^rW_i$, where each $W_i$ has dimension two
over $\F_2$, and let $0\ne f\in W^*$. The nonzero pure tensors
annihilated by $f$ span $\ker f$. Thus if pure tensors $e,t$ satisfy
$f(e)=f(t)=1$, then $t$ is $e$ plus a sum of pure zeros of $f$.
\end{lemma}

\begin{proof}
It suffices to show that a functional $g$ vanishing on the pure zeros
of $f$ is either $0$ or $f$. Induct on $r$, the case $r=0$ being
immediate. Slice at the three nonzero vectors of the last factor,
numbering so that $f_3=f_1+f_2$ and $g_3=g_1+g_2$.
If $f_i=0$, then $g_i=0$ because the smaller pure tensors span.
Otherwise induction gives $g_i=\epsilon_i f_i$ for a binary scalar
$\epsilon_i$. If all three $f_i$ are nonzero, $f_1,f_2$ are independent,
so the sum identities force all $\epsilon_i$ to agree. If one $f_i$
vanishes, the other two, and their $g_i$, agree. All three cannot vanish.
Hence $g=0$ or $g=f$.
\end{proof}

Take $f=\det M$ and let $e$ be the all-absent state. Every invertible
state is the empty state plus a sum of singular pure states.
Moreover, on any selected set of active sites, some choice of pair
toggles makes $M$ invertible: the iterated determinant difference over
all selected pairs is the empty determinant, one. These facts turn the
Selmer nullity bound into a constraint on symbol evaluations.

\subsection{Making one address singular}

The final construction singles out one nonzero vector of a binary
cube. Every position is activated by a linear form, as required for
interpolation; auxiliary split sites repair the matrix at all other
nonzero vectors.

\begin{lemma}[A block address program]\label{lem:lm:address-program}
Let $M_*$ be the singular matrix of a finite site state, and let $D$
be a sum of its pair toggles such that $N_*=M_*+D$ is invertible.
For $b\geq2$ and $0\ne x_*\in\F_2^b$, one can add finitely many
primary split sites, each activated by a linear form, so that:
\begin{itemize}
\item all sites present in the original state have activation $l_0$,
with $l_0(x_*)=1$;
\item no site is active at zero, and at $x_*$ precisely the original
state $M_*$ remains;
\item the matrix is invertible at every other nonzero address;
\item the active auxiliary matrix is always invertible, and its Schur
complement contributes $D$ exactly when the original sites are active
and $x\ne x_*$.
\end{itemize}
All entries use the freedoms of
Lemma~\ref{lem:lm:arithmetic-realization}.
\end{lemma}

\begin{proof}
Choose a basis $y=(y_1,\ldots,y_{b-1})$ for the forms annihilating
$x_*$ and complete it by $l_0$. Then $y=0$ exactly at $0,x_*$.
For each required toggle and each nonzero $\eta\in\F_2^{b-1}$,
choose a basis of forms all equal to one at $\eta$. Such a basis is
obtained from a form $u$ with $u(\eta)=1$ by adding $u$ to a basis
of the annihilator of $\eta$. Its unique simultaneous all-one
solution is $\eta$.

Construct a path with $r=b-1$ steps, using two new split positions at
each step, both activated by the corresponding basis form. In scalar
block notation, where a scalar denotes that scalar times $J$, put
\begin{equation}\label{eq:lm:path-matrix}
 P=\begin{pmatrix}0&T\\T^t&B_e\end{pmatrix},
 \qquad T=I+\text{lower consecutive shift}.
\end{equation}
The near and far endpoints are steps $1$ and $r$. The matrix $B_e$
has one nonzero diagonal entry at the far endpoint. Restrict both sides
to active steps, retaining only the original consecutive edges. The
restricted $T$ is triangular
with diagonal one, so $P$ is invertible, with scalar inverse
\[
 P^{-1}=\begin{pmatrix}
 T^{-t}B_eT^{-1}&T^{-t}\\T^{-1}&0
 \end{pmatrix}.
\]
Couple the addressed coordinate pair by $J$ to the first side at the
near endpoint. Such a pairing is available also for two simple lines.
Solving $Tv=e_{\rm near}$ propagates one to the far endpoint exactly
when every step is active; a missing step interrupts the consecutive
shift. The Schur contribution is consequently $J$ on a complete path
and zero otherwise, including when the near endpoint is absent.

For each nonzero $y$, exactly one path, the one indexed by $\eta=y$,
is complete. Their direct sum supplies the toggle when $y\ne0$ and
supplies zero when $y=0$. Repeat for each toggle in $D$, with zero
mutual cross-$M$ entries. The actual matrices are the scalar matrices
tensored with $J$, so the same inverse calculation gives the asserted
binary Schur complements. All blocks lie in the permitted spans, and
primary diagonals are $0$ or $J$.

At $x_*$ the auxiliary sites are absent and $l_0=1$, leaving $M_*$.
At any other nonzero address with $l_0=1$, eliminating the auxiliary
sites leaves $M_*+D=N_*$. If $l_0=0$, only the invertible auxiliary
matrix remains. Every activation is linear and vanishes at zero.
\end{proof}

\section{The converse when \texorpdfstring{$E[2]$}{E[2]} is irreducible}
\label{sec:sa:saturation}

The local matrices now supply the families needed for interpolation.
We shall arrange that a coefficient symbol is a unit at every nonzero
vertex of a binary cube, with all added primes absent at zero.
The key observation is that a unit evaluation has only boundedly many
singular matrix blocks.  A maximal list of such blocks will therefore
force all the required unit evaluations.  In the odd case, a second
construction supplies simultaneous nonvanishing for an auxiliary even
symbol.

\begin{theorem}[The irreducible converse]
\label{thm:irreducible-converse}
Let $E/\Q$ be an elliptic curve for which $E[2]$ is irreducible as a
$G_{\Q}$-module.  If
\[
 c_2(E)=\corank_{\Ztwo}\Sel_{2^\infty}(E/\Q)\in\{0,1\},
\]
then $\ord_{s=1}L(E,s)=c_2(E)$.
\end{theorem}

The curve is fixed throughout.  Every constant may depend on it and on
the finite choices specified below, all made before the cube dimension
grows.  The result applies pointwise to any quadratic twist by taking
that twist as the initial curve.

\subsection{Finite evaluation tables and ordered profiles}

A block is a finite set of the three-state sites of
Section~\ref{sec:lm:local-matrices}, together with their prescribed
matrix entries. A state chooses one of the three alternatives at each
site. Its evaluated matrix $M_s$ deletes the coordinates of absent sites
and applies the toggles of the chosen active alternatives. The empty state has determinant $1$.
We call a state singular when $\det M_s=0$.

The matrix and the coefficient symbol give different information about
such a state. The raw matrix $B$ records the active-prime Selmer equations;
$M$ is the symmetric matrix obtained after the fixed-support correction.
The symbol still depends on the arithmetic realization, even when those
matrix entries have been prescribed. What links them is
Proposition~\ref{prop:selmer-matrix-bound}: a symbol value $1$ permits
only boundedly many singular blocks. We shall first stabilize the
remaining arithmetic dependence and then use this bound.

Fix a normalized coefficient symbol $A$ on the filter of the parameter
$1$. For the odd source retain its condition at the moving detecting
prime $\ell$. Consider an actual finite array of blocks, numbered in
increasing order, with $M=0$ between different blocks. For chosen
indices $i_1<\cdots<i_j$ and states $s_1,\ldots,s_j$, write
\[
 \mathcal E(i_1,\ldots,i_j;s_1,\ldots,s_j)\in\F_2
\]
for its coefficient-symbol value. Each table entry uses its own
aggregate-root probe, which restores the total support label of $1$
and gives the chosen diagonals at active split sites. Root substitution
makes its value independent of the probe with these data. The calibration
below will realize all entries used by a cube with one common set of
linearly activated root primes. The total label, rather than the aggregate root's own label,
is fixed: deleting a site transfers that site's labels to the root.
The substitution rules therefore preserve this value when an unused
site or an empty block is deleted. In particular, the empty table entry
is $A(1)$.

There are two independent indices in these tables. Inside each block,
site slots specify where a state and its addressing paths can be put.
The indices $i_1<\cdots<i_j$ specify ordered copies of the whole block.
We will use a common countable library of site slots in every copy,
closed under the finite addressing constructions of
Lemma~\ref{lem:lm:address-program}. A state in this library uses only
finitely many sites, with every other site absent. A \emph{word} is an ordered finite
list of states from this library, with empty states omitted. An
\emph{ordered profile} assigns a value to each such word, consistently
with the three-state relations in each block. The next lemma constructs
one by stabilizing finite tables: its value depends on the ordered
states, but not on which increasing list of block indices carries them.
It asserts no invariance under permuting the states.

\begin{lemma}[A finitely realizable ordered profile]
\label{lem:sa:profile}
For the fixed symbol $A$, there is a countable library of site slots in each block
and an ordered profile with the following
properties.
\begin{enumerate}[label=\textup{(\roman*)}]
\item Every prescribed finite part is realized by an array of distinct
primes with the data of Lemma~\ref{lem:lm:arithmetic-realization}.
For the odd symbol, it is realized at arbitrarily advanced stages of
the prescribed precision stabilization.
\item A word evaluation depends on the states and on their order, and
not on the particular distinct block indices used to place them.
The tensor relations hold separately in all of its blocks.
\item The site library contains a transferred minimal-weight unit witness and
the primary split sites needed for every finite collection of requests
in
Lemma~\ref{lem:lm:address-program}.
\item For a constant $C$ independent of the finite array, a unit
evaluation contains at most $C$ singular blocks.
\end{enumerate}
Couplings to finitely many fixed primes, kept outside the address
symbol, may be prescribed in advance.  In the even case, all primes
may also satisfy a bound on $v_2(p^2-1)$ independent of the finite
array.
\end{lemma}

We first derive the consequence of this lemma. Its proof follows the
saturation argument and explains how the ordered values retain finite
arithmetic realizations, including for the moving odd source.

\begin{lemma}[Saturation of a punctured cube]
\label{lem:sa:punctured-cube}
For the profile in Lemma~\ref{lem:sa:profile}, either the empty word
has value~$1$, or there is a fixed finite word $W$ of singular blocks
with value~$1$ such that, for every sufficiently large $b$, an addressed
family has symbol~$1$ at every nonzero point of $\F_2^b$ and has all
sites absent at zero.  Every such finite cube is arithmetically
realizable.
\end{lemma}

\begin{proof}
In the tensor space of a block, every nonsingular state is the empty
state plus a sum of singular states, by
Lemma~\ref{lem:binary-tensor-span}.  Expand the transferred unit
witness this way.  Some word consisting entirely of singular blocks,
possibly the empty word, has value~$1$.

The lengths of all such words are bounded by
Lemma~\ref{lem:sa:profile}\textup{(iv)}.  Choose one of maximal
length $m$.  If $m=0$, the empty word has value~$1$.  Otherwise fix
one particular maximizing word
\[
 W=(W_1,\ldots,W_m).
\]
Each $W_i$ uses finitely many slots.  These slots, and thus their total
number, are fixed before $b$ grows.  A bound on the number of blocks
alone would not give this conclusion.

For each $a\in\F_2^b\setminus\{0\}$, place a copy of $W$ in its
original block order.  Address its $i$th block to $a$ using
Lemma~\ref{lem:lm:address-program}.  At $a$ the state is exactly
$W_i$, with every path site absent.  At every other nonzero address
the evaluated block is nonsingular, and at zero it is empty.  All
these finitely many programs are available in the recursive slot list.

At a nonzero address $x$, expand every block outside the copy
designated for $x$ as empty plus singular states.  Choosing the empty
state in every such block leaves $W$ in its original order, hence
gives~$1$.  Every other term contains its $m$ singular blocks and at
least one additional singular block, so maximality forces its value
to be zero.  This includes alternatives involving path slots:
maximality was taken over the whole slot profile.  Multilinearity
therefore gives~$1$ at $x$.  All evaluations use one finite restriction
of the profile, which is realizable by Lemma~\ref{lem:sa:profile}.
\end{proof}

\begin{proof}[Proof of Lemma~\ref{lem:sa:profile}]
\emph{The witness and available slots.}
Begin with the nonroot positions of a minimal-weight unit witness.
The minimum-transfer rule in Theorem~\ref{thm:symbol-rules} moves them
to fresh primes while retaining their fixed-support Frobenius labels.
Neither prescribed internal couplings nor couplings to stationary
primes change this transfer: the latter primes are outside the
coefficient index.  The total simple parity is even, so pair its
simple positions into sites.  Fix admissible generator templates,
including support squareclasses and self types, for this pattern.

At each finite stage place a fresh copy of the pattern in every block
to be colored. First apply Lemma~\ref{tc:minimum-transfer} with
matched actual root probes. The nonroot support labels are unchanged,
so retaining the aggregate root's support label restores the same total
parameter label. Adjust only its bits at active vertices to give the
prescribed diagonals. By Proposition~\ref{tc:irreducible-rules}\textup{(ii)}, a split
root-bit derivative deletes weight one and vanishes below the minimal
weight, while a simple root-bit derivative is zero. Thus every singleton
witness still has value~$1$, even after all internal and cross-block
entries have been prescribed.  For the odd symbol, make the transfer
at the actual coefficient level and retain the witness's labels at
the detecting prime.  Thus the witness entry of every singleton color
is~$1$ before taking homogeneous restrictions.

We now enlarge the site library inside each block.  For every finite list of earlier
site slots and every finite list of addressing requests on it, allocate
fresh primary split site slots realizing all those requests.  Programs
used together receive distinct path slots.  Set every undesignated
coupling to earlier slots equal to zero in $M$.  This gives a
countable formal list, since there are countably many finite requests.
Any finite set of requested slots, together with its finitely many
predecessors, is realizable by
Lemma~\ref{lem:lm:arithmetic-realization}.
The two simple-pair variants have identical outside bits and the
prescribed internal toggle; the aggregate root supplies the split
diagonal alternatives.  Corresponding types have the same reciprocity
correction $L$, so a forward entry $B=M+L$ gives the required reverse
entry as well.  In every alternative impose $M=0$ between blocks.

\emph{Order-invariant evaluations.}
Fix the first $q$ sites of the common library. There are finitely many
states on these sites, and their matrix entries have been prescribed.
For each $j\le q$, color an increasing $j$-tuple of block indices by
the finite table of values
$\mathcal E(i_1,\ldots,i_j;s_1,\ldots,s_j)$ on these states.  There are finitely
many colors.  Apply finite Ramsey successively for these finitely
many arities.  Starting with sufficiently many blocks gives a
homogeneous list of any prescribed finite length.

Let the lengths, arities, and slot cutoffs tend to infinity.
At each fixed cutoff there are only finitely many possible color
vectors.  Successive subsequences and a diagonal subsequence make
every finite vector constant.  These compatible values define the
ordered profile.  Every finite restriction occurs in an original
finite array.  The tensor identities, deletion rules, and unit
witness survive because each involves only such a restriction.

For the odd symbol, carry out this construction along its fixed
nonprincipal ultrafilter of coefficient precisions.  The witness has
bounded weight, so its fixed-$S$ generator squareclasses and binary
self data have only finitely many possibilities; first stabilize these
data. Finer ray-neighborhood depths are realization data and need not
stabilize. Impose the extra labels at $\ell$ at each actual level.  For every fixed weight
bound, Lemma~\ref{cs:fixed-weight-uniformity} gives an
ultrafilter-large set on which all required lower bounds and
smaller-weight vanishing assertions hold.  Otherwise one could choose
a violating index of that bounded weight at each failing level,
contradicting the definitions of the least depth or the least witness
weight.  More explicitly, let $G_q$ be an ultrafilter-large set on
which the witness, cutoff lower bounds, and minimal-weight vanishing
assertions needed for the $q$th finite coloring hold.  Choose
precisions $M_q\to\infty$ in
$\bigcap_{j\le q}G_j$, and apply finite Ramsey there.  Stabilize the
finitely many colors and take the diagonal construction above.  The
extracted subsequence need not itself be ultrafilter-large: every
finite restriction occurs at arbitrarily high precisions, which is
the required conclusion.  We assert no single level valid at all
weight bounds, and choose no infinite array of primes.

\emph{The bound on singular blocks.}
Proposition~\ref{prop:selmer-matrix-bound} and the uniform unit-symbol
bound in Proposition~\ref{prop:unit-symbol-detection} bound the sum
of the block nullities.  Each singular block contributes at least one.
The type alphabet and its correction cost depend on the fixed support
$S$.  The moving place $\ell$ and any bounded stationary environment
contribute only the bounded additional local cost allowed in
Proposition~\ref{prop:selmer-matrix-bound}; they do not enlarge that
alphabet as the array grows.  This proves~\textup{(iv)} in every
finite realization, and therefore in the profile.

For the final even-symbol assertion, choose fixed templates at
sufficient $2$-adic precision for the witness and its simple variants,
and a fixed sufficiently primary template for the path sites.
Refine each rational residue class to a finite precision determining
a finite value of $v_2(p^2-1)$.  Only finitely many templates occur.
Cross-bit conditions at new odd primes do not change these residues.
The same argument will apply to the root primes used below.
\end{proof}

\subsection{Local calibration and the even case}

The preceding argument determines the symbol values.  We next realize
all its local twist classes simultaneously, while keeping each prime's
activation linear in the cube coordinates.

\begin{lemma}[Calibration by root primes]
\label{lem:sa:root-calibration}
The cubes in Lemma~\ref{lem:sa:punctured-cube} admit realizations by
signed squarefree parameters
\[
 h_x=\prod_p(p^*)^{f_p(x)},\qquad x\in\F_2^b,
 \qquad h_0=1,
\]
where each $f_p:\F_2^b\to\F_2$ is linear.  The prescribed local
classes and diagonal choices hold at every assignment.  Moreover,
$h_x$ may be required to be a local square at any prescribed finite
set of stationary primes.  Calibration adds only root primes, whose
invariant dimension $t_p$ is zero.  In the even case, all added primes
may satisfy a bound on $v_2(p^2-1)$ independent of $b$.
\end{lemma}

\begin{proof}
The address and path programs already give a linear activation form
for each nonroot prime.  Both primes of a simple site have the same
form, so their total simple parity vanishes coefficient by coefficient.
For each coordinate $x_j$, add a fresh root prime $q_j$, active when
$x_j=1$.  Choose its rational radical label at the support to cancel
the coefficient of $x_j$ in the total nonroot label.  The paired-simple
condition makes the residual discriminant obstruction to a rotation
label zero.  Lemma~\ref{lem:lm:arithmetic-realization} and root
substitution in Theorem~\ref{thm:symbol-rules} therefore allow these
choices.

The same roots supply the diagonal choices.  At a nonroot split prime
$p$, the unit twist from the other nonroot primes is a linear function
of $x$.  A desired constant diagonal bit on the active set of $p$ is
represented by that constant times $f_p$.  The correction needed is
therefore linear, say $\sum_j c_{pj}x_j$.  Give $q_j$ the Legendre
bit $c_{pj}$ at $p$.  These bits are independently prescribable for
the different $p$; simple positions need no diagonal correction.
At each stationary prime, choose the root bits to cancel the existing
unit twist coefficient by coefficient.  To check compatibility, let $L$ be the compositum of the rational
quadratic fields recording these finite tests. The fresh-prime radical
classes are independent of the support classes by their valuations.
Moreover, $F\cap L=\Q$ for $G=C_3$; for $G=S_3$ it is at most
the residual discriminant field. Thus agreement on the already matched
discriminant coordinate is the only additional condition for a rotation
in $\Gal(F/\Q)$. Chebotarev supplies the required roots.

Using signed prime discriminants gives the prescribed quadratic
characters at every place, including infinity and~$2$. At each $x$,
the nonroot states and their internal data are those of its table entry.
The product of the active root primes has the support label and the
bits at active split primes of the aggregate-root probe in the
corresponding table entry. Its bits at simple primes do not change
the symbol. Root substitution in Theorem~\ref{thm:symbol-rules}
therefore identifies $A(|h_x|)$ with that profile value. The same
comparison includes $x=0$, where the root collection is empty and all
added primes are absent.  Finally, root support labels lie in a
fixed finite set, so fixed rational $2$-adic residue refinements
determine $v_2(q_j^2-1)$.  The extra conditions, ramified at fresh
odd primes, are compatible with these refinements as in
Lemma~\ref{lem:sa:profile}.
\end{proof}

\begin{proof}[Proof of Theorem~\ref{thm:irreducible-converse} when
$c_2(E)=0$]
By Lemma~\ref{pr:invariance}, the functional equation has
sign~$+1$.  We use the even coefficient source on the filter of~$1$.
To obtain its lower bound, choose an auxiliary twist of analytic rank
one, with opposite infinity sign and matching local classes at $2N$,
by nonvanishing in these fixed local classes
(Proposition~\ref{prop:rough-nonvanishing}).  Exclude its prime
support and the exceptional quadratic fields as in
Proposition~\ref{prop:even-coefficients}.  The reverse genus
comparison and Theorem~\ref{thm:heegner-genus-bound} give the lower
bound required there.  Fix this partner, the exclusions, and the
normalized even symbol before choosing the profile or cube dimension.

If the empty word is a unit, Proposition~\ref{prop:unit-symbol-detection}
gives $L(E,1)\ne0$.  Otherwise
Lemmas~\ref{lem:sa:punctured-cube} and~\ref{lem:sa:root-calibration}
give arbitrarily large cubes with unit symbol at all $x\ne0$.
At each such point,
\[
 L(E^{(h_x)},1)\ne0,
 \qquad
 v(\mathcal L(h_x))-2w(h_x)-s(h_x)\le C_A,
\]
where $C_A$ depends only on the fixed normalized symbol.  The local
classes are fixed and the added-prime valuations $v_2(p^2-1)$ are
bounded independently of $b$.  These are precisely the remaining
hypotheses of Theorem~\ref{thm:cyclotomic-interpolation}, which yields
$L(E,1)\ne0$.
\end{proof}

\subsection{A bounded companion for the odd cube}

Assume $c_2(E)=1$. Parity gives functional sign $-1$, so the address
symbol $A$ is odd and its filter is the filter of $1$. Choose its fixed
genus partner of analytic rank zero, with opposite infinity sign and
matching classes at $2N$, by
Proposition~\ref{prop:rough-nonvanishing}. The irreducible genus bound
supplies the coefficient lower estimate. Choose also an even symbol
$A_0$ on a fixed filter of negative odd fundamental discriminants $k$
for which $2N$ splits in $\Q(\sqrt{k})$. It has no condition at the
moving prime $\ell$. Its complementary genus partner has analytic
rank one and is chosen by the same nonvanishing proposition.
Fix these partners, all excluded prime supports and exceptional fields,
and both coefficient normalizations.

If $A(1)=1$, unit-symbol detection already gives the desired analytic
rank. Otherwise the following construction supplies every arithmetic
hypothesis needed for varying-field Heegner interpolation.

\begin{proposition}[An odd cube with a bounded even companion]
\label{prop:sa:bounded-companion}
For the fixed symbols $A,A_0$ just chosen, suppose $A(1)=0$.
There is a constant $C$ such that, for every sufficiently large $b$,
one can find an odd negative fundamental discriminant $k$ and a finite
family
\[
 h_x=\prod_p(p^*)^{f_p(x)},\qquad x\in\F_2^b,
 \qquad f_p:\F_2^b\longrightarrow\F_2\ \text{linear},
\]
with the following properties.
\begin{enumerate}[label=\textup{(\roman*)}]
\item The primes are distinct good odd primes outside the fixed
exclusions; $h_0=1$, and every $h_x$ has the local squareclasses of
$1$ at the fixed support and at infinity. Each $h_x$ is coprime to $k$.
The field $K=\Q(\sqrt{k})$ satisfies the standing Heegner conditions.
\item The prime support of $k$ has cardinality at most $C$. At every
nonroot prime $p\mid\prod_xh_x$, the parameter $k$ is a local square.
At every nonroot prime $q\mid k$, every $h_x$ is a local square.
\item The simultaneous symbol values are
\[
 A_0(|k|)=1,\qquad
 A(|h_x|)=A_0(|kh_x|)=1\quad(x\ne0).
\]
Consequently $E^{(k)}$ has analytic rank zero and $s(k)\le C$.
\end{enumerate}
For the odd source, these assertions mean simultaneous actual tests at
arbitrarily high admissible precisions after $b$ and its finite set of
tests have been fixed. The constant $C$ is independent of $b$, that
precision, and all moving primes. The parameter $k$ may vary with the
finite realization.
\end{proposition}

The construction uses the even symbol twice: at $k$ to control the
base Selmer group over $K$, and at $kh_x$ to make the product have a
simple zero at every nonzero vertex. We first condition $A_0$ by a
fixed singular prefix. Only after this conditioning will its later
values become determinants.

\begin{lemma}[A determinant-conditioned prefix]
\label{lem:sa:prefix}
For the fixed even symbol $A_0$, there are a finite list $a$ of
singular blocks and an environment consisting of at most one simple
vertex with the following property.  Fix these vertices and their
split diagonal entries.  Let $z$ be any further block of tensor sites
whose alternatives all have $B=L$ to every block of $a$.  Choose
aggregate root data preserving the prefix diagonals and the fixed
total label.  Then
\begin{equation}
 A_0(a,z)=\det M_z.
 \label{eq:sa:conditioned-determinant}
\end{equation}
The left side includes the fixed environment.  Its couplings need
only be compatible with the tensor alternatives.
\end{lemma}

\begin{proof}
If the filter of $A_0$ has odd simple parity, retain one simple
vertex of a unit witness as the environment; otherwise take the
environment empty.  The remaining simple positions have even parity,
so pair them into sites and put all variable positions of the witness
in one block.  Allow aggregate roots with the prescribed total class
throughout.

Consider all actual finite configurations with this environment and
a list of singular blocks, with $B=L$ between blocks, at which
$A_0=1$.  This collection, allowing an empty list, is nonempty.
Indeed, if the witness block is nonsingular, expand it as empty plus
singular states by Lemma~\ref{lem:binary-tensor-span}; some resulting
term has value~$1$.  Proposition~\ref{prop:selmer-matrix-bound}
bounds the lengths of these lists, since the environment has fixed
dimension cost.  Choose an actual maximizing list $a$ and fix all
its primes, singular matrices, and diagonal choices.  Only this one
finite list is retained; no bound on the sizes of arbitrary maximizing
blocks is asserted.

For a proposed further block, fixing the prefix and environment gives
a multilinear symbol functional $f$.  It vanishes on every singular
pure state $z$, for otherwise appending $z$ contradicts maximality.
Also $f(\varnothing)=1$, because its aggregate root preserves the
total label and the prefix diagonal data of the maximizing evaluation.
Singular pure states span the kernel of the determinant functional
by Lemma~\ref{lem:binary-tensor-span}, and the empty state has
determinant~$1$.  Hence $f=\det$, proving
Equation~\eqref{eq:sa:conditioned-determinant}.  Maximality here is
over actual singular states.  It requires no alternative choice for
any fixed prefix prime against a later prime.
\end{proof}

\subsubsection{The number of surviving copies}

Fix a nonempty maximal singular unit word
$W=(W_1,\ldots,W_m)$ in an odd profile, as in
Lemma~\ref{lem:sa:punctured-cube}. Write
\[
 M=\bigoplus_{i=1}^mM_i,\qquad
 D=\bigoplus_{i=1}^mD_i,\qquad N=M+D,
\]
where $D_i$ is a sum of site toggles making $M_i+D_i$ invertible.
The matrix $N$ is invertible and $D$ is alternating. Their common
coordinate space $U$ has even dimension, fixed with the word.

\begin{lemma}[Address copies with even total activation]
\label{lem:sa:balanced-copies}
For every $b\ge2$, the address construction for $W$ may be chosen so
that its symbol is $1$ at every nonzero $x$ and all sites are absent
at zero. At a nonzero $x$, eliminating its invertible path sites
leaves the designated copy of $M$ and an odd number of copies of $N$.
The different copies have zero mutual $M$-entries.
\end{lemma}

\begin{proof}
For every $a\ne0$ choose a linear form $l_a$ with
\begin{equation}
 l_a(a)=1,\qquad \sum_{a\ne0}l_a=0.
 \label{eq:sa:activation-parity}
\end{equation}
To obtain these forms, begin with arbitrary choices satisfying the
first condition. Allowed changes at $a$ form its annihilator.
The annihilators of two distinct nonzero addresses are distinct
hyperplanes and span the entire dual space. Changes at those two
addresses can therefore cancel the initial sum.

Activate all word sites in the copy designated for $a$ by $l_a$, and
use a basis of the annihilator of $a$ for its path coordinates in
Lemma~\ref{lem:lm:address-program}. The saturation proof is unchanged.
At a nonzero $x$, the designated copy leaves $M$. Each other active
copy leaves $N$, and a copy with $l_a(x)=0$ leaves only invertible
paths. Equation~\eqref{eq:sa:activation-parity} says that the total
number of active copies is even. The designated copy is active, so
an odd number of other copies remain.
\end{proof}

\subsubsection{A bounded interface}

We will couple every surviving copy to the same fixed collection of
split sites. When two $N$-copies are eliminated, their contributions
to that collection cancel in characteristic two. Thus the calculation
reduces to one $M$-copy, one $N$-copy, and the fixed collection. The
following lemma chooses its couplings; its dimension depends on $W$
and not on the number of addresses.

\begin{lemma}[Alternating interface repair]
\label{lem:sa:interface}
Let $U$ be an even-dimensional vector space over $\F_2$, let $N$
be symmetric and invertible on $U$, and let $D$ be alternating.
There are at most $\dim U$ linear maps $b_i:U\to\F_2^2$ such that,
on writing
$J=\left(\begin{smallmatrix}0&1\\1&0\end{smallmatrix}\right)$ and
$H=\sum_i b_i^tJb_i$, the matrix
\begin{equation}
 S(H)=\begin{pmatrix}D&N\\N&N+H\end{pmatrix}
 \label{eq:sa:interface-matrix}
\end{equation}
is invertible.
If $U$ has the rotation-field structure $\F_4^d$ occurring in the
cyclic cubic case, the maps may be chosen $\F_4$-linear from
$U$ to the field line $\F_4=\F_2^2$.
\end{lemma}

\begin{proof}
Start with $H=0$, and write $\pi_2:U\oplus U\to U$ for projection
to the second coordinate.  For every $H$, this projection is injective
on $\ker S(H)$: if $(u,0)$ belongs to the kernel, then $Nu=0$,
so $u=0$.  The update $H\mapsto H+b^tJb$ adds the alternating
form $(b\pi_2)^tJ(b\pi_2)$ to $S=S(H)$.  We show how to choose
an update reducing every nonzero kernel.

If $\dim\ker S\ge2$, choose $b$ of rank~$2$ on
$\pi_2(\ker S)$.  Take a complement of $\ker S$ and modify its
basis vectors by kernel vectors so that $b\pi_2$ vanishes on the
new complement.  The restriction of $S$ to that complement remains
nondegenerate.  The update is zero there and in its cross pairings,
and has rank~$2$ on $\ker S$.  Hence it decreases the nullity by
exactly~$2$.

Suppose instead that $\ker S=\langle k\rangle$ is one-dimensional.
Let $d_S$ denote the diagonal vector.  Since $z^tSz=d_S^tz$ for
every binary symmetric matrix, $d_S$ annihilates the kernel and lies
in the image of $S$.  Choose $s$ with $Ss=d_S$.  We shall use
\begin{equation}
 d_S^ts=\rank S\pmod2.
 \label{eq:sa:diagonal-rank}
\end{equation}
To prove it, restrict to a nondegenerate complement of the radical,
with invertible symmetric matrix $A$ and diagonal vector $d_A$.
The left side is $d_A^tA^{-1}d_A$.  Off-diagonal terms cancel in
pairs, leaving $\sum_i(A^{-1})_{ii}A_{ii}$.  Applying the same
cancellation to $\tr(A^{-1}A)$ identifies this with $\dim A$
modulo~$2$, proving Equation~\eqref{eq:sa:diagonal-rank}.

Since the ambient dimension is even and the kernel has dimension one,
$d_S^ts=1$.  The alternating matrices $D$ and $H$ have zero
diagonal, so $d_S$ is supported on the second coordinates.
Consequently $\pi_2(s)$ and $\pi_2(k)$ are binary-independent:
otherwise the former would be zero or equal to the latter, forcing
$d_S^ts=0$.  Choose $b$ sending them to independent vectors of
$\F_2^2$.

Let $z$ belong to the kernel after this update.  Self-pairing gives
$d_S^tz=0$.  Pairing the updated equation with $k$ gives
\[
 (b\pi_2(k))^tJ(b\pi_2(z))=0.
\]
Pairing with $s$ gives the same equation with $s$ in place of $k$,
since $s^tSz=d_S^tz=0$.  Their independent images force
$b\pi_2(z)=0$.  The updated equation then reduces to $Sz=0$, so
$z$ is a multiple of $k$.  As $b\pi_2(k)\ne0$, this multiple is
zero.  The update has made $S$ invertible.

Both choices of $b$ remain possible under the field-linear restriction.
Two binary-independent vectors $u,v\in\F_4^d$ are either
field-independent, in which case a field-linear map can send them to
any binary basis of $\F_4$, or satisfy $v=\alpha u$ with
$\alpha\in\F_4\setminus\F_2$.  In the latter case any field-linear
map nonzero on $u$ gives binary-independent images.  These maps are
allowed by the rotation-field off-diagonal spans, and the preceding
argument applies to the underlying binary forms.

Finally, injectivity of $\pi_2$ bounds every kernel dimension by
$\dim U$.  Each update strictly decreases it, so at most $\dim U$
maps are required.
\end{proof}

\begin{proof}[Proof of Proposition~\ref{prop:sa:bounded-companion}]
Choose the prefix and possible one-coordinate environment from
Lemma~\ref{lem:sa:prefix} and fix all their primes. Build the odd
address profile disjoint from them, with $B=L$ to every prefix block.
The alternatives at a simple pair have identical compatible outside
bits against the environment. Include these stationary primes among
the local-square conditions in Lemma~\ref{lem:sa:root-calibration}.

Since $A(1)=0$, Lemma~\ref{lem:sa:punctured-cube} supplies a nonempty
maximal singular unit word $W$. Fix it once. It determines $U,M,D,N$
before $b$ is chosen.

Choose the maps of Lemma~\ref{lem:sa:interface} for this fixed pair
$(N,D)$.  Let $J_{\mathrm{int}}$ be the direct sum of one standard
$J$-block for each map, and let $G$ be the matrix obtained by stacking
the maps $Jb_i$.  Then
\[
 G^tJ_{\mathrm{int}}^{-1}G=\sum_i b_i^tJb_i=H.
\]
For a given $b$, use
Lemma~\ref{lem:sa:balanced-copies} to obtain the address programs and
then calibrate their support labels and diagonals by root primes.

Lemma~\ref{lem:lm:arithmetic-realization} realizes these couplings by
a fixed number of primary split interface positions, with the
field-linear restriction when required.

Only now, for each actual finite address array, construct the auxiliary
parameter $k$.  Its support consists of the stationary prefix, its
possible environment, fresh primes for the fixed number of interface
positions, and one fresh root prime.  Prescribe $B=L$ between the
interface and the prefix blocks.  In the matrix $M$ of the further
block, couple the interface by zero to the path sites and by the same
matrix $G$ to the word coordinates of every address copy.  Keep all
other bits of the simple variants identical.

Require $h_x$ to be a local square at every interface prime for every
$x$.  This is compatible with the matrix entries: the rational norm
bits against address primes are independently prescribable, so their
coefficients against the activation forms can be made zero.  Include
the already chosen address root primes in these conditions; they need
only rational norm bits. Treat them as additional residue places in
Lemma~\ref{lem:lm:arithmetic-realization}. The independent matrix and
norm prescriptions concern the nonroot primes, and the further norm
conditions occur at disjoint root primes. Their quadratic radicals,
ramified at those fresh primes, create no relation with the preceding
extensions, which are unramified there. Root calibration has already made $h_x$ a local
square at the prefix and environment primes.

Choose the last root prime of $k$ to satisfy four requirements: restore
the fixed total support class; make $k$ negative and split at $2N$;
make $k$ a local square at every nonroot address prime; and give the
prescribed diagonals at the prefix and interface.  These are
simultaneous finite rational radical conditions.  Their residual
discriminant coordinate is compatible with a rotation, because the
environment accounts for any odd simple parity of the filter and the
remaining variable simple positions are paired.  The arithmetic
realization lemma therefore supplies the prime.  All new primes can
be taken fresh, so $k$ is squarefree, coprime to every $h_x$, and
avoids the excluded fields.  Address root primes need no splitting
condition, since their invariant dimension is zero.

These choices preserve all matrices needed for the two symbols.
At every actual assignment, the address matrices for $A_0$ equal
those for $A$, since the extra twist by $k$ is a local unit square
at every nonroot address prime.  The prefix and interface matrices
also stay fixed as $x$ varies.  In the tensor evaluations used in
Equation~\eqref{eq:sa:conditioned-determinant}, aggregate roots
preserve these same diagonal choices.

At zero, the further block is just the interface and has invertible
matrix $J_{\mathrm{int}}$.  Equation~\eqref{eq:sa:conditioned-determinant}
therefore gives
\begin{equation}
 A_0(|k|)=1.
 \label{eq:sa:auxiliary-unit}
\end{equation}
At a nonzero address, eliminate the invertible paths.  There remains
one copy of $M$ and an odd number of copies of $N$, each coupled
to the interface by $G$.  Eliminate the $N$-copies in pairs:
each pair has invertible matrix $N\oplus N$ and contributes
$GN^{-1}G^t+GN^{-1}G^t=0$ to the interface.  Retain the one
unpaired copy.  Eliminating the interface then reduces invertibility
of the further block to invertibility of
\[
 \begin{pmatrix}M+H&H\\H&N+H\end{pmatrix}.
\]
The substitution $(u,v)=(u',u'+v')$ transforms this matrix by
congruence to
\[
 \begin{pmatrix}M+N&N\\N&N+H\end{pmatrix}
 =S(H),
\]
which is invertible by Lemma~\ref{lem:sa:interface}.  A second
application of Equation~\eqref{eq:sa:conditioned-determinant} yields
\begin{equation}
 A_0(|k h_x|)=1\quad\text{for every }x\ne0.
 \label{eq:sa:paired-units}
\end{equation}
The odd symbol is also $1$ by saturation. The support of $k$ consists
of the fixed prefix and environment, at most $\dim U$ interface
primes, and one root prime. Its cardinality is bounded independently
of $b$. The symbol $A_0(|k|)=1$ gives analytic rank zero and a uniform
bound for $s(k)$ by Proposition~\ref{prop:unit-symbol-detection}.

For each fixed finite profile restriction, first fix its weight bound
and the finite list of symbol and operator comparisons. Choose a
sufficiently high actual odd-source level, and realize the address
array there as above. Keeping all these address primes fixed, place
the required Galois data in a common finite quotient and choose only
the fresh interface primes and final root prime of $k$ to meet the
remaining matrix and norm prescriptions. All the unit-symbol tests
then hold at once. All constants were fixed
with the symbols, prefix and word. This proves the proposition.
\end{proof}

\subsection{Completion of the odd case}

\begin{proof}[Completion of the proof of
Theorem~\ref{thm:irreducible-converse}]
Apply Proposition~\ref{prop:sa:bounded-companion}. It gives
$\omega(k)$ and $s(k)$ bounded independently of $b$, with
$E^{(k)}$ of analytic rank zero. For $K=\Q(\sqrt{k})$,
Lemma~\ref{lem:h:base-selmer-bound} gives
\[
 \corank_{\Ztwo}\Sel_{2^\infty}(E/K)=1
\]
and bounds the length of its finite quotient by the divisible part
uniformly in $k$.  Its hypotheses hold because $E$ is the fixed
corank-one curve, the other twist has rank zero with bounded $s(k)$,
and $\omega(k)$ is bounded.  This step does not assume finiteness
of the Tate--Shafarevich group of $E$.

For $x\ne0$, the two unit symbols give analytic rank one for
$E^{(h_x)}$ and analytic rank zero for $E^{(kh_x)}$.  Their
coefficient comparisons also give uniform bounds
\[
 j(P_1(h_*))\le w(h_x)+C_1,
 \qquad
 v(\mathcal L(kh_x))\le 2w(kh_x)+C_2.
\]
Here $h_*$ is the fixed genus partner of the odd symbol. The genus point
$P_1(h_*)$ is formed over $\Q(\sqrt{h_xh_*})$, so the point varies with
$x$ while $h_*$ remains fixed. Coprimality
gives $w(kh_x)=w(k)+w(h_x)$. Proposition~\ref{prop:heegner-comparison}
first yields $2j(P_{|h_x|}(h_x))\le4w(h_x)+2w(k)+O(1)$.
Since $w(k)\le\omega(k)$ is uniformly bounded, this becomes
\begin{equation}
 2j(P_{|h_x|}(h_x))\le 4w(h_x)+C_3.
 \label{eq:sa:heegner-upper}
\end{equation}
Subtracting the nonnegative finite lengths $s(h_x)+s(kh_x)$ gives
the uniform upper bound in Theorem~\ref{thm:heegner-interpolation}.

We have also imposed every other hypothesis of that theorem:
$K$ is imaginary and split at $2N$; each $h_x$ is coprime to $k$;
$h_0=1$; signed prime discriminants are activated by linear forms;
the local classes at $2N$ are fixed; and each added prime either
splits in $K$ or is a root prime with $t_p=0$.
For each fixed finite cube, take a sufficiently high precision from
its realization set so that all its unit-symbol tests hold at once.
The constants depend only on the already fixed normalized symbols,
prefix, and word, not on $b$.

The varying-field uniformity in
Theorem~\ref{thm:heegner-interpolation} therefore applies to these
cubes of arbitrarily large dimension.  It gives a simple zero at~$1$
for $L(E,s)L(E^{(k)},s)$.  The second factor is nonzero by
Equation~\eqref{eq:sa:auxiliary-unit}, so $L(E,s)$ has a simple
zero.  This proves the remaining case $c_2(E)=1$.
\end{proof}

\section{Proof of auxiliary nonvanishing}
\label{sec:auxiliary-nonvanishing}

We prove Proposition~\ref{prop:rough-nonvanishing}. The even-sign
assertion follows from Hoffstein--Luo after one preliminary twist.
For the odd-sign assertion we retain the prescribed local conditions
throughout a derivative first-moment calculation.

\begin{proof}[Proof of Proposition~\ref{prop:rough-nonvanishing}\textup{(i)}]
An odd fundamental discriminant is squarefree and congruent to $1$
modulo $4$, so $c_2$ is represented by $u_2\in\{1,5\}$ modulo $8$.
For odd $p\in S$, choose a nonzero residue $u_p$ representing $c_p$.
The Chinese remainder theorem and Dirichlet's theorem give a prime
$q\notin S\cup B$ satisfying
\[
 q\equiv\sigma u_2\pmod8,\qquad
 \sigma q\equiv u_p\pmod p\quad(p\in S,\ p\ne2).
\]
Then $D_0=\sigma q$ is an odd fundamental discriminant in the prescribed
filter. Put $A=E^{(D_0)}$ and let $N_A$ be its conductor. Thus
$\varepsilon(A)=1$.

The theorem of Hoffstein--Luo
\citep[pp.~435--436]{HoffsteinLuo1997} gives infinitely many odd
fundamental discriminants $h$ such that
\[
 L(A^{(h)},1)\ne0,\qquad \omega(|h|)\le4,
\]
and $\chi_h(p)=1$ at every prime in the fixed finite set consisting of
$S\cup B\cup\{q\}$ and the prime divisors of $N_A$.
The twisting formula for the root number gives
\[
 1=\varepsilon(A^{(h)})=\varepsilon(A)\chi_h(-N_A)
 =\operatorname{sign}(h),
\]
so $h>0$. For odd $p$ the condition $\chi_h(p)=1$ says that $h$ is a
unit square in $\Q_p$. At $2$, odd fundamentality and $\chi_h(2)=1$
give $h\equiv1\pmod8$, which has the same meaning.
Since $q\nmid h$, the product $D=D_0h$ is squarefree and congruent to
$1$ modulo $4$. It retains the prescribed sign and squareclasses and
avoids $B$. Finally,
\[
 E^{(D)}\simeq A^{(h)},\qquad \omega(|D|)=1+\omega(h)\le5.
\]
Distinct choices of $h$ give distinct $D$.
\end{proof}

\subsection{An odd-sign witness in a fixed progression}
\label{odd:fixed-progression}

We retain the modulus of the local conditions throughout the first-moment
calculation.  Iwaniec's derivative moment
\citep[Theorem on p.~367 and Sections~6--10]{Iwaniec1990} treats the
negative-discriminant split family.  The following argument allows any
fixed admissible progression.  It makes no uniformity assertion as that
progression varies.

\begin{lemma}[Odd first moment]\label{nv:odd-first-moment}
Let $f$ be the weight-two newform of an elliptic curve of conductor $N$,
with its $L$-function normalized to have center $1/2$.
Fix $H$ divisible by $8$ and $N$, a reduced class $a\pmod H$, and
$\sigma\in\{1,-1\}$ with $\sigma a\equiv1\pmod4$.
Suppose the twists with $n>0$ squarefree and $n\equiv a\pmod H$
have root number $-1$.
For every nonnegative, nonzero $\Phi\in C_c^\infty((1,2))$, there are
$c>0$, $c_1\in\mathbb R$, and an absolute $\eta>0$ such that
\begin{equation}\label{odd:moment}
 \sum_{\substack{n>0\ \mathrm{squarefree}\\n\equiv a\ (H)}}
 L'(1/2,f\otimes\chi_{\sigma n})\Phi(n/X)
 =cX\log X+c_1X+O_{f,H,a,\sigma,\Phi}(X^{1-\eta}).
\end{equation}
Consequently this progression contains arbitrarily large $n$ whose
twists have a simple central zero.
\end{lemma}

\begin{proof}
Write $L(s,f)=\sum_{m\ge1}\lambda(m)m^{-s}$, so
$|\lambda(m)|\le\tau(m)$, where $\tau$ is the divisor function, and put $Q=\sqrt N/(2\pi)$.
In estimates below, constants may depend on the fixed data in the Lemma;
all unspecified power exponents are absolute.
The proof separates two issues: preserving the fixed progression while
removing the squarefree condition, and controlling the frequencies after
Poisson summation. The zero frequency will give a positive main term of
order $X\log X$; the other frequencies and the discarded square divisors
must together have a power-saving bound. We keep the modulus $H$ fixed
throughout these steps.
Set
\[
 V(y)=\frac1{2\pi i}\int_{(2)}
 e^{z^2}\Gamma(1+z)y^{-z}\frac{dz}{z^2}.
\]
The conductor of the twist is $Nn^2$.  Its odd functional equation gives
\begin{equation}\label{odd:afe}
 L'(1/2,f\otimes\chi_{\sigma n})
 =2\sum_{m\ge1}\frac{\lambda(m)\chi_{\sigma n}(m)}{\sqrt m}
           V\left(\frac m{Qn}\right).
\end{equation}
Indeed, shift the defining completed integral across its double pole
and use the functional equation.  The central value is zero, so the
derivative of the completion contributes no additional term.
Use the right side of Equation~\eqref{odd:afe} to define $\mathcal A(n)$
also for nonsquarefree $n\equiv a\pmod H$.

\paragraph{Large square divisors.}
We first record the precise large-sieve input.  For every $\rho>0$ and
fixed $0<\delta<1/4$,
\begin{equation}\label{odd:second-moment}
 \sum_{\substack{0<t\le Z\ \mathrm{squarefree}\\
             \sigma t\ \mathrm{fundamental},\ (t,2N)=1}}
 |L(1/2+\delta+i v,f\otimes\chi_{\sigma t})|^2
 \ll_{f,\rho,\delta} Z^{1+\rho}(1+|v|)^{C_{\rho,\delta}}.
\end{equation}
To justify this consequence of the quadratic large sieve, its form in
\citep[Theorem~1]{HeathBrown1995} is
\[
 \sum_{\substack{t\le Z\ \mathrm{odd, squarefree}}}
 \left|\sum_{\substack{u\le U\ \mathrm{odd, squarefree}}}
 b_u\left(\frac tu\right)\right|^2
 \ll_\epsilon (ZU)^\epsilon(Z+U)\sum_u|b_u|^2.
\]
Quadratic reciprocity permits either orientation of the symbol after
separating classes modulo $4$.  In a Dirichlet polynomial with coefficients
$\lambda(m)m^{-1/2}$, write $m=ur^2$, $u$ squarefree.  For each $r$,
$\chi_{\sigma t}(r^2)$ restricts $t$ to $(t,r)=1$; discard that
restriction only after taking absolute squares.  Since
$|\lambda(ur^2)|\le\tau(u)\tau(r^2)$, Minkowski's inequality bounds the
square root of the moment by
\begin{equation}\label{odd:square-index-bound}
 (ZU)^{O(\epsilon)}
 \sum_{r\le\sqrt U}\frac{\sqrt{Z+U/r^2}}r
 \ll (ZU)^{O(\epsilon)}(\sqrt Z+\sqrt U).
\end{equation}
The powers of primes dividing $2N$ contribute convergent geometric sums.
On a dyadic interval $t\asymp T$, the approximate functional equation
has length $U\ll_f T(1+|v|)^C$, apart from rapidly decreasing tails.
The main sum's additional factor $m^{-\delta}$ improves partial
summation.  The dual sum loses at most $U^\delta$ in coefficients, but
its functional-equation multiplier contributes $T^{-2\delta}$, up to
a fixed power of $1+|v|$.  Thus it incurs no positive power of $T$.
Mellin inversion separates the smooth conductor weights; the gamma
factors give integrable decay in the Mellin variables.  Applying
Equation~\eqref{odd:square-index-bound} with sufficiently small
$\epsilon$, and summing dyadic intervals, proves
Equation~\eqref{odd:second-moment}.

Expand $\mu^2(n)=\sum_{d^2\mid n}\mu(d)$ and retain $d\le Y$, where
$1\le Y\le\sqrt X$.  The omitted contribution is
\begin{equation}\label{odd:square-tail}
 O_\rho(X^{1+\rho}/Y).
\end{equation}
Here are the details for the nonsquarefree terms.  Write $n=tv^2$ with
$t$ squarefree.  An omitted term has $v>Y$, with at most $\tau(v)$ choices
of $d$.  Also $\sigma t$ is fundamental, $(tv,2N)=1$, and
\[
 \chi_{\sigma n}(m)
   =\chi_{\sigma t}(m)\mathbf1_{(m,v)=1}.
\]
On a Mellin line $\Re z=\delta>0$, the series defining $\mathcal A(tv^2)$
therefore uses
\[
 L(1/2+z,f\otimes\chi_{\sigma t})
 \prod_{p\mid v}
 \left(1-\frac{\lambda(p)\chi_{\sigma t}(p)}{p^{1/2+z}}
             +\frac{\chi_{\sigma t}(p)^2}{p^{1+2z}}\right).
\]
The product is $O_\epsilon(v^\epsilon)$.  By Cauchy's inequality and
Equation~\eqref{odd:second-moment}, the sum of the absolute $L$-values
over $t\le2X/v^2$ is
$O((X/v^2)^{1+\epsilon}(1+|\Im z|)^C)$.
The remaining Mellin factors are bounded by $O(X^\delta)$ times a
rapidly decreasing function of $\Im z$.  Choose $\delta$ and $\epsilon$
sufficiently small in terms of $\rho$.  Absorbing the divisor count and
these small powers gives
$X^{1+\rho}\sum_{v>Y}v^{-2}\ll X^{1+\rho}/Y$, as claimed.
The double pole in the kernel is harmless on this fixed positive line.

\paragraph{Poisson summation.}
For $(d,H)=1$ put $J=d^2$ and $c_J\equiv aJ^{-1}\pmod H$.
Let $T(J)$ be the sum of $\mathcal A(n)\Phi(n/X)$ over
$n\equiv a\pmod H$ with $J\mid n$, without a squarefree restriction.
Writing $m=m_1q$ with $m_1\mid H^\infty$ and $(q,H)=1$, put
$\kappa(m_1)=\chi_{\sigma a}(m_1)$ and
\[
 G_\nu(q)=\sum_{y\bmod q}\left(\frac yq\right)e(\nu y/q),\qquad
 W(z,R)=\int_0^\infty\Phi(x)V\left(\frac z{Qx}\right)e(-Rx)\,dx,
\]
where $e(x)=e^{2\pi ix}$. Poisson summation gives
\begin{equation}\label{odd:poisson}
 \begin{split}
 T(J)={}&\frac{2X}{HJ}
 \sum_{m_1\mid H^\infty}\frac{\lambda(m_1)\kappa(m_1)}{\sqrt{m_1}}
 \sum_{(q,H)=1}\frac{\lambda(q)}{q^{3/2}}
                  \left(\frac{\sigma JH}{q}\right)\\
 &\hspace{9mm}\cdot\sum_{\nu\in\mathbb Z}
 G_\nu(q)e(\nu c_J\bar q/H)
 W\left(\frac{m_1q}{X},\frac{\nu X}{HJq}\right),
 \end{split}
\end{equation}
where $q\bar q\equiv1\pmod H$ and $G_\nu(1)=1$.
For clarity, the Chinese remainder theorem evaluates the finite sum as
\[
 \sum_{\substack{x\bmod Hq\\x\equiv c_J\ (H)}}
 \left(\frac xq\right)e(\nu x/(Hq))
 =e(\nu c_J\bar q/H)\left(\frac Hq\right)G_\nu(q).
\]
Multiplication by $(\sigma J/q)$ yields Equation~\eqref{odd:poisson};
this factor vanishes when $(q,J)>1$.

\paragraph{The zero frequency.}
The Gauss sum $G_0(q)$ is zero unless $q$ is a square, when it is
$\varphi(q)$.  Define
\[
 H_p(z)=1+(1-p^{-1})\sum_{r\ge1}
                    \lambda(p^{2r})p^{-r(1+2z)},\qquad
 F_J(z)=\prod_{p\nmid HJ}H_p(z),
\]
\[
 L_H(z)=\sum_{m_1\mid H^\infty}
          \lambda(m_1)\kappa(m_1)m_1^{-1/2-z},\qquad
 \widehat\Phi_0(z)=\int_0^\infty\Phi(x)x^z\,dx.
\]
Mellin inversion expresses the zero-frequency contribution as
\begin{equation}\label{odd:zero-integral}
 T_0(J)=\frac{2X}{HJ}\frac1{2\pi i}\int_{(2)}
 e^{z^2}\Gamma(1+z)(QX)^z\widehat\Phi_0(z)L_H(z)F_J(z)
 \frac{dz}{z^2}.
\end{equation}
At a good prime, let $\alpha_p\beta_p=1$ and
$|\alpha_p|=|\beta_p|=1$.  The identity
\[
 \sum_{r\ge0}\lambda(p^{2r})x^r
 =\frac{1+x}{(1-\alpha_p^2x)(1-\beta_p^2x)}
\]
shows that $H_p(z)$ is the symmetric-square Euler factor at $1+2z$
times a polynomial
$1+O(p^{-2-4\Re z}+p^{-2-2\Re z})$.
The correction product converges absolutely near $\Re z\ge-1/16$.
To omit a prime dividing $J$, omit its correction factor and insert an
inverse symmetric-square Euler polynomial; this introduces no poles at
zeros of $H_p$.  These omitted factors cost at most $J^B$ on that strip.
The finite Euler product $L_H$ is holomorphic there.
Symmetric-square continuation, its functional equation, and polynomial
strip bounds now allow the shift to $\Re z=-1/16$, with error
\begin{equation}\label{odd:zero-error}
 T_0(J)=R_J(X)+O(J^B X^{1-1/16}),
\end{equation}
where $R_J(X)$ is the residue of the integrand in
Equation~\eqref{odd:zero-integral}, including its factor $2X/(HJ)$.
For this analytic input see
\citep[Theorem~9.3 and Remark~9.9]{GelbartJacquet1978}; in the CM case
the untwisted symmetric square is a product of nontrivial Dirichlet and
Hecke-character $L$-functions, so the same continuation holds near $1$.

The sum of the residues, with $d$ extended to infinity, has coefficient
function
\[
 M(z)=e^{z^2}\Gamma(1+z)Q^z\widehat\Phi_0(z)L_H(z)
              \prod_{p\nmid H}\bigl(H_p(z)-p^{-2}\bigr).
\]
The product is interpreted by its symmetric-square factorization and
an absolutely convergent correction.  It is holomorphic near zero.
Moreover $M(0)>0$.  Indeed
\[
 F_p^{\rm ev}:=\frac{1+p^{-1}}
 {(1-\alpha_p^2/p)(1-\beta_p^2/p)}>0,\qquad
 H_p(0)=p^{-1}+(1-p^{-1})F_p^{\rm ev}>p^{-1}.
\]
All the fixed local factors in $L_H(0)$ are positive, and
$L(1,\operatorname{sym}^2 f)>0$ by the Rankin--Selberg residue formula.
The remaining correction factors are positive and form a convergent
product.  The full residue is therefore
\begin{equation}\label{odd:residue}
 \frac{2X}{H}\bigl(M(0)\log X+M'(0)\bigr),
 \qquad c=\frac{2M(0)}H>0.
\end{equation}
Extending the residue sum costs $O_\rho(X^{1+\rho}/Y)$:
its terms contain $d^{-2}\prod_{p\mid d}H_p(0)^{-1}$, and this
product, as well as its differentiated version, is
$O_\epsilon(d^\epsilon(1+\log d))$.
The resulting $d^{-2+\epsilon}$ tail also justifies differentiating
the residue sum.

\paragraph{The nonzero frequencies.}
The zero frequency has now supplied the positive main term. We next
bound the remaining frequencies uniformly in the square divisor $J$
by a fixed power of $J$, so that summation over the retained square
divisors will preserve a power saving in $X$.
We prove the bound
\begin{equation}\label{odd:offdiagonal}
 T_{\ne0}(J)\ll J^B X^{7/8}.
\end{equation}
For odd $q$, set $\epsilon_q=1$ or $i$ according as $q\equiv1$ or
$3\pmod4$, and $\widetilde G_\nu(q)=\epsilon_q^{-1}G_\nu(q)$.
Quadratic reciprocity makes $\widetilde G_\nu$ multiplicative in $q$.
As $4\mid H$, the phase $\epsilon_q e(\nu c_J\bar q/H)$ is a
function on $(\mathbb Z/H\mathbb Z)^\times$.  Expand it in Dirichlet
characters $\psi\pmod H$; every expansion coefficient has absolute
value at most one.  For $\nu\ne0$, consider
\[
 D_{\psi,J,\nu}(s)=\sum_{(q,H)=1}
 \frac{\lambda(q)}{q^s}\frac{\widetilde G_\nu(q)}{\sqrt q}
                 \left(\frac{\sigma JH}{q}\right)\psi(q).
\]
At $p\nmid2HJ\nu$, its local factor is
$1+\lambda(p)\xi(p)p^{-s}$, where
$\xi(q)=\psi(q)(\sigma JH\nu/q)$ is a Dirichlet character.
Multiplication by the inverse $\mathrm{GL}_2$ Euler factor leaves
\[
 1+(1-\lambda(p)^2)\xi(p)^2p^{-2s}
     +\lambda(p)\xi(p)^3p^{-3s}.
\]
These factors converge absolutely for $\Re s>1/2$.

The exceptional primes cause only a fixed power loss.  For an odd prime
$p$, put $v=v_p(\nu)$.  Direct summation over residue lifts gives
\[
 G_\nu(p^r)=
 \begin{cases}
 \varphi(p^r),&r\text{ even},\ r\le v,\\
 -p^{r-1},&r\text{ even},\ r=v+1,\\
 \epsilon_p p^{r-1/2}
       \left(\dfrac{\nu/p^{r-1}}p\right),
          &r\text{ odd},\ r=v+1,\\
 0,&\text{otherwise}.
 \end{cases}
\]
For even $r$ this is the Ramanujan sum.  For odd $r$, the sum of the
lifts vanishes unless $p^{r-1}\mid\nu$, and the primitive quadratic
Gauss sum then gives the formula.  Thus the local sum is a polynomial;
on $\Re s\ge3/4$ its absolute value is at most
$\sum_{r\ge0}(r+1)p^{-r/4}\le C_0$.
The inverse Euler polynomial is also bounded absolutely.  Primes dividing
$J$ remove the local sum altogether.  Passing to the primitive character
underlying $\xi$ introduces only such inverse Euler polynomials.  Hence
\[
 D_{\psi,J,\nu}(s)=L(s,f\otimes\xi^{\rm prim})C_{\psi,J,\nu}(s),
\]
where $C_{\psi,J,\nu}$ is holomorphic for $\Re s>1/2$ and bounded
on $\Re s\ge3/4$ by $(2J|\nu|)^C$.  Indeed the exceptional-prime
cost is at most $C_1^{\omega(J\nu)}$, and the other corrections have
summable $O(p^{-3/2})$ errors there.
The conductor of $\xi^{\rm prim}$ is at most $4HJ|\nu|$.
The twisted cusp-form $L$-function is entire; its functional equation
and polynomial strip bounds therefore give
\begin{equation}\label{odd:gauss-strip}
 |D_{\psi,J,\nu}(3/4+i t)|
 \ll_{f,H}(J|\nu|)^C(1+|t|)^C.
\end{equation}
The exponents are absolute because the degree and weight are fixed.

It remains to sum these bounds without losing the power saving.
The kernel and integration by parts in $x$ give, for fixed $r,s$ and
every $K>0$,
\begin{equation}\label{odd:weight-decay}
 |(z\partial_z)^r(R\partial_R)^s W(z,R)|
 \ll_{r,s,K}(1+z)^{-K}(1+|R|)^{-K}(1+|\log z|).
\end{equation}
The logarithm accounts for the double pole when $z$ tends to zero.
On the actual sums $z\ge X^{-1}$.
For a fixed small $\rho>0$, insert smooth cutoffs
$z\ll X^\rho$ and $|R|\ll X^\rho$.
The discarded part is negligible: use $|G_\nu(q)|\le q$ and
\[
 \sum_{\nu\in\mathbb Z}
 (1+|\nu|X/(HJq))^{-K}\ll_H1+Jq/X
\]
to obtain a polynomial absolute majorant in $X,J$ for
Equation~\eqref{odd:poisson}.  More explicitly, its absolute value is
$O_\epsilon(X^{3/2+\epsilon}\log X)$, uniformly for $J\ge1$:
the $q$-sums contribute
$(X/m_1)^{1/2+\epsilon}+(J/X)(X/m_1)^{3/2+\epsilon}$,
and the remaining $m_1\mid H^\infty$ sums converge.
For $z>1$, absorb the positive logarithm into an arbitrarily small
power of $z$ in this calculation.  On either discarded tail, increasing
$K$ in Equation~\eqref{odd:weight-decay} supplies an arbitrarily large
factor $X^{-\rho K}$.
The retained nonzero-frequency range satisfies
\begin{equation}\label{odd:ranges}
 m_1q\ll X^{1+\rho},\qquad
 0<|\nu|\ll_H JX^{2\rho},\qquad
 q\gg_H X^{1-\rho}/J.
\end{equation}

Partition the $q$-sum into smooth segments $q\asymp T$.
By Equation~\eqref{odd:weight-decay}, the finitely many derivative
norms needed for Mellin inversion of the segment's weight are
$O(X^{C'\rho})$, uniformly in $m_1,J,\nu,T$.
Mellin inversion uses $D_{\psi,J,\nu}(1+s)T^s$.
Shift its contour to $\Re s=-1/4$; the factorization above shows that
no pole is crossed.  Equation~\eqref{odd:gauss-strip} and the vertical
decay of the weight's Mellin transform give
\begin{equation}\label{odd:dyadic}
 \sum_q\frac{\lambda(q)}q\frac{\widetilde G_\nu(q)}{\sqrt q}
 \left(\frac{\sigma JH}{q}\right)\psi(q)W_*(q/T)
 \ll (J|\nu|)^C X^{C'\rho}T^{-1/4}.
\end{equation}
Here $W_*$ includes the dyadic and cutoff factors and both arguments
of the original weight; their logarithmic derivatives are controlled
by Equation~\eqref{odd:weight-decay}.
The sum over $m_1$ is bounded since
\[
 \sum_{m_1\mid H^\infty}\frac{|\lambda(m_1)|}{\sqrt{m_1}}
 \le\prod_{p\mid H}(1-p^{-1/2})^{-2}<\infty.
\]
Equation~\eqref{odd:ranges} bounds the dyadic sum of $T^{-1/4}$ by
$O_H(J^{1/4}X^{-1/4+\rho/4})$, even if its formal lower endpoint
is below one.  Also
\[
 \sum_{0<|\nu|\ll_H JX^{2\rho}}(J|\nu|)^C
 \ll_H J^{2C+1}X^{2(C+1)\rho}.
\]
Multiplying these estimates by $X/J$ from Equation~\eqref{odd:poisson}
proves $T_{\ne0}(J)\ll J^B X^{3/4+C''\rho}$.
Choose $\rho$ absolutely small enough that $C''\rho\le1/8$.
This proves Equation~\eqref{odd:offdiagonal}.

\paragraph{Conclusion.}
Sum the zero- and nonzero-frequency errors over $d\le Y$, with $J=d^2$.
Together with the large-square and residue tails, the total error is
\[
 O\left(Y^{2B+1}X^{1-\kappa}+X^{1+\rho}/Y\right),
 \qquad \kappa=1/16,
\]
after enlarging $B$.  Choose an absolute $\delta_0>0$ with
$(2B+1)\delta_0<\kappa/2$, put $Y=X^{\delta_0}$, and choose the
tail exponent $0<\rho<\delta_0/4$.
Both errors are $O(X^{1-\eta})$ for an absolute $\eta>0$.
Equation~\eqref{odd:residue} proves Equation~\eqref{odd:moment}.
Its leading coefficient is positive, so some derivative is nonzero in
every sufficiently large dyadic interval.  The odd functional equation
already forces the central value to vanish, proving the final assertion.
\end{proof}

\begin{proof}[Proof of Proposition~\ref{prop:rough-nonvanishing}\textup{(ii)}]
Choose $H$ divisible by $8N$ and by all primes in $S\cup B$.
The prescribed squareclasses and sign determine a reduced class
$a\pmod H$ with $\sigma a\equiv1\pmod4$, by the Chinese remainder
theorem. At a prime in $B\setminus S$, choose any nonzero residue.
Then every positive squarefree $n\equiv a\pmod H$ gives an odd
fundamental discriminant $D=\sigma n$ in the prescribed filter and
prime to every member of $B$. Lemma~\ref{nv:odd-first-moment} supplies
arbitrarily large such $n$ with $L'(E^{(D)},1)\ne0$.
\end{proof}

\section{Transfer on a binary cube}\label{sec:binary-algebra}
\label{sec:h:binary-algebra}

Both analytic arguments will produce one determinant coordinate at each
vertex of a binary cube. The following algebra transfers a bounded
valuation at every nonzero vertex to the zero vertex. Its essential
feature is that the required dimension depends on a bounded number of
coefficients, even when the family involves arbitrarily many primes.

Put $\Lambda=\Ztwo[[t]]$ and
$O=\widehat{\Lambda_{(2)}}$. Let $G\simeq(\Z/2\Z)^b$ and
$R=\Lambda[G]$. Index the sign homomorphisms $G\to\{\pm1\}$
by $x\in\F_2^b$, with $x=0$ the augmentation. For an element or
matrix over $O[G]$, a subscript $x$ denotes evaluation at this character.
The proofs first control Laurent coefficients, then clear determinant
denominators, and finally use the parity of the number of solutions of
low-degree equations on the cube.

\begin{lemma}[Laurent expansions]\label{lem:h:laurent-expansions}
For $M\ge1$,
\[
 O/2^MO\simeq(\Z/2^M\Z)((t)).
\]
Thus an element of $O$ has a unique expansion
$\sum_{n\in\Z}a_nt^n$ with $a_n\in\Ztwo$ and
$v(a_n)\to\infty$ as $n\to-\infty$. Its $O$-valuation is
$\min_n v(a_n)$.
\end{lemma}

\begin{proof}
If $s\in\Lambda\setminus(2)$, write
$s=t^h v_1+2w_1$, with $v_1\in\Lambda^\times$ and
$w_1\in\Lambda$. Modulo $2^M$ its inverse is a geometric expansion
of length $M$ in $2w_1/(t^h v_1)$. It is a Laurent series with finite
negative support. Conversely every Laurent series with finite negative
support lies in the localization of $\Lambda$ modulo $2^M$, since
$t$ is invertible there. Taking inverse limits proves the expansion
statement. Reduction modulo successive powers of $2$ proves the
valuation formula.
\end{proof}

\begin{lemma}[A uniform pole bound for Schur elimination]
\label{lem:h:bounded-poles}
Let $A_1,\ldots,A_m$ be square matrices over $\Lambda$ with nonzero
residual determinants and
$\ord_t(\overline{\det A_j})\le h$. Put
$A=\diag(A_1,\ldots,A_m)$. If $B=A+2E_1$ with $E_1$ a matrix over
$\Lambda$, then $B$ is invertible over $O$. Modulo $2^M$, every entry
of $B^{-1}$ has no powers below $-Mh$, independently of $m$ and of
the matrix sizes. The same bound holds for a Schur complement
$D=P-QB^{-1}S$ when $P,Q,S$ have entries in $\Lambda$.

In particular, suppose a block over $R$ augments to $A$. It is
invertible over $O[G]$, and the assertion applies uniformly at all sign
evaluations. If the lower off-diagonal block $S$ augments to zero, the
Schur complement augments to the old block $P$.
\end{lemma}

\begin{proof}
Write
$\det A_j=t^{h_j}v_j+2w_j$, where $h_j\le h$,
$v_j\in\Lambda^\times$, and $w_j\in\Lambda$. Put
\[
 C_j=\frac{\operatorname{adj}(A_j)}{t^{h_j}v_j},
 \qquad C=\diag(C_1,\ldots,C_m).
\]
Every entry of $C$ has pole at most $h$. We have $BC=1+2F_1$, with
every entry of $F_1$ having pole at most $h$. Consequently
\[
 B^{-1}\equiv C\sum_{j=0}^{M-1}(-2F_1)^j\pmod{2^M}.
\]
The $j$th summand has pole at most $(j+1)h\le Mh$. The number of
summands in a matrix product does not increase a valuation or pole
bound. Multiplication by power-series matrices likewise introduces no
negative powers. This proves the first assertion.

The ring $O[G]$ is local with maximal ideal generated by $2$ and the
augmentation ideal of $G$. The residual determinant of the proposed
block is that of $A$, a unit over $\F_2((t))$, so the block is
invertible. Each sign evaluation differs from augmentation by a
matrix divisible by $2$ with entries in $\Lambda$. Apply the first
assertion to it. The final assertion is the Schur formula with
$S_0=0$.
\end{proof}

\begin{lemma}[A bounded clearing factor]\label{lem:h:pivot-clearing}
Let $D$ be a based bounded complex of finite free $O[G]$-modules, of
total graded rank at most $r$, whose augmentation $D_0$ is a complex
over $\Lambda$. Suppose that the rational specialization of $D_0$ at
$t=0$ has one of the following cohomology patterns:
\begin{enumerate}[label=\textup{(\roman*)}]
 \item dimension one in degrees one and two, and zero elsewhere;
 \item dimension two in degree one, and zero elsewhere.
\end{enumerate}
Suppose that the total torsion length in the integral specialized
cohomology is at most $h_1$.

In case \textup{(i)}, take an integral degree-one cycle $z$ and a closed
functional $\lambda\in2^{-a_0}(D^2)^*$, and suppose each character
constituent over the fraction field of $O$ is either acyclic or has
pattern \textup{(i)}. Let $u_x$ be the determinant tensor of their
cohomology classes, with value zero in the acyclic case. In case
\textup{(ii)}, take two cycles whose denominator exponents sum to at
most $a_0$, suppose each constituent has pattern \textup{(ii)}, and
let $u_x$ be their wedge determinant coordinate. One may multiply these
coordinates by any one element $\eta\in O[G]$.

Then there are $f,d\in O[G]$ with $d_x=f_xu_x$ at every character,
such that $f$ is $2^{a_0}$ times a polynomial of bounded degree in
the differential entries of $D$, the degree depending only on $r$.
Moreover
\[
 f_0\in\Lambda,\qquad f_0(0)\ne0,\qquad
 v(f_0(0))\le C(r,h_1,a_0).
\]
If, at every fixed precision, the entries of $D_x$ have a uniform bound
on negative Laurent support, then so do the $f_x$. No such assertion is
required for $d_x$, for the cycles, or for $\eta$.
\end{lemma}

\begin{proof}
Over $\Ztwo$, put the specialized complex into a sum of free
cohomology terms and one-dimensional two-term complexes with
differentials $2^{a_j}$, by elementary-divisor decomposition. The sum
of the positive $a_j$ is at most $h_1$. Lift the constant integral
basis changes to $D$. Cancel the same pivot positions formally, retaining
their pivot numerators as denominators. The number of operations is bounded
by $r$. Each pivot and basis change is rational in the differential
entries. Let $P_1$ be a product of sufficiently high powers of the
successive pivot numerators so that every denominator in these changes
and in the determinant comparison divides $P_1$. Its degree is bounded
in terms of $r$; its specialization $P_{1,0}(0)$ is nonzero with
valuation bounded in terms of $r,h_1$. Increase the powers once more
if needed to give all final expressions the same denominator.

Over $O[G][1/P_1]$, case \textup{(i)} leaves a complex
$[O[G]\xrightarrow{\delta}O[G]]$ in degrees one and two. Let
$\alpha,\beta$ be the projected cycle and closed functional. Their
determinant tensor, including the cancelled-pivot determinant factor,
is a rational expression in $D,z,2^{a_0}\lambda$. If a field
constituent has $\delta\ne0$, the cycle identity
$\delta\alpha=0$ gives $\alpha=0$, so the same expression yields
the stipulated zero coordinate. In case \textup{(ii)} the remaining
complex is $O[G]^2[-1]$ and its coordinate is the determinant of the
two projected cycles times the pivot determinant factor. In either
case, after denominator clearing, for every constituent on which
$P_{1,x}\ne0$ we have
\[
 u_x=2^{-a_0}\eta_x Q_{1,x}/P_{1,x},\qquad Q_1\in O[G].
\]
Set $f=2^{a_0}P_1^2$ and $d=\eta P_1Q_1$. This gives $d_x=f_xu_x$
on those constituents. Where a pivot numerator first vanishes,
$P_{1,x}=0$, and both sides of the cleared identity are zero. Thus the
identity holds on every character. Since $f$ involves only the
bounded number of differential entries in the pivot calculation,
its degree, augmentation valuation and pole bounds are as asserted.
The factor $\eta$ has been placed in $d$, never in $f$.
\end{proof}

The preceding lemmas isolate the denominator problem: after eliminating
new blocks, the clearing factor uses a complex of bounded size.
The remaining task is to match finitely many coefficients at zero and
at some nonzero vertex. The next lemma gives that match once the cube
has sufficiently many coordinates.

\begin{lemma}[Finite coefficient congruences]
\label{lem:h:boolean-coefficients}
Let $A_1(x),\ldots,A_q(x)\in\Ztwo$ each be one specified Laurent
coefficient of the evaluation of an element of $O[G]$. If
\[
 b>q(2^{M-1}-1),
\]
there is $x\ne0$ such that
$A_i(x)\equiv A_i(0)\pmod{2^M}$ for every $i$.
\end{lemma}

\begin{proof}
For one coefficient write
\[
 A(x)=\sum_{g\in G}a_g(-1)^{g\cdot x}
     =\sum_{S\subseteq\{1,\ldots,b\}}c_Sx_S,
 \qquad x_S=\prod_{i\in S}x_i,
 \qquad 2^{|S|}\mid c_S.
\]
The last divisibility follows by expanding each sign as a product of
$1-2x_i$. The $j$th binary digit of $z\in\Ztwo$ is
$\delta_j(z)=\binom z{2^j}\bmod2$; this holds on nonnegative integers
by binary expansion and then on $\Ztwo$ by continuity. In the ring
of functions on the binary cube with coefficients in $\F_2[[Z]]$,
\[
 (1+Z)^{A(x)}=(1+Z)^{c_\varnothing}
  \prod_{S\ne\varnothing}
    \left(1+x_S\big((1+Z)^{c_S}-1\big)\right).
\]
The nonconstant factor indexed by $S$ starts at degree at least
$2^{|S|}$ in $Z$. A term contributing to $Z^{2^j}$, for $j\ge1$,
therefore has Boolean degree at most
\[
 \sum_S|S|\le\sum_S2^{|S|-1}\le2^{j-1}.
\]
The zeroth digit is constant. The congruences in the statement are
thus equations over $\F_2$ of total degree at most
$q\sum_{j=1}^{M-1}2^{j-1}=q(2^{M-1}-1)$, all vanishing at zero.
For completeness, the number of their common zeros is even: modulo
two it is the sum over the cube of the product of $1+Q$ over all
defining equations $Q$. Every monomial has degree below $b$, hence
omits a variable and sums to zero. This is the
Chevalley--Warning argument \citep{Chevalley1935,Warning1935}. The zero solution therefore has a
nonzero companion.
\end{proof}

We now apply these coefficient congruences to the cleared identity
$d_x=f_xu_x$. Only the bounded negative support of $f_x$ is prescribed;
for $d_x$ we will match its constant coefficient alone. This distinction
keeps the required number of congruences independent of the numerator's
other coefficients.

\begin{lemma}[Transfer through a bounded clearing factor]
\label{lem:h:boolean-transfer}
Let $f,d\in O[G]$ and let $u_x\in\Lambda[1/2]$ satisfy
$d_x=f_xu_x$ at each character. Suppose
$f_0\in\Lambda$ and $c_0=v(f_0(0))<\infty$. Fix $M>c_0$, and
suppose every $f_x\bmod2^M$ has no powers below $-L(M)$. If
\begin{equation}\label{eq:h:cube-threshold}
 b>(L(M)+2)(2^{M-1}-1),
\end{equation}
some $x\ne0$ satisfies
\begin{equation}\label{eq:h:central-transfer}
 v\bigl(u_x(0)-u_0(0)\bigr)\ge M-2c_0.
\end{equation}
Consequently, when $c_0$ and each $L(M)$ have bounds independent of
the family, a uniform upper bound $v(u_x(0))\le B_1$ at all nonzero
vertices forces $v(u_0(0))\le B_1$ for all sufficiently large $b$.
The base value itself is allowed to vary with the family.
\end{lemma}

\begin{proof}
Apply Lemma~\ref{lem:h:boolean-coefficients} to the coefficients of
$f_x$ in degrees $-L(M),\ldots,0$ and to the constant coefficient
of $d_x$. At the supplied nonzero solution, every negative coefficient
of $f_x$ is divisible by $2^M$, and its constant coefficient is
congruent to $f_0(0)$. Therefore its constant coefficient has valuation
$c_0$, and $v_O(f_x)\le c_0$. Since $d_x\in O$, the identity
$d_x=f_xu_x$ gives $v_O(u_x)\ge-c_0$; the same bound holds for
$u_0$. By Lemma~\ref{lem:h:laurent-expansions}, all their coefficients
are in $2^{-c_0}\Ztwo$. Both have zero negative coefficients because
they lie in $\Lambda[1/2]$.

The convergent constant-term convolution gives
\[
 [t^0]d_x\equiv([t^0]f_x)u_x(0)\pmod{2^{M-c_0}},
 \qquad [t^0]d_0=f_0(0)u_0(0).
\]
Combine these identities with the imposed coefficient congruences and
divide by $f_0(0)$. This proves
Equation~\eqref{eq:h:central-transfer}. For the consequence, choose
one integer $M>\max\{c_*,B_1+2c_*\}$, where $c_*$ bounds $c_0$,
and then take $b$ above Equation~\eqref{eq:h:cube-threshold}. If the
base valuation exceeded $B_1$, the transferred nonzero valuation would
also exceed $B_1$, a contradiction. The pivot valuation affects the
required precision and dimension, but not the transferred upper bound.
\end{proof}

\subsection{Finite cochain models}
\label{sec:common-finite-models}

The arithmetic applications also need finite free complexes that retain
maps to the original cochains. The following algebraic construction
provides these models and transports their comparison diagrams.
It applies before any determinant coordinate is chosen.

\begin{lemma}[Finite Artin models]\label{rl:finite-model}
\label{lem:h:finite-models}
Let $(A,\mathfrak m,k)$ be an Artin local ring, and let $D$ be a complex of
flat $A$-modules.  Suppose $H^*(D\otimes_A k)$ is finite dimensional and
supported in $[a,b]$.  There is a bounded finite free complex $C$, with
\[
 \operatorname{rank}_A C^j=\dim_k H^j(D\otimes_A k),
\]
and a strong deformation retraction
\[
 i:C\longrightarrow D,\qquad p:D\longrightarrow C,\qquad
 pi=1,\qquad 1-ip=dh+hd.
\]
Thus maps to the original complex, as well as maps between finite models,
are available.  Any fixed finite collection of chain maps, homotopies,
triangles, and perfect pairings can be transported to these models.
\end{lemma}

\begin{proof}
Flat modules over an Artin local ring are free, including modules of infinite
rank.  Here is a useful version of the argument.  Lift a basis of $M/\mathfrak
mM$ to a flat module $M$, and map the corresponding free module to $M$.
Its cokernel is zero because $\mathfrak m$ is nilpotent.  Flatness shows that
the kernel has zero reduction, and nilpotence again makes the kernel zero.

Split the residual complex as its cohomology, with zero differential, plus
contractible two-term disks.  Lift all the graded bases.  These lifts are
graded isomorphisms: an inverse modulo $\mathfrak m$ lifts to an inverse by
a finite geometric series.  Let $d_0$ be the lifted disk differential and
write $d=d_0+\epsilon$.  All coefficients of $\epsilon$ belong to
$\mathfrak m$.  For the disk contraction $(i_0,p_0,h_0)$, arrange
$h_0^2=h_0i_0=p_0h_0=0$ and
$1-i_0p_0=d_0h_0+h_0d_0$.  The operators
$1+h_0\epsilon$ and $1+\epsilon h_0$ have finite geometric inverses.
The homological perturbation lemma gives explicit contraction maps;
we record the formulas to retain their compatibility with the
arithmetic diagrams \citep[Sections~2.2--2.4 and~3.2]{Crainic2004}.
The perturbation identities give the differential
\[
 d_C=p_0\epsilon(1+h_0\epsilon)^{-1}i_0
\]
on the lifted residual cohomology, and give
\[
 i=(1+h_0\epsilon)^{-1}i_0,\quad
 p=p_0(1+\epsilon h_0)^{-1},\quad
 h=h_0(1+\epsilon h_0)^{-1}.
\]
Substitution, using $(d_0+\epsilon)^2=0$, proves the asserted contraction
identities.  These calculations are valid for infinite disk sums: each
operator is a module homomorphism and every geometric series terminates
at the nilpotence exponent of $\mathfrak m$.

Transport a map $f$ as $p'fi$, retaining the contraction homotopies
for compositions. A quasi-isomorphism between bounded free complexes
has a contractible cone and a homotopy inverse. This transports finite
diagrams, including triangles and perfect dualities.
\end{proof}

\section{Cyclotomic interpolation at analytic rank zero}
\label{sec:cy:interpolation}

The cyclotomic argument transfers a nonzero central value across a binary
family. Its input is a uniform upper bound for the normalized valuations
at every nonzero vertex. Its output is nonvanishing at a corank-zero
base. The class used here is Kato's zeta class, and the additional variable
is the real cyclotomic $\Ztwo$-extension of $\Q$.

Retain $T=T_2(E)$ and put $V_2=T\otimes_{\Ztwo}\Qtwo$.
The local weights, finite Sha lengths and period normalization are those
of Section~\ref{sec:preliminaries}, and $D(h)$ is defined in
Equation~\eqref{eq:transfer-normalized-value}. We prove
Theorems~\ref{thm:cyclotomic-bound} and~\ref{thm:cyclotomic-interpolation}
with their stated distinction between the uniform specialization error
and the dimension threshold for a fixed base.

For a two-dimensional representation over $\F_2$, reducibility is
equivalent to being an extension of trivial modules. These are therefore
the two exhaustive residual cases. The proof has three stages: construct
a class in a fixed integral lattice, identify its central determinant,
and apply the binary congruences. Residual reducibility makes the
cyclotomic complex free after localization at $(2)$; the irreducible
case instead uses the bounded clearing factor of
Lemma~\ref{lem:h:pivot-clearing}.

\subsection{The positive cyclotomic complex}

At the coefficient prime $2$, the real place must be included in the
cochain construction. The positive complex defined below removes the
real cochains and retains their connecting invariant line. We will
combine this fixed line with the zeta class to form the determinant
whose central value detects nonvanishing at a corank-zero base.

Put
\[
 \Gamma=\Gal(\Q_\infty/\Q),\qquad
 \Lambda=\Ztwo[[\Gamma]]=\Ztwo[[t]],\qquad
 \mathcal O=\widehat{\Lambda_{(2)}},\qquad
 \Omega=\mathcal O/2\mathcal O=\F_2((t)),
\]
where a generator of $\Gamma$ corresponds to $1+t$.
The completion is $2$-adic. An element of $\mathcal O$ has a finite
negative Laurent part modulo every $2^M$, but need not have a value at
$t=0$. Every central evaluation below is made on a rational Iwasawa
class or on an element of $\Lambda[1/2]$ before extending coefficients
to $\mathcal O$.

For a family in Equation~\eqref{eq:cy:cube}, let
$G=(\Z/2\Z)^b$ have characters indexed by $x\in\F_2^b$, and let
$\psi:G_{\mathbb Q}\to G$ be its universal
quadratic character: evaluation at $x$ is $\chi_{h_x/h_0}$.  Put
\[
 R=\Lambda[G],\qquad
 \mathbb T=T\otimes_{\mathbb Z _2}R,
\]
with the fixed twist $\chi_{h_0}$, the universal twist $\psi$, and the
inverse cyclotomic action.  One may replace $G$ by the actual image of
$\psi$ and subsequently extend scalars.  In a family in
Theorem~\ref{thm:cyclotomic-interpolation}, a nonzero character trivial on that image
already gives the conclusion directly.

Let $S$ be the set of primes dividing $2N h_0\prod_{q\in\mathcal Q}q$.
For $q\in\mathcal Q$, write $g_q\in G$ for its inertia element, so
$x(g_q)=(-1)^{\lambda_q(x)}$. Use continuous cochains with their completed
coefficient topology, and set
\begin{equation}\label{eq:cy:positive-complex}
 C_S^+(\mathbb T)=
 \operatorname{Cone}\bigl(R\Gamma(G_{\mathbb Q,S},\mathbb T)
              \longrightarrow R\Gamma(\mathbb R,\mathbb T)\bigr)[-1].
\end{equation}
The real complex here uses ordinary cochains in nonnegative degrees.
There are no imposed local conditions at the finite places. We write
$C_S(\mathbb T)=R\Gamma(G_{\Q,S},\mathbb T)$ for the unmodified complex.
Throughout this section, \emph{ordinary global cohomology} means this
unmodified cohomology; it imposes no ordinary-reduction hypothesis at~$2$.

\begin{lemma}\label{lem:cy:positive-model}
The complex in Equation~\eqref{eq:cy:positive-complex} is perfect over $R$ and computes derived
coefficient specializations.  There is a class $a\in H^1(C_S^+(\mathbb T))$
whose rational span is the real-place summand.  Its integral index in
that summand depends only on $E$ and the fixed real sign.
\end{lemma}
\begin{proof}
Here are the finiteness and coefficient issues in the use of global
cochains.  For a finite coefficient quotient of $R$, continuous cochains
are filtered unions of finite products of free coefficient modules.
Refinement of finite partitions makes these modules flat; over an Artin
local coefficient ring they are free.  For finite residual coefficients,
global cohomology is finite in each degree, and restriction to the real
places is an isomorphism in degrees greater than two, by the real-place
clause of global duality \cite[\S5.7.1.7]{Nekovar2006}.
Thus the residual fiber has bounded
amplitude and finite-dimensional cohomology.  The same amplitude bound
holds for every finite coefficient module, by a composition series and
the cohomology long exact sequence.

This gives a bounded finite free model as follows.  Over the residue
field split off the acyclic disks, leaving its finitely many cohomology
spaces.  Lift each invertible disk differential over an Artin coefficient
quotient and cancel it.  Invertibility of a lifted inverse follows from
nilpotence of the maximal ideal.  Equivalently one can perform this on a
free resolution; a minimal free model cannot have an additional term
whose reduction has zero cohomology, because its residual differential
is zero.  Hence the remaining ranks are the finite residual ranks and
are bounded independently of the coefficient quotient.  Minimal models
identified by a homotopy equivalence over an Artin local ring are
isomorphic.  Lifting their changes of basis therefore makes these models
compatible along $R/\mathfrak m^n$.  Taking the inverse limit gives a
bounded finite free model over~$R$.  The cochain construction with finite
free coefficients and the same cancellation also proves derived base
change.  This is a fixed-$S$ argument; it does not assert uniform bounds
on the number of generators as $S$ varies.

All the characters of the family have the same real sign, so the action
at the real place is the fixed action on $T\chi_{h_0}$ extended to $R$.
Choose a nonzero vector in its integral invariant line.  Its connecting
image in the fiber is $a$.  There are no global invariant vectors in the
rational elliptic representation, and consequently this connecting map
is injective rationally.  The choice and its index are fixed.  Notice also
that $H^1(\mathbb R,\mathbb T)$ is killed by~$2$; thus twice any ordinary
cohomology class lifts to the positive complex.  The real-place global duality input is its corrected form at~$2$
\cite[\S\S5.7.1--5.7.5]{Nekovar2006}; the invariant line and the lifting
assertion follow from the long exact sequence of this ordinary-real cone.
Complete real Tate cochains, used in the duality formulation, become
acyclic after inverting~$2$ and would not retain this invariant line.
\end{proof}

\subsection{Fixed-form zeta classes and the integral coefficient map}

To control denominators as the tame support grows, we construct every
class in the cohomological lattice of one fixed modular form.  Let $f$ be the fixed weight-two form of $E$, and choose
its fixed integral cohomological lattice, with the Tate twist making its
rational representation~$V_2$.  Kato's classes (8.1.2)--(8.1.3), for a fixed
symbol $\xi\in\mathrm{SL}_2(\mathbb Z)$, exist for arbitrary $m$ and any
$S$ containing the prime factors of $mN$.  The trace in Section~8.9 is
a homomorphism into the lattice
\[
 \mathbb Z _2[\operatorname{Gal}(\mathbb Q(\mu_m)/\mathbb Q)]
       \otimes V_{2,\mathbb Z _2}(Y_1(N))(1).
\]
Section~8.11 then projects to the lattice of the fixed form.
Proposition~8.12 supplies its norm relations.  These assertions hold at
the coefficient prime~$2$ \cite[Sections~5.1, 8.1, 8.9--8.12]{Kato2004}.
In particular, the auxiliary modular-curve level used to define a trace
does not become the level of the projected newform.

Choose an isomorphism of the resulting fixed rational representation with
$V_2$.  A single power $2^{c_E}$ sends its integral lattice into~$T$.
Choose fixed symbols with nonzero projections on the required period
lines; their existence and the construction with both signs are made
explicit in \cite[Sections~13.6 and 13.9]{Kato2004}.  Taking a fixed linear
combination, if needed, costs another fixed denominator.

The following integral coefficient map encodes the binary family. Choose
$m=|h_0|\prod_{q\in\mathcal Q}q$, which is odd and divisible by the
conductors of all the tame characters in the family, including
$\chi_{h_0}$.  At a sufficiently large
cyclotomic level write $A_{m,n}=\Gal(\Q(\mu_{m2^n})/\Q)$
and let $\Gamma_n$ be its quotient in the real cyclotomic tower.  Shapiro identifies the cohomology of the cyclotomic field
with global cohomology with induced group-ring coefficients.  Push the
class along the ring homomorphism
\begin{equation}\label{eq:cy:coefficient-map}
 \mathbb Z _2[A_{m,n}]\longrightarrow
       \mathbb Z _2[\Gamma_n\times G],\qquad
 [\sigma]\longmapsto
       \chi_{h_0}(\sigma)[\sigma|_{\Gamma_n}][\psi(\sigma)].
\end{equation}
Use inverse actions on both sides if this is required by the chosen
Shapiro convention.  This is a homomorphism of integral coefficient
modules and acts on actual cohomology classes.  It requires neither
division by $|G|$ nor surjectivity of a cohomological specialization map.
The prime $2$ belongs to the fixed omitted support, so Proposition~8.12
gives exact norm compatibility as $n$ increases.  The pushed classes
therefore give an integral smoothed Iwasawa class with coefficients in
$\mathbb T$, up to the single fixed power~$2^{c_E}$.

We have placed all character constituents in the lattice of the same
modular form. The next normalization issue is smoothing: its
removal must be harmless both after localization at $(2)$ and at the
central value $t=0$. These are distinct requirements, checked next.

We next remove smoothing.  In weight two and with $r=r'=1$, its factors
are
\begin{equation}\label{eq:cy:smoothing}
 (c^2-c\sigma_c)(d^2-d\sigma_d).
\end{equation}
This is the specialization of
\cite[Theorem~6.6 and (4.2.4)]{Kato2004}.  Choose a fixed
$e\geq\max(2,v_2(N))$, increasing it to accommodate the finitely many
fixed symbol levels.  The Chinese remainder theorem permits $c,d>1$,
prime to $6Nm$, such that
\[
 c,d\equiv1\pmod{3N_{\mathrm{odd}}m},\qquad
 c,d\equiv1+2^e\pmod{2^{e+1}}.
\]
Their tame actions in \eqref{eq:cy:coefficient-map} are trivial.  If
$\sigma_c$ acts in $\Gamma$ as $(1+t)^{\kappa(c)}$, then
$v_2(\kappa(c))=e-2$.  Modulo the maximal ideal of $\mathcal O[G]$ its
smoothing factor is
\[
 1-(1+t)^{\kappa(c)}\ne0\quad\hbox{in }\mathbb F _2((t)).
\]
It is therefore a unit in $\mathcal O[G]$.  Its central value is
$c(c-1)$, of valuation exactly~$e$.  The same statements hold for $d$
and for inverse-translate conventions.  Thus unsmoothing introduces
neither a $2$-adic denominator at $(2)$ nor a growing central denominator.
The integers $c,d$ vary with $m$; the congruences fix the needed valuations,
not their exact cyclotomic group elements.

Denote the resulting rational ordinary class by $z_S$.  \citet[Theorem~9.7]{Kato2004} applies to $k=2,r=r'=1$.  In the character constituent
$x$, its central reciprocity formula reads
\begin{equation}\label{eq:cy:reciprocity-center}
 v_2\bigl(\exp_2^*z_{S,x}(0)\bigr)
 =v_2(\mathcal L(h_x))+
  \sum_{\substack{q\in\mathcal Q\,,\ q\nmid h_x}}
       v_2\bigl(\#E^{(h_x)}(\mathbb F _q)\bigr)+O_E(1)
\end{equation}
when the value is nonzero, and a zero central value gives a zero dual
exponential.  The omitted primes dividing $2N$ are included in the error.
To check the normalization, the character sum in
\eqref{eq:cy:coefficient-map} is the unnormalized sum, as in Kato's trace.
Its tame quadratic Gauss sum supplies $\sqrt{|h_x|}$ up to a root of
unity.  Tame conductors and their Gauss sums are $2$-adic units.  The
period vector is the fixed symbol vector of $f$.  More explicitly,
in the required sign line write the period represented by
$\delta(f,1,\xi)^{\pm}$ as $c_\xi^{\pm}\Omega_E^{\pm}$.
The nonzero scalar $c_\xi^{\pm}$ is fixed: the character-sum formula
of \cite[Theorem~6.6]{Kato2004} uses this same $\delta(f,1,\xi)^{\pm}$
for every tame character.  The tame character enters through its
$L$-value and Gauss sum, not through a newly chosen period vector.
The comparison with the minimal differential of $E^{(h_x)}$ at $2$ has bounded valuation,
because there are only finitely many local twisting classes there.
At a good unused prime the omitted factor is
\[
 1-\frac{a_q(E)\chi_{h_x}(\operatorname{Fr}_q)}q+\frac1q
       =\frac{\#E^{(h_x)}(\mathbb F _q)}q;
\]
at a good ramified prime the local $L$-factor is~$1$.  These account for
every varying factor in \eqref{eq:cy:reciprocity-center}.

\subsection{Rational divisibility and the local Euler corrections}

We now exclude poles of the zeta determinant away from $(2)$ and
account for every omitted Euler factor.  Write $A=\Lambda[1/2]$.  Theorem~12.5(3) of
\cite{Kato2004} is a height-one divisibility over~$A$; its possible local
$H^2$ correction comes from the arithmetic place~$2$. It need not be
supported at the Iwasawa ideal~$(2)$. The global groups
there use $j_*$ outside that prime.  The stronger integral assertion of
Theorem~12.5(4) is not used.  For a non-CM elliptic curve the characteristic-zero
image condition is supplied by the open image theorem, independently of
the residual image.  For CM the required rational statement at every
prime, including~$2$, is supplied by
\cite[Theorem~2.6 and Remark~2.7]{BurungaleTian2026}.  Neither statement
asserts uniform integrality in a quadratic group ring. The rational
cohomology statements, including \citet[Theorem~12.4]{Kato2004}, give rank one in degree one and torsion in degree
two, with no other positive cohomological rank. The real-place
modification adds its invariant line. Hence every generic character
constituent of the positive complex has rank two in degree one and no
other positive rank.

Fix a character constituent $h$ and write $V_h=V_2\chi_h$.  Let $C_{j_*}$
be the rational Iwasawa complex in Kato's formulation and $C_S$ the
rational full global complex.  At an odd prime put
\[
 W_q=H^1(I_q,V_h).
\]
The localization triangle has local cone
\begin{equation}\label{eq:cy:local-block}
 Q_q=R\Gamma\bigl(G_{\mathbb F _q},
                   W_q\otimes A(\Gamma^\iota)\bigr)[-1]
   =\left[W_q\otimes A
       \xrightarrow{\ 1-\gamma_q\operatorname{Fr}_{W_q}\ }
             W_q\otimes A\right]
\end{equation}
in degrees one and two.  Here $\gamma_q$ denotes the scalar cyclotomic
Frobenius in the chosen convention.  Its exponent is nonzero.  Consequently
$D_q=\det(1-\gamma_q\operatorname{Fr}_{W_q})$ is a nonzero element of~$A$,
and the map in \eqref{eq:cy:local-block} is injective.  This remains true
after localization at a height-one prime dividing $D_q$: localization
does not mean reduction modulo that prime.  Thus
\[
 H^1(Q_q)=0,\qquad
 \operatorname{length}_{A_{\mathfrak p}}H^2(Q_q)_{\mathfrak p}
       =v_{\mathfrak p}(D_q).
\]
The long exact sequence gives
\begin{equation}\label{eq:cy:local-sequence}
 \begin{gathered}
 H^1(C_{j_*})\simeq H^1(C_S),\\
 0\longrightarrow H^2(C_{j_*})\longrightarrow H^2(C_S)
       \longrightarrow\bigoplus_{q\in S\setminus\{2\}}H^2(Q_q)
       \longrightarrow H^3(C_{j_*}).
 \end{gathered}
\end{equation}
In particular, the length of $H^2(C_S)$ is at most the length of
$H^2(C_{j_*})$ plus the sum of the local lengths.  No assertion that
$H^3(C_{j_*})$ vanishes is required here.

At a good unused prime, $W_q=V_h(-1)$, so, on including the quadratic
Frobenius sign in $Z$,
\begin{equation}\label{eq:cy:local-euler}
 D_q=P_q(Z),\qquad P_q(Z)=1-\frac{a_q(E)}q Z+\frac1q Z^2.
\end{equation}
At a good active prime, inertia acts by $-1$ on $V_h$.  Its rational
inertia cohomology is zero, so $W_q=0$ and $D_q=1$.

At fixed bad primes the general expression is $H^1(I_q,V_h)$ in
\eqref{eq:cy:local-block}.  Replacing it indiscriminately by $V_h^{I_q}(-1)$
would be incorrect.  For example, at split multiplicative reduction,
\[
 V_h^{I_q}=\mathbb Q _2(1),\qquad
 H^1(I_q,V_h)=\mathbb Q _2(-1),
\]
and the correct determinant is $1-\gamma_q/q$, with nonzero central
value.  The incorrect replacement would give $1-\gamma_q$.
In general $D_q(0)\ne0$ follows from local duality and the Weil pairing:
$H^2(\mathbb Q _q,V_h)=H^0(\mathbb Q _q,V_h)^*=0$, since elliptic torsion
over a local field is finite.  The fixed bad-prime corrections therefore
have nonzero central values for every reduction type.

Let $z_{\mathrm{BK},h}$ denote a generator of the rational Kato zeta
line in its real cyclotomic component.  The rational comparison with
the fixed-form class constructed above is made only after taking the
character constituent.  Both global Iwasawa $H^1$ modules have rank one.
Their reciprocity laws at cyclotomic finite-order characters agree up to
the omitted Euler factors, a nonzero constant, and a group-like unit.
This follows by coprime Gauss-sum factorization and the fixed period
comparison; restricting conductor exponents to either fixed parity
removes any factor $\chi_h(2)^n$.  Infinitely many of these tests have
nonzero central value by cyclotomic nonvanishing, as used in
\cite[Theorem~13.5(2)]{Kato2004}.  After clearing the finitely many
denominators of the scalar ratio, the difference is a power series
vanishing at infinitely many finite-order characters.  Weierstrass
preparation forces that difference to be zero.  This proves the generic
comparison without identifying the integral lattices of varying
twisted newforms.  Its nonzero scalar in $\mathbb Q _2$ may depend on
$h$ and need not have a uniformly bounded $2$-valuation: every such
scalar is a unit of $A$ and does not affect height-one divisibility
over~$A$.  The uniform integral and central bounds instead concern
the actual fixed-form class and the fixed period vector in
\eqref{eq:cy:reciprocity-center}.

Equations~\eqref{eq:cy:local-sequence}--\eqref{eq:cy:local-euler} now show exactly
how the imprimitive factors enter divisibility.  At any height-one
$\mathfrak p$ of $A$, enlarging the global complex increases its
degree-two length by at most $\sum_qv_{\mathfrak p}(D_q)$, whereas
multiplying the rank-one zeta generator by $\prod_qD_q$ adds exactly
that quantity to its index.  To cover Kato's coefficient-prime correction, multiply the
class also by a fixed element $F_2\in\Lambda$ divisible over $A$ by the
characteristic ideal of the relevant local $H^2_{\mathrm{Iw}}$.
It may be chosen with $F_2(0)\ne0$.  Indeed derived local control gives
\[
 H^2_{\mathrm{Iw}}(\mathbb Q _2,T\chi_h)_\Gamma\otimes\mathbb Q _2
       =H^2(\mathbb Q _2,V_h)=0,
\]
so the characteristic ideal has no factor $t$.  Local twisting classes
at $2$ are fixed, and one element covers their finite set.

Likewise any additional fixed Euler factors in the rational comparison
can be cleared by a fixed product with nonzero central value, using the
actual blocks \eqref{eq:cy:local-block}.  Denote the total fixed multiplier
by~$F_E$.  Multiplying by it only changes the fixed constants in
\eqref{eq:cy:reciprocity-center} and the integrality bound.  For the corrected
class, still written $z_S$, we have at every height-one $\mathfrak p$ of~$A$
\begin{equation}\label{eq:cy:divisibility}
 \operatorname{length}_{A_{\mathfrak p}}H^2(C_S)_{\mathfrak p}
 \leq
 \operatorname{length}_{A_{\mathfrak p}}
       \bigl(H^1(C_S)/Az_{S,h}\bigr)_{\mathfrak p}.
\end{equation}

There is an integral convention check for the varying unused primes.
In $\mathcal O[G]$, both $P_q(Z)$ and $Z^2P_q(Z^{-1})$ are units: their
reductions are nonzero polynomials evaluated at the nonconstant
cyclotomic Frobenius.  The identity
\[
 P_q(Z)-Z^2P_q(Z^{-1})=\frac{q-1}{q}(1-Z^2)
\]
shows that their ratio $r_q$ is $1$ modulo~$2$.  If $g_q$ is the universal
inertia element, then
\begin{equation}\label{eq:cy:reciprocal-correction}
 U_q=1+\frac{r_q-1}{2}(1+g_q)\in\mathcal O[G]^\times
\end{equation}
is integral.  At an active character it is $1$; at an unused character
it is $r_q$.  Its unused central value is $1$, since $Z=\pm1$ there.
Thus inverse Frobenius conventions can be reconciled before character
evaluation, without a denominator for each prime.  Products of these
units have the same property.  The resulting rational character
coordinates have the Euler convention in \eqref{eq:cy:divisibility}.

\begin{lemma}[Fixed-lattice zeta class]\label{lem:cy:kato-class}
For every binary family and its ramification set $S$, the corrected
class $z_S$ satisfies the central reciprocity formula \eqref{eq:cy:reciprocity-center} and the rational divisibility
\eqref{eq:cy:divisibility}. For some $c_E\ge0$ depending only on $E$,
\[
 2^{c_E}z_S\in
 H^1\!\left(C_S(\mathbb T)\otimes_R^{\mathbf L}\mathcal O[G]\right).
\]
The same exponent works for all $S$, binary families, and fixed factors
$h_0$ with the specified local classes. Here the displayed cohomology
is derived scalar extension of the Iwasawa complex.
\end{lemma}
\begin{proof}
The integral coefficient map uses the fixed-form lattice and no
character averages. Unsmoothing and every varying-prime convention
correction are units over $\mathcal O[G]$. The multiplier $F_E$ is fixed.
The preceding construction therefore gives the displayed assertion
with a fixed exponent, as well as the two stated comparison formulas.
\end{proof}

\subsection{The determinant and its central specialization}

The zeta class now has two controls: rational divisibility will exclude
poles of its determinant coordinate away from $(2)$, and the fixed
lattice bounds the class's integral denominator at $(2)$. We next
identify its determinant coordinate at
$t=0$. This is the step that converts those cohomological controls
into a statement about $D(h)$.

Choose a basis of the invertible $R$-module
$\det_R^{-1}C_S^+(\mathbb T)$; $R$ is local, so such a basis exists.
Lift $2z_S$ to rational positive cohomology and take its wedge with $a$.
The factor $2$ is fixed, and two lifts differ by a real-place class,
which does not change this wedge.  Let $u$ be its generic determinant
coordinate, including the fixed multiplier $F_E$ above.  The real-place
summand introduces no height-one torsion over~$A$.  Thus
\eqref{eq:cy:divisibility} says that $u_x$ has no pole at any height-one
prime of $\Lambda$ except possibly $(2)$.  Since $\Lambda$ is normal,
\begin{equation}\label{eq:cy:rational-integrality}
 u_x\in\Lambda[1/2]\qquad\hbox{for every }x.
\end{equation}
The two cycles have a common denominator bounded by a fixed power
$2^{C_{\mathrm{int}}}$ after extension to $\mathcal O[G]$. This follows
from Lemma~\ref{lem:cy:kato-class} and the factor $2$ needed for the lift.

\begin{proposition}[Central determinant formula]\label{prop:cy:specialization}
There is a constant $C_{\mathrm{loc}}$, depending only on the fixed curve
and local classes, with the following property. Suppose
$c_2(E^{(h_x)})=0$. If $L(E^{(h_x)},1)=0$, then $u_x(0)=0$; if
$L(E^{(h_x)},1)\ne0$, then
\begin{equation}\label{eq:cy:determinant-center}
 \bigl|v_2(u_x(0))-D(h_x)\bigr|\leq C_{\mathrm{loc}}.
\end{equation}
The rational positive complex at $t=0$ has cohomology only in degree one,
of dimension two, so these are specializations of the same determinant.
\end{proposition}
\begin{proof}
Abbreviate $T_h=T_2E^{(h_x)}$ and $A_h=E^{(h_x)}[2^\infty]$.
Use the integral Kummer subgroup
\[
 K_v=\operatorname{im}\bigl(E^{(h_x)}(\mathbb Q _v)^{\wedge}_2
                              \longrightarrow H^1(\mathbb Q _v,T_h)\bigr)
\]
at every finite place.  For $v\ne2$ it is all of $H^1(\mathbb Q _v,T_h)$.
One way to see this integrally is the inverse-limit Kummer exact sequence:
its quotient is $T_2H^1(\mathbb Q _v,E^{(h_x)})$, which is zero because
local Tate duality identifies the $2$-primary local Weil--Ch\^atelet group
with the dual of the finite $2$-completion of the local point group.
At $2$ the quotient
\[
 L_2=H^1(\mathbb Q _2,T_h)/K_2
\]
is a free rank-one $\mathbb Z _2$-module.  The same Kummer sequence and
local Tate duality identify it with the $\mathbb Z _2$-dual of the free
part of $E^{(h_x)}(\mathbb Q _2)^{\wedge}_2$.

The dual exponential identifies $L_2\otimes\mathbb Q _2$ with the
one-dimensional de Rham target. To specify its integral comparison, put
\[
 B_2=E^{(h_x)}(\mathbb Q _2)^{\wedge}_2/\mathrm{tors}
\]
and choose a minimal differential $\omega_h$.
The Bloch--Kato exponential agrees with the Kummer image of the classical
elliptic exponential, which is the inverse of the formal logarithm on a
sufficiently small open subgroup. This comparison holds for every
reduction type \cite[Definition~3.10 and Example~3.11]{BlochKato1990}.
The dual exponential is the adjoint of that exponential under local
Tate duality \cite[\S II and Proposition~II.5]{Berger2003}.
Extend the formal logarithm to the rationalized local point group and write
$\log_{\omega_h}(B_2)=\lambda_h\mathbb Z _2$.
The integral identification
$L_2=\operatorname{Hom}_{\mathbb Z _2}(B_2,\mathbb Z _2)$ then gives
\[
 \exp_2^*(L_2)=\lambda_h^{-1}\mathbb Z _2\omega_h,
\]
using the Weil-pairing identification of the de Rham target with
$\mathbb Q _2\omega_h$. Thus the relative lattice valuation is fixed
by the local curve and differential. Since only finitely many local
twist classes occur, these valuations are uniformly bounded.

The assumption on the Selmer corank implies that the compact finite
Selmer group is torsion and that the discrete finite Selmer group has
length $s(h_x)$.  Rational Poitou--Tate consequently identifies ordinary
global $H^1$ with the singular line at $2$, and gives ordinary global
$H^2=0$ rationally.  Adding the real invariant line gives the claimed
dimension two in degree one.  This equals the generic rank, so derived
base change of the perfect complex specializes its determinant without
a change of rational cohomological rank.

Here is the integral length calculation.  Let $J$ be the image of the
free part of global $H^1(G_{\mathbb Q,S},T_h)$ in $L_2$, and put
$i=\operatorname{length}(L_2/J)$.  The compact/discrete Poitou--Tate
sequence for the mutually orthogonal Kummer conditions contains
\[
\begin{aligned}
 0\longrightarrow L_2/J
 &\longrightarrow \operatorname{Sel}_{2^\infty}(E^{(h_x)}/\mathbb Q)^\vee
 \longrightarrow H^2(G_{\mathbb Q,S},T_h)\\
 &\longrightarrow \bigoplus_{q\in S}H^2(\mathbb Q _q,T_h)
 \longrightarrow H^0(\mathbb Q,A_h)^\vee\longrightarrow0,
\end{aligned}
\]
with the usual real Tate-cohomology terms inserted.  Those terms are
killed by $2$ and have dimensions bounded by the fixed dimension of~$T$.
Taking lengths and using local duality gives
\begin{equation}\label{eq:cy:pt-length}
 \operatorname{length}H^2(G_{\mathbb Q,S},T_h)
   =s(h_x)-i+\sum_{q\in S}\operatorname{length}H^0(\mathbb Q _q,A_h)
      +O_E(1).
\end{equation}
This sequence follows from integral Poitou--Tate duality with the finite
Kummer conditions and the real Tate-cohomology correction
\cite[\S\S5.7.1.6, 5.7.4--5.7.5 and 6.9.1]{Nekovar2006}.
The free real invariant line is already accounted for by $a$; only the
bounded real torsion terms enter the error above.
No unramified condition replaces
the Kummer condition at a bad prime.  Global torsion is bounded by its
injection into $E^{(h_x)}(\mathbb Q _2)[2^\infty]$, whose order is bounded
by the fixed local twisting classes.  The finite terms at $2$ and at
fixed bad primes are bounded for the same reason.  Passing between
ordinary and positive cohomology changes only these bounded real and
global torsion terms.

Suppose first that the central value is nonzero.  The coordinate of
$z_{S,x}(0)$ in the free global line has valuation
\[
 v_2(\exp_2^*z_{S,x}(0))-i+O_E(1).
\]
For the inverse determinant, torsion in degree two subtracts its length
and torsion in degree one adds its length.  This sign can be checked on
the complex $[\mathbb Z _2\xrightarrow{2^a}\mathbb Z _2]$ in degrees one
and two: its inverse determinant lattice maps to $2^a\mathbb Z _2$ in
the rational determinant.  A vector's coordinate thus loses~$a$.
Using \eqref{eq:cy:pt-length}, the two appearances of $i$ cancel, leaving
\[
 v_2(u_x(0))=
 v_2(\exp_2^*z_{S,x}(0))-s(h_x)
       -\sum_{q\in S}\operatorname{length}H^0(\mathbb Q _q,A_h)+O_E(1).
\]

At an active good prime, quadratic inertia acts by $-1$ on the Tate
module.  Its invariants on $A_h$ are therefore exactly $E[2]$, and taking
Frobenius invariants gives length~$t_q$.  At an unused good prime, reduction
identifies the local $2$-primary torsion with
$E^{(h_x)}(\mathbb F _q)[2^\infty]$, of length
$v_2(\#E^{(h_x)}(\mathbb F _q))$.  These unused terms cancel exactly
those in \eqref{eq:cy:reciprocity-center}.  The active terms, including the primes
of $h_0$, sum to $2w(h_x)$.  All other terms have already been bounded
independently of the growing support.  This proves
\eqref{eq:cy:determinant-center}.

If the central value is zero, reciprocity makes the dual exponential
zero.  The rational ordinary specialization then satisfies the finite
condition at $2$, and it satisfies the finite conditions elsewhere
automatically.  It lies in the rational finite Selmer group, which is
zero.  Its wedge with the real-place class is therefore zero.  The
constant-rank determinant specialization proves $u_x(0)=0$.
\end{proof}

\subsection{Residual freeness and the lower bound}

The central formula expresses $D(h)$ through the determinant valuation.
When the residual representation is an extension of trivial modules,
we can bound that valuation directly: the positive complex becomes a
free rank-two module in degree one after extension to $\mathcal O[G]$.
The fixed-lattice class then supplies the lower bound without any loss
from a growing ramification support.

\begin{lemma}[Residual concentration]\label{lem:cy:residual-free}
If $E[2]$ is an extension of two trivial $\F_2$-modules, then there is an
isomorphism in the derived category
\[
 C_S^+(\mathbb T)\otimes_R\mathcal O[G]
       \simeq \mathcal O[G]^2[-1].
\]
\end{lemma}
\begin{proof}
The ring $\mathcal O[G]$ is local, with residue field $\Omega$ obtained
by setting $2=0$ and every $g=1$.  Thus the quadratic characters disappear.
Since $E[2]$ has a filtration with trivial factors, it is enough first to
compute the positive residual complex for the trivial representation.

Let $K_n$ be the real cyclotomic layer of degree $d_n=2^n$.  The number
of finite primes of $K_n$ above a fixed $S$ is bounded with $n$: at $2$
there is total ramification, and the Frobenius of each fixed odd prime
has nonzero image in $\Gamma$.  The $2$-ranks of the class groups of
$K_n$ are also bounded. Let $X$ be the inverse limit, under norm maps, of
their $2$-primary class groups. The unique prime above~$2$ is totally
ramified in this tower, so class-field-theoretic control gives
\[
 X_{\Gamma^{2^n}}\simeq\Cl(K_n)[2^\infty]
\]
\cite[Proposition~3.2.11 and Remark~3.2.12]{SharifiIwasawa}.
Ferrero--Washington gives $\mu(X)=0$ for this abelian tower
\cite{FerreroWashington1979}. The Iwasawa structure theorem then makes
$X$ finitely generated over $\mathbb Z _2$, so all its coinvariants have
uniformly bounded $2$-rank. Passing to the quotient $S$-class groups
preserves this upper bound.

The signature map of units in $K_n$ is surjective.  For $n\geq1$, let
$\zeta$ be a primitive $2^{n+2}$th root and put
$\epsilon=1+\zeta+\zeta^{-1}$.  If $f_1(X)=X^2-2$ and
$f_{n+1}(X)=f_n(X^2-2)$, then $f_n$ is the minimal polynomial of
$\zeta+\zeta^{-1}$ and $f_n(-1)=-1$.  Hence
$N_{K_n/\mathbb Q}(\epsilon)=-1$.  Its signature vector has odd
augmentation.  The embeddings form a regular set for the cyclic
$2$-group $\operatorname{Gal}(K_n/\mathbb Q)$, whose group algebra over
$\mathbb F _2$ is local.  A vector with odd augmentation is a unit in
this group algebra, so its translates span every signature.  For $n=0$
the assertion follows from the unit~$-1$.

The signature calculation supplies the real-place surjectivity needed
for the positive complex. The class-group and support terms are bounded
along the tower. We can therefore identify the one cohomological degree
whose dimension grows with the layer degree, and then recover the rank
after inverting $t$.

Kummer theory now determines the cohomological growth.
Writing $s_n$ for the number of finite support places and $c_n$ for the
$2$-rank of the $S$-class group, one has
\[
 \dim H^1(G_{K_n,S},\mu_2)=d_n+s_n+c_n.
\]
The real degree-one restriction is surjective by the signature result.
In degree two, the Kummer sequence has class-group term of dimension
$c_n$ and Brauer term.  The Brauer invariants at the real and finite
support places have their single sum relation.  Since a finite support
place above $2$ is present, restriction of the Brauer term to the real
places is surjective and has kernel of dimension $s_n-1$.  Thus the
degree-two kernel of restriction to infinity is bounded.  Restriction
is an isomorphism in higher degrees by global duality.  The degree-zero
restriction is injective.  The cohomology of the positive fiber is
therefore bounded in every degree except one; its degree-one dimension
is $d_n+O_S(1)$.

Apply Lemma~\ref{lem:cy:positive-model} modulo $2$.  Derived specialization to the
$n$th layer is specialization modulo $t^{2^n}$.  For a perfect complex
over $\mathbb F _2[[t]]$, the universal-coefficient exact sequence shows
that a rank-$r$ cohomology module contributes $r2^n+O(1)$ to these
dimensions; the torsion terms contribute only $O(1)$.  It follows that
after inverting $t$ the trivial-constituent fiber has dimension one in
degree one and vanishes in every other degree.  Extensions preserve
this concentration, so the fiber for $E[2]$ has dimension two in degree
one.  A minimal finite free complex over the local ring $\mathcal O[G]$
has zero residual differentials.  The concentration just proved forces
all its terms except the rank-two term in degree one to vanish.
No denominator is introduced by this exact freeness statement.
\end{proof}

\begin{proof}[Proof of Theorem~\ref{thm:cyclotomic-bound}]
Apply Lemma~\ref{lem:cy:residual-free} to the single discriminant $h$.
The fixed-lattice construction and the free complex give
$2^{C_{\mathrm{int}}}u_h\in\mathcal O$ with $C_{\mathrm{int}}$ depending
only on $E$. Equation~\eqref{eq:cy:rational-integrality} also gives
$u_h\in\Lambda[1/2]$. Since
$\mathcal O\cap\Lambda[1/2]=\Lambda$, it follows that
$v(u_h(0))\ge-C_{\mathrm{int}}$. Analytic rank zero supplies Selmer
corank zero by the forward theorem of Gross--Zagier--Kolyvagin.
Proposition~\ref{prop:cy:specialization} now gives
$D(h)\ge-C_{\mathrm{int}}-C_{\mathrm{loc}}$.
All constants come from the fixed form, the fixed real period line,
or the finitely many local twisting classes at $2N$.
\end{proof}

\subsection{Transfer to the base}

It remains to prove the missing-vertex theorem. Residual freeness gives
a single integral family in the reducible case. In the irreducible
case we instead need a bounded clearing factor after eliminating the
new local blocks. The bound on $v_2(q^2-1)$ enters precisely here: the
next lemma controls the negative Laurent powers in those inverse
blocks at each fixed precision.

\begin{lemma}\label{lem:cy:cyclotomic-poles}
Suppose that $q$ is an added prime and $v_2(q^2-1)\le L$.  At
augmentation $G\to1$, the inverse singular Frobenius block contributed
by $q$ has, modulo each fixed $2^M$, negative Laurent support bounded
in terms of $L$ and $M$.  This bound is independent of $q$.
\end{lemma}

\begin{proof}
Identify $\Gamma$ with $1+4\Ztwo$ using a fixed generator.  If
$\gamma_q=(1+t)^{\alpha_q}$ is the image of arithmetic Frobenius at
$q$, then
\[
 v_2(\alpha_q)=v_2(q^2-1)-3
\]
when the generator is $5$; changing the generator multiplies
$\alpha_q$ by a unit.  In particular $\alpha_q\ne0$.  The singular
quotient is represented by a block of bounded size whose determinant
is the Euler polynomial in Equation~\eqref{eq:cy:local-euler}.
Modulo $2$, with $z=(1+t)^{\alpha_q}$, it is
\[
 1+\overline{a_q}z+z^2.
\]
If $\overline{a_q}=1$, its constant coefficient is $1$.  If
$\overline{a_q}=0$, it is $(1+z)^2$, whose $t$-order is
$2^{v_2(\alpha_q)+1}\le2^{L-2}$.  Thus the inverse determinant has
uniformly bounded pole order modulo $2$.  Lift its inverse modulo
$2^M$ by a geometric expansion of length at most $M$.  Each term
contains only a bounded number of inverse residual factors, and all
numerators have nonnegative powers of $t$.  The adjugate formula gives
the same bound for every entry of the inverse block.
\end{proof}

\begin{proof}[Proof of Theorem~\ref{thm:cyclotomic-interpolation}]
At every nonzero vertex, analytic rank zero gives
$c_2(E^{(h_x)})=0$ by the forward rank-zero theorem. Thus
Proposition~\ref{prop:cy:specialization} applies at every vertex.

First assume $E[2]$ is an extension of trivial modules.
Lemma~\ref{lem:cy:residual-free} and the fixed-lattice construction give
$2^{C_{\mathrm{int}}}u\in\mathcal O[G]$.  Moreover, the character coordinates of $u$
are in $\Lambda[1/2]$.  Fourier inversion, performed rationally, shows
$u\in\Lambda[1/2][G]$; intersecting coefficientwise with
$\mathcal O[G]$ gives
\begin{equation}\label{eq:cy:universal-integrality}
 2^{C_{\mathrm{int}}}u\in\Lambda[G].
\end{equation}
The coefficients at $t=0$ therefore obey the bounded-degree Boolean
congruences of Lemma~\ref{lem:h:boolean-coefficients}. For any fixed
precision and sufficiently large $b$, there is a nonzero $x$ for which
$u_x(0)$ agrees with $u_0(0)$ to that precision.  The fixed multiplier
in Equation~\eqref{eq:cy:universal-integrality} affects the required
precision but has no dependence on $b$.

Now suppose $E[2]$ is irreducible.  At augmentation, represent the
positive complex as the old complex, with ramification set excluding
$\mathcal Q$, extended by the direct sum of singular blocks at the
new primes.  The differential is upper triangular with the old
complex as subcomplex. The finite-model argument of
Lemma~\ref{lem:h:finite-models} lifts this model over $R$: stabilize the
finite free models and lift their graded changes of basis. These lifts
are invertible because the augmentation ideal lies in the Jacobson
radical. Cancel all additional lifted unit disks before retaining the
old-plus-singular grading; their augmented cross terms are zero, so the
remaining augmentation is exactly the prescribed block model. Thus no
stabilization disk enlarges the old complex.
Lemma~\ref{lem:h:bounded-poles} then applies to this block model.  All the new
blocks are invertible over $\mathcal O[G]$: modulo its maximal ideal
their Frobenius determinants are nonzero in $\Omega$.
Schur elimination leaves a complex of bounded size. Its augmentation
is the fixed old complex, and its entries have bounded negative
Laurent support modulo each fixed $2^M$. Under the determinant-line
identification, write
\[
 u=\eta u_{\mathrm{bd}},\qquad \eta\in\mathcal O[G]^\times,
\]
where $u_{\mathrm{bd}}$ is the coordinate on this bounded complex and
$\eta$ is the determinant factor of the eliminated blocks. No uniform
Laurent bound on $\eta$ is needed.

The pole bound remains uniform as primes are added.  At augmentation the inverse new block is a direct
sum of the inverses in Lemma~\ref{lem:cy:cyclotomic-poles}.  At any
character its perturbation has entries divisible by $2$ and no
negative powers of $t$.  Modulo $2^M$, the inverse uses a geometric
series of length at most $M$, independently of the number of blocks.
Each product therefore uses a bounded number of inverse entries.
Sums of an arbitrary number of such terms cannot decrease the
nonarchimedean valuation or the common lower Laurent bound.  This is
the required bound for Schur elimination.

At $t=0$ the old positive complex has rational cohomology only in
degree one, of dimension two, by
Proposition~\ref{prop:cy:specialization} at $h_0$.  Its finite torsion
and all old local data are fixed.  Choose boundary pivots there,
leaving those two free coordinates.  Their sizes and valuations may
depend on $h_0$ but not on $b$.  Apply the determinant clearing argument
of Lemma~\ref{lem:h:pivot-clearing} to the wedge of the two
bounded-integral cycles $a,2z_S$ on the remaining complex. Its allowance
for an arbitrary factor $\eta$ puts the eliminated-block factor into
the numerator alone. It yields
\begin{equation}\label{eq:cy:cleared-determinant}
 d=fu\in\mathcal O[G],\qquad f\in\mathcal O[G],\qquad
 f_0\in\Lambda,\quad f_0(0)\ne0.
\end{equation}
For each fixed $M$, the negative Laurent support of $f_x\bmod2^M$
is bounded independently of $b$ and $x$.  The clearing factor is a
polynomial of bounded degree in the remaining bounded-size matrix,
with the pivot denominators cleared by adjugates. Neither the
eliminated-block unit, the zeta class nor $d$ is required to have a uniform bound on
its negative Laurent support.

Equation~\eqref{eq:cy:rational-integrality} ensures that every $u_x$
has no negative $t$-powers.  Thus all hypotheses of
Lemma~\ref{lem:h:boolean-transfer} hold for
Equation~\eqref{eq:cy:cleared-determinant}.  In particular, for each
fixed precision and sufficiently large $b$, some nonzero $x$ satisfies
the corresponding central congruence between $u_x(0)$ and $u_0(0)$.
The dependence on $h_0$ enters the required precision and dimension,
through $v(f_0(0))$, and not the eventual valuation comparison.

In either residual case let $C$ bound the absolute error in
Equation~\eqref{eq:cy:determinant-center}.  The nonzero vertices have
$v(u_x(0))\le B+C$.  If $L(E^{(h_0)},1)=0$, the central coordinate at
zero is zero, and a congruence of precision larger than $B+C$ is
impossible.  Hence the central value is nonzero.  If
$D(h_0)>B+2C$, then
$v(u_0(0))>B+C$, giving the same contradiction.  This proves
Equation~\eqref{eq:cy:transferred-bound} with $C_E=2C$.
Only the specialization error determines this final loss; all pivot
and pole bounds determine how large a cube is needed.
\end{proof}

\section{Heegner points and transfer}
\label{sec:heegner-interpolation}

We prove the Heegner comparison, the genus and ring bounds, and the
interpolation criterion stated in Section~\ref{sec:transfer-criteria}.
The point and its paired functional define one determinant coordinate.
Its central value measures the point index after the Selmer and local
contributions have been subtracted. Local switches then control its
divisors. An outer limit supplies the uniform ring bound when residual
irreducibility is unavailable; a universal binary family supplies
interpolation.

\subsection{Height and base Selmer comparisons}

Retain the genus and ring data, character sums $P_c(e)$, and lattice
indices of Section~\ref{sec:transfer-criteria}. We first justify the two
comparisons used there.

\begin{proof}[Proof of Proposition~\ref{prop:heegner-comparison}]
The explicit Gross--Zagier formula for ring classes
\citep[Theorem~1.1, pp.~2524--2525]{CaiShuTian2014} applies. Automorphic induction of the
quadratic character gives the product in Equation~\eqref{eq:h:height}. Discriminant and conductor are prime to $N$, and the
classical split Heegner hypothesis holds; consequently the shared-prime
corrections in the general formula are absent. The unit-index convention
is constant in the fields retained here. With a character sum, solving
the formula for the height gives exactly $c\sqrt{|k|}$, as in
Equation~\eqref{eq:h:height}. The forward Gross--Zagier--Kolyvagin theorem
gives the rank and finiteness assertion
\citep{GrossZagier1986,Kolyvagin1990}.

For Equation~\eqref{eq:h:index-comparison}, both free points are compared
with the free line on $E^{(h)}(\Q)$. Its regulator cancels from the ratio
of their heights, since absolute canonical height is unchanged by twist
identification. This height invariance follows from the canonical-height
limit and the degree formula for heights under morphisms
\citep[Chapter~IV, Lemmas~4.1 and~4.5--4.6, and Remark~4.12]{MilneEC2021}.
The remaining variable factor is
$\sqrt{|hk|}L(E^{(hk)},1)$; its infinity sign is that of $h_*$. Period
and lattice changes have bounded $2$-valuation. If the rank-one line is
instead that of the fixed curve $E^{(h_*)}$, its regulator is fixed and
the same calculation gives Equation~\eqref{eq:h:reverse-comparison}.
\end{proof}

\begin{proof}[Proof of Lemma~\ref{lem:h:base-selmer-bound}]
Restriction, twisting, and addition or subtraction of conjugates give
maps in both directions between the Selmer group over $K$ and
\[
 \Sel_{2^\infty}(E^{(h_0)}/\Q)
 \oplus\Sel_{2^\infty}(E^{(h_0k)}/\Q)
\]
whose composites are multiplication by $2$. They induce isogenies on
divisible subgroups. On the finite quotients the kernel of either induced
map is killed by $2$: subtract an element of the divisible subgroup to
lift a divisible image, and then use the composite. The generator
dimensions are uniformly bounded. Indeed the mod-$2$ Selmer group over
$K$ is bounded by Kummer cohomology over $K(E[2])$; this is a quadratic
extension of a fixed field with a bounded number of ramified primes.
Kummer theory bounds the generator number of the corresponding pro-$2$
group over the fixed field, and the index-two subgroup bound preserves
a uniform bound. The fixed local support includes that of $h_0$.
The finite quotient on the displayed product is bounded by hypothesis;
therefore its counterpart over $K$ has bounded length as well. For a
finite collection of fields, finiteness of each quotient is enough.
\end{proof}

\subsection{Integral models and ring-class deformation}
\label{rl:section}\label{sec:h:finite-models}

The determinant argument uses characters of increasing conductor.
We construct their cochain limits together with maps to the original
cochains. These maps will let us detect limiting Selmer classes by
Galois evaluation and choose local switches. Later we retain the same
maps through the outer limit needed for the ring bound.

Write $T=T_2E$ and $\Lambda=\mathbf Z_2[[t]]$. For a quadratic
squareclass $h$, put $T^h=T\otimes\chi_h$.  The imaginary quadratic field
$K$ satisfies the split Heegner hypothesis at $2N$.  We omit the two fields
with extra units and, when $E$ has complex multiplication, its CM field.
A quadratic character $\chi$ comes from a quadratic twist over $\mathbf Q$.
We use
\[
 A_i=(\mathbf Z/2^i\mathbf Z)[t]/(t^i).
\]
The same constructions work over $A_i[G]$ and $\Lambda[G]$, where $G$ is a
fixed finite elementary abelian $2$-group and the additional character has
values in $G$.  Statements uniform in the number of generators of $G$ will
be identified separately; finite-diagram compactness by itself is a statement
for fixed $G$.

\subsubsection{Finite models with their cochain maps}

The general finite-model construction is Lemma~\ref{rl:finite-model}.
We apply it to the following arithmetic cochains and retain its actual
contractions throughout.

For a profinite group and a finite free coefficient module over an Artin local ring $A$,
continuous inhomogeneous cochains are flat: locally constant functions are
the filtered union of the modules of functions on finite clopen partitions,
and these modules are finite free.  Their reduction is the cochain complex
with reduced coefficients.  The same observations apply to restrictions,
mapping cones, and local conditions given by finite free complexes.  In the
applications below the groups are Galois groups of totally imaginary number
fields or nonarchimedean local fields.  Their $2$-cohomological dimension and
the usual finiteness theorems give the bounded residual cohomology required
in Lemma~\ref{rl:finite-model}.

For an inner limit, the quadratic field $K$, the active conductor support,
the finite group $G$, and the number of derivative primes are fixed.
The residual representation is $E[2]$: all the scalar characters and
group-like coefficients become trivial in the residue field. Put
$F=K(E[2])$. The number of places of $F$ in the permitted set $S_i$ is
bounded even when its auxiliary primes move. Kummer theory bounds
$H^1(G_{F,S_i},\mathbf F_2)$ by the $S_i$-units modulo squares and
$\operatorname{Cl}(F)[2]$; their dimensions are bounded by the unit
rank, $\#S_i$, and the fixed class-group dimension. Hochschild--Serre
for the fixed group $\operatorname{Gal}(F/K)$ then bounds the residual
global $H^1$ for $E[2]$. The global Euler characteristic bounds $H^2$,
and $H^0$ is already bounded. Local cohomology and the specified local
conditions have bounded residual dimensions, since the local degrees
of $K$ are at most two and the number of places is bounded. Thus the
target Selmer cones over $K$ have uniformly bounded finite-model ranks
in this inner limit. This bound may depend on the fixed active support;
the stronger outer bound is proved in Lemma~\ref{rl:residual-bound}.
No rank bound is asserted for cochains over the growing ring-class
fields, or for the induced modules in their Shapiro descriptions.

\begin{lemma}[Limits of finite diagrams]\label{rl:finite-limits}
Suppose diagrams of bounded finite free complexes over $A_i$, or
over $A_i[G]$, have a common bounded cohomological degree interval
and bounded ranks in every degree. Include the chain homotopies
and inverses for every asserted identity or equivalence. Along a
cofinal subsequence, or a nonprincipal ultrafilter, their matrices
converge to the same diagram of perfect complexes over $\Lambda$
or $\Lambda[G]$.

Moreover, the contractions into the original cochains may be retained for
the evaluation construction in Lemma~\ref{rl:evaluation}.  Adding finitely
many new objects or maps does not require changing already chosen models
for the old objects. Finite-stage arithmetic operations with unbounded
source complexes are permitted when only finitely many of their outputs,
and the witnesses to their identities, are retained in the bounded
target complexes.
\end{lemma}

\begin{proof}
There are finitely many possible graded ranks.  Fix them on an
ultrafilter-large set.  At a fixed precision $j$, all matrix entries belong
to the finite ring $A_j$ (or $A_j[G]$), so have unique ultrafilter limits.
These limits are compatible as $j$ varies.  Matrix identities, including
the homotopy-inverse identities, are finite polynomial identities and pass
to the limit.  New models and maps may be mapped to old models by their
already fixed contractions.  Passing to any slower cofinal precision has
the same old limits.  For the all-sequence evaluation statements below,
fix this ultrafilter once and for all, and keep the original contraction
at every stage.  Later arithmetic diagrams must extend these same
stages.  Restriction to an ultrafilter-large set and passage to slower
cofinal coefficient precision preserve every old evaluation on every
old group sequence. Thus later operations use the same stage diagrams
and ultrafilter, preserving the arithmetic maps already constructed.

Here is the last assertion in cochain terms. Let $D_i$ be an original
target cochain complex, with contraction
$(\iota_i,\pi_i,h_i)$ onto its bounded finite model $C_i$.
Perform transfer and scalar weighting at stage $i$ before applying
this contraction. A resulting cycle $z_i'\in D_i^1$ gives
$z_i=\pi_i z_i'\in C_i^1$, and
\[
 z_i'-\iota_i z_i=d(h_i z_i').
\]
Thus its relation to the original arithmetic cocycle is retained by
the original contraction. If an identity in the target is
$u_i'-v_i'=d w_i'$, retain the vectors
$\pi_i u_i',\pi_i v_i',\pi_i w_i'$; the corresponding identity is
still exact. For identities involving maps between target complexes,
use the transported maps $\pi_i'f_i\iota_i$ and their contraction
homotopies, so that the same assertion includes the required boundary
corrections. Local lifts, closed degree-two functionals, and witnesses
to specialization, duality and switch identities are treated in the
same way. They are finitely many vectors or matrices of bounded size.
Their entries, as well as the differentials, therefore have simultaneous
limits. Neither a free model of the arithmetic source nor the entire
transfer map from that source is included in the limiting diagram.
Any precision lost in a finite-stage calculation is accounted for
before retaining these vectors; passing to the same slower cofinal
precision for the whole diagram preserves all its identities.
\end{proof}

The finite models now retain both the original cochains and their
comparison maps. We next specify the local conditions on those cochains.
The point is to retain the actual Kummer images at every place dividing $2N$;
the determinant comparisons will use their integral orthogonality, not
just the rational Selmer ranks.

\subsubsection{Integral local conditions and duality}

\begin{lemma}[Integral Kummer local condition]\label{rl:kummer}
\label{lem:h:duality}
Let $F$ be a finite extension of a nonarchimedean local field of
characteristic zero, and let $B/F$ be any elliptic curve.  Put
$M_F=B(F)^\wedge_2=\varprojlim_n B(F)/2^nB(F)$.  There is a perfect local
condition
\[
 U_F=M_F[-1]\longrightarrow R\Gamma(F,T_2B).
\]
For every $n\geq1$, derived reduction has
\[
 H^0(U_F\otimes^{\mathbf L}\mathbf Z/2^n)=B(F)[2^n],\qquad
 H^1(U_F\otimes^{\mathbf L}\mathbf Z/2^n)=B(F)/2^nB(F).
\]
Its map in degree zero is the isomorphism onto
$H^0(F,B[2^n])$, and its map in degree one is the finite Kummer injection.
The condition is its own exact orthogonal complement under local Tate
duality, for every reduction type and with no hypothesis on $B[2]$.
\end{lemma}

\begin{proof}
The local analytic description of $B(F)$ shows that $M_F$ is a finitely
generated $\mathbf Z_2$-module.  Its torsion is the finite group
$B(F)[2^\infty]$.  Consequently $M_F$ has a free resolution of length at
most one.  Also $H^0(F,T_2B)=0$, since $B(F)[2^\infty]$ has bounded
exponent.  The integral Kummer injection into $H^1(F,T_2B)$ therefore
defines the displayed morphism using the truncation in degrees at most
one of $R\Gamma(F,T_2B)$.

Tensor the free resolution of $M_F$ with $\mathbf Z/2^n$.  Its two
cohomology groups are $M_F[2^n]$ and $M_F/2^nM_F$. The inverse limits
of the ordinary Kummer sequences identify their maps as asserted.

For orthogonality, compose cup product and the Weil pairing with the local
invariant map.  A morphism
\[
 M_F[-1]\otimes^{\mathbf L}M_F[-1]
       \longrightarrow\mathbf Z_2[-2]
\]
is determined by its homomorphism $M_F\otimes M_F\to\mathbf Z_2$:
the possible $\operatorname{Tor}_1$ is in degree one and gives no extra
degree-zero morphism to this target.  The indicated homomorphism vanishes,
because the finite Kummer image annihilates itself at every level by local
Tate duality.  Choose a nullhomotopy in free models.  The induced morphism
\[
 R\Gamma(F,T_2B)/U_F
       \longrightarrow R\operatorname{Hom}_{\mathbf Z_2}
                         (U_F,\mathbf Z_2)[-2]
\]
is a quasi-isomorphism.  This can be checked after reduction modulo $2$:
degree one is the perfect pairing of the finite Kummer quotient with the
Kummer image, and degree two is dual to $B(F)[2]$.  Its perfect cone is
therefore zero by derived Nakayama.  This proves exact orthogonality and
also its compatibility with derived reduction.  This is precisely the
finite-module duality input of \cite[Proposition~5.2.4]{Nekovar2006}.
\end{proof}

Where the scalar deformation is trivial, extend this local condition by
derived scalar extension to $A_i$ or $A_i[G]$.  At an unramified good prime
use the inflated residue-field complex
\[
 U_v=[\,M\xrightarrow{\operatorname{Fr}_v-1}M\,]
 \quad\text{in degrees }0,1.
\]
Its quotient in full local cochains is the singular local complex.
The full local condition is the identity on local cochains, and its
orthogonal strict condition is zero.  At a derivative prime the
conditions used later are the following specified truncations instead.
In its split local model write
\[
 L_v=D^0\oplus(F_v\oplus S_v)[-1]\oplus D^2[-2],
\]
where the differential is zero, $D^0$ is dual to $D^2$, and the finite
and transverse planes $F_v,S_v$ are each isotropic.  Put
\[
 U_{\rm str}=D^0,\qquad
 U_{\rm rel}=D^0\oplus(F_v\oplus S_v)[-1],\qquad
 U_F=D^0\oplus F_v[-1],\qquad U_S=D^0\oplus S_v[-1].
\]
Then $U_{\rm str}^{\perp}=U_{\rm rel}$ and each plane condition is
its own orthogonal complement.  Every displayed condition maps
isomorphically onto local $H^0$ and injectively into local $H^1$.
Thus they satisfy precisely the detection hypotheses used below.
Their finite arithmetic realization and the isotropic-plane
nullhomotopies are part of the derivative construction; these
truncated conditions are not identified with the full local condition.

For a finite permitted set $S$ the Selmer complex is, by definition,
\begin{equation}\label{rl:selmer-definition}
 C_S=\operatorname{Cone}\left(
 R\Gamma(G_{K,S},M)\oplus\bigoplus_{v\in S}U_v
 \longrightarrow\bigoplus_{v\in S}R\Gamma(K_v,M)
 \right)[-1].
\end{equation}
At $2N$ take the Kummer conditions just constructed.  Restriction and
localization triangles are the triangles of this cone.  The finite rings
$A_i$ and $A_i[G]$ are Frobenius coefficient rings; alternatively the
duality maps may be checked after reduction to their residue field.
Exact local orthogonality and global duality therefore give
\begin{equation}\label{rl:global-duality}
 C_S\simeq R\operatorname{Hom}(C_S^{\perp},A_i)[-3],
\end{equation}
with the appropriate dual representation.  This application uses the
orthogonality hypotheses of
\cite[Theorem~6.3.4]{Nekovar2006}; the theorem is not being used to assert
the existence of a varying-prime limit.  Complex places contribute zero
Tate complexes.  If the character is inverted by the nontrivial
automorphism of $K/\mathbf Q$, conjugate transport followed by the Weil
pairing identifies the two representations over the same scalar ring.
All maps and nullhomotopies in this paragraph are finite diagrams to
which Lemma~\ref{rl:finite-limits} applies.

\subsection{The moving conductor character}
\label{sec:h:conductor-direction}

The conductor character must preserve the local conditions at $2N$
while making the new local terms invertible after reduction modulo $2$
and inversion of $t$. The following proposition supplies those local
values. Its separate bounds on Tate Frobenius control the contribution
of the moving prime when we later specialize to $t=0$.

\begin{proposition}[Ring-class characters with prescribed local values]
\label{rl:characters}\label{lem:h:conductor-primes}
Fix finitely many split primes of $K$ above odd rational primes, outside
$2N$, and orient one prime in each conjugate pair.  There are distinct
split rational primes $r_i$ avoiding all the fixed data and characters
\[
 \xi_i:\operatorname{Gal}(H_{r_i}/K)\longrightarrow
                 \mathbf Z/2^{n_i}\mathbf Z,
 \qquad n_i\longrightarrow\infty,
\]
such that $\xi_i$ is surjective on the inertia quotient at $r_i$, is zero
at the primes above $2N$, and has odd value at each prescribed oriented
prime.  We can require $r_i\to1$ in $\mathbf Z_2$ and
\begin{equation}\label{eq:h:r-avoidance}
 v_2(2+a_{r_i}(E)),\quad v_2(2-a_{r_i}(E)),\quad
 v_2\det(1\pm\operatorname{Fr}_{r_i}\mid T)\ \leq c_E.
\end{equation}
The constant $c_E$ is independent of the list and of $K$.
\end{proposition}

\begin{proof}
We give the compatibility argument before imposing the Tate condition.
Choose prime ideals $\mathfrak p_1,\ldots,\mathfrak p_s$ with pairwise
disjoint conjugate pairs, including the oriented prescribed primes and
the primes above $2N$, and enough additional split prime ideals to
generate $\operatorname{Cl}(K)$.  Such additional ideals may avoid any
fixed finite set.  Put weight $w_j=1$ at a prescribed odd prime and
$w_j=0$ at every additional prime and every prime above $2N$.
The conjugate weight is $-w_j$.  Let
\[
 R=\ker\bigl(\mathbf Z^s\longrightarrow\operatorname{Cl}(K)\bigr),
\]
choose a basis $m^{(a)}$ of $R$, and choose $\beta_a\in K^\times$ with
$(\beta_a)=\prod_j\mathfrak p_j^{m_j^{(a)}}$.

Work first over $F_\infty=K(\mu_{2^\infty})$ and adjoin $2^n$-th roots
of the finitely many $\beta_a/\overline\beta_a$.  Prescribe on these
roots the multipliers
\begin{equation}\label{rl:kummer-prescription}
 \zeta_{2^n}^{\,-\sum_j m_j^{(a)}w_j}.
\end{equation}
The minus sign fixes the choice of connecting map in the ring-class
sequence and could be reversed throughout.  These prescriptions respect
every Kummer relation.  Indeed, if
\[
 \prod_a(\beta_a/\overline\beta_a)^{b_a}
                         \in(F_\infty^\times)^{2^n},
\]
take its valuation at a prime above each $\mathfrak p_j$ of nonzero
weight.  These odd primes are unramified in $F_\infty/K$.  Thus
$2^n$ divides $\sum_a b_am_j^{(a)}$ for each such $j$, proving the
required compatibility of~\eqref{rl:kummer-prescription}.  Primes over
$2$ have weight zero and impose no division by a ramification index.
Units impose no additional weighted relation.  Kummer duality supplies
the prescribed automorphism of the radical extension over $F_\infty$.

At a rational prime $r$ split in a sufficiently large finite cyclotomic
subfield, the power-residue character of $\mathbf F_r^\times$ is
surjective onto $\mu_{2^n}$.  The split ring-class exact sequence is
\[
 1\longrightarrow
  (\mathcal O_K/r\mathcal O_K)^\times/(\mathbf Z/r\mathbf Z)^\times
 \longrightarrow\operatorname{Pic}(\mathbf Z+r\mathcal O_K)
 \longrightarrow\operatorname{Cl}(K)\longrightarrow1.
\]
The excluded extra-unit fields are the only source of a unit
correction.  The first group is $\mathbf F_r^\times$, with a principal
relation measured by $\beta/\overline\beta$ at the oriented prime over
$r$.  Consequently~\eqref{rl:kummer-prescription} says exactly that its
primitive power-residue character extends to the class group with the
specified values $w_j$.  More explicitly, the inertia group and the
chosen ideal classes present the ring-class group, and the basis of $R$
lists its additional relations.  We have imposed those relations on the
inertia character itself; we have not divided by a class-group order
to extend a previously fixed character.  Thus this construction introduces
no denominator, even when the class number has a large $2$-part.  The
class group and the number of prescribed primes change the finite
Chebotarev extension, but not the integrality of the character.
Since the ring-class extension is unramified at $2N$, value zero on
Frobenius means that the resulting local character is trivial there.
Conjugation acts by inversion.

We now show that imposing these finitely many radical conditions cannot
force either forbidden Tate trace.  In the non-CM case the image of
$G_{K(\mu_{2^\infty})}$ is an open subgroup of
$\operatorname{SL}_2(\mathbf Z_2)$, by the fixed-prime open-image theorem
and the cyclotomic determinant identity
\cite[Chapter~IV, \S2.2, theorem, p.~IV-11, and Chapter~I, \S1.2,
Example~2, p.~I-4]{Serre1968}: the image is the intersection of the
open image of $G_K$ with the determinant-one subgroup.
Its index in the corresponding fixed-$E$ image changes
by at most two when $K$ changes.  Its commutator subgroup contains a
fixed open subgroup $U_E$ of $\operatorname{SL}_2(\mathbf Z_2)$.
One may see uniformity by starting with a principal congruence subgroup
in the fixed image: every index-two subgroup contains its squares, and
the commutators of the resulting fixed congruence subgroup are open.
Because the radical extension over $F_\infty$ is abelian, these
commutators act trivially on it.  Hence every permitted Tate coset
contains a translate of $U_E$.

For completeness, in the CM case let $M$ be the CM field.  The image over
$KM(\mu_{2^\infty})$ contains a fixed open subgroup of the norm-one
Cartan; the index changes by a bounded factor, as $[KM:\mathbf Q]\leq4$.
Choose $\gamma\in G_{KM}$ with
$\chi_{\rm cyc}(\gamma)=5^4=625$.  This is possible uniformly: the
cyclotomic image of $G_{KM}$ has index at most four in
$\mathbf Z_2^\times$, so contains fourth powers.  On the radical Galois
group, conjugation by $\gamma$ is raising to $625$, because all the
radicands lie in $K$.  On the Tate image over $KM$, conjugation is
trivial.  Therefore
\[
 \gamma u\gamma^{-1}u^{-625},\qquad
 u\in G_{KM(\mu_{2^\infty})},
\]
fixes every prescribed radical and acts on $T$ as $\rho_T(u)^{-624}$.
Raising to the fixed nonzero power $624$ carries an open subgroup of the
one-dimensional norm-one Cartan onto an open subgroup.  This proves the
required uniform kernel statement also in the CM case.  The open Cartan
image used here is the CM part of the usual $\ell$-adic image theorem
\cite[Chapter~IV, \S2.2, Remark, p.~IV-13]{Serre1968}.

On a coset of an open subgroup of $\operatorname{SL}_2$, trace is not
identically $2$ or $-2$.  The same holds on a coset of an open norm-one
Cartan subgroup; on the other Cartan-normalizer component trace is zero.
The relevant determinant-one sets are compact and have finitely many
cosets modulo the fixed open subgroup.  Choose in each coset an element
whose trace is neither forbidden value.  The largest of the finitely
many resulting valuations is a constant depending only on $E$.

Finally restrict the prescribed radical automorphism and such a Tate
lift to a finite Galois extension, including as many roots of unity as
required.  Chebotarev gives infinitely many rational primes with this
Frobenius class; choose the place over the prime to realize the chosen
representative.  They split in $K$ and have arbitrarily high prescribed
congruence to $1$.  Since
$\det(1\pm\operatorname{Fr}_r)=1\pm a_r+r$, a congruence precision
larger than the fixed trace bound gives all of~\eqref{eq:h:r-avoidance}.
All these conditions concern a finite extension at each stage.
\end{proof}

Choose $n_i$ large enough that $1+t$ has order dividing $2^{n_i}$ in
$A_i^\times$, and put $\Psi_i=(1+t)^{\xi_i}$.  For example $n_i\geq2i$
is sufficient. Orient the inertia generator to act by $1+t$;
the prescribed split Frobenius exponents remain odd in the limit.
Use global cochains unramified outside $2N$, $r_i$, the active conductor
primes and the derivative primes already retained. At $2N$ use the
Kummer conditions above and elsewhere outside the conductors use the
unramified condition. Reduction modulo $t$, with the unramified
condition restored at $r_i$, gives the ordinary fixed-field complex;
retain its comparison map at every finite stage.

\begin{lemma}[Acyclic relaxed local complexes]\label{rl:local-acyclic}
\label{lem:h:relaxed-local}
For a primitive character constituent, let $C$ denote the limiting
Selmer complex with full conditions at $r$ and at its active good
conductor primes, and the prescribed Kummer, unramified or derivative
conditions elsewhere.  The limiting full
local complexes are acyclic over every height-one residue field of
$\Lambda$ of characteristic zero.  At $(2)$, with residue field
$\Omega=\mathbf F_2((t))$, the $r$-complex is acyclic, as are the
complexes at active split primes with odd prescribed Frobenius exponent
of the deformation character and
at active primes $q$ with $t_q=0$. At each of these acyclic relaxed
places, the full and orthogonal strict conditions agree on the
indicated test. If every relaxed place is acyclic at that test and
the remaining local conditions are self-dual, global duality makes
the Selmer complex self-dual. The Selmer complex
has no global invariants on these tests.
\end{lemma}

\begin{proof}
Use tame inertia and unramified Frobenius, or the Hochschild--Serre
complex for these two procyclic directions.  At $r$, inertia minus one
is multiplication by $t$ times a unit.  It is invertible on all the
indicated fields except the characteristic-zero fiber $t=0$.  There,
$r\to1$, and the two remaining Frobenius differentials have nonzero
determinant by~\eqref{eq:h:r-avoidance}, including either sign of the
quadratic twist.  The bound also ensures that the integral local
torsion at this central fiber has bounded length.

At an active odd conductor prime, inertia of the quadratic character
acts by $-1$.  Its differential is thus $-2$ and is invertible in
characteristic zero.  In characteristic two the quadratic character
disappears.  At a split prime the Frobenius action is $zB$, where
$B\in\operatorname{GL}_2(\mathbf F_2)$ and
$z=(1+t)^a$ has odd $a\in\mathbf Z_2$.  Its coefficient of $t$ is one,
so $z$ is nonconstant.  No nonzero polynomial over the finite constant
field vanishes at $z$.  Thus both Frobenius complexes, with operators
$zB-1$ and $q^{-1}zB-1$, are invertible over $\Omega$; here $q=1$ in
the residue field.  No irreducibility assumption is used.

If $t_q=0$, the rational residual Frobenius $B$ has order three.
At a split place the operator is again $zB-1$, invertible whether
$z$ is nonconstant or $z=1$. At an inert place the ring-class
Frobenius is trivial, since the ideal $(q)$ is principal, and
the Tate Frobenius is $B^2$, which also has no eigenvalue one.
These observations give the additional $t_q=0$ assertion.

The orthogonal full and strict conditions agree at each place where
full local cochains are acyclic. When this holds at every relaxed
place and the remaining local conditions are self-dual,
Equation~\eqref{rl:global-duality} gives global self-duality.
Global invariants in characteristic two are excluded by the nontrivial
$r$-inertia scalar.  In characteristic zero the same argument works
unless the scalar is trivial; in that case $T\otimes\mathbf Q_2$ over
$G_K$ is absolutely irreducible.  For non-CM curves this follows from
the open image, and for CM curves from the open Cartan and its
nontrivial normalizer component, since $K$ is not the CM field.
Lemma~\ref{rl:evaluation} below realizes these actions on the limiting
complex as well. On the self-dual tests, duality pairs degrees
one and two.
\end{proof}

\subsection{The Heegner determinant and its central value}\label{sec:h:specialization}

The point index is not itself the quantity transferred across a growing
family. Each new conductor prime changes the local Selmer terms as well
as the point. We combine these changes in a determinant coordinate.
Its central valuation will be twice the point index, less the finite
Selmer length and the conductor contribution. This formula reduces the
index estimates to lower bounds for the determinant coordinate.

Retain a primitive genus or ring instance from Section~\ref{sec:transfer-criteria},
with point $P_c(e)$ over $K$. Use the conductor characters just
constructed, taking $n_i\ge2i$, and the representations $T\chi_e\Psi_i$.
Retain the same auxiliary primes and the same fixed inner ultrafilter,
now denoted $\mathcal U_{\mathrm{in}}$.

The Selmer complex $C$ is a coefficientwise limit of the corresponding
finite free cochain models along $\mathcal U_{\mathrm{in}}$. It has the integral Kummer condition at
$2N$, the full condition at $r_i$ and the active good conductor primes,
and the unramified condition at all other permitted finite places.
Here the full condition means that no restriction is imposed on local
cochains. Section~\ref{sec:h:finite-models} and Lemma~\ref{rl:local-acyclic}
constructed the local models, their duality maps, and the comparison
at $t=0$. We use these concrete complexes and maps;
cohomologically equivalent complexes without the comparison maps would
not suffice for the later prime choices. Over
$\operatorname{Frac}(\Lambda)$, the modified full local complexes
are acyclic, so the full conditions agree with their orthogonal strict
conditions. Global duality, followed by conjugate transport and the
Weil pairing, therefore gives a perfect pairing
\[
 H^1(C\otimes_\Lambda\operatorname{Frac}(\Lambda))\times H^2(C\otimes_\Lambda\operatorname{Frac}(\Lambda))
 \longrightarrow\operatorname{Frac}(\Lambda).
\]
Here conjugation in $K/\Q$ inverts the ring-class character. These are
the pairing and identification supplied by
Equation~\eqref{rl:global-duality} and Lemma~\ref{rl:local-acyclic}.

Fix the single integer multiplier $L$ of
Lemma~\ref{lem:h:derivative-classes}. It depends on the curve,
parametrization and fixed local types, and clears descent and local
Kummer ambiguities for the complete conductor. In particular its
valuation does not grow with the number of primes. The following
construction and formulas include that multiplier throughout.

At each finite stage transfer the Kummer class of $y_{cr_i}$ from
$H_{cr_i}$ to $K$ with the specified scalar coefficients, using
$Ly_{cr_i}$. Complete the finite coset sum before contracting its
target cocycle and local Kummer lifts into the models over $K$.
The resulting cycles $P_i$, with their contraction witnesses, give
$P\in H^1(C)$ by Lemma~\ref{rl:finite-limits}.
Local Kummer conditions at $2N$ are preserved because $\Psi_i$ is
trivial there. At the center of the $i$th finite model, the norm relation
has the factor
\begin{equation}\label{eq:h:trace-center}
 \begin{aligned}
 b_i&=a_{r_i}-\chi_e(\Fr_{\mathfrak r_i})
            -\chi_e(\Fr_{\bar{\mathfrak r}_i}),
 \qquad r_i=\mathfrak r_i\bar{\mathfrak r}_i,\\
 P_i\big|_{t=0}&=\varepsilon_i b_i\kappa(LP_c(e)).
 \end{aligned}
\end{equation}
where $\kappa$ denotes the Kummer class in the central comparison
and $\varepsilon_i$ is a unit; the equality is taken modulo the finite
coefficient precision. Since $\chi_e$ is defined over $\Q$, its values
at the conjugate Frobenius elements agree; explicitly,
$b_i=a_{r_i}-2\chi_e(r_i)$.
To retain the central identity at the cochain level, let $B_i$ be the
reduction of a fixed finite model of the ordinary Kummer complex over
$K$, and let $\beta_i:B_i\to C_i/(t)$ be its comparison map to the
relaxed central complex. Choose a cycle $q_i$ representing
$\kappa(LP_c(e))$. The finite norm relation gives a boundary witness
$v_i$ in the target model such that
\[
 P_i\bmod t-\varepsilon_i b_i\beta_i(q_i)=d v_i.
\]
Retain $\beta_i,q_i,v_i$, and the target contraction homotopies on
the same stages and at the same coefficient precision as $P_i$.
Retain these units in the same ultrafilter
$\mathcal U_{\mathrm{in}}$ as all the comparison maps. Then
$b=\lim_{\mathcal U_{\mathrm{in}}}b_i$ has bounded valuation by
Equation~\eqref{eq:h:r-avoidance}, and
$\varepsilon=\lim_{\mathcal U_{\mathrm{in}}}\varepsilon_i$ is a unit. Thus
$P|_{t=0}=\varepsilon b\kappa(LP_c(e))$, with $b\ne0$.
The special complex differs from the ordinary Kummer complex only by
the active-prime relaxed quotients and the limit of the singular
quotient at the single moving prime $r_i$, the latter of bounded length.
Inflation from the complexes omitting $r_i$ gives this comparison. All modified
local complexes are rationally acyclic at the center, so pairings agree
there as well.

If the product has a simple zero, its rational special fiber has
one-dimensional cohomology in degrees one and two, and $P$ specializes
nontrivially by Gross--Zagier. The generic dimensions are at most these
special dimensions. Moreover, the long exact sequence for multiplication
by $t$ shows that $H^1(C[1/2])[t]=0$, since the special $H^0$ vanishes.
A torsion class that disappeared over the generic field would be killed
by a power of $t$ after localization at $(t)$; hence the nonzero special
class cannot disappear. Thus the generic ranks are one in degrees one
and two, $P$ spans $H^1$, and the rational rank is constant at $t=0$.
If only base Selmer corank one is known, the same dimension bound gives
generic ranks at most one; a zero generic point is allowed.

\begin{definition}\label{def:h:determinant-coordinate}
Choose an integral basis of $\det^{-1}_{\Lambda}C$. In generic rank one,
the element formed by $P$ in $H^1$ and the functional
$\langle P,\cdot\rangle$ on $H^2$ determines a scalar
$u(C,P)\in\operatorname{Frac}(\Lambda)$ in that basis. Put $u=0$ if
$P=0$ or the generic rank is zero. Integral changes of determinant basis
multiply $u$ by a unit.
\end{definition}

In the simple-zero case, constancy of the rational cohomology rank at
$t=0$ makes $u$ regular in the local ring $\Lambda_{(t)}$. Its value
$u(0)\in\Qtwo$ is therefore already defined. This specialization takes
place at $(t)$; the stronger assertion $u\in\Lambda[1/2]$ will follow
from the divisor estimates below.

The determinant has a simple integral interpretation. On a rank-one
free Selmer line, the point and its paired functional each contribute
the point's index. Torsion in degree two subtracts its length from the
coordinate in the inverse determinant lattice. For example, the complex
$[\Ztwo\xrightarrow{2^a}\Ztwo]$ in degrees one and two has inverse
determinant lattice mapping to $2^a\Ztwo$ in
the rational determinant. The coordinate of a fixed rational tensor
therefore loses $a$ in valuation. The next lemma carries out this
calculation with the actual local Kummer conditions. It also compares
the Selmer length over $K$ with those of the two rational twists.

\begin{lemma}[Central determinant valuation]
\label{lem:h:central-valuation}
For a simple-zero instance,
\begin{equation}\label{eq:h:central-determinant}
 \begin{aligned}
 v(u(0))&=2j(P_c(e))
   -\length\Sha(E^{(e)}/K)[2^\infty]
   -\sum_{q\mid c}\sum_{v\mid q}t_v+O_E(1),\\
 t_v&=\dim_{\F_2}E[2]^{\Fr_v}.
 \end{aligned}
\end{equation}
Furthermore,
\begin{equation}\label{eq:h:weil-restriction-length}
 \length\Sha(E^{(e)}/K)[2^\infty]
       +\sum_{q\mid c}\sum_{v\mid q}t_v
 =s(e)+s(ek)+\sum_{q\mid k}t_q
                         +2\sum_{q\mid c}t_q+O_E(1).
\end{equation}
The constants include the one fixed class multiplier and the bounded
local contribution at the moving prime $r_i$.
\end{lemma}

\begin{proof}
The ordinary integral Kummer complex over $K$ is self-dual. Its $H^1$ is
the compact Selmer module, and the torsion in $H^2$ has the length of the
finite quotient of the discrete Selmer group by its divisible subgroup.
This follows by derived base change to $\Qtwo/\Ztwo$: the local
complexes give exactly the discrete Kummer conditions. In the present
simple-zero case that finite quotient is
$\Sha(E^{(e)}/K)[2^\infty]$. Duality pairs the free parts in degrees
one and two perfectly. Torsion in degrees one and three costs a bounded
amount, by global torsion and duality. The uniform torsion bound here comes
from a fixed nonidentity Tate scalar acting trivially on every ring-class
field in the construction: its scalar minus one annihilates that torsion.
Lemma~\ref{rl:kernel} constructs this scalar, and the descent argument
of Lemma~\ref{lem:h:derivative-classes} uses the same bound.

At an active odd conductor place inertia acts by minus identity.
All integral $H^1$ is Kummer, and the relaxed quotient contributes only
$H^2(K_v,T^e)$ in degree two. Its length is $t_v$, by local duality
and the fixed $2$-torsion. At the moving prime $r_i$ the contribution
is bounded by
Equation~\eqref{eq:h:r-avoidance}. The generic pairings specialize to
the ordinary Kummer pairing because the modified local complexes are
rationally acyclic at the center. Thus the free determinant tensor has
valuation twice the point index, whereas torsion in degree two
contributes negative length to the inverse determinant. Together with
Equation~\eqref{eq:h:trace-center}, this proves
Equation~\eqref{eq:h:central-determinant}. The specialization is at
constant rational cohomology rank, so no specialization torsion is
being substituted for the point.

For Equation~\eqref{eq:h:weil-restriction-length}, apply the isogeny
invariance of the BSD quotient to quadratic Weil restriction and
$E^{(e)}\times E^{(ek)}$
\citep[Theorem~2.1 and the preceding discussion]{Tate1966}.
All the groups in this
comparison are finite apart from the known rank-one lattices; no BSD
equality is assumed. The forward and reverse isogenies have composites
multiplication by $2$. Consequently the rank-one lattice indices,
height ratios, torsion ratios and period ratios contribute bounded
$2$-powers. For periods, the pullback indices of N\'eron differentials
at $2$ are bounded by those composites; real kernels and component
ratios are bounded as well.

The remaining terms are the Tamagawa numbers. A ramified quadratic
twist of good reduction at an odd prime $q$ has type $I_0^*$ and
component number one plus the number of roots of its reduced auxiliary
cubic \citep[Theorem~4.1, Table~4 and proof, Case~3]{BarriosEtAl2025}.
To identify this number, write a good integral model as $y^2=f(x)$,
retaining the quadratic term, and the twisting parameter as $d=\pi u$.
The auxiliary cubic is $\bar u^3\bar f(T/\bar u)$; it is separable
and its roots correspond to the nonzero rational two-torsion points.
Thus the component number is $\#E[2]^{\mathrm{Fr}_q}$ and its
$2$-valuation is $t_q$. This argument also applies in residue
characteristic $3$. Its connected reduction is unipotent. Some inertia element acts
by $-1$ on the twisted $2$-adic representation, so its local
$2$-primary torsion is precisely its Frobenius-fixed $2$-torsion.
Over an unramified extension the same formula holds with the
Frobenius of the enlarged residue field. At unramified
or split places, Weil-restriction component numbers are the products
over $v\mid q$. At a prime $q\mid k$ outside $2N$, the curve over
$K$ has good reduction and its Weil restriction has connected special
fiber. Indeed, Weil restriction of its N\'eron model is the N\'eron
model of the restriction of the curve
\citep[Section~7.6, Proposition~6]{BoschLutkebohmertRaynaud1990}.
Its special fiber has the original connected elliptic fiber as quotient,
with connected vector-group layers from ramification. Exactly one rational twist
has the ramified component factor. At a conductor prime $q\mid c$,
the two rational twists give a combined valuation $2t_q$, and the
corresponding contribution over $K$ is $\sum_{v\mid q}t_v$. All
places above $2N$ have bounded local contributions. These are precisely
the terms of Equation~\eqref{eq:h:weil-restriction-length}.
\end{proof}


In a simple-zero genus instance, the fixed partner absorbs the terms involving $h_*$,
while $c=1$ and $k=hh_*$. The displayed formulas therefore identify
\[
 v(u(0))=2j(P_1(h_*))-2w(h)-s(h)+O_{E,h_*}(1).
\]
In a simple-zero ring instance with $\omega(k)\le a$, the discriminant term
is bounded and the conductor term is $4w(h)$, giving
\[
 v(u(0))=2j(P_{|h|}(h))-4w(h)-s(h)-s(hk)+O_{E,a}(1).
\]
The ring identity also applies to the interpolation families in which
an added nonsplit prime has $t_q=0$; such a prime contributes no weight.
The all-split hypothesis is needed for the uniform ring lower bound.
These formulas explain both the lower-bound statements and the
normalization in Heegner interpolation. To prove the lower bounds,
we must show that $u$ has no negative divisor away from $(2)$ and
has a uniformly bounded negative divisor at $(2)$ in the stated
cases. We will construct integral local switches that preserve its determinant
valuation, and use inner detection to choose switches lowering excess
Selmer rank.
The uniform ring estimate will require an additional characteristic-zero
limit; the universal binary family will instead require a bounded
clearing factor after the new local blocks have been removed.

The main intermediate estimate is
\begin{equation}\label{eq:h:divisor-bounds}
 \begin{aligned}
 v_{\mathfrak p}(u)&\ge0
       &&\bigl(\operatorname{ht}\mathfrak p=1,
                         \ \mathfrak p\ne(2)\bigr),\\
 v_{(2)}(u)&\ge-O_{E,h_*}(1)
       &&\text{in genus data with irreducible }E[2],\\
 v_{(2)}(u)&\ge-O_{E,a}(1)
       &&\text{in ring data with }\omega(k)\le a
                         \text{ and all active primes split.}
 \end{aligned}
\end{equation}
To construct the switches used in these estimates, we first retain
enough Galois evaluation data to select their derivative primes.

\subsection{Evaluation and the exact kernel}

The determinant estimates require switches that detect the excess
Selmer dimension at a height-one test. To choose
switching primes, we must still detect a nonzero limiting Selmer class
by Galois evaluation. The next two lemmas first retain that detection
under coefficient specialization, then restrict it to sequences acting
trivially on the full Tate module, all ring-class fields, and the
quadratic characters in use.

\begin{lemma}[Evaluation on arbitrary Galois sequences]\label{rl:evaluation}
\label{lem:h:evaluation}
Fix $K$ and a sequence of the finite cochain constructions above.  Keep
the contraction maps into the original global cochains.  Let
$\Gamma_K=\prod_i G_K$ and
$\Gamma=\Gamma_K\rtimes\langle\tau\rangle\subset\prod_iG_{\mathbf Q}$,
with simultaneous complex conjugation $\tau$.  Thus $\Gamma_K$ has
index two in $\Gamma$; the larger unrestricted product of rational
Galois groups is not used for induction.  The matrix limit defines a
coefficient action of $\Gamma_K$ on $\Lambda^2$.  For every field $L$
receiving a map from $\Lambda$, there is a natural injection
\[
 H^1(C\otimes_\Lambda L)\longrightarrow
 H^1_{\rm abs}(\Gamma_K,L^2),
\]
where the right side denotes abstract crossed-homomorphism classes.
Here $C$ may be the global complex or a Selmer complex whose local
conditions contain all local degree-zero invariants and inject into
local $H^1$.  The same statement holds for the induced representation
of $\Gamma$.

The evaluation maps and their degree-two cochain identities exist
for every group sequence and commute with coefficient extension.
\end{lemma}

\begin{proof}
Let $C_i$ be a finite model and let $\iota_i$ be its chosen map to the
original global cochains, using the projection from the Selmer cone
when appropriate.  For every $g=(g_i)$ and $h=(h_i)$ define matrices
\[
 e^0_i:C_i^0\to A_i^2,\quad
 e^1_{g,i}:C_i^1\to A_i^2,\quad
 e^2_{g,h,i}:C_i^2\to A_i^2
\]
by evaluating $\iota_i$ in degrees zero, one, and two.  Their source
ranks are fixed and their entries belong to $A_i$.  Thus each has a
coefficientwise ultrafilter limit, for every choice of $g,h$.
The same fixed ultrafilter defines these limits for all $g,h$.  Likewise
$\rho(g)=\lim_i\rho_i(g_i)$ is a representation.  The original
cochain identities imply
\begin{align}
 e^1_g d&=(\rho(g)-1)e^0,\label{rl:eval-boundary}\\
 e^2_{g,h}d&=\rho(g)e^1_h-e^1_{gh}+e^1_g.
                         \label{rl:eval-cocycle}
\end{align}
All identities pass to the limit because their matrix expressions
have finitely many terms.  A cycle $x\in C^1\otimes L$ consequently
gives the crossed homomorphism $g\mapsto e^1_g(x)$, and a boundary
gives a coboundary.

We prove injectivity in a way that also controls base change.  At each
finite stage let $k=\mathbf F_2$ be the residual field.  Choose a
space of representatives for residual $H^1$.  Evaluate pairs $(c,m)$,
where $c$ is in this space and $m\in k^2$, by
$b_c(g)-(\rho(g)-1)m$.  The intersection of the kernels of all
these evaluations consists precisely of pairs $(0,m)$ with $m$ an
invariant.  Since the vector space is finite dimensional, a list of
at most $\dim H^1+2$ elements already has this intersection.  For
these elements evaluation is a chain map
\[
 C_i\longrightarrow D_i=
 [\,A_i^2\xrightarrow{((\rho_i(g_j)-1))_j}(A_i^2)^d\,]
\]
in degrees zero and one, inducing an isomorphism on residual $H^0$
and an injection on residual $H^1$.  For Selmer complexes this uses
the exact sequence of~\eqref{rl:selmer-definition}: the indicated
local degree-zero isomorphisms give $H^0(C_i)=H^0_{\rm glob}$ and
$H^1(C_i)\hookrightarrow H^1_{\rm glob}$ residually.

Retain these matrices and their chosen group-element sequences.
The residual global and local complexes, and the local conditions,
have no negative cohomology.  The cone definition therefore gives
$H^n(C\otimes k)=0$ for $n<0$.  Put $Q=\operatorname{Cone}(C\to D)$.
The long exact sequence gives $H^n(Q\otimes k)=0$ for $n<-1$;
its groups in degrees $-1$ and $0$ vanish because the detecting map is
an isomorphism on $H^0$ and injective on $H^1$.  Any $H^3(C\otimes k)$
contributes in degree two of this cone, so causes no negative-degree
obstruction.  Successively cancel the unit differential blocks of
$Q$.  This gives a homotopy-equivalent bounded free complex starting
in degree one.  For any $\Lambda$-algebra $L$, its ordinary tensor
product computes $Q\otimes^{\mathbf L}_\Lambda L$ and still starts
in degree one; no flatness is needed.  Thus the long exact sequence
makes $H^1(C\otimes L)\to H^1(D\otimes L)$ injective.  A crossed cocycle that is a coboundary
on all of $\Gamma_K$ is in particular a coboundary on the detecting
list, proving the desired injection.  For induction the same proof
has coefficient rank four and the detecting bound is $\dim H^1+4$.
\end{proof}

\begin{lemma}[The exact common kernel]\label{rl:kernel}
In the notation of Lemma~\ref{rl:evaluation}, let $H$ be the group of
sequences each of whose entries acts trivially on the full Tate module,
on all ring-class fields of $K$, and on the quadratic characters in
use.  Restriction from the crossed classes in that Lemma to $H$ is
injective at every characteristic-zero test.  It is also injective
at the residual $(2)$ test.  The characteristic-zero annihilator
needed for this assertion can be chosen with a fixed finite
$2$-valuation depending only on $E$, also as $K$ varies.
\end{lemma}

\begin{proof}
The Tate image contains an open subgroup of scalar homotheties, both
in the CM and non-CM cases.  Choose once and for all
$g\in G_{\mathbf Q}$ with Tate action $u\,1$, where
$u\in1+4\mathbf Z_2$ and $u\ne1$.  Its square belongs to every
quadratic subgroup.  Put
\[
 z=g^2\tau g^2\tau^{-1}.
\]
The generalized dihedral property of every ring-class extension
implies that $z$ is ring-class trivial.  It is also trivial on every
quadratic character, and its Tate action is the fixed scalar $u^4$.
For every sequence in $\Gamma_K$, its commutator with the constant
sequence $z$ is in the exact group $H$, since its Tate commutator and
all its abelian character images are trivial.  Thus $z$ is central in
$\Gamma_K/H$ and acts on the coefficients by $u^4$.

Here the elementary central-scalar argument suffices.  If a crossed
cocycle $b$ is defined on a group with such a central element, comparison
of $b(zg)$ and $b(gz)$ gives
\[
 (u^4-1)b(g)=(g-1)b(z).
\]
In characteristic zero this makes $b$ a coboundary.  Inflation--restriction
then gives injectivity on $H$.  The valuation of $u^4-1$ is fixed.

In characteristic two choose, at each stage, inertia at $r_i$ mapping
onto the chosen generator of $\xi_i$. Its Tate action is trivial and
its ring-class images commute, so the resulting sequence is central
in $\Gamma_K/H$ and acts on $M$ by $1+t$. Its scalar minus one is
invertible in $\mathbf F_2((t))$.

For the induced representation $W$, first restrict to $\Gamma_K$.
There $W=M\oplus M^\tau$, and this inertia sequence acts on the
summands by $1+t$ and $(1+t)^{-1}$. Hence $W^{\Gamma_K}=0$.
Inflation--restriction for the index-two inclusion gives an injection
$H^1(\Gamma,W)\to H^1(\Gamma_K,W)$, even in characteristic two.
Apply the central-scalar calculation separately to the two summands
to obtain injection on restriction to $H$. The same argument works
in characteristic zero with the fixed scalar $u^4$.
Throughout, $H$ is defined by exact actions at every stage.
\end{proof}

\begin{lemma}[Irreducibility on the characteristic-zero tests]
\label{rl:irreducibility}
On an inner characteristic-zero test where $\Psi^2\ne1$,
the four-dimensional representation induced from $T\chi\Psi$ is
absolutely irreducible.  If $\Psi^2=1$ on an inner test, it splits
as two inequivalent absolutely irreducible quadratic twists.
\end{lemma}

\begin{proof}
The two induced summands are distinguished by $r$-inertia, whose
scalars are $\Psi$ and $\Psi^{-1}$.  On either summand, multiplying
the Tate matrices by scalars does not alter invariant subspaces.
The Tate representation restricted to $G_K$ is absolutely
irreducible under the exclusions already specified.  These facts
give absolute irreducibility of the induction when the two inertia
characters differ.

If $\Psi^2=1$, extension across conjugation gives the two quadratic
twists.  They are inequivalent, since a quadratic self-twist of the
Tate representation here would make its restriction to $G_K$
reducible over an algebraic closure.  In the CM case this would
force $K$ to be the CM field, which was excluded.
\end{proof}

\begin{lemma}[Simultaneous finite Chebotarev realization]
\label{rl:chebotarev-evaluation}
For any element $u$ of an exact kernel in Lemma~\ref{rl:kernel},
derivative primes can be chosen, at increasing finite precision,
whose rational Frobenius is $\tau u$ on all retained data.  Their
finite localization is evaluation on $(\tau u)^2$.  The assertion
is compatible with any fixed finite list of previous derivative
operations.
\end{lemma}

\begin{proof}
At one Artin stage, a coordinate of a cochain in the image of a
finite model is a locally constant function.  The finitely many
coordinates being used factor through one finite quotient of the
absolute Galois group.  Enlarge it to contain the Tate action at
the required precision, the finitely many ring-class fields and
quadratic characters currently in use, and all old finite data.
Chebotarev realizes the conjugacy class of $\tau u$ on this finite
quotient, with a chosen place realizing the desired representative
\citep[Chapter~V, Theorem~3.23]{MilneCFT2020}.
It gives a rational prime inert in $K$, outside the old finite set.
Because $u$ is exactly Tate-trivial, its Frobenius tends to
$\rho_T(\tau)$; in particular $a_\ell\to0$ and $\ell\to-1$.
Because $u$ is exactly ring-class trivial, the square Frobenius on
the $K$-prime is the identity on every retained ring-class field.
The unramified local coordinate is therefore the stated
Frobenius-square evaluation.

Repeat at each finite stage with increasing precision.  All
conditions remain finite, and no density bound or prime-size
bound is required.
Previously imposed finite diagrams are unchanged.
\end{proof}

In induced-module notation, restriction to the exact kernel is additive.
For a crossed cocycle $b$ and an exact-kernel element $u$,
\[
 b((\tau u)^2)=(1+\tau)b(u).
\]
Indeed $\tau^2=1$, $u$ acts trivially on the coefficients, and the
cocycle identity cancels the two occurrences of $b(\tau)$.  Projection
to either induced summand identifies $\operatorname{im}(1+\tau)$
with a two-dimensional coefficient module, also in characteristic two.

\subsection{The local switch and its integral determinant}
\label{rd:local-subsection}

Fix a complex conjugation $\tau$ and write $C_\tau$ for its action on $T$.
The next two Lemmas construct the integral local switches and the classes
that pass through them. Their finite-level formulas are needed before
applying the limit constructions above.

\begin{lemma}[The split local model]\label{rd:local-model}\label{lem:h:switching-pairing}
Let a sequence of inert good primes \(\ell\) have rational Frobenius
lifts tending on \(T\) to a fixed complex conjugation \(C_\tau\).
At the unique prime of \(K\) above \(\ell\), the limiting local complex
has
\begin{equation}
 H^0=M,\qquad H^1=F\oplus S=M\oplus M,\qquad H^2=M,
 \label{rd:local-cohomology}
\end{equation}
with zero differentials in a free split model.  Here \(M=T\otimes\Lambda\)
with the relevant scalar coefficient structure.  The two degree-one
coordinates of a cocycle are
\begin{equation}
 f(c)=c(\gamma^2),\qquad s(c)=c(\sigma),
 \label{eq:h:local-coordinates}
\end{equation}
where \(\sigma\) is an oriented pro-\(2\) tame generator.  The plane
conditions \(F=\{s=0\}\) and \(S=\{f=0\}\) include \(H^0\) and exclude
\(H^2\).  Each is its own orthogonal complement for conjugate-Weil
duality.  The strict condition \(H^0\) and the relaxed condition
\(H^0\oplus H^1[-1]\) are mutually dual.  All four inclusions can be
equipped with compatible integral duality triangles.
\end{lemma}

\begin{proof}
The prime-to-\(\ell\) local cochains are computed by the tame inertia and
residue-field directions.  In these two directions the differentials
are formed from inertia minus one and Frobenius minus one, with the
usual Frobenius action on tame inertia.  On \(M\), inertia is trivial at
the derivative prime.  The square of a rational Frobenius is trivial
on every ring-class character.  In the limit it is also the identity on
\(T\), while \(\ell^2\) tends to one.  Thus both differentials become zero,
giving \eqref{rd:local-cohomology}.

Here is a finite-precision check of the pairing, needed at \(2\).
At a retained quotient killed by \(2^a\), impose also
\(q=\ell^2\equiv1\pmod{2^{a+1}}\), as well as trivial
Frobenius-square action on the coefficient module at this quotient.
The tame two-generator complex then has zero differentials,
with the coordinates in \eqref{eq:h:local-coordinates}.  Its degree-one
scalar generators are the unramified and tame characters.
More explicitly, use the presentation
\(\phi\sigma\phi^{-1}\sigma^{-q}=1\).
After augmentation with trivial coefficient action its Fox
resolution has terms of ranks \(1,2,1\), with
\(d^0=0\) and \(d^1=(0,1-q)\), in the ordered finite and
tame coordinates.  The diagonal from the same resolution
has zero pure-finite term, unit mixed terms, and pure-tame
coefficient \(\pm\binom q2\).  The latter coefficient is
the augmented second derivative of \(\sigma^{-q}\)
together with its adjacent \(\sigma\) in the relation.
Both \(1-q\) and \(\binom q2\) vanish modulo \(2^a\)
under the imposed congruence.  Thus this is an actual
split complex and an actual pure-plane zero pairing
at the retained precision.
The field contains the \(2^{a+1}\)-st roots of unity, so
\(-1\) is a \(2^a\)-th power.  The norm-residue identity
\((x,x)=(x,-1)\) for the \(2^a\)-Hilbert symbol therefore
makes the square of each scalar generator zero.
The two generators form a basis; local Tate duality makes
their cross pairing a unit.  The same calculation applies
coefficientwise to the free Artin coefficient module, including
its nilpotent \(t\) and quadratic variables.  Thus the pure
plane restrictions are already zero at every retained quotient.
All these congruences can be imposed simultaneously with the
derivative Frobenius conditions by increasing their precision.
Taking the compatible limit proves integral vanishing of the
pure squares.  It is not deduced from skew symmetry in
characteristic two.

Equivalently, in the limiting tame presentation
\(\phi\sigma\phi^{-1}=\sigma^q\) with \(q\to1\), the
rank-\((1,2,1)\) resolution specializes to the exterior
resolution for two procyclic directions.  The finite
Hilbert-symbol calculation specifies the cup product on
its split cohomology and rules out a hidden square term.
Conjugation transports the finite Tate coordinate by \(C_\tau\) and the
singular coordinate by \(-C_\tau\), with possible scalar units coming from
the orientations.  Consequently the cross pairing is
\begin{equation}
 \beta(x,y)=\epsilon\,\langle x,C_\tau y\rangle,\qquad
 \epsilon\in\Lambda^\times,\quad x\in F,\ y\in S,
 \label{eq:h:cross-pairing}
\end{equation}
where \(\langle\, ,\,\rangle\) is the perfect alternating Weil form.
This proves the assertions about orthogonal complements.  Since the
restrictions of this explicit pairing to either pure plane are zero as
maps of the split complexes, their isotropy homotopies may be chosen to
be zero.  Only a product of two degree-one terms could map to the
degree-two local invariant target; the degree-zero terms introduce
no additional restriction.  These zero homotopies agree on the
strict subcomplex.  These are the actual
orthogonality data required in the Selmer duality construction
\cite[Theorem~6.3.4]{Nekovar2006}; the compatibility assertion is therefore
an assertion about these displayed complexes and maps.  Lemma~\ref{rl:finite-limits} is used to retain this finite diagram together
with the unchanged local data.
\end{proof}

For a finite set $I$ of derivative primes, let $C_I$ be obtained
from $C$ by replacing the finite plane condition $F$ by $S$ at each
prime of $I$, retaining all other local conditions. The classes below
use the same multiplier $L$ already included in $P$; in particular,
$P_{\varnothing}=P$.

\begin{lemma}[Descent and the switch at \(2\)]\label{rd:derivative}\label{lem:h:derivative-classes}
Suppose that the finite ring-class transfers are performed before
retaining their target cycles and witnesses as in
Lemma~\ref{rl:finite-limits}. There is one
integer \(L\ne0\), depending on the fixed curve, parametrization and
fixed local types, such that the classes obtained from the Heegner
derivatives of \(L\) times the point satisfy, for every finite switched
set \(I\),
\begin{equation}
 P_I\in H^1(C_I),\qquad
 f(P_{I\cup\{\ell\}})=0,\qquad
 s(P_{I\cup\{\ell\}})=C_\tau f(P_I).
 \label{eq:h:finite-singular}
\end{equation}
The same \(L\) works for every number of active and derivative primes.
The previously imposed singular conditions are retained.
\end{lemma}

\begin{proof}
We first descend the derivative classes using one multiplier for the
complete conductor, and verify their Kummer conditions at $2N$. We then
compute both local coordinates in Equation~\eqref{eq:h:finite-singular}.
The last step retains these classes and identities in the same bounded
target cochain diagrams used for evaluation.

Choose the orientations coherently.  The relative ring-class group at
an inert prime \(\ell\) is cyclic of order \(\ell+1\).  For distinct new
primes these groups form a product: the ring-class exact sequence has
kernel the product of the local residue-unit quotients, and the only
global units here are \(1,-1\), whose images in these quotients are
trivial.  Set
\[
 D_\ell=\sum_{i=1}^{\ell} i\sigma_\ell^i,\qquad
 D_I=\prod_{\ell\in I}D_\ell.
\]
Direct multiplication gives
\begin{equation}
 (\sigma_\ell-1)D_\ell=(\ell+1)-
             \sum_{i=0}^{\ell}\sigma_\ell^i.
 \label{rd:derivative-identity}
\end{equation}
The Hecke correspondence and reduction of its CM isogenies give, for
every class-group translate,
\begin{equation}
 \operatorname{Tr}(y_{c r n})=a_\ell y_{c r n/\ell},
 \qquad
 \widetilde y_{c r n}=\operatorname{Fr}_\ell
                       \widetilde y_{c r n/\ell}.
 \label{eq:h:heegner-norm-reduction}
\end{equation}
These are the point relations, not an application of a residual
irreducibility theorem.  A fixed multiple of the parametrization removes
cuspidal terms.

There is a uniform bound on the \(2\)-primary torsion over the union of
the ring-class fields in use.  To see this, choose a fixed nonidentity
scalar homothety of infinite order in the \(2\)-adic Tate image.
Square a lift so that it
belongs to every quadratic subgroup under consideration.  Multiply it
by its conjugate under complex conjugation.  Its action on each
ring-class field is trivial, because conjugation inverts the abelian
ring-class action, while its action on \(T\) is still a fixed
nonidentity scalar.  If \(2^e\) annihilates its scalar minus one up to
a unit, then \(2^e\) kills the torsion in question.  The existence of
this fixed scalar, including the CM case, is the Tate-image input
used in Lemma~\ref{rl:kernel}.

Work first with coefficients modulo \(2^a\).  Require \(a_\ell\) and
\(\ell+1\) to vanish to a precision substantially larger than \(a\).
Equations \eqref{rd:derivative-identity} and
\eqref{eq:h:heegner-norm-reduction} then make the derivative Kummer class
invariant under the relative ring-class group.  In the
inflation--restriction sequence, both the obstruction to descent and
the ambiguity of descent are killed by \(2^e\), since their
coefficient module is killed by \(2^e\).  Multiplication by \(2^{2e}\)
therefore gives a descent independent of this ambiguity: descend
\(2^e\) times the class and multiply its lift by \(2^e\).
There is no repetition of this factor for individual primes; descent
is performed once from the field with the complete conductor.

The global descent has cost only the fixed factor $2^{2e}$, regardless
of the number of derivative primes. It remains to ensure that the
descended class satisfies the Kummer condition at every fixed bad place;
this is a separate local descent issue.

At a place over \(2N\) these descent extensions are unramified.
Here is the local descent argument, valid also for even extension degree
and for toric or additive reduction.  Let $F'/F$ be a finite unramified
extension with group $G$, and write $k'/k$ for its residue extension.
The successive formal-group quotients $E_j(F')/E_{j+1}(F')$ are additive
copies of $k'$.  The normal basis theorem makes each a cohomologically
trivial $G$-module, without a restriction on $|G|$.  Every finite
quotient $E_1(F')/E_j(F')$ is therefore cohomologically trivial.
The inverse systems of these finite groups and their invariant groups
satisfy the Mittag--Leffler condition.  Taking inverse limits in finite
$G$-cochains consequently gives $H^1(G,E_1(F'))=0$.
Let $\mathcal E/\mathcal O_F$ be the N\'eron model and put
$\widetilde E^0=\mathcal E_k^0$, a smooth geometrically connected
commutative algebraic group over the finite field $k$.
Lang's finite-extension cocycle theorem
\citep[Proposition~3, p.~559]{Lang1956} gives
$H^1(G,\widetilde E^0(k'))=0$.
The smooth reduction sequence then gives $H^1(G,E^0(F'))=0$.
Let $\Phi=\mathcal E_k/\widetilde E^0$ be the finite \'etale
component-group scheme. Each fiber over a point of $\Phi(k')$
is a torsor under $\widetilde E^0_{k'}$, and hence has a $k'$-point
by \citet[Theorem~2, pp.~557--558]{Lang1956}.
Smooth lifting therefore makes $E(F')\to\Phi(k')$ surjective. Hence
\[
 H^1(G,E(F'))\lhook\joinrel\longrightarrow
 H^1(G,\Phi(k')).
\]
Existence of the smooth finite-type model and smooth lifting follow from
\citet[Section~1.3, Corollary~2, and Section~2.3, Proposition~5]{BoschLutkebohmertRaynaud1990}.
The right side is killed by the exponent of $\Phi(\overline k)$.
By inflation--restriction, the left side is exactly the kernel of
$H^1(F,E)\to H^1(F',E)$.  Thus the same fixed exponent makes a
class in $H^1(F,E[2^a])$ Kummer whenever its restriction is Kummer.
N\'eron models commute with unramified base extension
\citep[Section~7.2, Theorem~1 and Corollary~2]{BoschLutkebohmertRaynaud1990},
so their geometric component groups are unchanged. Only the finitely many fixed local twisting
types at $2N$ occur, so the product of these component exponents is
bounded independently of the conductor and of the extension degree.
Include this product, the single descent factor $2^{2e}$, and the
parametrization factor in \(L\).  Thus the descended class has the
integral Kummer condition at \(2N\).
No good-ordinary or semistability condition was used.  Transfer with
the scalar character preserves this condition because \(\Psi=1\)
there.

The single multiplier $L$ now gives descent with the required Kummer
conditions. We next prove the finite-to-singular identity itself. Both
coordinates matter: the singular coordinate must become $C_\tau f$,
while the new finite coordinate must vanish at the retained precision.

Here is the local calculation, including the division by two.  Apply
the derivatives at all primes other than \(\ell\), and write \(X,Y\)
for the resulting lower and upper points.  Put \(F_\ell=\operatorname{Fr}_\ell\).
The lower field splits at the \(K\)-prime \((\ell)\), since that ideal
is principal.  Reduction therefore gives \(F_\ell^2X=X\).
The following Frobenius-polynomial identities are taken on the reduction;
good odd-prime reduction is injective on the $2$-primary torsion
coordinates being compared.
If \(Q\) is a \(2^a\)-division point of \(LX\), the finite coordinate
of its Kummer class is
\begin{equation}
 f=(F_\ell^2-1)Q=
       \bigl(a_\ell F_\ell-(\ell+1)\bigr)Q.
 \label{rd:finite-formula}
\end{equation}
We give the descent-cocycle calculation of the other coordinate.
Choose \(Z\) with \(2^aZ=LD_\ell Y\), and choose the descended
cocycle \(b\) so that on the upper-field Galois group it is
\(b(g)=(g-1)Z\).  A change by the already bounded descent ambiguity
has no effect on the eventual limit.  The crossed identity gives
\[
 U=(\sigma_\ell-1)Z-b(\sigma_\ell)
                  \in E(H_{crn}).
 \]
For \(g\) in the upper-field Galois group, compare the two
expansions of \(b(g\sigma_\ell)\); they give \(gU=U\).
Moreover,
\[
 2^aU=L\bigl((\ell+1)Y-a_\ell X\bigr).
 \]
Both integers \(\ell+1,a_\ell\) are divisible by \(2^a\).
Consequently
\[
 U_0=L\frac{\ell+1}{2^a}Y-
                      L\frac{a_\ell}{2^a}X
 \]
is an actual upper-field point with the same \(2^a\)-multiple
as \(U\).  Thus \(U-U_0\) is upper-field torsion of bounded
exponent.  This comparison with an actual integer quotient
is what bounds the ambiguity; arbitrary division points
would not give this bound.  Reduction makes
\(\sigma_\ell\) act trivially and
\(\widetilde Y=F_\ell\widetilde X\).  Therefore
reduction of \(b(\sigma_\ell)=-U\) gives its singular coordinate
\begin{equation}
 s=\bigl(a_\ell-(\ell+1)F_\ell\bigr)Q
 \label{rd:singular-formula}
\end{equation}
up to the uniformly bounded descent torsion already described.
The Frobenius polynomial yields the exact operator identity
\begin{equation}
 (F_\ell-a_\ell)
     \bigl(a_\ell F_\ell-(\ell+1)\bigr)
       =a_\ell-(\ell+1)F_\ell.
 \label{rd:exact-frobenius}
\end{equation}
Thus \(s=(F_\ell-a_\ell)f\); when \(a_\ell\) vanishes at the
retained torsion precision, this is \(s=F_\ell f\).  Taking the
Frobenius limit gives \(s=C_\tau f\).

If $\gamma$ is the chosen rational Frobenius lift, its image in the
generalized-dihedral ring-class Galois group lies outside the abelian
subgroup, so $\gamma^2$ acts trivially on the full upper ring-class field.
The same lift therefore computes the descended finite Kummer coordinate
from the reduced point over the common residue field $\mathbf F_{\ell^2}$.
For the new finite coordinate, reduction of the derivative point is
\begin{equation}
 \sum_{i=1}^{\ell}iF_\ell X
       =\frac{\ell(\ell+1)}2F_\ell X.
 \label{rd:half-factor}
\end{equation}
The finite Kummer coordinate consequently vanishes modulo \(2^a\)
if \(v_2(\ell+1)\ge a+1\).  The additional \(1\) is essential at
\(2\).  It is imposed in the choice of the next precision, and does
not multiply the class by a new denominator.  Bounded torsion
ambiguities disappear on passing from \(E[2^a]\) to \(E[2^m]\) with
\(a-m\ge e\).  Choose \(a\) and the Frobenius precision successively
faster than \(m\).

The calculation is compatible with every translate in the transfer:
transport the place and its Frobenius before applying that translate.
The scalar weight on the lower and upper class is identical.  Applying
the same calculation at each \(\ell\in I\) proves the earlier
singular conditions as well.

The local formulas now give the switch. To use it in the rank-reduction
argument, we must retain more than its cohomology class: the transferred
target cycles, local lifts and degree-two cochain identities must extend
the original finite diagrams. We record this compatibility next.

Here the transfer itself is integral at the cochain level.
For a finite Artin coefficient ring \(A\), let
\(\Delta=\operatorname{Gal}(H_{cr}/K)\), through which the
finite scalar and quadratic characters factor.  Shapiro's
cochain map and the coefficient homomorphism
\[
 A[\Delta]\longrightarrow A,\qquad
 [\sigma]\longmapsto\chi_e(\sigma)\Psi(\sigma)
 \]
send the descended class to the desired twisted class.
For a universal quadratic family, replace the quadratic
value by its group element in \(A[G]\), also including the
fixed base character.  Choose inverse actions on both
sides when using that Shapiro convention.  Explicitly,
these are the finite coset sums with the indicated scalar
weights; there is no division by \(\#\Delta\).
These maps and their degree-zero and degree-two compatibility formulas
are used at this finite stage, before passage to a bounded target model.
Naturality of restriction and Kummer corestriction
supplies the local lifts. At \(2N\) the scalar character
is trivial, so these are precisely the Kummer conditions
of the indicated quadratic twist.

More explicitly, in the universal case let
$\nu:\Delta\to G$ be the additional quadratic character.
Perform the coefficient map
\[
 A_i[\Delta]\longrightarrow A_i[G],\qquad
 [\sigma]\longmapsto
 \chi_{h_0}(\sigma)\Psi_i(\sigma)[\nu(\sigma)]
\]
and the finite coset sum before any character evaluation. This gives
one cycle over $A_i[G]$, with its local lifts. Every sign evaluation
$A_i[G]\to A_i$ sends this cycle and these lifts to the corresponding
weighted construction, because it commutes with finite sums and
differentials. After contraction, compare an evaluated target model
with its already chosen scalar model using their maps to the same
original target cochains. Retain these comparison matrices and their
contraction homotopies for all the finitely many characters of this
fixed $G$. No character idempotent or division by a class-group
order is used. At $t=0$ the split-prime norm multiplier on this universal
cycle is
\[
 a_{r_i}-2\chi_{h_0}(r_i)[\nu(\Fr_{\mathfrak r_i})].
\]
The two conjugate Frobenius values of $\nu$ agree because each
quadratic character comes from $\Q$. Evaluation at the character
corresponding to $h$ therefore gives exactly
$a_{r_i}-2\chi_h(r_i)$, with the central boundary witness described
above.

Finally contract only the resulting target cycles, local lifts and
boundary witnesses into the models over $K$. Retain the target
comparison maps, the split local Fox complexes, their zero isotropy
homotopies, and their duality and switch identities in degrees zero,
one and two. The proof of Lemma~\ref{rl:finite-limits} shows explicitly
how to transport these identities and their witnesses. In particular
the unbounded cochains over $H_{cr}$ and the full Shapiro maps are not
objects or maps of this limiting diagram.
When a further derivative prime is selected, extend the
same original finite-stage diagram and its contraction
maps, using the same inner ultrafilter.  Do not construct
an unrelated cohomologically equivalent old complex.
The earlier evaluations, including all degree-two
cocycle identities, are consequently retained by the
same contraction maps after each new switch. For a fixed finite set of
switches and fixed $G$, the target rank bounds above apply.
Lemma~\ref{rl:finite-limits} therefore applies to these extended
target diagrams at their common retained precision.
Their limiting classes give
\eqref{eq:h:finite-singular}.  The displayed local formulas agree with the
calculation underlying \cite[Proposition~1.7.4]{Howard2012}; Howard's
odd-prime global hypotheses are not being asserted at \(2\).
\end{proof}

\begin{lemma}[Preservation of the determinant coordinate]
\label{lem:h:determinant-switch}
Let $\mathcal O$ be a height-one discrete valuation ring at which the
complexes and switching conditions are self-dual. Suppose the generic
Selmer group for a set $I$ has rank one and is spanned by $P_I$, and
$v_0=f(P_I)$ is nonzero at the next derivative prime. Then the switched
generic group has rank one and is spanned by $P_{I\cup\{l\}}$, and
\begin{equation}\label{eq:h:switch-equality}
 v_{\mathcal O}(u_I)=v_{\mathcal O}(u_{I\cup\{l\}}).
\end{equation}
This assertion does not require a primitive localization in either
integral plane.
\end{lemma}

\begin{proof}
Compare the strict condition and the relaxed condition of
Lemma~\ref{lem:h:switching-pairing}. Their complexes are dual. Over the
fraction field, the image of relaxed $H^1$ in $F\oplus S$ is
self-annihilating by Poitou--Tate, hence two-dimensional. Strict $H^1$ is
zero because the original rank-one group localizes nontrivially. The
relaxed image meets $F$ in the line spanned by $v_0$ and, by
Equation~\eqref{eq:h:finite-singular}, meets $S$ in the nonzero line
spanned by $C_\tau v_0$. Its intersection with either plane is therefore
exactly a line. This proves the rank and spanning assertions.

For the integral determinant comparison let $A=P_I$ and
$B=P_{I\cup\{l\}}$ be a field basis of relaxed $H^1$. The dual basis
$A^\vee,B^\vee$ belongs to strict $H^2$. In the triangle adding $F[-1]$
to the strict complex, the map $F\to H^2_{\mathrm{str}}$ pairs a vector
$x$ against the relaxed localization. By
Equations~\eqref{eq:h:cross-pairing} and
\eqref{eq:h:finite-singular}, it is a unit times
$\langle x,v_0\rangle B^\vee$. The remaining $H^2$ quotient maps to
$A^\vee$. Choose $x$ mapping to $B^\vee$. Although $v_0$ and $x$ need
not be integral separately, the volume of the ordered basis $(v_0,x)$
of $F$ is a unit of $\mathcal O$, precisely because its symplectic
pairing is a unit. For $S$ use $C_\tau v_0$ and a preimage of
$A^\vee$; its volume is again a unit. The two determinant tensors thus
map, up to units, to the same tensor in the strict determinant, using
$A^\vee\wedge B^\vee$. All triangles and plane lattices in this
comparison are integral. Their determinant isomorphisms consequently
give Equation~\eqref{eq:h:switch-equality}, including any torsion or
localization-index contribution.
\end{proof}

\subsection{Rank reduction and nonnegative divisors}\label{sec:h:rank-reduction}

The derivative-prime reduction follows the Kolyvagin-system argument of
\citet[Lemmas~1.5.3 and~1.6.4]{Howard2012}. We carry its switches
through integral Selmer complexes and preserve their determinant
coordinate. The local calculation and descent above supply the
characteristic-two pairing and one fixed multiplier; the argument
below supplies the uniformity as the conductor varies.

We now prove the first two bounds in Equation~\eqref{eq:h:divisor-bounds}.
A height-one specialization can have larger Selmer dimension than the
generic rank-one line. Local switches will remove that excess dimension
while preserving the determinant valuation. At test dimension one,
integrality is immediate. At the coefficient prime $(2)$ this argument
uses residual irreducibility and therefore applies here to genus data.

Use the sequence groups $\Gamma_K\subset\Gamma$ of
Lemma~\ref{rl:evaluation}, and write $\mathcal H$ for the exact
kernel of Lemma~\ref{rl:kernel}. Over the relevant test field, put
$W=\operatorname{Ind}_{\Gamma_K}^{\Gamma}(T\chi_e\Psi)$.
Lemma~\ref{rl:chebotarev-evaluation} realizes finite localization as
the projection of $(1+\tau)b(u')$, for $u'\in\mathcal H$, to one
induced summand. Restriction to $\mathcal H$ is injective by
Lemma~\ref{rl:kernel}, and evaluation there is additive and
Galois-equivariant. The image of $1+\tau$ is a copy of the
two-dimensional summand, also in characteristic two.

\begin{lemma}[Rank reduction at a height-one prime]
\label{lem:h:height-one-reduction}
Suppose a primitive instance has generic cohomology of rank at most one
in degrees one and two, zero elsewhere, and $P$ spans $H^1$ when it is
nonzero. Then the first line of Equation~\eqref{eq:h:divisor-bounds}
holds. These hypotheses hold in every simple-zero instance and also at
the base of Theorem~\ref{thm:heegner-interpolation}. The second line holds in
genus data with irreducible $E[2]$, with a bound independent of the
number of active or derivative primes.
\end{lemma}

\begin{proof}
There is nothing to prove if $u=0$. Otherwise the generic rank is one
and $P$ is nonzero. Localize, and if convenient complete, at a
height-one prime. Test its
residue field; at $(2)$ do this argument only in genus data with
irreducible residual representation. Suppose first that $\Psi^2\ne1$
on the field. Then $W$ is irreducible. Inertia at $r$ distinguishes its
two induced summands. On one summand, irreducibility of $T$ over $K$
comes from open image or the CM normalizer in characteristic zero, and
from the assumed residual irreducibility at $(2)$. In the latter case
we have omitted the fields meeting $\Q(E[2])$ nontrivially. The cyclic
cubic residual case remains irreducible over $\F_2((t))$, with a
possibly quadratic commuting field. Constant Tate matrix images are
realized on the sequence group; scalar character twists do not change
this irreducibility.

Let the test Selmer $H^1$ have dimension $s\ge2$. Include the reduction of
a primitive generator of the generic free $H^1$ lattice among two
independent classes. We can arrange that their localization has rank
two. To see this, if the classes are independent over the commuting
field of $W$, simplicity and the injective restriction in
Lemma~\ref{rl:kernel} imply that the field span of their joint
evaluations on $\mathcal H$ is the whole of $W\oplus W$. Otherwise
a nontrivial equivariant linear relation would vanish on restriction,
and hence on the original classes. Projection by $1+\tau$ therefore
spans all pairs in the two-dimensional localization space. Some actual
evaluation pair has nonzero determinant: a homogeneous quadratic
vanishing on an additive group vanishes on its field span, since its
polar form vanishes on every pair from the group. This argument is
valid in characteristic two as well. If the classes are dependent over
the commuting field but independent over the ground field, projection
commutes with that field action on the graph of $\tau$. At the residual
prime~$(2)$, in the cyclic cubic case, this commuting field is
$\F_4((t))$: rational complex
conjugation has trivial residual image, so the commuting-field action
is the same on the two induced summands. If the second class is $\alpha$ times the first, with
$\alpha$ outside $\F_2((t))$, a nonzero projected evaluation $v$ gives
the ground-field-independent pair $v,\alpha v$. Thus this case also
gives rank-two localization.

Choose the derivative prime realizing that evaluation. Its localization
is nonzero on the primitive generic generator. The relaxed image on the
test fiber is self-annihilating of dimension two, and contains all of
$F$ because the finite localization has rank two. It is therefore $F$.
Switching to $S$ lowers test $H^1$ dimension by two. Generically,
Lemmas~\ref{lem:h:derivative-classes} and
\ref{lem:h:determinant-switch} keep rank one, a nonzero class, and the
same determinant valuation.

It remains to treat $\Psi^2=1$, which occurs here only in
characteristic zero. Extend $\Psi$ across $\tau$ and split into the
two conjugation signs. Conjugation on the cochains includes
$1+t\mapsto(1+t)^{-1}$, which is trivial on this fiber. The fixed
Kummer conditions and relaxed conditions are preserved. At previously
switched places the split local model has diagonal conjugation action
on $f,s$, so its plane conditions are preserved too. The induced
representation now splits into two nonisomorphic absolutely irreducible
twists. Finite and singular localization each have one line per sign.
Within a sign the pairing is a symmetric hyperbolic pairing of those
two isotropic lines; the signs are orthogonal.

If both Selmer signs occur, the restriction argument just used gives
rank-two localization by choosing one class of each sign, with a
nonzero component of the primitive generator among them. If only one
sign occurs and $s>1$, choose rank-one localization nonzero on that
generator. The relaxed image has a self-annihilating line in each
sign. In the sign hit this is the finite line; in the other it must
be the singular line, since a symmetric hyperbolic plane in
characteristic zero has exactly these two isotropic lines. The switch
preserves the total dimension but introduces the other sign: the old
strict kernel has dimension $s-1\ge1$ and the new singular line
contributes the second sign. A subsequent rank-two switch reduces
dimension. Iteration reaches test rank one; it cannot pass below one,
because the generic rank-one class has been retained throughout.

At test rank one the minimal self-dual complex over the DVR has one
free term in each of degrees one and two and zero differential; there
are no invariants. An integral class and its integral dual functional
therefore give an integral determinant coordinate. Reversing the
switches gives the first assertion. At $(2)$ the initial fixed descent
and local multiples cost a bounded valuation, independent of every
subsequent switch. This proves the second assertion with the stated
uniformity.
\end{proof}

\subsection{Uniformity without residual irreducibility}
\label{sec:h:uniform-ring-bound}

Only the ring-data bound at $(2)$ remains. Residual irreducibility is
not assumed, so the preceding argument on that residue field does not
cover every instance. We first bound the model sizes and retain evaluation through a second
limit. These constructions turn a hypothetical
sequence of unbounded negative determinant valuations into a
characteristic-zero outer limit. There the retained evaluation maps
permit rank reduction and force a bounded denominator, giving a
contradiction.

\begin{lemma}[Uniform residual sizes]\label{rl:residual-bound}
Suppose $\omega(\operatorname{disc}K)\leq A$, the local twisting classes
at $2N$ belong to a fixed finite list, and all active conductor primes
are split with odd Frobenius exponent of the deformation character.  After localization and completion at
$(2)$, the minimal free ranks of the Selmer complexes are bounded in
terms of $E,A$, and the number of retained derivative primes.  They do
not depend on the number of active conductor primes.
\end{lemma}

\begin{proof}
Reduce first over $\mathbf F_2[[t]]$.  The quadratic character is now
trivial.  Removing the active conductor primes from the permitted
ramification set changes the global complex by their singular local
complexes.  These, and their unramified local complexes, are acyclic
after inverting $t$, by Lemma~\ref{rl:local-acyclic}.  The resulting
smaller-support complex, including its comparison maps, can be
constructed at the original Artin stages before taking the limit.
Its support consists of $2N,r$, and the retained derivative primes.
Its $\mathbf F_2((t))$ cohomology dimensions are bounded by those at
$t=0$: choose a minimal finite free model over $\mathbf F_2[[t]]$.

Here is a uniform bound for the latter residual cohomology.  Put
$F=\mathbf Q(E[2])$, a fixed field.  Then $FK/F$ is at most quadratic.
Enlarge the allowed set of primes of $F$ by those ramified in $K$.
Its cardinality is bounded by a constant depending on $E,A$ and the
number of derivative primes.  Kummer theory bounds the number of
generators of the maximal pro-$2$ Galois group over $F$ unramified
outside this set, using the fixed unit group, the fixed $2$-class group,
and the allowed valuations.  The subgroup defining $FK$ has index at
most two; the elementary Schreier bound bounds its generator number.
Every finite Galois $2$-extension of $FK$ with this ramification has
Galois closure of $2$-power degree over $F$, so this subgroup indeed
controls the required pro-$2$ extensions of $FK$.

Over $FK$ the residual Tate module is trivial.  Restriction and
inflation--restriction therefore bound the residual $H^1$ over $K$.
The $H^0$ bound is immediate, and the global Euler characteristic or
global duality bounds $H^2$.  The local terms at the remaining bounded
set have bounded dimensions, with degrees of the local fields at most
two over the fixed rational local fields.  At $2N$ there are only
finitely many local quadratic twisting classes.  The same bounds
apply to the Selmer cone.  Lemma~\ref{rl:finite-model} now bounds the
minimal ranks. Determinant valuations will require a separate argument.
\end{proof}

\subsubsection{Evaluation through the outer limit}

The preceding evaluation argument concerns one inner construction over a
fixed $K$. The ring bound must also be uniform when $K$ and the conductor
vary. When all active primes split and $\omega(k)$ is bounded,
Lemma~\ref{rl:residual-bound} bounds the model sizes after any fixed
number of derivative operations. The next proposition shows that passing to a second
limit of such models preserves detection, including classes represented
only by approximate cycles before that limit.

\begin{proposition}[Integral evaluation through the outer limit]
\label{rl:outer-evaluation}
Let $C_j$ be any family of the preceding limiting complexes, completed
over $\mathcal O_j=\widehat{\Lambda_{(2)}}$, possibly for different
quadratic fields $K_j$.  Suppose their minimal ranks are bounded.
Assume a common bounded cohomological degree interval. Keep the
actual evaluation matrices and transport them to the minimal models
through integral contractions.  For a nonprincipal
ultrafilter $\mathcal U$, put
\[
 \mathcal O_\infty=
 \left(\prod_{\mathcal U}\mathcal O_j\right)/\mathfrak a,
 \qquad
 \mathfrak a=\{(x_j):
    \{j:v_2(x_j)\geq n\}\in\mathcal U\text{ for every }n\}.
\]
Then $\mathcal O_\infty$ is a DVR with uniformizer $2$ and a
characteristic-zero fraction field $L_\infty$. The outer $H^1$ over
$L_\infty$ injects into crossed classes on the inner quadratic sequence
group product, or its extension by one simultaneous complex conjugation
for the induced representation.  Its restriction to the product
of their exact common kernels is injective.  All arbitrary-sequence
degree-two evaluation identities survive.  These conclusions also
hold after retaining any fixed finite collection of additional
integral diagrams.
\end{proposition}

\begin{proof}
A nonzero element of the quotient has a representative whose
valuations are bounded on an ultrafilter-large set.  A finite
partition then makes its valuation equal to a single integer $n$
on such a set.  It is $2^n$ times a unit.  Products add these finite
valuations, so the quotient is a domain, and every nonzero ideal has
an element of smallest valuation.  It is consequently a DVR.  The
integer $2^n$ is nonzero for every $n$, so its fraction field has
characteristic zero.

For each $j$, all-sequence evaluation maps originally have entries
in $\Lambda$.  Extension to $\mathcal O_j$ and composition with an
integral contraction still give integral matrices.  Thus evaluation
on any sequence of inner group elements has an entrywise image in
$\mathcal O_\infty$.  Equations~\eqref{rl:eval-boundary} and
\eqref{rl:eval-cocycle} survive for all such sequences.  In particular,
if an outer cycle is represented by vectors $x_j$ for which
$d_jx_j$ has valuation tending to infinity, the evaluation of
$d_jx_j$ does too: the degree-two evaluation matrix is integral.
Thus an outer cycle gives an exact crossed cocycle even when its
representatives are only approximate inner cycles.

Over the inner residue field $\Omega_j$, choose a detecting list as
in Lemma~\ref{rl:evaluation}; its length is at most
$\dim_{\Omega_j}H^1(C_j\otimes\Omega_j)+\operatorname{rank}M$.
These lengths are bounded.  Include in the integral diagrams the
evaluation map to its two-term complex and a contraction of its cone
to a complex starting in degree one.  Such a contraction uses only
unit pivots over $\mathcal O_j$ and has bounded graded sizes.
All its matrices, including inverses and homotopies, are integral
and pass to the quotient.  Its cone still starts in degree one
after extension to $L_\infty$.  This proves injectivity just as in
Lemma~\ref{rl:evaluation}.

Let $H_j$ be the inner exact kernel. In the product of the quadratic
sequence groups, the central elements of Lemma~\ref{rl:kernel} give
a central element modulo $\prod_jH_j$ with scalar $u^4$. Since
$u^4-1$ remains nonzero in $\mathcal O_\infty$, the same calculation
gives restriction injectivity. For the induced representation first
restrict to the quadratic subgroup, as in Lemma~\ref{rl:kernel}.
Additional finite diagrams pass through the same entrywise operations.
\end{proof}

\begin{lemma}[Irreducibility in the outer limit]
\label{rl:outer-irreducibility}
On the outer characteristic-zero test, $\Psi$ has infinite order and
the four-dimensional representation induced from $T\chi\Psi$ is
absolutely irreducible.
\end{lemma}

\begin{proof}
The inner irreducibility argument requires a uniform finite
witness to survive this limit.  The images of the quadratic subgroups have index at most
two in the fixed compact $2$-adic Tate image.  A compact $2$-adic
analytic group is topologically finitely generated and hence has
only finitely many closed subgroups of index at most two.  Choose
one on an ultrafilter-large set.  Its absolutely irreducible
representation has finitely many matrices spanning
$M_2(\mathbf Q_2)$; the determinant expressing this spanning
property is a fixed nonzero $2$-adic number.  Every one of these
Tate matrices is realized by choices of Galois elements in each
component.  Their accompanying character scalars are units and
do not change the spanning assertion.  Its determinant remains
nonzero in $L_\infty$.  This proves absolute irreducibility in
that field.  Also $1+t$ has infinite order already in every
$\Omega_j$: for a nonzero integer $n$, the first nonzero term of
$(1+t)^n-1$ occurs in finite degree.  It remains a unit modulo
$2$ in every component, so no such power becomes one in the
outer field.

The two induced summands therefore have different inertia scalars,
$\Psi$ and $\Psi^{-1}$, and the induction is absolutely irreducible
as in Lemma~\ref{rl:irreducibility}.
\end{proof}

\begin{lemma}[Prime realization through the outer limit]
\label{rl:outer-chebotarev}
The finite Chebotarev realization of Lemma~\ref{rl:chebotarev-evaluation}
is compatible with Proposition~\ref{rl:outer-evaluation} and with any
fixed finite list of previous derivative operations.
\end{lemma}

\begin{proof}
Apply Lemma~\ref{rl:chebotarev-evaluation} at each finite stage.
For an outer sequence do this separately
for each inner sequence, retaining the relevant finite evaluation
matrices and degree-two identities before taking the outer
quotient.  Approximate inner cycles cause no problem, by the
integrality calculation in Proposition~\ref{rl:outer-evaluation}.
Previously imposed finite diagrams are unchanged.
\end{proof}

\begin{theorem}[Compatibility of the limit constructions]
\label{rl:limit-evaluation-package}
The finite free models, ring-class characters, local conditions and
pairings constructed above can be retained simultaneously. Their
inner limits have the stated height-one acyclicity and all-sequence
cohomology detection. For all-split ring data with bounded
$\omega(\operatorname{disc}K)$, the residual model sizes are bounded
independently of the active conductor support. Their outer limits
retain arbitrary evaluations, exact-kernel injectivity and finite
Chebotarev realization after any fixed number of additional
integral arithmetic operations.
\end{theorem}

\begin{proof}
Use Lemma~\ref{rl:finite-model} at each finite stage and keep the
same contractions and ultrafilter as in Lemma~\ref{rl:finite-limits}.
Perform the finite ring-class transfer and weighting first, retaining
only the resulting bounded target cycles and witnesses as in that
lemma; no limit of the growing Shapiro source is taken.
The local assertions and uniform residual sizes are
Lemmas~\ref{rl:kummer}, \ref{rl:local-acyclic} and
\ref{rl:residual-bound}; the required characters come from
Proposition~\ref{rl:characters}. Lemmas~\ref{rl:evaluation} and
\ref{rl:kernel} give inner detection. Proposition~\ref{rl:outer-evaluation}
and Lemma~\ref{rl:outer-chebotarev} preserve it through the
outer limit and later finite operations. Every retained identity is
an identity of actual integral matrices. The finite derivative classes and local identities of
Lemma~\ref{lem:h:derivative-classes} are among the integral arithmetic
target diagrams retained in these operations, including the central
norm relation before and after sign evaluation.
\end{proof}

\begin{lemma}[Uniform determinant bound in ring data]
\label{lem:h:uniform-ring-determinant}
For simple-zero ring data with every active conductor prime split in
$K$ and $\omega(k)\le a$, the coordinate of
Definition~\ref{def:h:determinant-coordinate} satisfies
\[
 v_{(2)}(u)\ge -C_{E,a},
\]
where $C_{E,a}$ is independent of $h$ and $K$.
There is no residual irreducibility assumption. This is the last
line of Equation~\eqref{eq:h:divisor-bounds}.
\end{lemma}

\begin{proof}
Put $O=\widehat{\Lambda_{(2)}}$. By
Lemma~\ref{rl:residual-bound}, after any fixed number of derivative
switches the minimal $O$-models have bounded size in terms of $E$, $a$
and that number, independently of $h$. The full local complexes at
$r$ and the active split primes are contractible over $O$, and the
switched planes are integrally self-dual. Thus the class and its paired
closed functional have a common bounded clearing exponent, supplied by
the fixed descent multiplier and the retained global duality map.

Suppose the asserted lower bound fails. Choose instances with
$v_O(u)\to-\infty$. Cancel disks over $O$ and take an outer limit of
their bounded free models. Transport the class, closed functional,
evaluations and the retained comparison and switch maps through these
integral contractions, applying the same contractions to their boundary
witnesses. Only the resulting bounded target vectors and matrices enter
the outer diagram; the eliminated list of active-place complexes does
not. Concretely, take the ultraproduct of their
valuation rings and quotient by the ideal of sequences whose valuation
tends to infinity. The resulting ring $\mathcal O_\infty$ is a DVR
with uniformizer $2$: a nonzero sequence has a finite valuation on an
ultrafilter-large set and is that power of $2$ times a unit. Its
fraction field has characteristic zero. All integral matrix identities,
duality diagrams, and the local switch diagrams pass to this limit.

Proposition~\ref{rl:outer-evaluation} gives injective evaluation
on the outer sequence group and on its exact common kernel. This
includes outer cycles represented by approximate inner cycles, since
their coboundaries are evaluated by integral degree-two matrices.
The fixed central scalar remains different from one in the outer
field. Lemma~\ref{rl:outer-irreducibility} gives absolute irreducibility
of the induced representation there: the scalar $1+t$ has infinite
order and distinguishes the two summands.

Each individual minimal $O$-complex is a square complex in degrees
one and two, by self-duality and the absence of residual invariants.
Choose a primitive vector in its generic kernel. Such vectors survive
the outer limit, so outer $H^1$ has rank at least one. If its rank exceeds
one, choose rank-two localization by the absolutely irreducible
evaluation argument in Lemma~\ref{lem:h:height-one-reduction}, making
it nonzero on one of these primitive vectors. The inner Chebotarev construction realizes the corresponding
sequences of derivative primes. At each original Artin stage retain
the old contractions and evaluations, adjoin the new local and
derivative diagrams, and pass through the same inner and outer
ultrafilters. Slower cofinal coefficient precision is permitted;
independent reconstruction of the old diagrams is not. For an ultrafilter-large set of instances the localization is
generically nonzero. Lemma~\ref{lem:h:determinant-switch} preserves
the valuations tending to minus infinity, while the outer rank falls
by two. The new complexes and integral classes obey the same bounds,
with only the now fixed finite number of added derivative primes.

The process must reach outer rank one: ranks are finite and cannot
fall below one while the primitive kernel vector is retained. At this
stage some maximal boundary minor is nonzero over the outer fraction
field. Its valuation is consequently bounded in the individual DVRs
on an ultrafilter-large set. Split that minor over their fraction
fields. The entries of the class and duality maps are integral and
their common initial multipliers are bounded. The formula for their
determinant coordinate, or the elementary-divisor decomposition, then
has a lower valuation bound depending on this minor and the bounded
complex size. The Heegner classes themselves need not survive the outer
limit: a numerator tending to zero only increases its valuation. It is
the primitive kernel vector that preserves rank, and the boundary minor
that bounds the denominators. This contradicts $v_O(u)\to-\infty$, and
proves the uniform estimate.
\end{proof}

\begin{proof}[Proof of Theorems~\ref{thm:heegner-genus-bound} and
\ref{thm:heegner-ring-bound}]
By Lemmas~\ref{lem:h:height-one-reduction} and
\ref{lem:h:uniform-ring-determinant}, the element $u$ has nonnegative
valuation away from $(2)$ and a uniformly bounded negative valuation
at $(2)$. Normality of $\Lambda$ implies that a fixed power of $2$
times $u$ belongs to $\Lambda$. Its central value therefore has the
same lower valuation bound.

In genus data, $s(h_*)$ and the weight of $h_*$ are fixed, $c=1$,
and $k=hh_*$. Equations~\eqref{eq:h:central-determinant} and
\eqref{eq:h:weil-restriction-length} give
Equation~\eqref{eq:h:genus-bound}. The excluded intersection fields
are finite in number, since an irreducible two-division field has
at most one quadratic subfield. In ring data, the sum over $q\mid k$
is bounded by $2\omega(k)$, whereas
$2\sum_{q\mid h}t_q=4w(h)$. The same equations give
Equation~\eqref{eq:h:ring-bound}.
\end{proof}

\subsection{The universal Heegner class}
\label{sec:h:universal-class}

For interpolation we combine the non-$(2)$ divisor estimate and the
central valuation formula with a universal binary family. The
characterwise coordinates must have a uniformly controlled clearing
factor. We first remove the
norm and local determinant factors belonging to primes unused at a
character, then eliminate the new-prime blocks before choosing that
clearing factor.

Recall the binary coefficient rings:
\[
 G\simeq(\Z/2\Z)^b,\qquad R=\Lambda[G],\qquad
 O=\widehat{\Lambda_{(2)}}.
\]
We work with the universal quadratic character over $R$ for the
family in Theorem~\ref{thm:heegner-interpolation}. At a new prime $q$
let $g_q\in G$ be its inertia bit: at a character $x$, it evaluates
to $-1$ when $q$ is active and to $1$ when it is unused. Choose the
conductor character of Proposition~\ref{rl:characters} with odd
exponents at all the new split primes. In the universal complex relax
at every new prime and at the old active primes, and take the trace
class $p$ from the conductor containing all new primes, together with
$r$ and the old conductor. As in the proof of
Lemma~\ref{lem:h:derivative-classes}, form its $A_i[G]$-weighted
target cycle at each finite stage, together with its local lifts and
boundary witnesses, before evaluating any sign character. Apply the
limit construction for this fixed $G$ to those target data; the source
ring-class cochains do not enter the limiting diagram. Choose an integral determinant basis and
write $u_{\mathrm{rel}}$ for the characterwise determinant coordinate.

\begin{lemma}[Removing the unused-prime factors]
\label{lem:h:unused-prime-correction}
There is $\eta\in O[G]$ such that
$u=\eta u_{\mathrm{rel}}$ agrees at each character with the primitive
determinant coordinate up to a $\Lambda$-unit. In particular,
$u_x\in\Lambda[1/2]$. If the base product derivative is zero, then
$u_0(0)=0$. At the nonzero vertices,
\begin{equation}\label{eq:h:normalized-vertex-bound}
 v(u_x(0))\le B+O_{E,h_0,a}(1),
\end{equation}
where $a$ bounds $\omega(k)$.
\end{lemma}

\begin{proof}
At an unused prime, the norm relation multiplies the primitive point
class by
\[
 e_q=\begin{cases}
 a_q,&q\text{ inert in }K,\\
 a_q-z-z^{-1},&q\text{ split in }K,
 \end{cases}
\]
where $z=\chi_{h_0}\Psi$ times the universal quadratic scalar at an
oriented Frobenius. The relaxed singular quotient contributes
\[
 l_q=\prod_{v\mid q}
  \det\left(1-q_v^{-1}\eta_v(\Fr_v)\Fr_v\mid T\right),
\]
where $q_v$ is the residue cardinality and $\eta_v$ denotes the full
scalar character. Thus $u_{\mathrm{rel}}$ differs from the primitive
coordinate by $e_q^2/l_q$, up to a $\Lambda$-unit. This follows from
the localization triangle with its integral determinant isomorphism;
the pairing comparison is made over the fraction field. Choose
conjugate-compatible Frobenius lifts in these formulas, taking the
square of rational Frobenius at an inert place, and extend the
expressions universally over $O[G]$.

Both $e_q$ and $l_q$ are units in $O[G]$. At a split prime, the odd
Frobenius exponent gives a nonconstant scalar in $\F_2((t))$; at an
inert new prime, the hypothesis $t_q=0$ gives an odd trace. Moreover,
already in the group ring,
$l_q\equiv e_q^2\pmod2$. In the split case the residual identity is
\[
 (1+a_qz+z^2)(1+a_qz^{-1}+z^{-2})
       =a_q^2+z^2+z^{-2}=(a_q+z+z^{-1})^2.
\]
In the inert case the Frobenius-square polynomial gives $a_q^2$.
Here the group-like quadratic scalars have square one, so the identities
hold before sign evaluation. Consequently
\begin{equation}\label{eq:h:group-ring-correction}
 \eta=\prod_{q\ \mathrm{new}}
  \left(1+\frac{l_q/e_q^2-1}{2}(1+g_q)\right)
       \quad\text{belongs to }O[G].
\end{equation}
Its factor is $l_q/e_q^2$ at unused characters and $1$ at active
characters. It makes the required correction without an idempotent
denominator for each prime.

For all the primitive constituents, generic cohomology has rank at
most one. At nonzero vertices this follows from the simple-zero
assumption. At the base it follows from Selmer corank one; a nonzero
generic class then spans the possible rank-one group. Hence
Lemma~\ref{lem:h:height-one-reduction} applies to every nonzero
primitive determinant coordinate. Its nonnegative valuations away
from $(2)$ imply membership in $\Lambda[1/2]$, by normality. The
same assertion is automatic for a zero coordinate.

If the base derivative vanishes, Equation~\eqref{eq:h:height} makes
the specialized point torsion. Either its generic determinant
coordinate is zero or the specialization is at constant rational
rank one and its degree-one class is zero there. In either case
$u_0(0)=0$. At a nonzero vertex use
Equations~\eqref{eq:h:central-determinant} and
\eqref{eq:h:weil-restriction-length}. Subtracting the two Sha lengths
and $4w(h_x)$ leaves precisely the expression bounded in
Equation~\eqref{eq:h:interpolation-hypothesis}, together with bounded
terms at $k$, the old support and $r$. This gives
Equation~\eqref{eq:h:normalized-vertex-bound}.
\end{proof}

The corrected coordinates have the required upper bounds at nonzero
vertices and vanish at the center if the base derivative vanishes.
To apply binary transfer, we still need $f,d\in O[G]$ with $d_x=f_xu_x$ at every character.
The required controls are a bound on $v(f_0(0))$ and, at each fixed
precision, a bound on the negative Laurent support of every $f_x$. The next lemma obtains those bounds from the
old complex, whose size is bounded, rather than from the full growing
conductor.

\begin{lemma}[A uniform clearing factor for the Heegner family]
\label{lem:h:family-clearing}
For the element $u$ of Lemma~\ref{lem:h:unused-prime-correction},
there are $f,d\in O[G]$ with $d_x=f_xu_x$ at every character,
$f_0\in\Lambda$, and $v(f_0(0))$ uniformly bounded. For each fixed
$M$, the negative Laurent support of every $f_x\bmod2^M$ has a
bound independent of $b$ and of the number of new primes.
\end{lemma}

\begin{proof}
First make the pairing functional integral up to a bounded multiplier.
The class $p$ is integral. Lift $2^{a_0}p$ to the complex strict at
the relaxed places and use integral duality against the relaxed
complex. At $r$ and every new prime the full local complex over
$O[G]$ is contractible: its reduction at the maximal ideal is acyclic,
by the local calculation in Lemma~\ref{lem:h:relaxed-local} and the
root-type alternative. At the old active primes its cohomology is
killed by a bounded power of $2$, because inertia is minus identity.
The old support is fixed. Thus $a_0$ is bounded independently of the
new primes. After inverting $2$, lifting is canonical. We obtain a
closed representative of the paired functional over
$2^{-a_0}O[G]$.

At augmentation, present the complex as the old complex, with no new
primes, extended by the direct sum of their singular complexes.
Each singular complex is a bounded-size block in degrees one and two,
with differential $A_q$ and determinant $l_{q,0}$. The differential is
upper triangular with the old complex as a subcomplex. This graded
free model lifts to $R$: stabilize the finite free models, identify their
augmentations, and lift the corresponding graded changes of basis.
The lifts are invertible because the augmentation ideal lies in the
Jacobson radical. Cancel the additional lifted unit disks before
retaining the old-plus-singular grading. Their augmented cross terms
are zero, so cancellation leaves precisely the displayed augmented
model; the added disks do not enlarge its old part.
The old model has bounded size, since $h_0$ is fixed and
$\omega(k)$ is bounded.
Indeed, the fixed-field Kummer and Schreier argument in
Lemma~\ref{rl:residual-bound}, now retaining the fixed primes of $h_0$,
bounds the residual global and local dimensions uniformly for
$\omega(k)$ bounded. This size bound is separate from the central
torsion-length bound used below.

Over $O[G]$ the entire new-prime block is invertible. Schur elimination
leaves a complex $D$ of the bounded old size, with augmentation exactly
the old complex. For each $A_q$, the order of its residual determinant
at $t=0$ is bounded uniformly. At a split prime this follows from the
odd Frobenius exponent and the degree-two Frobenius polynomial; at a
root-type prime the constant term is already nonzero. Every sign
evaluation of the lifted block is its augmentation plus a matrix
divisible by $2$ with power-series entries. Therefore
Lemma~\ref{lem:h:bounded-poles} bounds the negative Laurent support of
the entries of every $D_x$ modulo each fixed $2^M$, independently of
the number of blocks.

The integral central specialization of the old complex is controlled
by the ordinary base Kummer complex, the fixed old conductor primes,
and the bounded singular contributions at $r$. Its rational $H^1$ and
$H^2$ both have dimension one. Its torsion lengths are bounded by the
finite Selmer quotient assumed in
Theorem~\ref{thm:heegner-interpolation}, together with bounded global
torsion and these local contributions. Choose its elementary-divisor
bases at the center and lift. Apply case \textup{(i)} of
Lemma~\ref{lem:h:pivot-clearing} to $D$, the projected cycle and
paired functional. The eliminated large block contributes an
$O[G]$-unit to the determinant, and the multiplier in
Equation~\eqref{eq:h:group-ring-correction} belongs to $O[G]$.
Include both in the allowed factor $\eta$ of that lemma, so neither
enters the clearing factor $f$. The lemma gives $d=fu$, a bounded
central valuation for $f_0$, and the asserted pole bounds for $f_x$.
\end{proof}

\begin{proof}[Proof of Theorem~\ref{thm:heegner-interpolation}]
The primitive ring-class product has odd functional-equation sign
under the standing Heegner conditions. Suppose the base product does
not have a simple zero. Its first derivative then vanishes. By
Lemma~\ref{lem:h:unused-prime-correction}, $u_0(0)=0$, whereas every
nonzero vertex has central determinant valuation bounded by one
constant $B_1$ independent of the family. By
Lemma~\ref{lem:h:family-clearing}, the hypotheses of
Lemma~\ref{lem:h:boolean-transfer} hold with uniform central and pole
bounds. Choose its fixed precision larger than $B_1$ plus twice the
central bound, and then take the cube dimension sufficiently large.
It supplies a nonzero vertex whose central valuation exceeds $B_1$,
a contradiction. Thus the base product has a simple zero.
\end{proof}

\section{The unrestricted converse and Goldfeld's densities}
\label{sec:conclusions}

The two residual constructions now give the pointwise theorem. The
density statement then requires only Smith's Selmer distribution and
an elementary change of counting convention.

\begin{proof}[Proof of Theorem~\ref{thm:main-converse}]
By Lemma~\ref{pr:residual-cases}, either $E(\Q)[2]\ne0$ or $E[2]$
is irreducible. Theorem~\ref{thm:rational-converse} proves
$a(E)=c_2(E)$ in the first case, and
Theorem~\ref{thm:irreducible-converse} proves it in the second.
Both use the local Kummer conditions defining the corank in the
statement. The isogeny choices in the rational-torsion proof preserve
that corank and the $L$-function by Lemma~\ref{pr:invariance}.
Thus every curve in the theorem has analytic rank zero or one.
Lemma~\ref{pr:ranks} supplies the Mordell--Weil rank equality and
finiteness of the whole Tate--Shafarevich group.
\end{proof}

\begin{proof}[Proof of Corollary~\ref{cor:finite-descent}]
Lemma~\ref{pr:finite-descent} gives $c_2(E)=d_2(E)$ and shows that
$\Sha(E/\Q)[2^\infty]$ is divisible. Theorem~\ref{thm:main-converse}
gives the rank equalities and whole-Sha finiteness. A finite divisible
$2$-group is zero, so its $2$-primary part vanishes.
\end{proof}

\begin{lemma}[From integer to squarefree parameters]\label{lem:squarefree-density}
Let $\mathcal P$ be a property of nonzero signed integers unchanged by
multiplication by an integer square. Suppose
\[
 A(X)=\#\{n\in\Z:0<|n|\le X,\ \mathcal P(n)\}
      =2\delta X+o(X).
\]
Then
\[
 S(X)=\#\{d\in\mathcal D(X):\mathcal P(d)\}
      =\frac{2\delta}{\zeta(2)}X+o(X).
\]
In particular the property has density $\delta$ among signed
squarefree parameters ordered by absolute value.
\end{lemma}

\begin{proof}
Every nonzero integer has a unique expression $n=m^2d$ with $m>0$
and $d$ signed squarefree. Thus
\[
 A(X)=\sum_{m\le\sqrt X}S(X/m^2),\qquad
 S(X)=\sum_{m\le\sqrt X}\mu(m)A(X/m^2).
\]
The second equality follows by substituting the first and using
$\sum_{m\mid r}\mu(m)=\mathbf1_{r=1}$.
For a fixed integer $M$, the terms with $m\le M$, divided by $X$,
tend to $2\delta\sum_{m\le M}\mu(m)/m^2$. The remaining absolute
contribution is at most
\[
 2X\sum_{m>M}m^{-2},
\]
because $A(Y)\le2\lfloor Y\rfloor\le2Y$.
First let $X$ and then $M$ tend to infinity. The absolutely convergent
series $\sum\mu(m)/m^2=1/\zeta(2)$ proves the formula.
Applying it to the property always true gives
$\#\mathcal D(X)\sim2X/\zeta(2)$, proving the density assertion.
\end{proof}

\begin{proof}[Proof of Theorem~\ref{thm:main-density}]
Smith's theorem states that, for every fixed elliptic curve $E/\Q$
and $j=0,1$ \citep[Theorem~1.1]{Smith2025},
\[
 \#\{n\in\Z:0<|n|\le X,\ c_2(E^{(n)})=j\}=X+o(X)
 .
\]
Quadratic twisting depends only on the squareclass, so
Lemma~\ref{lem:squarefree-density} gives density $1/2$ for each
corank among the signed squarefree parameters.

For $j=0,1$, Theorem~\ref{thm:main-converse} gives
$c_2(E^{(d)})=j\Rightarrow a(E^{(d)})=j$, while
Lemma~\ref{pr:ranks} gives the converse implication.
The analytic-rank counts therefore equal the corresponding corank
counts exactly. The two densities sum to one. Their complement,
the twists of analytic rank at least two, has density zero, and
Lemma~\ref{pr:ranks} gives the final rank and finiteness assertions
on the density-one union.
\end{proof}

\bibliographystyle{plainnat}
\bibliography{references}
\end{document}
